% =====================================================================================================
% Chương 4, PHẦN A (mở chương, mục 4.1–4.5) — bản sửa v5 (2026-10-13), BẢN DÙNG ĐỂ GHÉP.
% Sửa có mục tiêu từ drafts/08-chuong4-phanA-v4.tex theo R34-CH4A-V4-VERIFY-CODEX-v2.md và brief R36.
% Quyết định trước: reviews/CH4A-REVISION-LOG-v2.md; ghi chú cũ: drafts/08-chuong4-phanA-v2-NOTES.md.
% Bản v4 và mọi bản trước giữ nguyên. Tệp tiếp tục R36: 08-chuong4-phanA-v5-v3.tex;
% từng thay đổi, bằng chứng mới và việc bàn giao: 08-chuong4-phanA-v5-CHANGES-v3.md. Phần B v10 nối sau tệp này.
% Mỗi dòng "% EV:" trỏ tới EVIDENCE-REGISTER.md (A01…K52), các phụ lục 1 (N-xx, L1a-xx), 2 (X3-xx),
% 4 (C4A-xx), 6 (N4B8-xx), hoặc sổ bổ sung R36-A01…03 trong tệp CHANGES-v3.md bên cạnh.
% Không dùng gói mới: TikZ và array đã nạp trong main.tex; chỉ sửa độ rộng các bảng được brief chỉ ra.
% =====================================================================================================

\chapter{Thực nghiệm và đánh giá}
\label{chuong:thucnghiem}

Chương này đo luật C được gì và mất gì trên mô phỏng FleetPy, so với điều phối không trần, luật hạn chót và cách đệm
lời hứa, tức cộng thêm một khoảng cố định vào mọi lời hứa. Chương cũng xét lời hứa theo dự báo và những chỗ dự báo
chưa đủ. Các mục đi theo ba phần của độ lệch giờ trả so với lời hứa đầu (mục~\ref{muc:so}). Phần tính cho điều phối,
gắn với việc ghép thêm khách, ở mục~\ref{muc:c4-ketqua}--\ref{muc:c4-vaitro}; phần được hồ sơ tắc đường lịch sử giải
thích, cùng cách đệm lời hứa, ở mục~\ref{muc:c4-dubao} và~\ref{muc:c4-dem}; phần còn lại ở mục~\ref{muc:c4-tuyet}.
Bảng~\ref{bang:c4-lotrinh} nối ba câu hỏi nghiên cứu ở mục~\ref{muc:c1-cauhoi} với các mục trả lời chúng.
% EV: — (đoạn dẫn, không nêu số); ba phần theo eq:so của chương 3 (X3-12); "đệm" theo GLOSSARY-VI mục B

\begin{table}[htbp]
\centering
\caption[Câu hỏi nghiên cứu và các mục trả lời]{Ba câu hỏi nghiên cứu và các mục của Chương~\ref{chuong:thucnghiem} trả lời chúng}
\label{bang:c4-lotrinh}
\small
\begin{tabular}{@{}p{6.3cm}p{2.4cm}p{6.3cm}@{}}
\toprule
\textbf{Câu hỏi} & \textbf{Mục} & \textbf{Đại lượng được đo}\\
\midrule
CH1. Phần dời do điều phối được đo thế nào, và trần trên phần đó giữ được gì, không giữ được gì?
& \ref{muc:c4-dieuphoi}, \ref{muc:c4-congdon}, \ref{muc:c4-hople}
& Số khách bị dời giờ trả vì ghép thêm khách, độ dời lớn nhất; yêu cầu bị từ chối có phương án bị trần chặn;
phép cộng các lần dời qua cả chuyến; các bộ đếm kiểm tra\\
CH2. Giữ trần đổi số lượt phục vụ, số lượt không trễ quá một và hai phút, cùng độ dài lời hứa thế nào so
với điều phối không trần, luật hạn chót và cách đệm lời hứa?
& \ref{muc:c4-ketqua}, \ref{muc:c4-chontran}, \ref{muc:c4-vaitro}, \ref{muc:c4-dem}
& Lượt phục vụ; lượt trễ và không trễ so với lời hứa đầu; quyết định bị ép; các mức trần; thứ tự tiêu chí
của bộ chọn; độ dài thêm của lời hứa\\
CH3. Lời hứa theo hồ sơ tắc đường lịch sử đổi số lượt không trễ quá một phút thế nào, và phần nào còn
chưa được giải thích?
& \ref{muc:c4-dubao}, \ref{muc:c4-tuyet}
& Lượt trễ khi có và không có dự báo; độ dài thêm của lời hứa; ba phần của độ lệch trên ngày thường và chiều
tuyết\\
\bottomrule
\end{tabular}
\end{table}
% EV: — (lộ trình, không có khẳng định; câu hỏi theo OUTLINE-V2 cuối §1.2)

% -----------------------------------------------------------------------------------------------------
\section{Thiết lập thực nghiệm}
\label{muc:c4-thietlap}

\subsection{Kịch bản thử}
\label{muc:c4-kichban}

FleetPy 1.0.2 mô phỏng xe và khách trên bộ dữ liệu chuẩn Manhattan; ở mỗi lượt quyết định, bộ chuyển đổi gửi tình
trạng hiện tại cho bộ chạy RideBound và nhận lại quyết định (mục~\ref{muc:kientruc}). Nhu cầu là các chuyến taxi New
York tuần 11--18/11/2018 trong bộ dữ liệu này~\cite{engelhardt2026}. Một kịch bản gồm một ngày, một khung giờ hai
tiếng và một trong bốn mẫu nhu cầu của ngày đó, mỗi mẫu gồm 10\% số chuyến của ngày: 108 yêu cầu trong khung giờ, phục
vụ bởi 8 xe 4 chỗ (Bảng~\ref{bang:c4-thamso}).
% EV: B01; X3-15; X3-21 (FC:395-429); B02 (FIX: 2018-11-14_sample_10_3.csv); J18 (mẫu 10% của một ngày); B08 (mẫu r1–r4); B03

Tắc đường là một hệ số chung cho cả thành phố, đổi theo khung 15 phút và lấy từ dữ liệu chuyến taxi, không phải tốc độ
từng đoạn đường; dưới đây gọi là hệ số tắc đường đo được. Mô phỏng không có nhiễu ngẫu nhiên, nên mỗi cặp
kịch bản và cấu hình chạy một lần và cho một kết quả. Ba kịch bản cùng ngày, cùng mẫu có chung 3 đến 20 trong 108 yêu
cầu, nên các kịch bản không độc lập với nhau.
% EV: B06 (sigma = shockRate = 0); K49; B11 (T38:90, 93); I05

\begin{table}[htbp]
\centering
\caption[Tham số của một kịch bản thử]{Tham số của một kịch bản thử. Các giá trị là cấu hình của thực nghiệm, không phải chuẩn ngành. Tập tắc
đường cố định (mục~\ref{muc:c4-taptu}) khác ở dòng tắc đường và chạy bằng bộ chạy cũ, có giới hạn bộ sinh và bộ
giải chưa được kiểm lại.}
\label{bang:c4-thamso}
\small
\begin{tabular}{@{}p{4.0cm}p{11.0cm}@{}}
\toprule
\textbf{Tham số} & \textbf{Giá trị}\\
\midrule
Khung giờ & 7:00--9:00, 8:00--10:00 hoặc 17:00--19:00, gọi là khung 7, 8 và 17 giờ; riêng chiều tuyết
là 14:00--16:00 và 15:00--17:00\\
Yêu cầu & 108 mỗi kịch bản, lấy trong khung giờ từ một mẫu 10\% số chuyến taxi\\
Xe & 8 xe, mỗi xe 4 chỗ\\
Đồ thị & 96 nút là các điểm đi, điểm đến lấy từ dữ liệu Manhattan; mỗi cặp nút có sẵn một thời gian đi. Đây là
đồ thị rút gọn, không phải mạng đường phố\\
Ràng buộc dịch vụ & Đón trong vòng 600 giây kể từ lúc gọi; ngồi xe không quá 1,5 lần thời gian đi
thẳng (Bảng~\ref{bang:rangbuoc})\\
Thời gian dừng đón, trả & 0 giây: xe không mất thời gian khi khách lên, xuống (kiểm trên 3 lần chạy mẫu)\\
Tắc đường & Một hệ số chung cho cả thành phố theo khung 15 phút; tốc độ gốc của mọi kịch bản lấy theo hệ số lúc
00:00\\
Bộ sinh phương án & Tối đa 100 phương án mỗi xe; mỗi lượt quyết định xét tối đa 2 yêu cầu mới\\
Bộ giải & OR-Tools 9.15, hạt giống ngẫu nhiên cố định bằng 7\\
\bottomrule
\end{tabular}
\end{table}
% EV (khung giờ, yêu cầu): B08; B20; B02; C4A-06 (chiều tuyết cũng 108 yêu cầu mỗi kịch bản, INV:14-15,17-22)
% EV (xe, đồ thị, ràng buộc dịch vụ): B03; K48 (PREF:22-26: mỗi cặp điểm đi–đến là một cạnh trực tiếp)
% EV (dừng, tắc đường): B05 (chỉ 3 lần chạy mẫu, SERV:2-4); B06; B07 (tốc độ nền lấy lúc 00:00)
% EV (bộ sinh, bộ giải): X3-17; X3-32 (giới hạn 2 yêu cầu áp cho cả lô của lượt, GEN:54-96); C4A-18 (tập tắc đường cố định: bộ chạy khác, NEW:13; giới hạn chưa kiểm)

\subsection{Các tập kịch bản và cấu hình được so}
\label{muc:c4-taptu}

Chương này dùng bốn tập kịch bản (Bảng~\ref{bang:c4-taptu}). Trên tập kiểm tra chính còn có phép so hai
thứ tự chọn phương án ở mục~\ref{muc:c4-vaitro}; phép so này dùng lại cùng kịch bản, không bổ sung ngày dữ liệu. Từ các ngày 14--18/11/2018 dựng được 49 kịch bản, chia
theo cặp ngày và mẫu để ba khung giờ của một cặp nằm cùng một phía. 34 kịch bản làm tập kiểm tra chính; 15 kịch bản
còn lại không dùng để tính các kết quả phục vụ và đúng giờ của chương này. Các lần chạy trên 15 kịch bản đó chỉ tham
gia phép kiểm tính nhất quán của sổ ở mục~\ref{muc:kiemtra}. Ngày 14--16/11 là ngày thường, 17 và 18/11 là cuối tuần.
% EV: B08 (SPLIT: chia theo nhóm ngày, mẫu); C17 (15 kịch bản hiệu chỉnh khoảng dự phòng γ); B09; B21 (11/11/2018 là Chủ nhật); A26; RUN38; T38:99-106; chương 3 v4-v3:725-727 (185 lần chạy gồm 15 lần hiệu chỉnh)

\begin{table}[htbp]
\centering
\caption[Các tập kịch bản dùng trong thực nghiệm]{Bốn tập kịch bản dùng trong chương. U: điều phối không trần; C-$B$: luật C trần $B$ giây;
M$^{+}$-$B$: luật hạn chót $B$ giây; V-30: luật trần độ dời khách thấy 30 giây; FREE: ghép tự do
(Bảng~\ref{bang:luat}). Cột cuối là số lần chạy, mỗi lần một cặp kịch bản và cấu hình.}
\label{bang:c4-taptu}
\footnotesize
\setlength{\tabcolsep}{2.5pt}
\begin{tabular}{@{}p{2.5cm}p{1.9cm}p{3.1cm}p{3.0cm}p{2.6cm}p{1.9cm}@{}}
\toprule
\textbf{Tập} & \textbf{Quy mô} & \textbf{Ngày, khung giờ} & \textbf{Cấu hình dùng ở chương này} &
\textbf{Bộ chạy} & \textbf{Lần chạy}\\
\midrule
Tập kiểm tra chính & 34 kịch bản, 3.672 yêu cầu & 14--18/11/2018; khung 7, 8, 17 giờ có 9, 14, 11
kịch bản & U, C-30, C-120, M$^{+}$-30, M$^{+}$-60, V-30 & Bản đã lưu, có luật phục hồi không tệ hơn &
170 (năm cấu hình), 34 (C-120); thêm 68 ưu tiên chi phí\\
% EV: B08; B09; B10 (170 = 34 × 5); L1a-01 (34 lần chạy C-120); A19; A28; X3-10 (R120:3-8)
Tập chỉnh trần buổi sáng & 8 kịch bản, 864 yêu cầu & 12--13/11/2018; khung 7 giờ & U; C và M$^{+}$ ở 30, 60,
90, 120, 150, 180, 240 và 300 giây & Như tập chính & 136 (trong 192 lần chạy đã lưu)\\
% EV: B15; B16; D01; D04; X3-10 (RW7:1-8; RW7X:6-8); X3-28; C4A-07 (17 cấu hình × 8 = 136; 88 + 104 = 192, RW7:33-34; RW7X:35-36)
Tập tắc đường cố định & 16 kịch bản, 1.728 yêu cầu & 12--13/11/2018; khung 8 và 17 giờ; hệ số không
đổi trong lần chạy & FREE, U, C-30, C-60, C-120 & Bản cũ, không có luật phục hồi & 16 mỗi cấu hình\\
% EV: B17; D08; D09; H13 (n = 16 mỗi cấu hình); LAD:2-6 (travelSnapshotVersion {1: 16}: một bảng giờ đi); C4A-04
Chiều tuyết 15/11/2018 & 8 kịch bản, 864 yêu cầu & 15/11/2018; khung 14 và 15 giờ & C-30 &
Như tập chính & 8\\
% EV: B20; N-26 (8 lần chạy C-30, 601 khách được phục vụ); C4A-06 (864 = 8 × 108, INV:14-15,17-22); N4B8-05 (CK8:65-72); R26 §Q8; agent-numeric-output-v2.txt (manifest, cấu hình phục hồi và DLL)
\bottomrule
\end{tabular}
\end{table}

Bản đã lưu là bản bộ chạy biên dịch và lưu lại, dùng chung cho tập kiểm tra chính, tập buổi sáng, các lần chạy trần
120 giây và chiều tuyết; bản cũ chưa có luật phục hồi (mục~\ref{muc:doichung}, \ref{muc:kiemtra}). Trừ ghép tự do, mọi
cấu hình trong từng thứ tự ưu tiên dùng cùng bộ sinh, bộ kiểm và bộ chọn với luật C; điều phối không trần chỉ
khác luật C trần 30 giây ở trần giờ trả. Các mục trước mục~\ref{muc:c4-vaitro}, cùng các phép so dự báo và đệm
sau đó, dùng thứ tự ưu tiên ổn định nếu không ghi khác. Ghép tự do không khóa gì và không qua bộ kiểm cam kết.
% EV: A28; X3-10 (RW7:1-8; RW7X:6-8; R120:3-8); N4B8-05 (CK8:65-72); R26 §Q8; B17; NEW:13 (bộ chạy cũ không có commitmentRecovery); A08; H01; A17; K50; X3-38 (HO:109, b1-rolling-cost)

Lần chạy hỏng vì lỗi hạ tầng (một lần ở tập kiểm tra chính) được thay bằng bản chạy lại hoàn tất (Phụ
lục~\ref{phuluc:C}). Bản kiểm kê ngày 08/10/2026 ghi 1.991 bản ghi lần chạy; kết quả của chương chỉ lấy từ bốn tập
trên và 68 lượt ưu tiên chi phí; các đợt chạy khác chỉ được nhắc khi nêu phạm vi kiểm tra (mục~\ref{muc:c4-congdon}).
% EV: C4A-17 (T38:12; ADD1:6: C-30-d20181114-s10-r3-w17-rerun1); ADDENDUM-1 mục 0.3; B19; C4A-20 (INV:1, bản kiểm kê ngày 08/10)

\subsection{Chỉ số}
\label{muc:c4-chiso}

Một lượt phục vụ là một yêu cầu, tức một khách, được nhận và chở tới nơi. Lượt trễ và lượt tới sớm được đếm trong nhóm
khách được phục vụ, so với lời hứa đầu (Bảng~\ref{bang:c4-chiso}). Số lượt không trễ bằng số lượt phục vụ trừ số lượt
trễ, tính trên toàn bộ yêu cầu, nên yêu cầu bị từ chối không được tính là đúng giờ. Ngưỡng chính là 60 giây; ngưỡng
120 giây được báo kèm vì có phép so đổi kết luận theo ngưỡng (mục~\ref{muc:c4-ketqua}).
% EV: C01; N-07 (khách được nhận = khách được phục vụ); C06; C08 (DIALPY:8-9: bị từ chối không tính là kịp; ở 120 s U hơn C-30 5 lượt)

Lời hứa đầu là giờ trả của kế hoạch lúc khách xác nhận đặt xe, tính với thời gian đi lúc đó. Tính từ lúc gọi xe, độ dài
trung bình của lời hứa là 14,15 phút trên 2.434 khách khi không trần và 14,12 phút trên 2.338 khách dưới luật C trần 30
giây, gộp theo khách được phục vụ. Hai giá trị trung bình gần nhau trên hai tập khách khác nhau; phép so này không nói
về độ dài lời hứa của từng khách.
% EV: A13; A05; B14 (BURD:5,11); R36-A03 (p0 trừ request.arrivalTimeMs, gộp theo khách); agent-numeric-output-v2.txt

Tỉ lệ \emph{gộp theo khách} chia tổng số khách thỏa điều kiện cho tổng số khách được phục vụ; \emph{trung bình theo
kịch bản} lấy trung bình của tỉ lệ từng kịch bản, mỗi kịch bản tính như nhau. Ví dụ dựng tay: với 5 trong 10 khách ở
một kịch bản và 5 trong 90 ở kịch bản kia, gộp theo khách cho 10\%, trung bình theo kịch bản cho 27,8\%. Chênh lệch là
cấu hình được thử trừ đối chứng; số dương nghĩa là nhiều hơn.
% EV: C27 (hai cách gộp cho hai số khác nhau); K24; ví dụ dựng tay (10/100 = 10%; (50% + 5,56%)/2 = 27,78%)

Kết quả ``có dự báo'' là phép phát lại trên các lần chạy đã ghi (mục~\ref{muc:dubao}): xe nào chở ai giữ nguyên, chỉ
các giờ báo cho khách được tính lại theo dự báo, và các chỉ số dùng $D_0$. Dự báo chỉ có ở 21 kịch bản ngày thường của
tập kiểm tra chính; 13 kịch bản cuối tuần và tập buổi sáng giữ lời hứa không dự báo.
% EV: C09; F03 (chỉ vế phát lại); X3-11 (D = max(p, S) áp cho mọi lời hứa đón và trả đã công bố); F02; N-28; D02 (12–13/11 không có dự báo)

\begin{table}[htbp]
\centering
\caption[Định nghĩa các chỉ số đánh giá]{Các chỉ số dùng trong chương. $a$ là giờ trả thực tế, $p_0$ là lời hứa đầu không dự báo, $D_0$
là lời hứa đầu theo dự báo (không có dự báo thì $D_0 = p_0$); $c$ là một lần dời giờ trả do điều phối, $z$ là phần
khách thấy của một lần cập nhật lời hứa (mục~\ref{muc:bamoc}). ``Quá'' nghĩa là lớn hơn hẳn: trễ đúng 60 giây
không tính là trễ quá 60 giây.}
\label{bang:c4-chiso}
\footnotesize
\setlength{\tabcolsep}{3pt}
\begin{tabular}{@{}p{2.7cm}p{5.3cm}p{3.5cm}p{3.5cm}@{}}
\toprule
\textbf{Chỉ số} & \textbf{Định nghĩa} & \textbf{Mẫu số, cách gộp} & \textbf{Vì sao đo}\\
\midrule
Lượt phục vụ & Yêu cầu được nhận và chở tới nơi & Số đếm; chia cho số yêu cầu thì ra tỉ lệ phục vụ & Giá của
luật tính bằng chuyến\\
Trễ quá $\theta$ & Khách được phục vụ có $a - D_0 > \theta$; $\theta$ = 30, 60, 120, 300 hoặc 600 giây & Đếm trong
khách được phục vụ & Lời hứa bị vượt ở phía muộn\\
Không trễ quá $\theta$ & Lượt phục vụ trừ lượt trễ quá $\theta$ & Chia cho số yêu cầu; yêu cầu bị từ chối không
tính là đúng giờ & Phục vụ và đúng giờ trên một mẫu số\\
Tới sớm quá 60 giây & Khách được phục vụ có $a - D_0 < -60$ giây & Đếm trong khách được phục vụ & Lời
hứa quá dài cũng sai\\
Lời hứa dài thêm & $D_0 - p_0$ & Trung bình gộp theo khách được phục vụ, kể cả khách không có dự báo & Giá của
dự báo hay đệm\\
Khách đi chung & Có khách khác trên cùng xe trong khoảng từ lúc được đón tới lúc được trả & Gộp theo
khách được phục vụ & Mức ghép chuyến còn lại\\
Khách bị điều phối dời giờ trả & Có ít nhất một lần $c \neq 0$; độ dời của khách là $\sum|c|$ &
Gộp theo khách được phục vụ & Đại lượng luật C giới hạn\\
Quyết định bị ép & Lượt quyết định có ít nhất một xe bị ép: luật lời hứa đã loại cả phương án giữ nguyên kế hoạch
của xe, và xe đi theo phương án được giữ nhờ luật phục hồi, có ít nhất một khách được miễn trừ
(mục~\ref{muc:doichung}) & Chia cho số lượt quyết định đã được bộ mô phỏng xác nhận áp dụng; gộp theo
lượt quyết định, mỗi lượt đếm một lần & Luật có tự giữ được nghĩa vụ không\\
Khách bị phá lời hứa & Khách mà luật đang dùng không giữ được điều nó hứa. Luật C: $\sum|c| > B$; luật hạn chót:
$a - p_0 > B$; luật trần độ dời khách thấy: $\sum|z| > B$ & Chia cho số khách được phục vụ; gộp theo
khách & Luật có giữ được điều nó hứa không\\
\bottomrule
\end{tabular}
\end{table}
% EV (phục vụ): C01; L1a-01 (2.387 hoàn tất + 1.285 bị từ chối = 3.672); N-07
% EV (trễ, không trễ, sớm): C06; C08 (DIALPY:8-9); N-04 (trễ > 60 s và sớm < −60 s rời nhau); C4A-01 (ngưỡng 30 s, cùng quy tắc sweep_analyze.py:107,128-132); C09; C10
% EV (dài thêm, đi chung, bị dời): F06; N-06 (gộp theo khách, ext_sum / served, ADD1:14,19,24); N-09 (định nghĩa đi chung); C20; N-07; A06; X3-05
% EV (bị ép): C4A-10 (RS:1012-1027; VALID:141-145; HVA:252-300; POL:443-463); X3-30; X3-09; C13 (trung bình theo kịch bản)
% EV (phá lời hứa): N-21 (định nghĩa của M⁺-30 và C-30; số gộp); A11; GLOSSARY-VI mục B "phá lời hứa"

\subsection{Mức trần 30 giây và tính chất thăm dò của kết quả}
\label{muc:c4-thamdo}

Trần chính $B$ = 30 giây là trần giờ trả của cấu hình có trần ở tập kịch bản tĩnh của đợt thử trước, nơi thời gian đi
không đổi (Phụ lục~\ref{phuluc:A}). Mức này được giữ nguyên cho tập kiểm tra chính, không chỉnh theo kết quả của tập
này. Nó không phải ngưỡng cảm nhận của khách hay chuẩn ngành, nên mỗi mức trần chỉ là một điểm trên đường đánh đổi
(mục~\ref{muc:c4-chontran}).
% EV: C4A-15 (TEX08:136-137: C1 = trần 30 giây cho tổng độ dời giờ trả; thay X3-18 đã bỏ); K29 (tập thời gian đi tĩnh); K33; B10 (năm cấu hình cố định trước lần chạy đầu, T38:3-6); D20; K52; K36

Các kết quả là thăm dò và dùng lại dữ liệu đã xem. Ngưỡng phân tích và cách chia tập được chọn sau khi đã xem tập buổi
sáng và các đợt chạy thử trước đó. Mức thử 120 giây được chọn sau khi đã xem kết quả của điều phối không trần và luật
C trần 30 giây trên chính tập kiểm tra chính, kể cả các phép tính thử ở 120 giây (mục~\ref{muc:c4-dieuphoi},
\ref{muc:c4-congdon}). Các con số vì vậy mô tả điều đã đo, không phải phép kiểm trên dữ liệu chưa từng xem.
% EV: I06; T38:8 (chọn sau khi xem Việc A, B, w07x); C4A-11 (BUD 07/10 = D14; CVP 10/10 = E06; C120R:10 chạy 10/10; T38:5 tập chính chạy 1–2/10); K42

% -----------------------------------------------------------------------------------------------------
\section{Kết quả chính khi chưa có dự báo}
\label{muc:c4-ketqua}

Với bộ chọn ưu tiên ổn định trên 3.672 yêu cầu của tập kiểm tra chính, luật C trần 30 giây không để phần dời giờ trả do điều phối của khách
nào vượt 30 giây: lớn nhất 29,2 giây, so với 590,1 giây (gần 10 phút) khi không trần. Không lượt nào của luật C trễ
quá 5 phút so với lời hứa đầu, so với 18 lượt khi không trần. Đổi lại, luật C phục vụ 2.338 lượt, ít hơn điều phối
không trần 96 lượt, và có ít hơn 242 khách đi chung. Số lượt không trễ quá một phút tăng 31, từ 2.231 lên 2.262, nhưng
số lượt không trễ quá hai phút giảm 5, từ 2.332 xuống 2.327.
% EV: C20; N-23; N-24; C07; C01; C02; N-20; C08

Mục~\ref{muc:c4-dem} cho thấy khi có dự báo, cộng thêm 30 giây vào mọi lời hứa của điều phối không trần cũng cho 2.328
lượt không trễ quá một phút, bằng luật C có dự báo. Vì vậy khóa luận không dùng chỉ số này làm lý do chọn luật C.
% EV: C10 (U + dự báo + 30 s: 2.328 = C-30 + dự báo: 2.328); K06; K51

Luật hạn chót chỉ mượn phần hạn chót của Schulz và Pfeiffer~\cite{schulz2026}: giờ trả mới không được muộn hơn lời hứa
đầu quá 30 hoặc 60 giây, thay vì 900 giây của họ. Nó chạy cùng bộ điều phối và luật phục hồi với luật C, nên đây không
phải phép so với hệ thống của hai tác giả.
% EV: J10; J11; A09; A19; H01; X3-28; X3-41

\begin{table}[htbp]
\centering
\caption[Kết quả chính khi chưa có dự báo]{Kết quả trên tập kiểm tra chính với lời hứa đầu không dự báo (34 kịch bản, 3.672 yêu cầu mỗi
cấu hình, với bộ chọn ưu tiên ổn định). Số đếm gộp trên mọi kịch bản; trễ tính so với lời hứa đầu. Hai cột ``không trễ quá'' là số lượt phục vụ
không trễ quá ngưỡng; chia cho 3.672 thì ra tỉ lệ trên yêu cầu. ``Chưa đếm'': chưa có số đếm chính xác.}
\label{bang:c4-chinh}
\small
\setlength{\tabcolsep}{4pt}
\begin{tabular}{@{}p{\dimexpr .12\linewidth-2\tabcolsep\relax}
>{\raggedleft\arraybackslash}p{\dimexpr .13\linewidth-2\tabcolsep\relax}
>{\raggedleft\arraybackslash}p{\dimexpr .16\linewidth-2\tabcolsep\relax}
>{\raggedleft\arraybackslash}p{\dimexpr .22\linewidth-2\tabcolsep\relax}
>{\raggedleft\arraybackslash}p{\dimexpr .23\linewidth-2\tabcolsep\relax}
>{\raggedleft\arraybackslash}p{\dimexpr .14\linewidth-2\tabcolsep\relax}@{}}
\toprule
\textbf{Cấu hình} & \textbf{Phục vụ} & \textbf{Trễ quá 60 s} & \textbf{Không trễ quá 60 s} &
\textbf{Không trễ quá 120 s} & \textbf{Khách đi chung}\\
\midrule
U & 2.434 & 203 & 2.231 & 2.332 & 555\\
C-30 & 2.338 & 76 & 2.262 & 2.327 & 313\\
M$^{+}$-30 & 2.338 & 74 & 2.264 & 2.327 & 307\\
% EV (U, C-30, M⁺-30): C01; C06; C08; N-09 (555 / 313 / 307, số đếm chính xác)
M$^{+}$-60 & 2.347 & 79 & 2.268 & 2.336 & chưa đếm\\
V-30 & 2.336 & 74 & 2.262 & 2.325 & chưa đếm\\
% EV (M⁺-60, V-30): C01; C4A-02 (CK4A check: trễ quá 60 / 120 s 79 / 11 và 74 / 11; EXP:5-6 trễ quá 60 s 79 và 74; đi chung chỉ có % làm tròn ở EXP)
\bottomrule
\end{tabular}
\end{table}

Cách đọc bảng: dưới trần 30 giây, 2.338 yêu cầu được phục vụ; 76 lượt trong số đó tới muộn hơn lời hứa đầu quá 60
giây, nên 2.338 $-$ 76 = 2.262 lượt không trễ quá 60 giây.
% EV: C01; C06; C08

Cái giá trước hết là số lượt phục vụ: bớt 96 lượt nghĩa là tỉ lệ phục vụ giảm 2,6 điểm phần trăm (96 trên 3.672 yêu
cầu), tức ít hơn 3,9\% số lượt của điều phối không trần (96 trên 2.434). Mức giảm nằm ở khung 8 giờ (1.014 xuống 975)
và 17 giờ (904 xuống 852); khung 7 giờ chỉ giảm 5. Số khách đi chung giảm từ 555 xuống 313, tức từ 22,8\% xuống 13,4\%
số khách được phục vụ (gộp theo khách). Khóa luận không có bằng chứng về giá trị của các chuyến đi chung này đối với
khách hay hãng xe.
% EV: C02; C03; N-09; N-20 (lưu ý: D20, I07, I10)

Số lượt không trễ quá một phút vẫn tăng vì phục vụ bớt 96 lượt nhưng trễ bớt 127 lượt: 127 $-$ 96 = 31. Hai cấu hình
cũng phục vụ những khách khác nhau: 2.207 khách chung, 227 khách chỉ được điều phối không trần và 131 khách chỉ được
luật C phục vụ. Vì vậy 31 là chênh lệch hai tổng, không phải 31 khách cụ thể được đưa tới đúng giờ.
% EV: C06; C08; C04 (câu về chờ, ngồi xe của khách chung, C05, cắt vì ngân sách trang; xem NOTES)

So với luật hạn chót 30 giây, luật C phục vụ đúng bằng số lượt (2.338) nhưng không hoàn toàn cùng những khách (2.332
khách chung, mỗi bên 6 khách riêng), và có 76 so với 74 lượt trễ quá một phút. Luật hạn chót 60 giây và luật trần độ
dời khách thấy cũng gần như vậy (Bảng~\ref{bang:c4-chinh}). Các luật lời hứa khác nhau ở chỗ mỗi luật có giữ được điều
chính nó hứa hay không (Bảng~\ref{bang:c4-nghiavu}).
% EV: C01; C04; C06; K05; K07; C32; C4A-02

\begin{table}[htbp]
\centering
\caption[Nghĩa vụ của chính từng luật lời hứa]{Nghĩa vụ của chính từng luật trên tập kiểm tra chính (34 kịch bản). Hai cột đầu ghi số đếm gộp: số quyết
định bị ép trên số quyết định được xác nhận, và số khách bị phá lời hứa trên số khách được phục vụ. Cột cuối ghi số
kịch bản có ít nhất một khách bị phá lời hứa trên 34 kịch bản. Định nghĩa ở Bảng~\ref{bang:c4-chiso}. Ở hàng luật C,
các số 0 là số đo trên tập đã chạy, phù hợp với tính chất có điều kiện ở Nhận xét~\ref{nx:giunguyen}.}
\label{bang:c4-nghiavu}
\small
\setlength{\tabcolsep}{4pt}
\begin{tabular}{@{}p{\dimexpr .12\linewidth-2\tabcolsep\relax}
>{\raggedleft\arraybackslash}p{\dimexpr .28\linewidth-2\tabcolsep\relax}
>{\raggedleft\arraybackslash}p{\dimexpr .31\linewidth-2\tabcolsep\relax}
>{\raggedleft\arraybackslash}p{\dimexpr .29\linewidth-2\tabcolsep\relax}@{}}
\toprule
\textbf{Cấu hình} & \textbf{Quyết định bị ép / được xác nhận} &
\textbf{Khách bị phá lời hứa / được phục vụ} &
\textbf{Kịch bản có khách bị phá / tổng kịch bản}\\
\midrule
C-30 & 0 / 20.101 & 0 / 2.338 & 0 / 34\\
M$^{+}$-30 & 7.710 / 20.103 & 289 / 2.338 & 28 / 34\\
M$^{+}$-60 & 4.031 / 20.163 & 79 / 2.347 & 22 / 34\\
V-30 & 10.442 / 20.090 & 456 / 2.336 & 34 / 34\\
\bottomrule
\end{tabular}
\end{table}
% EV: R36-A01 (sổ bổ sung trong CHANGES-v3.md; 170 tệp IVC và đối chiếu PAIRS); N-21; X3-30; A22

Cách đọc bảng: luật hạn chót 30 giây có 7.710 trong 20.103 lượt quyết định bị ép và 289 trong 2.338 khách tới muộn
hơn lời hứa đầu quá 30 giây; những khách này nằm ở 28 trong 34 kịch bản.
% EV: N-21; R36-A01

Trên cả 34 kịch bản đã chạy, luật C không có quyết định bị ép và không khách nào vượt trần. Theo Nhận
xét~\ref{nx:giunguyen}, nếu giả thiết (N) đúng ở lượt đang xét và mọi khách có lời hứa của xe có $0 \le L_r \le B$,
phương án giữ nguyên không thêm phần dời do điều phối nên vẫn qua kiểm tra trần. Khi (N) đúng ở mọi lượt, trần được
giữ từ lúc xác nhận. Các số đo ở đây phù hợp với tính chất đó; nhận xét không bảo đảm giờ trả thực tế. Điều phối không
trần cũng có 0 trong 20.760 quyết định bị ép. Luật hạn chót so giờ trả mới với lời hứa đầu, nên phần thay đổi khi giữ
nguyên kế hoạch vẫn có thể làm xe bị ép (mục~\ref{muc:doichung}).
% EV: A21; A22; N-21; X3-16 (U 0 / 20.760); X3-30; A10; K05; Ch3 v4-v3:371-390; R36-A01

Tính gộp, luật hạn chót 30 giây bị ép ở 7.710 trong 20.103 lượt quyết định (38,4\%) và phá lời hứa với 289 trong 2.338
khách (12,4\%). Gánh nặng dồn vào khung 7 giờ: cả 9 kịch bản có khách bị phá lời hứa, với 181 trong 509 khách được phục
vụ và 3.447 trong 4.564 lượt quyết định bị ép. Ở khung 17 giờ, 5 trong 11 kịch bản có khách bị phá, với 8 trong 852
khách và 328 trong 7.163 lượt quyết định bị ép. Các số này được gộp theo từng khung giờ. Ở mức 60 giây, các tỉ lệ toàn
tập giảm nhưng không về 0. Luật trần độ dời khách thấy bị ép nhiều nhất, phù hợp với định nghĩa: tổng độ lớn các lần
cập nhật mà nó kiểm không nhỏ hơn độ lùi so với lời hứa đầu (công thức~\eqref{eq:dongnhat}).
% EV: N-21 (gộp 7.710/20.103 và 289/2.338); R36-A02 (cặp đếm theo khung giờ, sổ bổ sung trong CHANGES-v3.md); C14; C13; R36-A01; A11; A12; X3-33 (eq:dongnhat: Σ|z| ≥ |Σz| = |p_n − p_0|)

Bảng~\ref{bang:c4-nghiavu} không so hai chuẩn đúng giờ: luật C hứa về phần dời do điều phối, còn luật hạn chót hứa về
toàn bộ độ trễ so với lời hứa đầu, một lời hứa lớn hơn. Đo bằng chính thước của luật hạn chót, tức tới muộn hơn lời
hứa đầu quá 30 giây, luật C có 292 trong 2.338 khách (12,5\%), so với 289 khách (12,4\%) của luật hạn chót. Cả hai
luật đều có khách tới muộn hơn lời hứa đầu quá 30 giây ở 28 trong 34 kịch bản. Vậy về độ đúng giờ, hai luật gần như
ngang nhau; luật C khác ở chỗ chỉ hứa phần mà bộ điều phối kiểm soát được, nên giữ được lời hứa đó.
% EV: C13; K05; C4A-01 (C-30 292/2.338 = 12,49%, 28/34); N-21; R36-A01

Bảng~\ref{bang:c4-phanbotre} cho thấy số lượt trễ giảm ở cả bốn ngưỡng đã đo, từ 1 tới 10 phút.
% EV: N-24

\begin{table}[htbp]
\centering
\caption[Phân bố trễ theo ngưỡng khi chưa có dự báo]{Phân bố trễ so với lời hứa đầu không dự báo trên tập kiểm tra chính: số lượt trễ quá mỗi ngưỡng
(trong số khách được phục vụ) và số lượt đó tính trên 100 yêu cầu (3.672 yêu cầu mỗi cấu hình).}
\label{bang:c4-phanbotre}
\small
\begin{tabular}{@{}lrrrr@{}}
\toprule
& \multicolumn{2}{c}{\textbf{U}} & \multicolumn{2}{c}{\textbf{C-30}}\\
\cmidrule(lr){2-3}\cmidrule(lr){4-5}
\textbf{Trễ quá} & \textbf{Lượt} & \textbf{Trên 100 yêu cầu} & \textbf{Lượt} & \textbf{Trên 100 yêu cầu}\\
\midrule
1 phút & 203 & 5,53 & 76 & 2,07\\
2 phút & 102 & 2,78 & 11 & 0,30\\
5 phút & 18 & 0,49 & 0 & 0\\
10 phút & 2 & 0,05 & 0 & 0\\
\midrule
Trễ lớn nhất (s) & \multicolumn{2}{c}{660,1} & \multicolumn{2}{c}{232,3}\\
\bottomrule
\end{tabular}
\end{table}
% EV: N-24 (203/102/18/2 và 76/11/0/0; 5,53/2,78/0,49/0,05 và 2,07/0,30/0/0); C07; C28; C4A-14 (CON:4,11: +660,1 s và +232,3 s; khớp C26 "11 phút" và "3,87 phút")

Cách đọc bảng: khi không trần, 18 lượt tới muộn hơn lời hứa đầu quá 5 phút, tức 0,49 lượt trên 100 yêu cầu.
% EV: N-24; C28

Khi không trần, lượt trễ nhất muộn 660,1 giây (khoảng 11 phút); dưới trần 30 giây là 232,3 giây (3,87 phút). Trong 76
khách còn trễ quá một phút dưới luật C, 72 người không có lần dời do điều phối nào được ghi và 61 người thuộc khung 7
giờ. Lời hứa theo dự báo (mục~\ref{muc:c4-dubao}) nhắm tới phần đoán được của loại trễ này.
% EV: N-24; C4A-14; C26; C21 (4/76 từng bị dời); K39 (không viết "trễ thuần vì kẹt xe"); N-11 (trễ quá 60 s của C-30 theo khung 61 / 15 / 0)

Tách theo kịch bản, luật C không có nhiều lượt trễ quá một phút hơn điều phối không trần ở cả 34 kịch bản và có số
lượt không trễ quá một phút bằng hoặc nhiều hơn ở 29 kịch bản. Ở khung 8 và 17 giờ, điều phối không trần nhiều lượt
trễ hơn ở cả 25 kịch bản; ở khung 7 giờ, luật C ít hơn ở 3 kịch bản và bằng ở 6. Vì 34 kịch bản thuộc 5 ngày và không
độc lập, đây là phép đếm mô tả, không phải kiểm định thống kê.
% EV: N-22 (34/34 trễ ≤ U; 29/34 không trễ ≥ U; bảng theo ngày × khung ở ADD1:173-186); C30; I05; B11

% -----------------------------------------------------------------------------------------------------
\section{Phần do điều phối: luật C chặn gì và mất gì}
\label{muc:c4-dieuphoi}

Dưới luật C trần 30 giây, 33 trong 2.338 khách được phục vụ bị điều phối dời giờ trả, mỗi khách nhiều nhất 29,2 giây;
khi không trần là 174 trong 2.434 khách, một khách bị dời tổng cộng 590,1 giây (Bảng~\ref{bang:c4-doidieuphoi}). Cái
giá là số yêu cầu bị từ chối tăng từ 1.238 lên 1.334; chỉ 148 trong số đó có phương án bị trần chặn được ghi lại.
% EV: C20; N-23; C01 (1.238 = 3.672 − 2.434; CON:3); D13 (1.334; 148)

\begin{table}[htbp]
\centering
\caption[Phần dời giờ trả do điều phối trên tập kiểm tra chính]{Phần dời giờ trả do điều phối trên tập kiểm tra chính, lời hứa đầu không dự báo. Mẫu số của
dòng đầu là số khách được phục vụ (2.434, 2.338 và 2.387); độ dời của một khách là tổng $|c|$ của khách
đó. Dòng cuối lấy tổng phần dời có dấu do điều phối chia cho tổng số giây trễ so với lời hứa đầu của chính các khách
trễ quá 60 s: 23.027 trên 31.353 giây, 82 trên 6.892 giây và 4.836 trên 12.450 giây. Cột trần 120 giây được dùng
lại ở mục~\ref{muc:c4-chontran}.}
\label{bang:c4-doidieuphoi}
\small
\begin{tabular}{@{}p{7.6cm}rrr@{}}
\toprule
\textbf{Chỉ số} & \textbf{U} & \textbf{C-30} & \textbf{C-120}\\
\midrule
Khách bị điều phối dời giờ trả & 174 (7,15\%) & 33 (1,41\%) & 103 (4,32\%)\\
Khách bị dời từ hai lần trở lên & 12 & 0 & 0\\
Tổng dời lớn nhất của một khách (s) & 590,1 & 29,2 & 118,9\\
Trung vị độ lớn mỗi lần dời (s) & 108,6 & 8,3 & 62,5\\
% EV: C20; C22 (trung vị theo lần: 108,6 s trên 187 lần, 8,3 s trên 33 lần); N-07 (C-120: 103, 0 khách ≥ 2 lần, nên trung vị theo lần = trung vị theo khách 62,54 s); N-23
Khách trễ quá 60 s từng bị dời & 134 trên 203 & 4 trên 76 & 60 trên 132\\
Phần tính cho điều phối trong tổng giây trễ của các khách trễ quá 60 s (\%) & 73,4 & 1,2 & 38,8\\
% EV: C21; C24; N-23 (cột C-120: 60/132, 38,8%); C4A-03 (tử số, mẫu số: 23.027,239 / 31.353,407; 81,742 / 6.891,526; 4.836,281 / 12.450,223 s)
\bottomrule
\end{tabular}
\end{table}

Cách đọc bảng: khi không trần, 174 khách bị điều phối dời giờ trả ít nhất một lần, 12 người trong số đó bị dời từ hai
lần, và tổng dời lớn nhất của một khách là 590,1 giây.
% EV: C20

Dưới trần 30 giây, mỗi khách bị dời chỉ một lần, trung vị 8,3 giây, so với 108,6 giây trên 187 lần dời khi không trần.
Trong 33 khách bị dời, 29 người bị dời muộn đi; 4 người bị kéo sớm hơn, mỗi người đúng 1 mili giây. Khi không trần,
134 trong 203 khách trễ quá một phút từng bị dời; con số chỉ cho biết hai nhóm trùng nhau, không cho biết việc bị dời
là nguyên nhân của độ trễ. Theo sổ theo từng khách (công thức~\eqref{eq:so}), phần tính cho điều phối bằng 73,4\% tổng
số giây trễ của 203 khách đó, còn dưới trần 30 giây chỉ 1,2\%: gần như mọi giây trễ còn lại nằm ở hai phần kia của sổ.
% EV: E02; C22; C23 (29 muộn, 4 sớm); N-05 (mỗi lần dời sớm đúng −1 ms; chỉ kiểm ở U, C-30, C-120, M⁺-30; nguyên nhân chưa kiểm, ADDENDUM-1 mục 0.2; C4A-19); C21 (giao hai tập); C24 (sổ mô hình, không nhân quả); K04; X3-12

Giới hạn này không dựa vào dự báo, vì dự báo không đi vào bộ điều phối. Ở 13 kịch bản cuối tuần không có dự báo, tổng
dời lớn nhất của một khách là 29,2 giây ở trần 30 giây và 116,5 giây ở trần 120 giây. Trên 8 kịch bản chiều tuyết,
8 trong 601 khách được phục vụ bị dời, lớn nhất 25,2 giây; không ai vượt trần và không quyết định nào bị ép trong
5.067 lượt.
% EV: F03; X3-11 (D chỉ tính sau khi chạy); N-28; N-26 (thay anchor F13 của RD §8 mục 4, ADDENDUM-1 mục 0.4); N4B8-05 (CK8:65-72); R26 §Q8

Một ví dụ trên tập buổi sáng cho thấy luật C và luật hạn chót khác nhau ở đâu. Trước một quyết định, giờ trả đang hứa
của một khách đã lùi 79,0 giây so với lời hứa đầu, toàn bộ là phần thay đổi khi giữ nguyên kế hoạch. Phương án ghép
thêm một khách mới sẽ dời giờ trả dự kiến của khách đó thêm 68,3 giây do điều phối. Ở mức 90 giây, luật hạn chót cộng
hai phần, khoảng 147 giây, nên loại phương án; luật C chỉ tính 68,3 giây nên cho qua. Trên tập này hai luật tách nhau
ở đúng ba chỗ cùng kiểu, ở mức 90 hoặc 120 giây; mỗi chỗ luật C nhận thêm một khách. Đây là phép so hai luật trên
giờ trả dự kiến tại cùng một trạng thái của mô hình, không phải phép đo nhân quả.
% EV: D05 (79,0 + 68,3 s; 60,9 + 61,2 s); C4A-12 (W07XM:6,11,16: "sum x 79.0, sum c 0.0, p_last − p_first 79.0", c = 68,3 s, x = 0,1 s; W07X:54-62: "khách đang có", giai đoạn chờ đón hay trên xe không ghi); K04

Khi luật C loại một phương án vì nó làm một khách đã có lời hứa vượt trần, bộ kiểm ghi một bản ghi chặn gồm khách đó,
phần trần đã dùng, lần dời mới và trần (Bảng~\ref{bang:kiemtran}). Dưới trần 30 giây, 148 trong 1.334 yêu cầu bị từ
chối có ít nhất một phương án mang bản ghi chặn; điều phối không trần phục vụ 116 trong số đó. 1.186 yêu cầu còn lại
không có bản ghi chặn nào được ghi, và điều phối không trần cũng từ chối 1.075 yêu cầu trong số này; điều đó không
chứng minh chúng bị từ chối chỉ vì ràng buộc dịch vụ. Bản ghi chặn cũng không phải khách: 845 bản ghi nằm ở 643 yêu
cầu, và 495 yêu cầu trong số đó cuối cùng vẫn được phục vụ.
Các số đếm theo yêu cầu trong đoạn này tính theo từng cặp kịch bản và yêu cầu.
% EV: X3-05 (nội dung bản ghi chặn); D13 (BUD:2-5: 148 = 116 + 32; 1.186 = 111 + 1.075); K10; C4A-09 (BUDPY:29-44: nhóm 1.186 là phần dư không có bản ghi chặn được ghi, không phải lý do cấu trúc đã chứng minh); D16 (SF:2: 643 cặp, 495 phục vụ, 148 từ chối); E01 (845)

Với mỗi yêu cầu trong 148 yêu cầu, phép phân tích chọn bản ghi có mức vượt nhỏ nhất, tức phần vượt quá 30 giây; trung
vị là 109,6 giây, lớn nhất 524,8 giây. Chỉ ở 2 trong 148 bản ghi được chọn, khách nêu trong bản ghi đã dùng một phần
trần trước đó; ở cả 148, riêng lần dời mới đã vượt 30 giây. Với trần 120 giây, mức vượt của 59 trong 148 yêu cầu nằm
trong phần nới thêm 90 giây. Không thể kết luận 59 yêu cầu đó sẽ được phục vụ: bộ kiểm dừng ở khách vi phạm gặp trước,
và khi trần đổi thì các quyết định trước cũng đổi. Mục~\ref{muc:c4-chontran} cho kết quả chạy thật.
% EV: D14 (BUD:6-8: trung vị 109,6 s; lớn nhất 524,8 s; 59/148 có mức vượt ≤ 90 s); D19 (BUD:10; FJC:4-6); C4A-09 (BUDPY:22,35-42,51-52: mức vượt = after − limit; 2/148 đếm theo bản ghi được chọn của từng yêu cầu, không phải số khách khác nhau); K11; A23; I08

% -----------------------------------------------------------------------------------------------------
\section{Chọn mức trần}
\label{muc:c4-chontran}

Trong ba điểm đã đo trên tập kiểm tra chính, không điểm nào tốt nhất ở mọi chỉ số. Trần 30 giây có ít lượt trễ quá một
phút nhất (76). Trần 120 giây có nhiều lượt không trễ quá hai phút nhất (2.366, so với 2.327 của trần 30 giây và 2.332
khi không trần) và phục vụ 2.387 lượt, nằm giữa hai điểm kia. Chọn mức nào tùy ngưỡng mà hãng xe coi trọng, điều thực
nghiệm không trả lời được.
% EV: N-06; L1a-02; C01; C06; C08; K12; K36

Mức trần được đo trên ba tập không cộng gộp (Hình~\ref{hinh:c4-tran}). Tập tắc đường cố định cho thấy trần làm gì khi
thời gian đi không đổi, tập buổi sáng dò tám mức trên cùng 8 kịch bản, tập kiểm tra chính thêm mức 120 giây.
% EV: B17; B15; B16; X3-28; D01; L1a-01; I12

\textbf{Tập tắc đường cố định.} Bảng~\ref{bang:c4-codinh} xếp các cấu hình từ lỏng nhất (ghép tự do) tới chặt nhất
(trần 30 giây), với hệ số tắc đường cố định suốt lần chạy.
% EV: B17; LAD:2-6

\begin{table}[htbp]
\centering
\caption[Các mức trần trên tập tắc đường cố định]{Các mức trần trên tập tắc đường cố định (16 kịch bản, 1.728 yêu cầu mỗi cấu hình; bộ chạy cũ,
không có luật phục hồi). Ba cột số đầu là số đếm; hai cột cuối là số khách được phục vụ thỏa điều kiện, trong ngoặc
là tỉ lệ gộp theo khách được phục vụ (\%).}
\label{bang:c4-codinh}
\small
\begin{tabular}{@{}lrrrrr@{}}
\toprule
\textbf{Cấu hình} & \textbf{Phục vụ} & \textbf{Trễ quá 60 s} & \textbf{Không trễ quá 120 s} &
\textbf{Đi chung} & \textbf{Bị dời giờ trả}\\
\midrule
FREE & 1.629 & 274 & 1.465 & 865 (53,1) & 388 (23,8)\\
U & 1.584 & 47 & 1.557 & 346 (21,8) & 62 (3,9)\\
C-120 & 1.563 & 23 & 1.563 & 303 (19,4) & 41 (2,6)\\
C-60 & 1.540 & 0 & 1.540 & 266 (17,3) & 20 (1,3)\\
C-30 & 1.534 & 0 & 1.534 & 241 (15,7) & 6 (0,4)\\
\bottomrule
\end{tabular}
\end{table}
% EV: D09 (số đếm, WP17); D08 (tỉ lệ đi chung, bị dời; gộp theo khách); C4A-04 (đi chung 865/346/303/266/241, CK4A fixed-shared; bị dời 388/62/41/20/6, WBU:4,7,10,13,16); B17

Cách đọc bảng: với trần 120 giây, 1.563 trong 1.728 yêu cầu được phục vụ, 23 lượt trễ quá 60 giây; 303 khách (19,4\%)
đi chung và 41 khách (2,6\%) bị dời giờ trả.
% EV: D08; D09; C4A-04

Ghép tự do khác điều phối không trần ở cả cách chọn phương án lẫn việc kiểm cam kết, nên chênh lệch giữa hai cấu hình
này không chỉ do trần. Trong nhóm cùng chính sách, siết từ không trần xuống trần 30 giây làm số lượt phục vụ giảm từ
1.584 xuống 1.534 và số lượt trễ quá 60 giây từ 47 xuống 0. Giữa hai mức đã chạy 60 và 120 giây, số lượt trễ tăng từ 0
lên 23. Các con số của bộ chạy cũ này không chuyển sang tập kiểm tra chính.
% EV: A17; K50; X3-38; NEW:13; D09; D11; B17; I12

\textbf{Tập chỉnh trần buổi sáng.} Trên 864 yêu cầu, số lượt phục vụ chỉ dao động trong 6 lượt, từ 660 ở trần 30 giây
tới 666 ở trần 120--180 giây (Bảng~\ref{bang:c4-sang}).
% EV: D01; B15; B16

\begin{table}[htbp]
\centering
\caption[Các mức trần trên tập chỉnh trần buổi sáng]{Tám mức trần trên tập chỉnh trần buổi sáng (8 kịch bản khung 7 giờ ngày 12--13/11/2018, 864 yêu cầu
mỗi cấu hình; lời hứa đầu không dự báo). Số đếm gộp; hai cột ``không trễ quá'' là số lượt trên 864 yêu cầu. Các mức
có cùng mọi số đếm được gộp vào một dòng.}
\label{bang:c4-sang}
\small
\begin{tabular}{@{}lrrrr@{}}
\toprule
\textbf{Cấu hình} & \textbf{Phục vụ} & \textbf{Trễ quá 60 s} & \textbf{Không trễ quá 60 s} &
\textbf{Không trễ quá 120 s}\\
\midrule
U & 665 & 50 & 615 & 656\\
C-30 & 660 & 47 & 613 & 658\\
C-60 & 661 & 47 & 614 & 659\\
C-90 & 664 & 50 & 614 & 660\\
C-120, C-150, C-180 & 666 & 51 & 615 & 659\\
C-240, C-300 & 665 & 50 & 615 & 656\\
\bottomrule
\end{tabular}
\end{table}
% EV: D01; D02 (U, C-30, C-90, C-120: CAP:3-7; DIAL17:3-7, hai bộ đọc độc lập); C4A-05 (C-60, C-150, C-180, C-240, C-300: CAP:5,8-11; DIAL17:5,8-11); D02 lưu ý (cột dự báo bằng ngây thơ ở 12–13/11)

Cách đọc bảng: với trần 120, 150 hoặc 180 giây, 666 yêu cầu được phục vụ, 51 lượt trễ quá 60 giây, 615 lượt không trễ
quá 60 giây và 659 lượt không trễ quá 120 giây.
% EV: D01; D02; C4A-05

Trần 120 giây phục vụ nhiều hơn điều phối không trần một lượt, nên điều phối không trần không phải mức phục vụ cao
nhất có thể đạt, chỉ là một điểm đo. Từ 30 lên 120 giây, trần thêm 6 lượt phục vụ, 4 lượt trễ quá 60 giây và 2 lượt
không trễ quá 60 giây. Luật C không có quyết định bị ép trong 64 lần chạy của tám mức trần này. Ở cả tám mức, hai
luật phục vụ chênh nhau không quá 2 lượt.
% EV: D01; K12; D03 (+6 / +4 / +2); D04; X3-28 (tám mức 30–300 s); E01 (64 lần C, 0 bị ép); agent-numeric-output-v2.txt (0 / 44.270 lượt quyết định)

\textbf{Tập kiểm tra chính.} Bảng~\ref{bang:c4-tran120} đặt trần 120 giây cạnh hai cấu hình của
mục~\ref{muc:c4-ketqua}, khi không có và khi có dự báo.

\begin{table}[htbp]
\centering
\caption[Trần 120 giây trên tập kiểm tra chính]{Trần 120 giây trên tập kiểm tra chính (34 kịch bản, 3.672 yêu cầu mỗi cấu hình). Số đếm gộp; các dòng
``không trễ quá'' là số lượt trên 3.672 yêu cầu. Khối dưới là phép phát lại lời hứa theo dự báo trên cùng các lần
chạy, bàn ở mục~\ref{muc:c4-dubao}; dự báo chỉ có ở 21 kịch bản ngày thường, và độ dài thêm của lời hứa là trung
bình gộp trên mọi khách được phục vụ, kể cả khách cuối tuần có độ dài thêm bằng 0.}
\label{bang:c4-tran120}
\small
\begin{tabular}{@{}p{8.4cm}rrr@{}}
\toprule
\textbf{Chỉ số} & \textbf{U} & \textbf{C-30} & \textbf{C-120}\\
\midrule
\multicolumn{4}{@{}l}{\textit{Lời hứa đầu không dự báo}}\\
Lượt phục vụ & 2.434 & 2.338 & 2.387\\
Lượt trễ quá 60 s & 203 & 76 & 132\\
Lượt trễ quá 120 s & 102 & 11 & 21\\
Lượt trễ quá 300 s & 18 & 0 & 0\\
Lượt tới sớm quá 60 s & 7 & 3 & 1\\
Không trễ quá 60 s & 2.231 & 2.262 & 2.255\\
Không trễ quá 120 s & 2.332 & 2.327 & 2.366\\
Khách đi chung & 555 & 313 & 428\\
% EV: L1a-02; L1a-04 (U, C-30 tái lập); N-06; C06; C07; C08; N-09 (L1a-07); C4A-13 (sớm quá 60 s 7 / 3 / 1, ADD1:12,17,22; C120R:14,16,18)
\midrule
\multicolumn{4}{@{}l}{\textit{Lời hứa đầu theo dự báo}}\\
Lượt trễ quá 60 s & 137 & 10 & 65\\
Lượt tới sớm quá 60 s & 8 & 4 & 2\\
Không trễ quá 60 s & 2.297 & 2.328 & 2.322\\
Không trễ quá 120 s & 2.347 & 2.338 & 2.379\\
Lời hứa dài thêm, trung bình gộp theo khách (s) & 7,18 & 7,37 & 7,27\\
% EV: L1a-03; N-03; N-06; N-16 (sớm quá 60 s 8 / 4 / 2); C09; F06
\bottomrule
\end{tabular}
\end{table}

Cách đọc bảng: với trần 120 giây và lời hứa đầu không dự báo, 2.387 yêu cầu được phục vụ, 132 lượt trễ quá 60 giây,
nên 2.255 lượt không trễ quá 60 giây trên 3.672 yêu cầu. Chỉ 1 lượt tới sớm quá 60 giây.
% EV: N-06; L1a-02; C4A-13

So với trần 30 giây, trần 120 giây phục vụ thêm 49 lượt và có thêm 56 lượt trễ quá một phút, nên ít hơn 7 lượt không
trễ quá một phút. Nhưng nó thêm 39 lượt không trễ quá hai phút, nhiều hơn cả trần 30 giây lẫn điều phối không trần.
Khi có dự báo, nó ít hơn trần 30 giây 6 lượt ở ngưỡng một phút và thêm 41 lượt ở ngưỡng hai phút. Độ dài thêm của lời
hứa gần như bằng nhau (7,27 so với 7,37 giây, trên hai tập khách khác nhau). Mục~\ref{muc:c4-dem} đặt điểm này cạnh
cách đệm lời hứa.
% EV: L1a-09 (+49 / +56 / −7 / +39; dự báo −6 / +41); N-06; L1a-02 (2.366 > U 2.332 và C-30 2.327); K12; L1a-03; N-16; L1a-08

Dưới trần 120 giây, 9 khách bị dời từ 108 giây trở lên (90\% của trần), không khách nào vượt trần và không quyết định
nào bị ép trong 20.435 lượt.
% EV: N-07 (≥ 0,9 B: 9); L1a-05; N-23

Hai tập khách khác nhau: 2.266 khách được cả hai mức trần phục vụ, 121 khách chỉ được trần 120 giây và 72 khách chỉ
được trần 30 giây phục vụ; 93 trong 121 khách đó cũng được phục vụ khi không trần. Vì vậy 49 lượt thêm là chênh lệch
hai tổng, không phải 49 trong 96 lượt mà trần 30 giây bỏ. Phần thêm dồn vào khung 8 và 17 giờ; so từng kịch bản, trần
120 giây phục vụ nhiều hơn ở 23, bằng ở 8 và ít hơn ở 3.
% EV: N-10 (lưu ý: không viết "lấy lại 49 trong 96 lượt"); L1a-08; N-11 (C-120 − C-30: phục vụ +4 / +26 / +19, trễ quá 60 s +2 / +34 / +20; 23 / 8 / 3 kịch bản); L1a-09

\begin{figure}[htbp]
\centering
\begin{tikzpicture}[font=\footnotesize,
  vuong/.style={draw,fill=white,rectangle,minimum size=4.2pt,inner sep=0pt},
  tamgiac/.style={draw,fill=white,regular polygon,regular polygon sides=3,minimum size=6pt,inner sep=0pt},
  tron/.style={draw,fill=white,circle,minimum size=4.6pt,inner sep=0pt},
  thoi/.style={draw,fill=black,diamond,minimum size=6.5pt,inner sep=0pt},
  nhan/.style={font=\scriptsize,inner sep=1.5pt}]
% (a) Tập tắc đường cố định. Trục ngang: phục vụ − 1.530, 0,06 cm mỗi lượt; trục dọc: trễ, 0,05 cm mỗi lượt.
\begin{scope}[shift={(0,0)}]
\draw[->] (0,0)--(3.95,0);
\draw[->] (0,0)--(0,3.35);
\foreach \vx/\so in {0/1.530,1.2/1.550,2.4/1.570,3.6/1.590}
  \draw (\vx,0)--(\vx,-0.07) node[below,nhan]{\so};
\foreach \vy/\so in {0/0,1/20,2/40,3/60}
  \draw (0,\vy)--(-0.07,\vy) node[left,nhan]{\so};
\draw[dashed] (0.24,0)--(0.60,0)--(1.98,1.15)--(3.24,2.35);
\node[vuong] at (0.24,0) {}; \node[nhan,above] at (0.24,0.05) {30};
\node[vuong] at (0.60,0) {}; \node[nhan,above] at (0.62,0.05) {60};
\node[vuong] at (1.98,1.15) {}; \node[nhan,above left] at (1.98,1.15) {120};
\node[vuong] at (3.24,2.35) {}; \node[nhan,above left] at (3.24,2.35) {U};
\node at (1.8,-0.72) {Lượt phục vụ};
\node[rotate=90] at (-0.85,1.6) {Lượt trễ quá 60 s};
\node[anchor=south] at (1.9,3.4) {(a) Tắc đường cố định};
\end{scope}
% (b) Tập chỉnh trần buổi sáng. Trục ngang: phục vụ − 658, 0,36 cm mỗi lượt; trục dọc: trễ, 0,05 cm mỗi lượt.
% C-120, C-150, C-180 trùng một điểm (666; 51); U, C-240, C-300 trùng một điểm (665; 50).
\begin{scope}[shift={(5.3,0)}]
\draw[->] (0,0)--(3.95,0);
\draw[->] (0,0)--(0,3.35);
\foreach \vx/\so in {0/658,0.72/660,1.44/662,2.16/664,2.88/666,3.6/668}
  \draw (\vx,0)--(\vx,-0.07) node[below,nhan]{\so};
\foreach \vy/\so in {0/0,1/20,2/40,3/60}
  \draw (0,\vy)--(-0.07,\vy) node[left,nhan]{\so};
\draw[densely dotted,thick] (0.72,2.35)--(1.08,2.35)--(2.16,2.50)--(2.88,2.55)--(2.52,2.50);
\node[tamgiac] at (0.72,2.35) {}; \node[nhan,below] at (0.72,2.25) {30};
\node[tamgiac] at (1.08,2.35) {}; \node[nhan,below] at (1.10,2.25) {60};
\node[tamgiac] at (2.16,2.50) {}; \node[nhan,below] at (2.16,2.40) {90};
\node[tamgiac] at (2.88,2.55) {}; \node[nhan,right] at (2.98,2.55) {120--180};
\node[tamgiac] at (2.52,2.50) {}; \node[nhan,above left] at (2.50,2.58) {U, 240, 300};
\node at (1.8,-0.72) {Lượt phục vụ};
\node[rotate=90] at (-0.85,1.6) {Lượt trễ quá 60 s};
\node[anchor=south] at (1.9,3.4) {(b) Buổi sáng};
\end{scope}
% (c) Tập kiểm tra chính. Trục ngang: phục vụ − 2.320, 0,03 cm mỗi lượt; trục dọc: trễ, 0,0125 cm mỗi lượt.
\begin{scope}[shift={(10.6,0)}]
\draw[->] (0,0)--(3.95,0);
\draw[->] (0,0)--(0,3.35);
\foreach \vx/\so in {0/2.320,1.2/2.360,2.4/2.400,3.6/2.440}
  \draw (\vx,0)--(\vx,-0.07) node[below,nhan]{\so};
\foreach \vy/\so in {0/0,1/80,2/160,3/240}
  \draw (0,\vy)--(-0.07,\vy) node[left,nhan]{\so};
\draw (0.54,0.95)--(2.01,1.65)--(3.42,2.54);
\node[tron] at (0.54,0.95) {}; \node[nhan,below right] at (0.54,0.95) {30};
\node[thoi] at (2.01,1.65) {}; \node[nhan,below right] at (2.06,1.62) {120};
\node[tron] at (3.42,2.54) {}; \node[nhan,above left] at (3.42,2.54) {U};
\node at (1.8,-0.72) {Lượt phục vụ};
\node[rotate=90] at (-0.85,1.6) {Lượt trễ quá 60 s};
\node[anchor=south] at (1.9,3.4) {(c) Tập kiểm tra chính};
\end{scope}
\end{tikzpicture}
\caption[Lượt trễ theo lượt phục vụ ở các mức trần]{Lượt trễ quá 60 giây so với lời hứa đầu (trục dọc) theo lượt phục vụ (trục ngang) ở các mức
trần, lời hứa đầu không dự báo. Mỗi khung là một tập với thang riêng: (a) tập tắc đường cố định, 1.728
yêu cầu; (b) tập chỉnh trần buổi sáng, 864 yêu cầu, đủ tám mức, các mức có cùng số đếm chung một điểm; (c) tập
kiểm tra chính, 3.672 yêu cầu. Số cạnh mỗi điểm là trần tính bằng giây, U là điều phối không trần. Đường nối các mức
của cùng một tập theo thứ tự trần tăng dần, chỉ giúp theo dõi các điểm đã đo và không biểu diễn kết quả ở các mức
chưa chạy; các tập không nối với nhau. Ở (a), ghép tự do (1.629 lượt phục vụ, 274 lượt trễ) nằm ngoài khung. Hình
thoi đen là trần 120 giây trên tập kiểm tra chính.}
\label{hinh:c4-tran}
\end{figure}
% EV (a): D09 (C-30 1.534/0; C-60 1.540/0; C-120 1.563/23; U 1.584/47; FREE 1.629/274)
% EV (b): D01; D02; C4A-05 (C-30 660/47; C-60 661/47; C-90 664/50; C-120 = C-150 = C-180 666/51; U = C-240 = C-300 665/50)
% EV (c): C01; C06; N-06 (C-30 2.338/76; C-120 2.387/132; U 2.434/203)
% Tọa độ: (a) x = (phục vụ − 1530) × 0,06, y = trễ × 0,05; (b) x = (phục vụ − 658) × 0,36, y = trễ × 0,05
% (C-60: x = 3 × 0,36 = 1,08; y = 47 × 0,05 = 2,35); (c) x = (phục vụ − 2320) × 0,03, y = trễ × 0,0125.

Trên cả ba tập, nới trần từ 30 lên 120 giây thêm lượt phục vụ và thêm lượt trễ quá một phút. Số lượt không trễ quá một
phút thêm 6, thêm 2 và bớt 7 lần lượt ở tập tắc đường cố định, tập buổi sáng và tập chính, còn ở ngưỡng hai phút trần
120 giây không kém trần 30 giây ở cả ba tập. Vì mỗi kịch bản chỉ chạy một lần và mức 120 giây được chọn sau khi đã xem
tập kiểm tra chính, trần 120 giây chỉ là một điểm đo thêm trên đường đánh đổi.
% EV: D10; D03; L1a-09; D09 (không trễ quá 120 s: 1.563 so với 1.534); D02 (659 so với 658); L1a-02; K36; D11; B11; I06; C4A-11; K42

% -----------------------------------------------------------------------------------------------------
\section{Phần cộng dồn của luật C: kiểm toàn bộ bản ghi chặn}
\label{muc:c4-congdon}

Luật C cộng mọi lần dời của một khách qua cả chuyến rồi mới so với trần. Trên các lần chạy đã kiểm, hiếm khi chỉ riêng
phép cộng ấy làm một phương án bị loại. Trong 98 lần chạy trần 30--300 giây của tập kiểm tra chính và tập buổi sáng,
cả 1.256 bản ghi chặn đều có một lần dời đơn lẻ đã vượt trần của cấu hình tương ứng. Trên 34 lần chạy trần 120 giây
của tập kiểm tra chính, 3 trong 524 bản ghi chỉ vượt trần khi cộng thêm phần trần đã dùng trước đó.
% EV: E01 (98 = 34 + 64; 1.256 = 845 + 411); K09; N-08; L1a-06

Một bản ghi chặn là \emph{chỉ do cộng dồn} khi lần dời mới không vượt $B$ nhưng phần đã dùng $L_r$ cộng lần dời mới
thì vượt (mục~\ref{muc:luatc}). Ví dụ, ở một bản ghi ngày 14/11 với trần 120 giây, khách đã dùng 25,844 giây và phương
án mới dời thêm 99,027 giây: mỗi phần dưới 120 giây, nhưng tổng 124,871 giây vượt trần. Bảng~\ref{bang:c4-congdon} đếm
các bản ghi này trên mọi bản ghi chặn của các lần chạy đã kiểm.
% EV: A06; E05; N-27 (25,844 + 99,027 = 124,871 s); N-08 (quy tắc riêng + cumulative_decisive_check)

\begin{table}[htbp]
\centering
\caption[Kiểm toàn bộ bản ghi chặn của luật C]{Kiểm toàn bộ bản ghi chặn do trần trên 140 lần chạy có luật C dùng hệ số tắc đường đo được đã được kiểm,
thuộc bốn nhóm dưới đây. Một bản ghi ứng với một phương án bị loại, không phải một khách.}
\label{bang:c4-congdon}
\small
\setlength{\tabcolsep}{4pt}
\begin{tabular}{@{}>{\raggedright\arraybackslash}p{\dimexpr .38\linewidth-2\tabcolsep\relax}
>{\raggedleft\arraybackslash}p{\dimexpr .11\linewidth-2\tabcolsep\relax}
>{\raggedleft\arraybackslash}p{\dimexpr .16\linewidth-2\tabcolsep\relax}
>{\raggedleft\arraybackslash}p{\dimexpr .16\linewidth-2\tabcolsep\relax}
>{\raggedleft\arraybackslash}p{\dimexpr .19\linewidth-2\tabcolsep\relax}@{}}
\toprule
\textbf{Lần chạy} & \textbf{Số lần chạy} & \textbf{Bản ghi chặn} & \textbf{Chỉ do cộng dồn} &
\textbf{Khách phục vụ bị dời từ hai lần}\\
\midrule
Tập kiểm tra chính, trần 30 giây & 34 & 845 & 0 & 0\\
Tập chỉnh trần buổi sáng, trần 30--300 giây & 64 & 411 & 0 & 0\\
% EV: E01 (98 = 34 + 64; 1.256 = 845 + 411; 0 chỉ do cộng dồn); E02 (repeated = 0 trên 98 lần chạy)
Tập kiểm tra chính, trần 120 giây & 34 & 524 & 3 & 0\\
% EV: N-08; L1a-06; N-07 (0 khách ≥ 2 lần)
Chiều tuyết 15/11/2018, trần 30 giây & 8 & 163 & 0 & 0\\
% EV: N-26 (163 bản ghi, 0 chỉ do cộng dồn, ≥ 2 lần 0); N4B8-05 (bộ chạy chiều tuyết đã kiểm, CK8:65-72); R26 §Q8
\midrule
Cộng & 140 & 1.943 & 3 & 0\\
% EV: N-15 (132 / 1.780 / 3; kể cả chiều tuyết 140 / 1.943 / 3); ADDENDUM-1 mục 0.3
\bottomrule
\end{tabular}
\end{table}

Cách đọc bảng: trên 34 lần chạy trần 120 giây có 524 phương án bị loại kèm bản ghi chặn do trần, 3 bản ghi trong số đó
chỉ do cộng dồn, và không khách được phục vụ nào bị dời từ hai lần.
% EV: N-08; N-07

Phạm vi kiểm là 132 lần chạy và 1.780 bản ghi (98 lần chạy với 64.371 lượt quyết định được xác nhận, cộng 34 lần chạy
trần 120 giây), hay 140 lần chạy và 1.943 bản ghi kể cả chiều tuyết. Chưa được kiểm: 15 lần chạy hiệu chỉnh, 16 lần
chạy trong đó bộ điều phối dùng dự báo khi quyết định (khác phép phát lại) và 12 lần chạy luật C dùng hệ số đo được
của một phép quét trước đó. Kết quả vì vậy không áp cho mọi lần chạy.
% EV: N-15 (SCOPE); ADDENDUM-1 mục 0.3 (15 lần chạy hiệu chỉnh, 16 thư mục C-30F); E01 (64.371); K09; F12; C4A-08 (SWEEP:13,70: C-x, x = 30…900 s, trên thế giới đo thật W5, 2 kịch bản mỗi mức; NEW:1 by_arm chỉ có t38 và w07)

Ba bản ghi chỉ do cộng dồn nằm ở ba kịch bản khung 8 giờ (14, 15 và 18/11), trong 3 lượt quyết định; khi không trần,
khách của bản ghi ngày 14/11 bị dời hai lần đúng bằng hai lượng ở ví dụ trên. Trong 3 yêu cầu mới liên quan, 2 cuối
cùng không được nhận dưới trần 120 giây (một trong hai cũng không dưới trần 30 giây). Vì bản ghi dừng ở lỗi gặp trước
và diễn biến của các cấu hình khác nhau, điều này không chứng minh phần cộng dồn làm mất hai yêu cầu đó.
% EV: L1a-06 (ba kịch bản w08 ngày 14, 15, 18/11); N-08 (3 quyết định; 2 không được nhận, 1 phục vụ); N-27; E07; A23; I08

Ngoài hệ số đo được, bản ghi chỉ do cộng dồn còn có ở hai nơi. Trên tập tắc đường cố định, nơi bộ kiểm ghi đủ mọi bản
ghi, có 2 bản ghi ở trần 60 giây và 4 ở trần 120 giây; mọi yêu cầu mới liên quan vẫn được phục vụ. Trên các kiểu tắc
đường tự tạo của phép quét trước đó, sau khi bỏ bản chạy trùng, có 4 bản ghi trong 3 lượt quyết định ở trần 120 giây;
1 yêu cầu mới liên quan cuối cùng bị từ chối.
% EV: E03; E04 (4/3 sau khi bỏ bản lặp; thô 6/5); A23 (NEW:13: tập tắc đường cố định bật chế độ ghi đủ); FJC:31-54; FJR:2-7; C4A-08 (thế giới tự tạo thuộc tier3-sweep-v1)

Về định nghĩa, luật C vẫn khác luật ``mỗi lần dời không quá 30 giây'': luật sau cho một khách bị dời ba lần, tổng 90
giây, còn luật C trần 30 giây thì không. Trên số liệu, khác biệt này hiếm: khi không trần chỉ 12 trong 2.434 khách bị
dời từ hai lần. Tính thử cả hai cách trên các lần chạy không trần đã ghi, chúng chỉ khác kết luận ở 4 khách với trần
120 giây và 0 khách với trần 30, 60 và 300 giây. Phép tính thử không cho biết kết quả nếu chạy thật với luật kia, vì
các quyết định sẽ khác; luật mỗi lần một trần chưa được cài đặt và chưa chạy.

Khi chuyển sang ưu tiên chi phí (mục~\ref{muc:c4-vaitro}), cấu hình có trần 30 giây có một bản ghi chỉ vượt
do cộng dồn: phần đã dùng 24,404 giây cộng lần dời 13,558 giây thành 37,962 giây. Yêu cầu mới liên quan cuối
cùng vẫn được phục vụ. Cấu hình này có hai khách bị sửa giờ trả từ hai lần. Các số 0 của bảng trên chỉ thuộc
phạm vi ưu tiên ổn định đã nêu, không áp cho mọi thứ tự chọn phương án.
% EV: S15-01 (5d--5e); S15-04. Một witness, không phải một yêu cầu cuối cùng bị từ chối.

Trên 2.502 lịch sử khách được phục vụ của cấu hình ưu tiên chi phí không trần, có 7 lịch sử mà mỗi lần dời
không quá 120 giây nhưng tổng độ lớn vượt 120 giây. Ví dụ minh họa, hai lần dời 70 và 60 giây đều dưới trần
từng lần 120 giây, nhưng tổng 130 giây vượt trần cộng dồn. Số 7 đếm lịch sử khách theo phép kiểm ở mức 120
giây, không đếm khách vi phạm trần 30 giây. Nhánh so trực tiếp với luật giới hạn từng lần không được thực
hiện vì phép sàng lọc chưa đạt mốc 10 lịch sử đã đặt để phân bổ nguồn lực. Mốc này không phải ngưỡng kiểm
định thống kê; việc dừng chưa chứng minh phép cộng dồn không có tác dụng.
% EV: S15-01 (H7); S15-02 (h7_served); STAGE1B-PLAN.md mục 6,9. Ví dụ 70+60 là dựng tay.

% EV: E08; C20; E06 (CVP:14-18: 0/0/4/0 trên 2.434 khách U); K42; E09; K08

Các mục trên so luật C với điều phối không trần khi giữ nguyên bộ chọn ưu tiên ổn định, tức bộ chọn xét mức dùng trần
và phần dời trước chi phí kế hoạch (Bảng~\ref{bang:rangbuoc}). Mục~\ref{muc:c4-vaitro} xét kết quả có phụ thuộc vào
thứ tự đó không.
% EV: A08; H01; H02; A14; X3-07

% =====================================================================================================
% Chương 4, PHẦN B (mục 4.6–4.11) — bản sửa v11 (2026-10-13), BẢN DÙNG ĐỂ GHÉP.
% Sửa có mục tiêu từ drafts/09-chuong4-phanB-v10.tex theo R37 và brief R39; các bản trước giữ nguyên.
% Quyết định trước: reviews/CH4B-REVISION-LOG.md; ghi chú cũ: drafts/09-chuong4-phanB-v8-NOTES.md.
% Từng thay đổi, phép kiểm số mới và việc bàn giao: 09-chuong4-phanB-v11-CHANGES-v2.md bên cạnh.
% Nối ngay sau 08-chuong4-phanA-v5-v3.tex trong cùng % =====================================================================================================
% Chương 4, PHẦN A (mở chương, mục 4.1–4.5) — bản sửa v5 (2026-10-13), BẢN DÙNG ĐỂ GHÉP.
% Sửa có mục tiêu từ drafts/08-chuong4-phanA-v4.tex theo R34-CH4A-V4-VERIFY-CODEX-v2.md và brief R36.
% Quyết định trước: reviews/CH4A-REVISION-LOG-v2.md; ghi chú cũ: drafts/08-chuong4-phanA-v2-NOTES.md.
% Bản v4 và mọi bản trước giữ nguyên. Tệp tiếp tục R36: 08-chuong4-phanA-v5-v3.tex;
% từng thay đổi, bằng chứng mới và việc bàn giao: 08-chuong4-phanA-v5-CHANGES-v3.md. Phần B v10 nối sau tệp này.
% Mỗi dòng "% EV:" trỏ tới EVIDENCE-REGISTER.md (A01…K52), các phụ lục 1 (N-xx, L1a-xx), 2 (X3-xx),
% 4 (C4A-xx), 6 (N4B8-xx), hoặc sổ bổ sung R36-A01…03 trong tệp CHANGES-v3.md bên cạnh.
% Không dùng gói mới: TikZ và array đã nạp trong main.tex; chỉ sửa độ rộng các bảng được brief chỉ ra.
% =====================================================================================================

\chapter{Thực nghiệm và đánh giá}
\label{chuong:thucnghiem}

Chương này đo luật C được gì và mất gì trên mô phỏng FleetPy, so với điều phối không trần, luật hạn chót và cách đệm
lời hứa, tức cộng thêm một khoảng cố định vào mọi lời hứa. Chương cũng xét lời hứa theo dự báo và những chỗ dự báo
chưa đủ. Các mục đi theo ba phần của độ lệch giờ trả so với lời hứa đầu (mục~\ref{muc:so}). Phần tính cho điều phối,
gắn với việc ghép thêm khách, ở mục~\ref{muc:c4-ketqua}--\ref{muc:c4-vaitro}; phần được hồ sơ tắc đường lịch sử giải
thích, cùng cách đệm lời hứa, ở mục~\ref{muc:c4-dubao} và~\ref{muc:c4-dem}; phần còn lại ở mục~\ref{muc:c4-tuyet}.
Bảng~\ref{bang:c4-lotrinh} nối ba câu hỏi nghiên cứu ở mục~\ref{muc:c1-cauhoi} với các mục trả lời chúng.
% EV: — (đoạn dẫn, không nêu số); ba phần theo eq:so của chương 3 (X3-12); "đệm" theo GLOSSARY-VI mục B

\begin{table}[htbp]
\centering
\caption[Câu hỏi nghiên cứu và các mục trả lời]{Ba câu hỏi nghiên cứu và các mục của Chương~\ref{chuong:thucnghiem} trả lời chúng}
\label{bang:c4-lotrinh}
\small
\begin{tabular}{@{}p{6.3cm}p{2.4cm}p{6.3cm}@{}}
\toprule
\textbf{Câu hỏi} & \textbf{Mục} & \textbf{Đại lượng được đo}\\
\midrule
CH1. Phần dời do điều phối được đo thế nào, và trần trên phần đó giữ được gì, không giữ được gì?
& \ref{muc:c4-dieuphoi}, \ref{muc:c4-congdon}, \ref{muc:c4-hople}
& Số khách bị dời giờ trả vì ghép thêm khách, độ dời lớn nhất; yêu cầu bị từ chối có phương án bị trần chặn;
phép cộng các lần dời qua cả chuyến; các bộ đếm kiểm tra\\
CH2. Giữ trần đổi số lượt phục vụ, số lượt không trễ quá một và hai phút, cùng độ dài lời hứa thế nào so
với điều phối không trần, luật hạn chót và cách đệm lời hứa?
& \ref{muc:c4-ketqua}, \ref{muc:c4-chontran}, \ref{muc:c4-vaitro}, \ref{muc:c4-dem}
& Lượt phục vụ; lượt trễ và không trễ so với lời hứa đầu; quyết định bị ép; các mức trần; thứ tự tiêu chí
của bộ chọn; độ dài thêm của lời hứa\\
CH3. Lời hứa theo hồ sơ tắc đường lịch sử đổi số lượt không trễ quá một phút thế nào, và phần nào còn
chưa được giải thích?
& \ref{muc:c4-dubao}, \ref{muc:c4-tuyet}
& Lượt trễ khi có và không có dự báo; độ dài thêm của lời hứa; ba phần của độ lệch trên ngày thường và chiều
tuyết\\
\bottomrule
\end{tabular}
\end{table}
% EV: — (lộ trình, không có khẳng định; câu hỏi theo OUTLINE-V2 cuối §1.2)

% -----------------------------------------------------------------------------------------------------
\section{Thiết lập thực nghiệm}
\label{muc:c4-thietlap}

\subsection{Kịch bản thử}
\label{muc:c4-kichban}

FleetPy 1.0.2 mô phỏng xe và khách trên bộ dữ liệu chuẩn Manhattan; ở mỗi lượt quyết định, bộ chuyển đổi gửi tình
trạng hiện tại cho bộ chạy RideBound và nhận lại quyết định (mục~\ref{muc:kientruc}). Nhu cầu là các chuyến taxi New
York tuần 11--18/11/2018 trong bộ dữ liệu này~\cite{engelhardt2026}. Một kịch bản gồm một ngày, một khung giờ hai
tiếng và một trong bốn mẫu nhu cầu của ngày đó, mỗi mẫu gồm 10\% số chuyến của ngày: 108 yêu cầu trong khung giờ, phục
vụ bởi 8 xe 4 chỗ (Bảng~\ref{bang:c4-thamso}).
% EV: B01; X3-15; X3-21 (FC:395-429); B02 (FIX: 2018-11-14_sample_10_3.csv); J18 (mẫu 10% của một ngày); B08 (mẫu r1–r4); B03

Tắc đường là một hệ số chung cho cả thành phố, đổi theo khung 15 phút và lấy từ dữ liệu chuyến taxi, không phải tốc độ
từng đoạn đường; dưới đây gọi là hệ số tắc đường đo được. Mô phỏng không có nhiễu ngẫu nhiên, nên mỗi cặp
kịch bản và cấu hình chạy một lần và cho một kết quả. Ba kịch bản cùng ngày, cùng mẫu có chung 3 đến 20 trong 108 yêu
cầu, nên các kịch bản không độc lập với nhau.
% EV: B06 (sigma = shockRate = 0); K49; B11 (T38:90, 93); I05

\begin{table}[htbp]
\centering
\caption[Tham số của một kịch bản thử]{Tham số của một kịch bản thử. Các giá trị là cấu hình của thực nghiệm, không phải chuẩn ngành. Tập tắc
đường cố định (mục~\ref{muc:c4-taptu}) khác ở dòng tắc đường và chạy bằng bộ chạy cũ, có giới hạn bộ sinh và bộ
giải chưa được kiểm lại.}
\label{bang:c4-thamso}
\small
\begin{tabular}{@{}p{4.0cm}p{11.0cm}@{}}
\toprule
\textbf{Tham số} & \textbf{Giá trị}\\
\midrule
Khung giờ & 7:00--9:00, 8:00--10:00 hoặc 17:00--19:00, gọi là khung 7, 8 và 17 giờ; riêng chiều tuyết
là 14:00--16:00 và 15:00--17:00\\
Yêu cầu & 108 mỗi kịch bản, lấy trong khung giờ từ một mẫu 10\% số chuyến taxi\\
Xe & 8 xe, mỗi xe 4 chỗ\\
Đồ thị & 96 nút là các điểm đi, điểm đến lấy từ dữ liệu Manhattan; mỗi cặp nút có sẵn một thời gian đi. Đây là
đồ thị rút gọn, không phải mạng đường phố\\
Ràng buộc dịch vụ & Đón trong vòng 600 giây kể từ lúc gọi; ngồi xe không quá 1,5 lần thời gian đi
thẳng (Bảng~\ref{bang:rangbuoc})\\
Thời gian dừng đón, trả & 0 giây: xe không mất thời gian khi khách lên, xuống (kiểm trên 3 lần chạy mẫu)\\
Tắc đường & Một hệ số chung cho cả thành phố theo khung 15 phút; tốc độ gốc của mọi kịch bản lấy theo hệ số lúc
00:00\\
Bộ sinh phương án & Tối đa 100 phương án mỗi xe; mỗi lượt quyết định xét tối đa 2 yêu cầu mới\\
Bộ giải & OR-Tools 9.15, hạt giống ngẫu nhiên cố định bằng 7\\
\bottomrule
\end{tabular}
\end{table}
% EV (khung giờ, yêu cầu): B08; B20; B02; C4A-06 (chiều tuyết cũng 108 yêu cầu mỗi kịch bản, INV:14-15,17-22)
% EV (xe, đồ thị, ràng buộc dịch vụ): B03; K48 (PREF:22-26: mỗi cặp điểm đi–đến là một cạnh trực tiếp)
% EV (dừng, tắc đường): B05 (chỉ 3 lần chạy mẫu, SERV:2-4); B06; B07 (tốc độ nền lấy lúc 00:00)
% EV (bộ sinh, bộ giải): X3-17; X3-32 (giới hạn 2 yêu cầu áp cho cả lô của lượt, GEN:54-96); C4A-18 (tập tắc đường cố định: bộ chạy khác, NEW:13; giới hạn chưa kiểm)

\subsection{Các tập kịch bản và cấu hình được so}
\label{muc:c4-taptu}

Chương này dùng bốn tập kịch bản (Bảng~\ref{bang:c4-taptu}). Trên tập kiểm tra chính còn có phép so hai
thứ tự chọn phương án ở mục~\ref{muc:c4-vaitro}; phép so này dùng lại cùng kịch bản, không bổ sung ngày dữ liệu. Từ các ngày 14--18/11/2018 dựng được 49 kịch bản, chia
theo cặp ngày và mẫu để ba khung giờ của một cặp nằm cùng một phía. 34 kịch bản làm tập kiểm tra chính; 15 kịch bản
còn lại không dùng để tính các kết quả phục vụ và đúng giờ của chương này. Các lần chạy trên 15 kịch bản đó chỉ tham
gia phép kiểm tính nhất quán của sổ ở mục~\ref{muc:kiemtra}. Ngày 14--16/11 là ngày thường, 17 và 18/11 là cuối tuần.
% EV: B08 (SPLIT: chia theo nhóm ngày, mẫu); C17 (15 kịch bản hiệu chỉnh khoảng dự phòng γ); B09; B21 (11/11/2018 là Chủ nhật); A26; RUN38; T38:99-106; chương 3 v4-v3:725-727 (185 lần chạy gồm 15 lần hiệu chỉnh)

\begin{table}[htbp]
\centering
\caption[Các tập kịch bản dùng trong thực nghiệm]{Bốn tập kịch bản dùng trong chương. U: điều phối không trần; C-$B$: luật C trần $B$ giây;
M$^{+}$-$B$: luật hạn chót $B$ giây; V-30: luật trần độ dời khách thấy 30 giây; FREE: ghép tự do
(Bảng~\ref{bang:luat}). Cột cuối là số lần chạy, mỗi lần một cặp kịch bản và cấu hình.}
\label{bang:c4-taptu}
\footnotesize
\setlength{\tabcolsep}{2.5pt}
\begin{tabular}{@{}p{2.5cm}p{1.9cm}p{3.1cm}p{3.0cm}p{2.6cm}p{1.9cm}@{}}
\toprule
\textbf{Tập} & \textbf{Quy mô} & \textbf{Ngày, khung giờ} & \textbf{Cấu hình dùng ở chương này} &
\textbf{Bộ chạy} & \textbf{Lần chạy}\\
\midrule
Tập kiểm tra chính & 34 kịch bản, 3.672 yêu cầu & 14--18/11/2018; khung 7, 8, 17 giờ có 9, 14, 11
kịch bản & U, C-30, C-120, M$^{+}$-30, M$^{+}$-60, V-30 & Bản đã lưu, có luật phục hồi không tệ hơn &
170 (năm cấu hình), 34 (C-120); thêm 68 ưu tiên chi phí\\
% EV: B08; B09; B10 (170 = 34 × 5); L1a-01 (34 lần chạy C-120); A19; A28; X3-10 (R120:3-8)
Tập chỉnh trần buổi sáng & 8 kịch bản, 864 yêu cầu & 12--13/11/2018; khung 7 giờ & U; C và M$^{+}$ ở 30, 60,
90, 120, 150, 180, 240 và 300 giây & Như tập chính & 136 (trong 192 lần chạy đã lưu)\\
% EV: B15; B16; D01; D04; X3-10 (RW7:1-8; RW7X:6-8); X3-28; C4A-07 (17 cấu hình × 8 = 136; 88 + 104 = 192, RW7:33-34; RW7X:35-36)
Tập tắc đường cố định & 16 kịch bản, 1.728 yêu cầu & 12--13/11/2018; khung 8 và 17 giờ; hệ số không
đổi trong lần chạy & FREE, U, C-30, C-60, C-120 & Bản cũ, không có luật phục hồi & 16 mỗi cấu hình\\
% EV: B17; D08; D09; H13 (n = 16 mỗi cấu hình); LAD:2-6 (travelSnapshotVersion {1: 16}: một bảng giờ đi); C4A-04
Chiều tuyết 15/11/2018 & 8 kịch bản, 864 yêu cầu & 15/11/2018; khung 14 và 15 giờ & C-30 &
Như tập chính & 8\\
% EV: B20; N-26 (8 lần chạy C-30, 601 khách được phục vụ); C4A-06 (864 = 8 × 108, INV:14-15,17-22); N4B8-05 (CK8:65-72); R26 §Q8; agent-numeric-output-v2.txt (manifest, cấu hình phục hồi và DLL)
\bottomrule
\end{tabular}
\end{table}

Bản đã lưu là bản bộ chạy biên dịch và lưu lại, dùng chung cho tập kiểm tra chính, tập buổi sáng, các lần chạy trần
120 giây và chiều tuyết; bản cũ chưa có luật phục hồi (mục~\ref{muc:doichung}, \ref{muc:kiemtra}). Trừ ghép tự do, mọi
cấu hình trong từng thứ tự ưu tiên dùng cùng bộ sinh, bộ kiểm và bộ chọn với luật C; điều phối không trần chỉ
khác luật C trần 30 giây ở trần giờ trả. Các mục trước mục~\ref{muc:c4-vaitro}, cùng các phép so dự báo và đệm
sau đó, dùng thứ tự ưu tiên ổn định nếu không ghi khác. Ghép tự do không khóa gì và không qua bộ kiểm cam kết.
% EV: A28; X3-10 (RW7:1-8; RW7X:6-8; R120:3-8); N4B8-05 (CK8:65-72); R26 §Q8; B17; NEW:13 (bộ chạy cũ không có commitmentRecovery); A08; H01; A17; K50; X3-38 (HO:109, b1-rolling-cost)

Lần chạy hỏng vì lỗi hạ tầng (một lần ở tập kiểm tra chính) được thay bằng bản chạy lại hoàn tất (Phụ
lục~\ref{phuluc:C}). Bản kiểm kê ngày 08/10/2026 ghi 1.991 bản ghi lần chạy; kết quả của chương chỉ lấy từ bốn tập
trên và 68 lượt ưu tiên chi phí; các đợt chạy khác chỉ được nhắc khi nêu phạm vi kiểm tra (mục~\ref{muc:c4-congdon}).
% EV: C4A-17 (T38:12; ADD1:6: C-30-d20181114-s10-r3-w17-rerun1); ADDENDUM-1 mục 0.3; B19; C4A-20 (INV:1, bản kiểm kê ngày 08/10)

\subsection{Chỉ số}
\label{muc:c4-chiso}

Một lượt phục vụ là một yêu cầu, tức một khách, được nhận và chở tới nơi. Lượt trễ và lượt tới sớm được đếm trong nhóm
khách được phục vụ, so với lời hứa đầu (Bảng~\ref{bang:c4-chiso}). Số lượt không trễ bằng số lượt phục vụ trừ số lượt
trễ, tính trên toàn bộ yêu cầu, nên yêu cầu bị từ chối không được tính là đúng giờ. Ngưỡng chính là 60 giây; ngưỡng
120 giây được báo kèm vì có phép so đổi kết luận theo ngưỡng (mục~\ref{muc:c4-ketqua}).
% EV: C01; N-07 (khách được nhận = khách được phục vụ); C06; C08 (DIALPY:8-9: bị từ chối không tính là kịp; ở 120 s U hơn C-30 5 lượt)

Lời hứa đầu là giờ trả của kế hoạch lúc khách xác nhận đặt xe, tính với thời gian đi lúc đó. Tính từ lúc gọi xe, độ dài
trung bình của lời hứa là 14,15 phút trên 2.434 khách khi không trần và 14,12 phút trên 2.338 khách dưới luật C trần 30
giây, gộp theo khách được phục vụ. Hai giá trị trung bình gần nhau trên hai tập khách khác nhau; phép so này không nói
về độ dài lời hứa của từng khách.
% EV: A13; A05; B14 (BURD:5,11); R36-A03 (p0 trừ request.arrivalTimeMs, gộp theo khách); agent-numeric-output-v2.txt

Tỉ lệ \emph{gộp theo khách} chia tổng số khách thỏa điều kiện cho tổng số khách được phục vụ; \emph{trung bình theo
kịch bản} lấy trung bình của tỉ lệ từng kịch bản, mỗi kịch bản tính như nhau. Ví dụ dựng tay: với 5 trong 10 khách ở
một kịch bản và 5 trong 90 ở kịch bản kia, gộp theo khách cho 10\%, trung bình theo kịch bản cho 27,8\%. Chênh lệch là
cấu hình được thử trừ đối chứng; số dương nghĩa là nhiều hơn.
% EV: C27 (hai cách gộp cho hai số khác nhau); K24; ví dụ dựng tay (10/100 = 10%; (50% + 5,56%)/2 = 27,78%)

Kết quả ``có dự báo'' là phép phát lại trên các lần chạy đã ghi (mục~\ref{muc:dubao}): xe nào chở ai giữ nguyên, chỉ
các giờ báo cho khách được tính lại theo dự báo, và các chỉ số dùng $D_0$. Dự báo chỉ có ở 21 kịch bản ngày thường của
tập kiểm tra chính; 13 kịch bản cuối tuần và tập buổi sáng giữ lời hứa không dự báo.
% EV: C09; F03 (chỉ vế phát lại); X3-11 (D = max(p, S) áp cho mọi lời hứa đón và trả đã công bố); F02; N-28; D02 (12–13/11 không có dự báo)

\begin{table}[htbp]
\centering
\caption[Định nghĩa các chỉ số đánh giá]{Các chỉ số dùng trong chương. $a$ là giờ trả thực tế, $p_0$ là lời hứa đầu không dự báo, $D_0$
là lời hứa đầu theo dự báo (không có dự báo thì $D_0 = p_0$); $c$ là một lần dời giờ trả do điều phối, $z$ là phần
khách thấy của một lần cập nhật lời hứa (mục~\ref{muc:bamoc}). ``Quá'' nghĩa là lớn hơn hẳn: trễ đúng 60 giây
không tính là trễ quá 60 giây.}
\label{bang:c4-chiso}
\footnotesize
\setlength{\tabcolsep}{3pt}
\begin{tabular}{@{}p{2.7cm}p{5.3cm}p{3.5cm}p{3.5cm}@{}}
\toprule
\textbf{Chỉ số} & \textbf{Định nghĩa} & \textbf{Mẫu số, cách gộp} & \textbf{Vì sao đo}\\
\midrule
Lượt phục vụ & Yêu cầu được nhận và chở tới nơi & Số đếm; chia cho số yêu cầu thì ra tỉ lệ phục vụ & Giá của
luật tính bằng chuyến\\
Trễ quá $\theta$ & Khách được phục vụ có $a - D_0 > \theta$; $\theta$ = 30, 60, 120, 300 hoặc 600 giây & Đếm trong
khách được phục vụ & Lời hứa bị vượt ở phía muộn\\
Không trễ quá $\theta$ & Lượt phục vụ trừ lượt trễ quá $\theta$ & Chia cho số yêu cầu; yêu cầu bị từ chối không
tính là đúng giờ & Phục vụ và đúng giờ trên một mẫu số\\
Tới sớm quá 60 giây & Khách được phục vụ có $a - D_0 < -60$ giây & Đếm trong khách được phục vụ & Lời
hứa quá dài cũng sai\\
Lời hứa dài thêm & $D_0 - p_0$ & Trung bình gộp theo khách được phục vụ, kể cả khách không có dự báo & Giá của
dự báo hay đệm\\
Khách đi chung & Có khách khác trên cùng xe trong khoảng từ lúc được đón tới lúc được trả & Gộp theo
khách được phục vụ & Mức ghép chuyến còn lại\\
Khách bị điều phối dời giờ trả & Có ít nhất một lần $c \neq 0$; độ dời của khách là $\sum|c|$ &
Gộp theo khách được phục vụ & Đại lượng luật C giới hạn\\
Quyết định bị ép & Lượt quyết định có ít nhất một xe bị ép: luật lời hứa đã loại cả phương án giữ nguyên kế hoạch
của xe, và xe đi theo phương án được giữ nhờ luật phục hồi, có ít nhất một khách được miễn trừ
(mục~\ref{muc:doichung}) & Chia cho số lượt quyết định đã được bộ mô phỏng xác nhận áp dụng; gộp theo
lượt quyết định, mỗi lượt đếm một lần & Luật có tự giữ được nghĩa vụ không\\
Khách bị phá lời hứa & Khách mà luật đang dùng không giữ được điều nó hứa. Luật C: $\sum|c| > B$; luật hạn chót:
$a - p_0 > B$; luật trần độ dời khách thấy: $\sum|z| > B$ & Chia cho số khách được phục vụ; gộp theo
khách & Luật có giữ được điều nó hứa không\\
\bottomrule
\end{tabular}
\end{table}
% EV (phục vụ): C01; L1a-01 (2.387 hoàn tất + 1.285 bị từ chối = 3.672); N-07
% EV (trễ, không trễ, sớm): C06; C08 (DIALPY:8-9); N-04 (trễ > 60 s và sớm < −60 s rời nhau); C4A-01 (ngưỡng 30 s, cùng quy tắc sweep_analyze.py:107,128-132); C09; C10
% EV (dài thêm, đi chung, bị dời): F06; N-06 (gộp theo khách, ext_sum / served, ADD1:14,19,24); N-09 (định nghĩa đi chung); C20; N-07; A06; X3-05
% EV (bị ép): C4A-10 (RS:1012-1027; VALID:141-145; HVA:252-300; POL:443-463); X3-30; X3-09; C13 (trung bình theo kịch bản)
% EV (phá lời hứa): N-21 (định nghĩa của M⁺-30 và C-30; số gộp); A11; GLOSSARY-VI mục B "phá lời hứa"

\subsection{Mức trần 30 giây và tính chất thăm dò của kết quả}
\label{muc:c4-thamdo}

Trần chính $B$ = 30 giây là trần giờ trả của cấu hình có trần ở tập kịch bản tĩnh của đợt thử trước, nơi thời gian đi
không đổi (Phụ lục~\ref{phuluc:A}). Mức này được giữ nguyên cho tập kiểm tra chính, không chỉnh theo kết quả của tập
này. Nó không phải ngưỡng cảm nhận của khách hay chuẩn ngành, nên mỗi mức trần chỉ là một điểm trên đường đánh đổi
(mục~\ref{muc:c4-chontran}).
% EV: C4A-15 (TEX08:136-137: C1 = trần 30 giây cho tổng độ dời giờ trả; thay X3-18 đã bỏ); K29 (tập thời gian đi tĩnh); K33; B10 (năm cấu hình cố định trước lần chạy đầu, T38:3-6); D20; K52; K36

Các kết quả là thăm dò và dùng lại dữ liệu đã xem. Ngưỡng phân tích và cách chia tập được chọn sau khi đã xem tập buổi
sáng và các đợt chạy thử trước đó. Mức thử 120 giây được chọn sau khi đã xem kết quả của điều phối không trần và luật
C trần 30 giây trên chính tập kiểm tra chính, kể cả các phép tính thử ở 120 giây (mục~\ref{muc:c4-dieuphoi},
\ref{muc:c4-congdon}). Các con số vì vậy mô tả điều đã đo, không phải phép kiểm trên dữ liệu chưa từng xem.
% EV: I06; T38:8 (chọn sau khi xem Việc A, B, w07x); C4A-11 (BUD 07/10 = D14; CVP 10/10 = E06; C120R:10 chạy 10/10; T38:5 tập chính chạy 1–2/10); K42

% -----------------------------------------------------------------------------------------------------
\section{Kết quả chính khi chưa có dự báo}
\label{muc:c4-ketqua}

Với bộ chọn ưu tiên ổn định trên 3.672 yêu cầu của tập kiểm tra chính, luật C trần 30 giây không để phần dời giờ trả do điều phối của khách
nào vượt 30 giây: lớn nhất 29,2 giây, so với 590,1 giây (gần 10 phút) khi không trần. Không lượt nào của luật C trễ
quá 5 phút so với lời hứa đầu, so với 18 lượt khi không trần. Đổi lại, luật C phục vụ 2.338 lượt, ít hơn điều phối
không trần 96 lượt, và có ít hơn 242 khách đi chung. Số lượt không trễ quá một phút tăng 31, từ 2.231 lên 2.262, nhưng
số lượt không trễ quá hai phút giảm 5, từ 2.332 xuống 2.327.
% EV: C20; N-23; N-24; C07; C01; C02; N-20; C08

Mục~\ref{muc:c4-dem} cho thấy khi có dự báo, cộng thêm 30 giây vào mọi lời hứa của điều phối không trần cũng cho 2.328
lượt không trễ quá một phút, bằng luật C có dự báo. Vì vậy khóa luận không dùng chỉ số này làm lý do chọn luật C.
% EV: C10 (U + dự báo + 30 s: 2.328 = C-30 + dự báo: 2.328); K06; K51

Luật hạn chót chỉ mượn phần hạn chót của Schulz và Pfeiffer~\cite{schulz2026}: giờ trả mới không được muộn hơn lời hứa
đầu quá 30 hoặc 60 giây, thay vì 900 giây của họ. Nó chạy cùng bộ điều phối và luật phục hồi với luật C, nên đây không
phải phép so với hệ thống của hai tác giả.
% EV: J10; J11; A09; A19; H01; X3-28; X3-41

\begin{table}[htbp]
\centering
\caption[Kết quả chính khi chưa có dự báo]{Kết quả trên tập kiểm tra chính với lời hứa đầu không dự báo (34 kịch bản, 3.672 yêu cầu mỗi
cấu hình, với bộ chọn ưu tiên ổn định). Số đếm gộp trên mọi kịch bản; trễ tính so với lời hứa đầu. Hai cột ``không trễ quá'' là số lượt phục vụ
không trễ quá ngưỡng; chia cho 3.672 thì ra tỉ lệ trên yêu cầu. ``Chưa đếm'': chưa có số đếm chính xác.}
\label{bang:c4-chinh}
\small
\setlength{\tabcolsep}{4pt}
\begin{tabular}{@{}p{\dimexpr .12\linewidth-2\tabcolsep\relax}
>{\raggedleft\arraybackslash}p{\dimexpr .13\linewidth-2\tabcolsep\relax}
>{\raggedleft\arraybackslash}p{\dimexpr .16\linewidth-2\tabcolsep\relax}
>{\raggedleft\arraybackslash}p{\dimexpr .22\linewidth-2\tabcolsep\relax}
>{\raggedleft\arraybackslash}p{\dimexpr .23\linewidth-2\tabcolsep\relax}
>{\raggedleft\arraybackslash}p{\dimexpr .14\linewidth-2\tabcolsep\relax}@{}}
\toprule
\textbf{Cấu hình} & \textbf{Phục vụ} & \textbf{Trễ quá 60 s} & \textbf{Không trễ quá 60 s} &
\textbf{Không trễ quá 120 s} & \textbf{Khách đi chung}\\
\midrule
U & 2.434 & 203 & 2.231 & 2.332 & 555\\
C-30 & 2.338 & 76 & 2.262 & 2.327 & 313\\
M$^{+}$-30 & 2.338 & 74 & 2.264 & 2.327 & 307\\
% EV (U, C-30, M⁺-30): C01; C06; C08; N-09 (555 / 313 / 307, số đếm chính xác)
M$^{+}$-60 & 2.347 & 79 & 2.268 & 2.336 & chưa đếm\\
V-30 & 2.336 & 74 & 2.262 & 2.325 & chưa đếm\\
% EV (M⁺-60, V-30): C01; C4A-02 (CK4A check: trễ quá 60 / 120 s 79 / 11 và 74 / 11; EXP:5-6 trễ quá 60 s 79 và 74; đi chung chỉ có % làm tròn ở EXP)
\bottomrule
\end{tabular}
\end{table}

Cách đọc bảng: dưới trần 30 giây, 2.338 yêu cầu được phục vụ; 76 lượt trong số đó tới muộn hơn lời hứa đầu quá 60
giây, nên 2.338 $-$ 76 = 2.262 lượt không trễ quá 60 giây.
% EV: C01; C06; C08

Cái giá trước hết là số lượt phục vụ: bớt 96 lượt nghĩa là tỉ lệ phục vụ giảm 2,6 điểm phần trăm (96 trên 3.672 yêu
cầu), tức ít hơn 3,9\% số lượt của điều phối không trần (96 trên 2.434). Mức giảm nằm ở khung 8 giờ (1.014 xuống 975)
và 17 giờ (904 xuống 852); khung 7 giờ chỉ giảm 5. Số khách đi chung giảm từ 555 xuống 313, tức từ 22,8\% xuống 13,4\%
số khách được phục vụ (gộp theo khách). Khóa luận không có bằng chứng về giá trị của các chuyến đi chung này đối với
khách hay hãng xe.
% EV: C02; C03; N-09; N-20 (lưu ý: D20, I07, I10)

Số lượt không trễ quá một phút vẫn tăng vì phục vụ bớt 96 lượt nhưng trễ bớt 127 lượt: 127 $-$ 96 = 31. Hai cấu hình
cũng phục vụ những khách khác nhau: 2.207 khách chung, 227 khách chỉ được điều phối không trần và 131 khách chỉ được
luật C phục vụ. Vì vậy 31 là chênh lệch hai tổng, không phải 31 khách cụ thể được đưa tới đúng giờ.
% EV: C06; C08; C04 (câu về chờ, ngồi xe của khách chung, C05, cắt vì ngân sách trang; xem NOTES)

So với luật hạn chót 30 giây, luật C phục vụ đúng bằng số lượt (2.338) nhưng không hoàn toàn cùng những khách (2.332
khách chung, mỗi bên 6 khách riêng), và có 76 so với 74 lượt trễ quá một phút. Luật hạn chót 60 giây và luật trần độ
dời khách thấy cũng gần như vậy (Bảng~\ref{bang:c4-chinh}). Các luật lời hứa khác nhau ở chỗ mỗi luật có giữ được điều
chính nó hứa hay không (Bảng~\ref{bang:c4-nghiavu}).
% EV: C01; C04; C06; K05; K07; C32; C4A-02

\begin{table}[htbp]
\centering
\caption[Nghĩa vụ của chính từng luật lời hứa]{Nghĩa vụ của chính từng luật trên tập kiểm tra chính (34 kịch bản). Hai cột đầu ghi số đếm gộp: số quyết
định bị ép trên số quyết định được xác nhận, và số khách bị phá lời hứa trên số khách được phục vụ. Cột cuối ghi số
kịch bản có ít nhất một khách bị phá lời hứa trên 34 kịch bản. Định nghĩa ở Bảng~\ref{bang:c4-chiso}. Ở hàng luật C,
các số 0 là số đo trên tập đã chạy, phù hợp với tính chất có điều kiện ở Nhận xét~\ref{nx:giunguyen}.}
\label{bang:c4-nghiavu}
\small
\setlength{\tabcolsep}{4pt}
\begin{tabular}{@{}p{\dimexpr .12\linewidth-2\tabcolsep\relax}
>{\raggedleft\arraybackslash}p{\dimexpr .28\linewidth-2\tabcolsep\relax}
>{\raggedleft\arraybackslash}p{\dimexpr .31\linewidth-2\tabcolsep\relax}
>{\raggedleft\arraybackslash}p{\dimexpr .29\linewidth-2\tabcolsep\relax}@{}}
\toprule
\textbf{Cấu hình} & \textbf{Quyết định bị ép / được xác nhận} &
\textbf{Khách bị phá lời hứa / được phục vụ} &
\textbf{Kịch bản có khách bị phá / tổng kịch bản}\\
\midrule
C-30 & 0 / 20.101 & 0 / 2.338 & 0 / 34\\
M$^{+}$-30 & 7.710 / 20.103 & 289 / 2.338 & 28 / 34\\
M$^{+}$-60 & 4.031 / 20.163 & 79 / 2.347 & 22 / 34\\
V-30 & 10.442 / 20.090 & 456 / 2.336 & 34 / 34\\
\bottomrule
\end{tabular}
\end{table}
% EV: R36-A01 (sổ bổ sung trong CHANGES-v3.md; 170 tệp IVC và đối chiếu PAIRS); N-21; X3-30; A22

Cách đọc bảng: luật hạn chót 30 giây có 7.710 trong 20.103 lượt quyết định bị ép và 289 trong 2.338 khách tới muộn
hơn lời hứa đầu quá 30 giây; những khách này nằm ở 28 trong 34 kịch bản.
% EV: N-21; R36-A01

Trên cả 34 kịch bản đã chạy, luật C không có quyết định bị ép và không khách nào vượt trần. Theo Nhận
xét~\ref{nx:giunguyen}, nếu giả thiết (N) đúng ở lượt đang xét và mọi khách có lời hứa của xe có $0 \le L_r \le B$,
phương án giữ nguyên không thêm phần dời do điều phối nên vẫn qua kiểm tra trần. Khi (N) đúng ở mọi lượt, trần được
giữ từ lúc xác nhận. Các số đo ở đây phù hợp với tính chất đó; nhận xét không bảo đảm giờ trả thực tế. Điều phối không
trần cũng có 0 trong 20.760 quyết định bị ép. Luật hạn chót so giờ trả mới với lời hứa đầu, nên phần thay đổi khi giữ
nguyên kế hoạch vẫn có thể làm xe bị ép (mục~\ref{muc:doichung}).
% EV: A21; A22; N-21; X3-16 (U 0 / 20.760); X3-30; A10; K05; Ch3 v4-v3:371-390; R36-A01

Tính gộp, luật hạn chót 30 giây bị ép ở 7.710 trong 20.103 lượt quyết định (38,4\%) và phá lời hứa với 289 trong 2.338
khách (12,4\%). Gánh nặng dồn vào khung 7 giờ: cả 9 kịch bản có khách bị phá lời hứa, với 181 trong 509 khách được phục
vụ và 3.447 trong 4.564 lượt quyết định bị ép. Ở khung 17 giờ, 5 trong 11 kịch bản có khách bị phá, với 8 trong 852
khách và 328 trong 7.163 lượt quyết định bị ép. Các số này được gộp theo từng khung giờ. Ở mức 60 giây, các tỉ lệ toàn
tập giảm nhưng không về 0. Luật trần độ dời khách thấy bị ép nhiều nhất, phù hợp với định nghĩa: tổng độ lớn các lần
cập nhật mà nó kiểm không nhỏ hơn độ lùi so với lời hứa đầu (công thức~\eqref{eq:dongnhat}).
% EV: N-21 (gộp 7.710/20.103 và 289/2.338); R36-A02 (cặp đếm theo khung giờ, sổ bổ sung trong CHANGES-v3.md); C14; C13; R36-A01; A11; A12; X3-33 (eq:dongnhat: Σ|z| ≥ |Σz| = |p_n − p_0|)

Bảng~\ref{bang:c4-nghiavu} không so hai chuẩn đúng giờ: luật C hứa về phần dời do điều phối, còn luật hạn chót hứa về
toàn bộ độ trễ so với lời hứa đầu, một lời hứa lớn hơn. Đo bằng chính thước của luật hạn chót, tức tới muộn hơn lời
hứa đầu quá 30 giây, luật C có 292 trong 2.338 khách (12,5\%), so với 289 khách (12,4\%) của luật hạn chót. Cả hai
luật đều có khách tới muộn hơn lời hứa đầu quá 30 giây ở 28 trong 34 kịch bản. Vậy về độ đúng giờ, hai luật gần như
ngang nhau; luật C khác ở chỗ chỉ hứa phần mà bộ điều phối kiểm soát được, nên giữ được lời hứa đó.
% EV: C13; K05; C4A-01 (C-30 292/2.338 = 12,49%, 28/34); N-21; R36-A01

Bảng~\ref{bang:c4-phanbotre} cho thấy số lượt trễ giảm ở cả bốn ngưỡng đã đo, từ 1 tới 10 phút.
% EV: N-24

\begin{table}[htbp]
\centering
\caption[Phân bố trễ theo ngưỡng khi chưa có dự báo]{Phân bố trễ so với lời hứa đầu không dự báo trên tập kiểm tra chính: số lượt trễ quá mỗi ngưỡng
(trong số khách được phục vụ) và số lượt đó tính trên 100 yêu cầu (3.672 yêu cầu mỗi cấu hình).}
\label{bang:c4-phanbotre}
\small
\begin{tabular}{@{}lrrrr@{}}
\toprule
& \multicolumn{2}{c}{\textbf{U}} & \multicolumn{2}{c}{\textbf{C-30}}\\
\cmidrule(lr){2-3}\cmidrule(lr){4-5}
\textbf{Trễ quá} & \textbf{Lượt} & \textbf{Trên 100 yêu cầu} & \textbf{Lượt} & \textbf{Trên 100 yêu cầu}\\
\midrule
1 phút & 203 & 5,53 & 76 & 2,07\\
2 phút & 102 & 2,78 & 11 & 0,30\\
5 phút & 18 & 0,49 & 0 & 0\\
10 phút & 2 & 0,05 & 0 & 0\\
\midrule
Trễ lớn nhất (s) & \multicolumn{2}{c}{660,1} & \multicolumn{2}{c}{232,3}\\
\bottomrule
\end{tabular}
\end{table}
% EV: N-24 (203/102/18/2 và 76/11/0/0; 5,53/2,78/0,49/0,05 và 2,07/0,30/0/0); C07; C28; C4A-14 (CON:4,11: +660,1 s và +232,3 s; khớp C26 "11 phút" và "3,87 phút")

Cách đọc bảng: khi không trần, 18 lượt tới muộn hơn lời hứa đầu quá 5 phút, tức 0,49 lượt trên 100 yêu cầu.
% EV: N-24; C28

Khi không trần, lượt trễ nhất muộn 660,1 giây (khoảng 11 phút); dưới trần 30 giây là 232,3 giây (3,87 phút). Trong 76
khách còn trễ quá một phút dưới luật C, 72 người không có lần dời do điều phối nào được ghi và 61 người thuộc khung 7
giờ. Lời hứa theo dự báo (mục~\ref{muc:c4-dubao}) nhắm tới phần đoán được của loại trễ này.
% EV: N-24; C4A-14; C26; C21 (4/76 từng bị dời); K39 (không viết "trễ thuần vì kẹt xe"); N-11 (trễ quá 60 s của C-30 theo khung 61 / 15 / 0)

Tách theo kịch bản, luật C không có nhiều lượt trễ quá một phút hơn điều phối không trần ở cả 34 kịch bản và có số
lượt không trễ quá một phút bằng hoặc nhiều hơn ở 29 kịch bản. Ở khung 8 và 17 giờ, điều phối không trần nhiều lượt
trễ hơn ở cả 25 kịch bản; ở khung 7 giờ, luật C ít hơn ở 3 kịch bản và bằng ở 6. Vì 34 kịch bản thuộc 5 ngày và không
độc lập, đây là phép đếm mô tả, không phải kiểm định thống kê.
% EV: N-22 (34/34 trễ ≤ U; 29/34 không trễ ≥ U; bảng theo ngày × khung ở ADD1:173-186); C30; I05; B11

% -----------------------------------------------------------------------------------------------------
\section{Phần do điều phối: luật C chặn gì và mất gì}
\label{muc:c4-dieuphoi}

Dưới luật C trần 30 giây, 33 trong 2.338 khách được phục vụ bị điều phối dời giờ trả, mỗi khách nhiều nhất 29,2 giây;
khi không trần là 174 trong 2.434 khách, một khách bị dời tổng cộng 590,1 giây (Bảng~\ref{bang:c4-doidieuphoi}). Cái
giá là số yêu cầu bị từ chối tăng từ 1.238 lên 1.334; chỉ 148 trong số đó có phương án bị trần chặn được ghi lại.
% EV: C20; N-23; C01 (1.238 = 3.672 − 2.434; CON:3); D13 (1.334; 148)

\begin{table}[htbp]
\centering
\caption[Phần dời giờ trả do điều phối trên tập kiểm tra chính]{Phần dời giờ trả do điều phối trên tập kiểm tra chính, lời hứa đầu không dự báo. Mẫu số của
dòng đầu là số khách được phục vụ (2.434, 2.338 và 2.387); độ dời của một khách là tổng $|c|$ của khách
đó. Dòng cuối lấy tổng phần dời có dấu do điều phối chia cho tổng số giây trễ so với lời hứa đầu của chính các khách
trễ quá 60 s: 23.027 trên 31.353 giây, 82 trên 6.892 giây và 4.836 trên 12.450 giây. Cột trần 120 giây được dùng
lại ở mục~\ref{muc:c4-chontran}.}
\label{bang:c4-doidieuphoi}
\small
\begin{tabular}{@{}p{7.6cm}rrr@{}}
\toprule
\textbf{Chỉ số} & \textbf{U} & \textbf{C-30} & \textbf{C-120}\\
\midrule
Khách bị điều phối dời giờ trả & 174 (7,15\%) & 33 (1,41\%) & 103 (4,32\%)\\
Khách bị dời từ hai lần trở lên & 12 & 0 & 0\\
Tổng dời lớn nhất của một khách (s) & 590,1 & 29,2 & 118,9\\
Trung vị độ lớn mỗi lần dời (s) & 108,6 & 8,3 & 62,5\\
% EV: C20; C22 (trung vị theo lần: 108,6 s trên 187 lần, 8,3 s trên 33 lần); N-07 (C-120: 103, 0 khách ≥ 2 lần, nên trung vị theo lần = trung vị theo khách 62,54 s); N-23
Khách trễ quá 60 s từng bị dời & 134 trên 203 & 4 trên 76 & 60 trên 132\\
Phần tính cho điều phối trong tổng giây trễ của các khách trễ quá 60 s (\%) & 73,4 & 1,2 & 38,8\\
% EV: C21; C24; N-23 (cột C-120: 60/132, 38,8%); C4A-03 (tử số, mẫu số: 23.027,239 / 31.353,407; 81,742 / 6.891,526; 4.836,281 / 12.450,223 s)
\bottomrule
\end{tabular}
\end{table}

Cách đọc bảng: khi không trần, 174 khách bị điều phối dời giờ trả ít nhất một lần, 12 người trong số đó bị dời từ hai
lần, và tổng dời lớn nhất của một khách là 590,1 giây.
% EV: C20

Dưới trần 30 giây, mỗi khách bị dời chỉ một lần, trung vị 8,3 giây, so với 108,6 giây trên 187 lần dời khi không trần.
Trong 33 khách bị dời, 29 người bị dời muộn đi; 4 người bị kéo sớm hơn, mỗi người đúng 1 mili giây. Khi không trần,
134 trong 203 khách trễ quá một phút từng bị dời; con số chỉ cho biết hai nhóm trùng nhau, không cho biết việc bị dời
là nguyên nhân của độ trễ. Theo sổ theo từng khách (công thức~\eqref{eq:so}), phần tính cho điều phối bằng 73,4\% tổng
số giây trễ của 203 khách đó, còn dưới trần 30 giây chỉ 1,2\%: gần như mọi giây trễ còn lại nằm ở hai phần kia của sổ.
% EV: E02; C22; C23 (29 muộn, 4 sớm); N-05 (mỗi lần dời sớm đúng −1 ms; chỉ kiểm ở U, C-30, C-120, M⁺-30; nguyên nhân chưa kiểm, ADDENDUM-1 mục 0.2; C4A-19); C21 (giao hai tập); C24 (sổ mô hình, không nhân quả); K04; X3-12

Giới hạn này không dựa vào dự báo, vì dự báo không đi vào bộ điều phối. Ở 13 kịch bản cuối tuần không có dự báo, tổng
dời lớn nhất của một khách là 29,2 giây ở trần 30 giây và 116,5 giây ở trần 120 giây. Trên 8 kịch bản chiều tuyết,
8 trong 601 khách được phục vụ bị dời, lớn nhất 25,2 giây; không ai vượt trần và không quyết định nào bị ép trong
5.067 lượt.
% EV: F03; X3-11 (D chỉ tính sau khi chạy); N-28; N-26 (thay anchor F13 của RD §8 mục 4, ADDENDUM-1 mục 0.4); N4B8-05 (CK8:65-72); R26 §Q8

Một ví dụ trên tập buổi sáng cho thấy luật C và luật hạn chót khác nhau ở đâu. Trước một quyết định, giờ trả đang hứa
của một khách đã lùi 79,0 giây so với lời hứa đầu, toàn bộ là phần thay đổi khi giữ nguyên kế hoạch. Phương án ghép
thêm một khách mới sẽ dời giờ trả dự kiến của khách đó thêm 68,3 giây do điều phối. Ở mức 90 giây, luật hạn chót cộng
hai phần, khoảng 147 giây, nên loại phương án; luật C chỉ tính 68,3 giây nên cho qua. Trên tập này hai luật tách nhau
ở đúng ba chỗ cùng kiểu, ở mức 90 hoặc 120 giây; mỗi chỗ luật C nhận thêm một khách. Đây là phép so hai luật trên
giờ trả dự kiến tại cùng một trạng thái của mô hình, không phải phép đo nhân quả.
% EV: D05 (79,0 + 68,3 s; 60,9 + 61,2 s); C4A-12 (W07XM:6,11,16: "sum x 79.0, sum c 0.0, p_last − p_first 79.0", c = 68,3 s, x = 0,1 s; W07X:54-62: "khách đang có", giai đoạn chờ đón hay trên xe không ghi); K04

Khi luật C loại một phương án vì nó làm một khách đã có lời hứa vượt trần, bộ kiểm ghi một bản ghi chặn gồm khách đó,
phần trần đã dùng, lần dời mới và trần (Bảng~\ref{bang:kiemtran}). Dưới trần 30 giây, 148 trong 1.334 yêu cầu bị từ
chối có ít nhất một phương án mang bản ghi chặn; điều phối không trần phục vụ 116 trong số đó. 1.186 yêu cầu còn lại
không có bản ghi chặn nào được ghi, và điều phối không trần cũng từ chối 1.075 yêu cầu trong số này; điều đó không
chứng minh chúng bị từ chối chỉ vì ràng buộc dịch vụ. Bản ghi chặn cũng không phải khách: 845 bản ghi nằm ở 643 yêu
cầu, và 495 yêu cầu trong số đó cuối cùng vẫn được phục vụ.
Các số đếm theo yêu cầu trong đoạn này tính theo từng cặp kịch bản và yêu cầu.
% EV: X3-05 (nội dung bản ghi chặn); D13 (BUD:2-5: 148 = 116 + 32; 1.186 = 111 + 1.075); K10; C4A-09 (BUDPY:29-44: nhóm 1.186 là phần dư không có bản ghi chặn được ghi, không phải lý do cấu trúc đã chứng minh); D16 (SF:2: 643 cặp, 495 phục vụ, 148 từ chối); E01 (845)

Với mỗi yêu cầu trong 148 yêu cầu, phép phân tích chọn bản ghi có mức vượt nhỏ nhất, tức phần vượt quá 30 giây; trung
vị là 109,6 giây, lớn nhất 524,8 giây. Chỉ ở 2 trong 148 bản ghi được chọn, khách nêu trong bản ghi đã dùng một phần
trần trước đó; ở cả 148, riêng lần dời mới đã vượt 30 giây. Với trần 120 giây, mức vượt của 59 trong 148 yêu cầu nằm
trong phần nới thêm 90 giây. Không thể kết luận 59 yêu cầu đó sẽ được phục vụ: bộ kiểm dừng ở khách vi phạm gặp trước,
và khi trần đổi thì các quyết định trước cũng đổi. Mục~\ref{muc:c4-chontran} cho kết quả chạy thật.
% EV: D14 (BUD:6-8: trung vị 109,6 s; lớn nhất 524,8 s; 59/148 có mức vượt ≤ 90 s); D19 (BUD:10; FJC:4-6); C4A-09 (BUDPY:22,35-42,51-52: mức vượt = after − limit; 2/148 đếm theo bản ghi được chọn của từng yêu cầu, không phải số khách khác nhau); K11; A23; I08

% -----------------------------------------------------------------------------------------------------
\section{Chọn mức trần}
\label{muc:c4-chontran}

Trong ba điểm đã đo trên tập kiểm tra chính, không điểm nào tốt nhất ở mọi chỉ số. Trần 30 giây có ít lượt trễ quá một
phút nhất (76). Trần 120 giây có nhiều lượt không trễ quá hai phút nhất (2.366, so với 2.327 của trần 30 giây và 2.332
khi không trần) và phục vụ 2.387 lượt, nằm giữa hai điểm kia. Chọn mức nào tùy ngưỡng mà hãng xe coi trọng, điều thực
nghiệm không trả lời được.
% EV: N-06; L1a-02; C01; C06; C08; K12; K36

Mức trần được đo trên ba tập không cộng gộp (Hình~\ref{hinh:c4-tran}). Tập tắc đường cố định cho thấy trần làm gì khi
thời gian đi không đổi, tập buổi sáng dò tám mức trên cùng 8 kịch bản, tập kiểm tra chính thêm mức 120 giây.
% EV: B17; B15; B16; X3-28; D01; L1a-01; I12

\textbf{Tập tắc đường cố định.} Bảng~\ref{bang:c4-codinh} xếp các cấu hình từ lỏng nhất (ghép tự do) tới chặt nhất
(trần 30 giây), với hệ số tắc đường cố định suốt lần chạy.
% EV: B17; LAD:2-6

\begin{table}[htbp]
\centering
\caption[Các mức trần trên tập tắc đường cố định]{Các mức trần trên tập tắc đường cố định (16 kịch bản, 1.728 yêu cầu mỗi cấu hình; bộ chạy cũ,
không có luật phục hồi). Ba cột số đầu là số đếm; hai cột cuối là số khách được phục vụ thỏa điều kiện, trong ngoặc
là tỉ lệ gộp theo khách được phục vụ (\%).}
\label{bang:c4-codinh}
\small
\begin{tabular}{@{}lrrrrr@{}}
\toprule
\textbf{Cấu hình} & \textbf{Phục vụ} & \textbf{Trễ quá 60 s} & \textbf{Không trễ quá 120 s} &
\textbf{Đi chung} & \textbf{Bị dời giờ trả}\\
\midrule
FREE & 1.629 & 274 & 1.465 & 865 (53,1) & 388 (23,8)\\
U & 1.584 & 47 & 1.557 & 346 (21,8) & 62 (3,9)\\
C-120 & 1.563 & 23 & 1.563 & 303 (19,4) & 41 (2,6)\\
C-60 & 1.540 & 0 & 1.540 & 266 (17,3) & 20 (1,3)\\
C-30 & 1.534 & 0 & 1.534 & 241 (15,7) & 6 (0,4)\\
\bottomrule
\end{tabular}
\end{table}
% EV: D09 (số đếm, WP17); D08 (tỉ lệ đi chung, bị dời; gộp theo khách); C4A-04 (đi chung 865/346/303/266/241, CK4A fixed-shared; bị dời 388/62/41/20/6, WBU:4,7,10,13,16); B17

Cách đọc bảng: với trần 120 giây, 1.563 trong 1.728 yêu cầu được phục vụ, 23 lượt trễ quá 60 giây; 303 khách (19,4\%)
đi chung và 41 khách (2,6\%) bị dời giờ trả.
% EV: D08; D09; C4A-04

Ghép tự do khác điều phối không trần ở cả cách chọn phương án lẫn việc kiểm cam kết, nên chênh lệch giữa hai cấu hình
này không chỉ do trần. Trong nhóm cùng chính sách, siết từ không trần xuống trần 30 giây làm số lượt phục vụ giảm từ
1.584 xuống 1.534 và số lượt trễ quá 60 giây từ 47 xuống 0. Giữa hai mức đã chạy 60 và 120 giây, số lượt trễ tăng từ 0
lên 23. Các con số của bộ chạy cũ này không chuyển sang tập kiểm tra chính.
% EV: A17; K50; X3-38; NEW:13; D09; D11; B17; I12

\textbf{Tập chỉnh trần buổi sáng.} Trên 864 yêu cầu, số lượt phục vụ chỉ dao động trong 6 lượt, từ 660 ở trần 30 giây
tới 666 ở trần 120--180 giây (Bảng~\ref{bang:c4-sang}).
% EV: D01; B15; B16

\begin{table}[htbp]
\centering
\caption[Các mức trần trên tập chỉnh trần buổi sáng]{Tám mức trần trên tập chỉnh trần buổi sáng (8 kịch bản khung 7 giờ ngày 12--13/11/2018, 864 yêu cầu
mỗi cấu hình; lời hứa đầu không dự báo). Số đếm gộp; hai cột ``không trễ quá'' là số lượt trên 864 yêu cầu. Các mức
có cùng mọi số đếm được gộp vào một dòng.}
\label{bang:c4-sang}
\small
\begin{tabular}{@{}lrrrr@{}}
\toprule
\textbf{Cấu hình} & \textbf{Phục vụ} & \textbf{Trễ quá 60 s} & \textbf{Không trễ quá 60 s} &
\textbf{Không trễ quá 120 s}\\
\midrule
U & 665 & 50 & 615 & 656\\
C-30 & 660 & 47 & 613 & 658\\
C-60 & 661 & 47 & 614 & 659\\
C-90 & 664 & 50 & 614 & 660\\
C-120, C-150, C-180 & 666 & 51 & 615 & 659\\
C-240, C-300 & 665 & 50 & 615 & 656\\
\bottomrule
\end{tabular}
\end{table}
% EV: D01; D02 (U, C-30, C-90, C-120: CAP:3-7; DIAL17:3-7, hai bộ đọc độc lập); C4A-05 (C-60, C-150, C-180, C-240, C-300: CAP:5,8-11; DIAL17:5,8-11); D02 lưu ý (cột dự báo bằng ngây thơ ở 12–13/11)

Cách đọc bảng: với trần 120, 150 hoặc 180 giây, 666 yêu cầu được phục vụ, 51 lượt trễ quá 60 giây, 615 lượt không trễ
quá 60 giây và 659 lượt không trễ quá 120 giây.
% EV: D01; D02; C4A-05

Trần 120 giây phục vụ nhiều hơn điều phối không trần một lượt, nên điều phối không trần không phải mức phục vụ cao
nhất có thể đạt, chỉ là một điểm đo. Từ 30 lên 120 giây, trần thêm 6 lượt phục vụ, 4 lượt trễ quá 60 giây và 2 lượt
không trễ quá 60 giây. Luật C không có quyết định bị ép trong 64 lần chạy của tám mức trần này. Ở cả tám mức, hai
luật phục vụ chênh nhau không quá 2 lượt.
% EV: D01; K12; D03 (+6 / +4 / +2); D04; X3-28 (tám mức 30–300 s); E01 (64 lần C, 0 bị ép); agent-numeric-output-v2.txt (0 / 44.270 lượt quyết định)

\textbf{Tập kiểm tra chính.} Bảng~\ref{bang:c4-tran120} đặt trần 120 giây cạnh hai cấu hình của
mục~\ref{muc:c4-ketqua}, khi không có và khi có dự báo.

\begin{table}[htbp]
\centering
\caption[Trần 120 giây trên tập kiểm tra chính]{Trần 120 giây trên tập kiểm tra chính (34 kịch bản, 3.672 yêu cầu mỗi cấu hình). Số đếm gộp; các dòng
``không trễ quá'' là số lượt trên 3.672 yêu cầu. Khối dưới là phép phát lại lời hứa theo dự báo trên cùng các lần
chạy, bàn ở mục~\ref{muc:c4-dubao}; dự báo chỉ có ở 21 kịch bản ngày thường, và độ dài thêm của lời hứa là trung
bình gộp trên mọi khách được phục vụ, kể cả khách cuối tuần có độ dài thêm bằng 0.}
\label{bang:c4-tran120}
\small
\begin{tabular}{@{}p{8.4cm}rrr@{}}
\toprule
\textbf{Chỉ số} & \textbf{U} & \textbf{C-30} & \textbf{C-120}\\
\midrule
\multicolumn{4}{@{}l}{\textit{Lời hứa đầu không dự báo}}\\
Lượt phục vụ & 2.434 & 2.338 & 2.387\\
Lượt trễ quá 60 s & 203 & 76 & 132\\
Lượt trễ quá 120 s & 102 & 11 & 21\\
Lượt trễ quá 300 s & 18 & 0 & 0\\
Lượt tới sớm quá 60 s & 7 & 3 & 1\\
Không trễ quá 60 s & 2.231 & 2.262 & 2.255\\
Không trễ quá 120 s & 2.332 & 2.327 & 2.366\\
Khách đi chung & 555 & 313 & 428\\
% EV: L1a-02; L1a-04 (U, C-30 tái lập); N-06; C06; C07; C08; N-09 (L1a-07); C4A-13 (sớm quá 60 s 7 / 3 / 1, ADD1:12,17,22; C120R:14,16,18)
\midrule
\multicolumn{4}{@{}l}{\textit{Lời hứa đầu theo dự báo}}\\
Lượt trễ quá 60 s & 137 & 10 & 65\\
Lượt tới sớm quá 60 s & 8 & 4 & 2\\
Không trễ quá 60 s & 2.297 & 2.328 & 2.322\\
Không trễ quá 120 s & 2.347 & 2.338 & 2.379\\
Lời hứa dài thêm, trung bình gộp theo khách (s) & 7,18 & 7,37 & 7,27\\
% EV: L1a-03; N-03; N-06; N-16 (sớm quá 60 s 8 / 4 / 2); C09; F06
\bottomrule
\end{tabular}
\end{table}

Cách đọc bảng: với trần 120 giây và lời hứa đầu không dự báo, 2.387 yêu cầu được phục vụ, 132 lượt trễ quá 60 giây,
nên 2.255 lượt không trễ quá 60 giây trên 3.672 yêu cầu. Chỉ 1 lượt tới sớm quá 60 giây.
% EV: N-06; L1a-02; C4A-13

So với trần 30 giây, trần 120 giây phục vụ thêm 49 lượt và có thêm 56 lượt trễ quá một phút, nên ít hơn 7 lượt không
trễ quá một phút. Nhưng nó thêm 39 lượt không trễ quá hai phút, nhiều hơn cả trần 30 giây lẫn điều phối không trần.
Khi có dự báo, nó ít hơn trần 30 giây 6 lượt ở ngưỡng một phút và thêm 41 lượt ở ngưỡng hai phút. Độ dài thêm của lời
hứa gần như bằng nhau (7,27 so với 7,37 giây, trên hai tập khách khác nhau). Mục~\ref{muc:c4-dem} đặt điểm này cạnh
cách đệm lời hứa.
% EV: L1a-09 (+49 / +56 / −7 / +39; dự báo −6 / +41); N-06; L1a-02 (2.366 > U 2.332 và C-30 2.327); K12; L1a-03; N-16; L1a-08

Dưới trần 120 giây, 9 khách bị dời từ 108 giây trở lên (90\% của trần), không khách nào vượt trần và không quyết định
nào bị ép trong 20.435 lượt.
% EV: N-07 (≥ 0,9 B: 9); L1a-05; N-23

Hai tập khách khác nhau: 2.266 khách được cả hai mức trần phục vụ, 121 khách chỉ được trần 120 giây và 72 khách chỉ
được trần 30 giây phục vụ; 93 trong 121 khách đó cũng được phục vụ khi không trần. Vì vậy 49 lượt thêm là chênh lệch
hai tổng, không phải 49 trong 96 lượt mà trần 30 giây bỏ. Phần thêm dồn vào khung 8 và 17 giờ; so từng kịch bản, trần
120 giây phục vụ nhiều hơn ở 23, bằng ở 8 và ít hơn ở 3.
% EV: N-10 (lưu ý: không viết "lấy lại 49 trong 96 lượt"); L1a-08; N-11 (C-120 − C-30: phục vụ +4 / +26 / +19, trễ quá 60 s +2 / +34 / +20; 23 / 8 / 3 kịch bản); L1a-09

\begin{figure}[htbp]
\centering
\begin{tikzpicture}[font=\footnotesize,
  vuong/.style={draw,fill=white,rectangle,minimum size=4.2pt,inner sep=0pt},
  tamgiac/.style={draw,fill=white,regular polygon,regular polygon sides=3,minimum size=6pt,inner sep=0pt},
  tron/.style={draw,fill=white,circle,minimum size=4.6pt,inner sep=0pt},
  thoi/.style={draw,fill=black,diamond,minimum size=6.5pt,inner sep=0pt},
  nhan/.style={font=\scriptsize,inner sep=1.5pt}]
% (a) Tập tắc đường cố định. Trục ngang: phục vụ − 1.530, 0,06 cm mỗi lượt; trục dọc: trễ, 0,05 cm mỗi lượt.
\begin{scope}[shift={(0,0)}]
\draw[->] (0,0)--(3.95,0);
\draw[->] (0,0)--(0,3.35);
\foreach \vx/\so in {0/1.530,1.2/1.550,2.4/1.570,3.6/1.590}
  \draw (\vx,0)--(\vx,-0.07) node[below,nhan]{\so};
\foreach \vy/\so in {0/0,1/20,2/40,3/60}
  \draw (0,\vy)--(-0.07,\vy) node[left,nhan]{\so};
\draw[dashed] (0.24,0)--(0.60,0)--(1.98,1.15)--(3.24,2.35);
\node[vuong] at (0.24,0) {}; \node[nhan,above] at (0.24,0.05) {30};
\node[vuong] at (0.60,0) {}; \node[nhan,above] at (0.62,0.05) {60};
\node[vuong] at (1.98,1.15) {}; \node[nhan,above left] at (1.98,1.15) {120};
\node[vuong] at (3.24,2.35) {}; \node[nhan,above left] at (3.24,2.35) {U};
\node at (1.8,-0.72) {Lượt phục vụ};
\node[rotate=90] at (-0.85,1.6) {Lượt trễ quá 60 s};
\node[anchor=south] at (1.9,3.4) {(a) Tắc đường cố định};
\end{scope}
% (b) Tập chỉnh trần buổi sáng. Trục ngang: phục vụ − 658, 0,36 cm mỗi lượt; trục dọc: trễ, 0,05 cm mỗi lượt.
% C-120, C-150, C-180 trùng một điểm (666; 51); U, C-240, C-300 trùng một điểm (665; 50).
\begin{scope}[shift={(5.3,0)}]
\draw[->] (0,0)--(3.95,0);
\draw[->] (0,0)--(0,3.35);
\foreach \vx/\so in {0/658,0.72/660,1.44/662,2.16/664,2.88/666,3.6/668}
  \draw (\vx,0)--(\vx,-0.07) node[below,nhan]{\so};
\foreach \vy/\so in {0/0,1/20,2/40,3/60}
  \draw (0,\vy)--(-0.07,\vy) node[left,nhan]{\so};
\draw[densely dotted,thick] (0.72,2.35)--(1.08,2.35)--(2.16,2.50)--(2.88,2.55)--(2.52,2.50);
\node[tamgiac] at (0.72,2.35) {}; \node[nhan,below] at (0.72,2.25) {30};
\node[tamgiac] at (1.08,2.35) {}; \node[nhan,below] at (1.10,2.25) {60};
\node[tamgiac] at (2.16,2.50) {}; \node[nhan,below] at (2.16,2.40) {90};
\node[tamgiac] at (2.88,2.55) {}; \node[nhan,right] at (2.98,2.55) {120--180};
\node[tamgiac] at (2.52,2.50) {}; \node[nhan,above left] at (2.50,2.58) {U, 240, 300};
\node at (1.8,-0.72) {Lượt phục vụ};
\node[rotate=90] at (-0.85,1.6) {Lượt trễ quá 60 s};
\node[anchor=south] at (1.9,3.4) {(b) Buổi sáng};
\end{scope}
% (c) Tập kiểm tra chính. Trục ngang: phục vụ − 2.320, 0,03 cm mỗi lượt; trục dọc: trễ, 0,0125 cm mỗi lượt.
\begin{scope}[shift={(10.6,0)}]
\draw[->] (0,0)--(3.95,0);
\draw[->] (0,0)--(0,3.35);
\foreach \vx/\so in {0/2.320,1.2/2.360,2.4/2.400,3.6/2.440}
  \draw (\vx,0)--(\vx,-0.07) node[below,nhan]{\so};
\foreach \vy/\so in {0/0,1/80,2/160,3/240}
  \draw (0,\vy)--(-0.07,\vy) node[left,nhan]{\so};
\draw (0.54,0.95)--(2.01,1.65)--(3.42,2.54);
\node[tron] at (0.54,0.95) {}; \node[nhan,below right] at (0.54,0.95) {30};
\node[thoi] at (2.01,1.65) {}; \node[nhan,below right] at (2.06,1.62) {120};
\node[tron] at (3.42,2.54) {}; \node[nhan,above left] at (3.42,2.54) {U};
\node at (1.8,-0.72) {Lượt phục vụ};
\node[rotate=90] at (-0.85,1.6) {Lượt trễ quá 60 s};
\node[anchor=south] at (1.9,3.4) {(c) Tập kiểm tra chính};
\end{scope}
\end{tikzpicture}
\caption[Lượt trễ theo lượt phục vụ ở các mức trần]{Lượt trễ quá 60 giây so với lời hứa đầu (trục dọc) theo lượt phục vụ (trục ngang) ở các mức
trần, lời hứa đầu không dự báo. Mỗi khung là một tập với thang riêng: (a) tập tắc đường cố định, 1.728
yêu cầu; (b) tập chỉnh trần buổi sáng, 864 yêu cầu, đủ tám mức, các mức có cùng số đếm chung một điểm; (c) tập
kiểm tra chính, 3.672 yêu cầu. Số cạnh mỗi điểm là trần tính bằng giây, U là điều phối không trần. Đường nối các mức
của cùng một tập theo thứ tự trần tăng dần, chỉ giúp theo dõi các điểm đã đo và không biểu diễn kết quả ở các mức
chưa chạy; các tập không nối với nhau. Ở (a), ghép tự do (1.629 lượt phục vụ, 274 lượt trễ) nằm ngoài khung. Hình
thoi đen là trần 120 giây trên tập kiểm tra chính.}
\label{hinh:c4-tran}
\end{figure}
% EV (a): D09 (C-30 1.534/0; C-60 1.540/0; C-120 1.563/23; U 1.584/47; FREE 1.629/274)
% EV (b): D01; D02; C4A-05 (C-30 660/47; C-60 661/47; C-90 664/50; C-120 = C-150 = C-180 666/51; U = C-240 = C-300 665/50)
% EV (c): C01; C06; N-06 (C-30 2.338/76; C-120 2.387/132; U 2.434/203)
% Tọa độ: (a) x = (phục vụ − 1530) × 0,06, y = trễ × 0,05; (b) x = (phục vụ − 658) × 0,36, y = trễ × 0,05
% (C-60: x = 3 × 0,36 = 1,08; y = 47 × 0,05 = 2,35); (c) x = (phục vụ − 2320) × 0,03, y = trễ × 0,0125.

Trên cả ba tập, nới trần từ 30 lên 120 giây thêm lượt phục vụ và thêm lượt trễ quá một phút. Số lượt không trễ quá một
phút thêm 6, thêm 2 và bớt 7 lần lượt ở tập tắc đường cố định, tập buổi sáng và tập chính, còn ở ngưỡng hai phút trần
120 giây không kém trần 30 giây ở cả ba tập. Vì mỗi kịch bản chỉ chạy một lần và mức 120 giây được chọn sau khi đã xem
tập kiểm tra chính, trần 120 giây chỉ là một điểm đo thêm trên đường đánh đổi.
% EV: D10; D03; L1a-09; D09 (không trễ quá 120 s: 1.563 so với 1.534); D02 (659 so với 658); L1a-02; K36; D11; B11; I06; C4A-11; K42

% -----------------------------------------------------------------------------------------------------
\section{Phần cộng dồn của luật C: kiểm toàn bộ bản ghi chặn}
\label{muc:c4-congdon}

Luật C cộng mọi lần dời của một khách qua cả chuyến rồi mới so với trần. Trên các lần chạy đã kiểm, hiếm khi chỉ riêng
phép cộng ấy làm một phương án bị loại. Trong 98 lần chạy trần 30--300 giây của tập kiểm tra chính và tập buổi sáng,
cả 1.256 bản ghi chặn đều có một lần dời đơn lẻ đã vượt trần của cấu hình tương ứng. Trên 34 lần chạy trần 120 giây
của tập kiểm tra chính, 3 trong 524 bản ghi chỉ vượt trần khi cộng thêm phần trần đã dùng trước đó.
% EV: E01 (98 = 34 + 64; 1.256 = 845 + 411); K09; N-08; L1a-06

Một bản ghi chặn là \emph{chỉ do cộng dồn} khi lần dời mới không vượt $B$ nhưng phần đã dùng $L_r$ cộng lần dời mới
thì vượt (mục~\ref{muc:luatc}). Ví dụ, ở một bản ghi ngày 14/11 với trần 120 giây, khách đã dùng 25,844 giây và phương
án mới dời thêm 99,027 giây: mỗi phần dưới 120 giây, nhưng tổng 124,871 giây vượt trần. Bảng~\ref{bang:c4-congdon} đếm
các bản ghi này trên mọi bản ghi chặn của các lần chạy đã kiểm.
% EV: A06; E05; N-27 (25,844 + 99,027 = 124,871 s); N-08 (quy tắc riêng + cumulative_decisive_check)

\begin{table}[htbp]
\centering
\caption[Kiểm toàn bộ bản ghi chặn của luật C]{Kiểm toàn bộ bản ghi chặn do trần trên 140 lần chạy có luật C dùng hệ số tắc đường đo được đã được kiểm,
thuộc bốn nhóm dưới đây. Một bản ghi ứng với một phương án bị loại, không phải một khách.}
\label{bang:c4-congdon}
\small
\setlength{\tabcolsep}{4pt}
\begin{tabular}{@{}>{\raggedright\arraybackslash}p{\dimexpr .38\linewidth-2\tabcolsep\relax}
>{\raggedleft\arraybackslash}p{\dimexpr .11\linewidth-2\tabcolsep\relax}
>{\raggedleft\arraybackslash}p{\dimexpr .16\linewidth-2\tabcolsep\relax}
>{\raggedleft\arraybackslash}p{\dimexpr .16\linewidth-2\tabcolsep\relax}
>{\raggedleft\arraybackslash}p{\dimexpr .19\linewidth-2\tabcolsep\relax}@{}}
\toprule
\textbf{Lần chạy} & \textbf{Số lần chạy} & \textbf{Bản ghi chặn} & \textbf{Chỉ do cộng dồn} &
\textbf{Khách phục vụ bị dời từ hai lần}\\
\midrule
Tập kiểm tra chính, trần 30 giây & 34 & 845 & 0 & 0\\
Tập chỉnh trần buổi sáng, trần 30--300 giây & 64 & 411 & 0 & 0\\
% EV: E01 (98 = 34 + 64; 1.256 = 845 + 411; 0 chỉ do cộng dồn); E02 (repeated = 0 trên 98 lần chạy)
Tập kiểm tra chính, trần 120 giây & 34 & 524 & 3 & 0\\
% EV: N-08; L1a-06; N-07 (0 khách ≥ 2 lần)
Chiều tuyết 15/11/2018, trần 30 giây & 8 & 163 & 0 & 0\\
% EV: N-26 (163 bản ghi, 0 chỉ do cộng dồn, ≥ 2 lần 0); N4B8-05 (bộ chạy chiều tuyết đã kiểm, CK8:65-72); R26 §Q8
\midrule
Cộng & 140 & 1.943 & 3 & 0\\
% EV: N-15 (132 / 1.780 / 3; kể cả chiều tuyết 140 / 1.943 / 3); ADDENDUM-1 mục 0.3
\bottomrule
\end{tabular}
\end{table}

Cách đọc bảng: trên 34 lần chạy trần 120 giây có 524 phương án bị loại kèm bản ghi chặn do trần, 3 bản ghi trong số đó
chỉ do cộng dồn, và không khách được phục vụ nào bị dời từ hai lần.
% EV: N-08; N-07

Phạm vi kiểm là 132 lần chạy và 1.780 bản ghi (98 lần chạy với 64.371 lượt quyết định được xác nhận, cộng 34 lần chạy
trần 120 giây), hay 140 lần chạy và 1.943 bản ghi kể cả chiều tuyết. Chưa được kiểm: 15 lần chạy hiệu chỉnh, 16 lần
chạy trong đó bộ điều phối dùng dự báo khi quyết định (khác phép phát lại) và 12 lần chạy luật C dùng hệ số đo được
của một phép quét trước đó. Kết quả vì vậy không áp cho mọi lần chạy.
% EV: N-15 (SCOPE); ADDENDUM-1 mục 0.3 (15 lần chạy hiệu chỉnh, 16 thư mục C-30F); E01 (64.371); K09; F12; C4A-08 (SWEEP:13,70: C-x, x = 30…900 s, trên thế giới đo thật W5, 2 kịch bản mỗi mức; NEW:1 by_arm chỉ có t38 và w07)

Ba bản ghi chỉ do cộng dồn nằm ở ba kịch bản khung 8 giờ (14, 15 và 18/11), trong 3 lượt quyết định; khi không trần,
khách của bản ghi ngày 14/11 bị dời hai lần đúng bằng hai lượng ở ví dụ trên. Trong 3 yêu cầu mới liên quan, 2 cuối
cùng không được nhận dưới trần 120 giây (một trong hai cũng không dưới trần 30 giây). Vì bản ghi dừng ở lỗi gặp trước
và diễn biến của các cấu hình khác nhau, điều này không chứng minh phần cộng dồn làm mất hai yêu cầu đó.
% EV: L1a-06 (ba kịch bản w08 ngày 14, 15, 18/11); N-08 (3 quyết định; 2 không được nhận, 1 phục vụ); N-27; E07; A23; I08

Ngoài hệ số đo được, bản ghi chỉ do cộng dồn còn có ở hai nơi. Trên tập tắc đường cố định, nơi bộ kiểm ghi đủ mọi bản
ghi, có 2 bản ghi ở trần 60 giây và 4 ở trần 120 giây; mọi yêu cầu mới liên quan vẫn được phục vụ. Trên các kiểu tắc
đường tự tạo của phép quét trước đó, sau khi bỏ bản chạy trùng, có 4 bản ghi trong 3 lượt quyết định ở trần 120 giây;
1 yêu cầu mới liên quan cuối cùng bị từ chối.
% EV: E03; E04 (4/3 sau khi bỏ bản lặp; thô 6/5); A23 (NEW:13: tập tắc đường cố định bật chế độ ghi đủ); FJC:31-54; FJR:2-7; C4A-08 (thế giới tự tạo thuộc tier3-sweep-v1)

Về định nghĩa, luật C vẫn khác luật ``mỗi lần dời không quá 30 giây'': luật sau cho một khách bị dời ba lần, tổng 90
giây, còn luật C trần 30 giây thì không. Trên số liệu, khác biệt này hiếm: khi không trần chỉ 12 trong 2.434 khách bị
dời từ hai lần. Tính thử cả hai cách trên các lần chạy không trần đã ghi, chúng chỉ khác kết luận ở 4 khách với trần
120 giây và 0 khách với trần 30, 60 và 300 giây. Phép tính thử không cho biết kết quả nếu chạy thật với luật kia, vì
các quyết định sẽ khác; luật mỗi lần một trần chưa được cài đặt và chưa chạy.

Khi chuyển sang ưu tiên chi phí (mục~\ref{muc:c4-vaitro}), cấu hình có trần 30 giây có một bản ghi chỉ vượt
do cộng dồn: phần đã dùng 24,404 giây cộng lần dời 13,558 giây thành 37,962 giây. Yêu cầu mới liên quan cuối
cùng vẫn được phục vụ. Cấu hình này có hai khách bị sửa giờ trả từ hai lần. Các số 0 của bảng trên chỉ thuộc
phạm vi ưu tiên ổn định đã nêu, không áp cho mọi thứ tự chọn phương án.
% EV: S15-01 (5d--5e); S15-04. Một witness, không phải một yêu cầu cuối cùng bị từ chối.

Trên 2.502 lịch sử khách được phục vụ của cấu hình ưu tiên chi phí không trần, có 7 lịch sử mà mỗi lần dời
không quá 120 giây nhưng tổng độ lớn vượt 120 giây. Ví dụ minh họa, hai lần dời 70 và 60 giây đều dưới trần
từng lần 120 giây, nhưng tổng 130 giây vượt trần cộng dồn. Số 7 đếm lịch sử khách theo phép kiểm ở mức 120
giây, không đếm khách vi phạm trần 30 giây. Nhánh so trực tiếp với luật giới hạn từng lần không được thực
hiện vì phép sàng lọc chưa đạt mốc 10 lịch sử đã đặt để phân bổ nguồn lực. Mốc này không phải ngưỡng kiểm
định thống kê; việc dừng chưa chứng minh phép cộng dồn không có tác dụng.
% EV: S15-01 (H7); S15-02 (h7_served); STAGE1B-PLAN.md mục 6,9. Ví dụ 70+60 là dựng tay.

% EV: E08; C20; E06 (CVP:14-18: 0/0/4/0 trên 2.434 khách U); K42; E09; K08

Các mục trên so luật C với điều phối không trần khi giữ nguyên bộ chọn ưu tiên ổn định, tức bộ chọn xét mức dùng trần
và phần dời trước chi phí kế hoạch (Bảng~\ref{bang:rangbuoc}). Mục~\ref{muc:c4-vaitro} xét kết quả có phụ thuộc vào
thứ tự đó không.
% EV: A08; H01; H02; A14; X3-07

% =====================================================================================================
% Chương 4, PHẦN B (mục 4.6–4.11) — bản sửa v11 (2026-10-13), BẢN DÙNG ĐỂ GHÉP.
% Sửa có mục tiêu từ drafts/09-chuong4-phanB-v10.tex theo R37 và brief R39; các bản trước giữ nguyên.
% Quyết định trước: reviews/CH4B-REVISION-LOG.md; ghi chú cũ: drafts/09-chuong4-phanB-v8-NOTES.md.
% Từng thay đổi, phép kiểm số mới và việc bàn giao: 09-chuong4-phanB-v11-CHANGES-v2.md bên cạnh.
% Nối ngay sau 08-chuong4-phanA-v5-v3.tex trong cùng % =====================================================================================================
% Chương 4, PHẦN A (mở chương, mục 4.1–4.5) — bản sửa v5 (2026-10-13), BẢN DÙNG ĐỂ GHÉP.
% Sửa có mục tiêu từ drafts/08-chuong4-phanA-v4.tex theo R34-CH4A-V4-VERIFY-CODEX-v2.md và brief R36.
% Quyết định trước: reviews/CH4A-REVISION-LOG-v2.md; ghi chú cũ: drafts/08-chuong4-phanA-v2-NOTES.md.
% Bản v4 và mọi bản trước giữ nguyên. Tệp tiếp tục R36: 08-chuong4-phanA-v5-v3.tex;
% từng thay đổi, bằng chứng mới và việc bàn giao: 08-chuong4-phanA-v5-CHANGES-v3.md. Phần B v10 nối sau tệp này.
% Mỗi dòng "% EV:" trỏ tới EVIDENCE-REGISTER.md (A01…K52), các phụ lục 1 (N-xx, L1a-xx), 2 (X3-xx),
% 4 (C4A-xx), 6 (N4B8-xx), hoặc sổ bổ sung R36-A01…03 trong tệp CHANGES-v3.md bên cạnh.
% Không dùng gói mới: TikZ và array đã nạp trong main.tex; chỉ sửa độ rộng các bảng được brief chỉ ra.
% =====================================================================================================

\chapter{Thực nghiệm và đánh giá}
\label{chuong:thucnghiem}

Chương này đo luật C được gì và mất gì trên mô phỏng FleetPy, so với điều phối không trần, luật hạn chót và cách đệm
lời hứa, tức cộng thêm một khoảng cố định vào mọi lời hứa. Chương cũng xét lời hứa theo dự báo và những chỗ dự báo
chưa đủ. Các mục đi theo ba phần của độ lệch giờ trả so với lời hứa đầu (mục~\ref{muc:so}). Phần tính cho điều phối,
gắn với việc ghép thêm khách, ở mục~\ref{muc:c4-ketqua}--\ref{muc:c4-vaitro}; phần được hồ sơ tắc đường lịch sử giải
thích, cùng cách đệm lời hứa, ở mục~\ref{muc:c4-dubao} và~\ref{muc:c4-dem}; phần còn lại ở mục~\ref{muc:c4-tuyet}.
Bảng~\ref{bang:c4-lotrinh} nối ba câu hỏi nghiên cứu ở mục~\ref{muc:c1-cauhoi} với các mục trả lời chúng.
% EV: — (đoạn dẫn, không nêu số); ba phần theo eq:so của chương 3 (X3-12); "đệm" theo GLOSSARY-VI mục B

\begin{table}[htbp]
\centering
\caption[Câu hỏi nghiên cứu và các mục trả lời]{Ba câu hỏi nghiên cứu và các mục của Chương~\ref{chuong:thucnghiem} trả lời chúng}
\label{bang:c4-lotrinh}
\small
\begin{tabular}{@{}p{6.3cm}p{2.4cm}p{6.3cm}@{}}
\toprule
\textbf{Câu hỏi} & \textbf{Mục} & \textbf{Đại lượng được đo}\\
\midrule
CH1. Phần dời do điều phối được đo thế nào, và trần trên phần đó giữ được gì, không giữ được gì?
& \ref{muc:c4-dieuphoi}, \ref{muc:c4-congdon}, \ref{muc:c4-hople}
& Số khách bị dời giờ trả vì ghép thêm khách, độ dời lớn nhất; yêu cầu bị từ chối có phương án bị trần chặn;
phép cộng các lần dời qua cả chuyến; các bộ đếm kiểm tra\\
CH2. Giữ trần đổi số lượt phục vụ, số lượt không trễ quá một và hai phút, cùng độ dài lời hứa thế nào so
với điều phối không trần, luật hạn chót và cách đệm lời hứa?
& \ref{muc:c4-ketqua}, \ref{muc:c4-chontran}, \ref{muc:c4-vaitro}, \ref{muc:c4-dem}
& Lượt phục vụ; lượt trễ và không trễ so với lời hứa đầu; quyết định bị ép; các mức trần; thứ tự tiêu chí
của bộ chọn; độ dài thêm của lời hứa\\
CH3. Lời hứa theo hồ sơ tắc đường lịch sử đổi số lượt không trễ quá một phút thế nào, và phần nào còn
chưa được giải thích?
& \ref{muc:c4-dubao}, \ref{muc:c4-tuyet}
& Lượt trễ khi có và không có dự báo; độ dài thêm của lời hứa; ba phần của độ lệch trên ngày thường và chiều
tuyết\\
\bottomrule
\end{tabular}
\end{table}
% EV: — (lộ trình, không có khẳng định; câu hỏi theo OUTLINE-V2 cuối §1.2)

% -----------------------------------------------------------------------------------------------------
\section{Thiết lập thực nghiệm}
\label{muc:c4-thietlap}

\subsection{Kịch bản thử}
\label{muc:c4-kichban}

FleetPy 1.0.2 mô phỏng xe và khách trên bộ dữ liệu chuẩn Manhattan; ở mỗi lượt quyết định, bộ chuyển đổi gửi tình
trạng hiện tại cho bộ chạy RideBound và nhận lại quyết định (mục~\ref{muc:kientruc}). Nhu cầu là các chuyến taxi New
York tuần 11--18/11/2018 trong bộ dữ liệu này~\cite{engelhardt2026}. Một kịch bản gồm một ngày, một khung giờ hai
tiếng và một trong bốn mẫu nhu cầu của ngày đó, mỗi mẫu gồm 10\% số chuyến của ngày: 108 yêu cầu trong khung giờ, phục
vụ bởi 8 xe 4 chỗ (Bảng~\ref{bang:c4-thamso}).
% EV: B01; X3-15; X3-21 (FC:395-429); B02 (FIX: 2018-11-14_sample_10_3.csv); J18 (mẫu 10% của một ngày); B08 (mẫu r1–r4); B03

Tắc đường là một hệ số chung cho cả thành phố, đổi theo khung 15 phút và lấy từ dữ liệu chuyến taxi, không phải tốc độ
từng đoạn đường; dưới đây gọi là hệ số tắc đường đo được. Mô phỏng không có nhiễu ngẫu nhiên, nên mỗi cặp
kịch bản và cấu hình chạy một lần và cho một kết quả. Ba kịch bản cùng ngày, cùng mẫu có chung 3 đến 20 trong 108 yêu
cầu, nên các kịch bản không độc lập với nhau.
% EV: B06 (sigma = shockRate = 0); K49; B11 (T38:90, 93); I05

\begin{table}[htbp]
\centering
\caption[Tham số của một kịch bản thử]{Tham số của một kịch bản thử. Các giá trị là cấu hình của thực nghiệm, không phải chuẩn ngành. Tập tắc
đường cố định (mục~\ref{muc:c4-taptu}) khác ở dòng tắc đường và chạy bằng bộ chạy cũ, có giới hạn bộ sinh và bộ
giải chưa được kiểm lại.}
\label{bang:c4-thamso}
\small
\begin{tabular}{@{}p{4.0cm}p{11.0cm}@{}}
\toprule
\textbf{Tham số} & \textbf{Giá trị}\\
\midrule
Khung giờ & 7:00--9:00, 8:00--10:00 hoặc 17:00--19:00, gọi là khung 7, 8 và 17 giờ; riêng chiều tuyết
là 14:00--16:00 và 15:00--17:00\\
Yêu cầu & 108 mỗi kịch bản, lấy trong khung giờ từ một mẫu 10\% số chuyến taxi\\
Xe & 8 xe, mỗi xe 4 chỗ\\
Đồ thị & 96 nút là các điểm đi, điểm đến lấy từ dữ liệu Manhattan; mỗi cặp nút có sẵn một thời gian đi. Đây là
đồ thị rút gọn, không phải mạng đường phố\\
Ràng buộc dịch vụ & Đón trong vòng 600 giây kể từ lúc gọi; ngồi xe không quá 1,5 lần thời gian đi
thẳng (Bảng~\ref{bang:rangbuoc})\\
Thời gian dừng đón, trả & 0 giây: xe không mất thời gian khi khách lên, xuống (kiểm trên 3 lần chạy mẫu)\\
Tắc đường & Một hệ số chung cho cả thành phố theo khung 15 phút; tốc độ gốc của mọi kịch bản lấy theo hệ số lúc
00:00\\
Bộ sinh phương án & Tối đa 100 phương án mỗi xe; mỗi lượt quyết định xét tối đa 2 yêu cầu mới\\
Bộ giải & OR-Tools 9.15, hạt giống ngẫu nhiên cố định bằng 7\\
\bottomrule
\end{tabular}
\end{table}
% EV (khung giờ, yêu cầu): B08; B20; B02; C4A-06 (chiều tuyết cũng 108 yêu cầu mỗi kịch bản, INV:14-15,17-22)
% EV (xe, đồ thị, ràng buộc dịch vụ): B03; K48 (PREF:22-26: mỗi cặp điểm đi–đến là một cạnh trực tiếp)
% EV (dừng, tắc đường): B05 (chỉ 3 lần chạy mẫu, SERV:2-4); B06; B07 (tốc độ nền lấy lúc 00:00)
% EV (bộ sinh, bộ giải): X3-17; X3-32 (giới hạn 2 yêu cầu áp cho cả lô của lượt, GEN:54-96); C4A-18 (tập tắc đường cố định: bộ chạy khác, NEW:13; giới hạn chưa kiểm)

\subsection{Các tập kịch bản và cấu hình được so}
\label{muc:c4-taptu}

Chương này dùng bốn tập kịch bản (Bảng~\ref{bang:c4-taptu}). Trên tập kiểm tra chính còn có phép so hai
thứ tự chọn phương án ở mục~\ref{muc:c4-vaitro}; phép so này dùng lại cùng kịch bản, không bổ sung ngày dữ liệu. Từ các ngày 14--18/11/2018 dựng được 49 kịch bản, chia
theo cặp ngày và mẫu để ba khung giờ của một cặp nằm cùng một phía. 34 kịch bản làm tập kiểm tra chính; 15 kịch bản
còn lại không dùng để tính các kết quả phục vụ và đúng giờ của chương này. Các lần chạy trên 15 kịch bản đó chỉ tham
gia phép kiểm tính nhất quán của sổ ở mục~\ref{muc:kiemtra}. Ngày 14--16/11 là ngày thường, 17 và 18/11 là cuối tuần.
% EV: B08 (SPLIT: chia theo nhóm ngày, mẫu); C17 (15 kịch bản hiệu chỉnh khoảng dự phòng γ); B09; B21 (11/11/2018 là Chủ nhật); A26; RUN38; T38:99-106; chương 3 v4-v3:725-727 (185 lần chạy gồm 15 lần hiệu chỉnh)

\begin{table}[htbp]
\centering
\caption[Các tập kịch bản dùng trong thực nghiệm]{Bốn tập kịch bản dùng trong chương. U: điều phối không trần; C-$B$: luật C trần $B$ giây;
M$^{+}$-$B$: luật hạn chót $B$ giây; V-30: luật trần độ dời khách thấy 30 giây; FREE: ghép tự do
(Bảng~\ref{bang:luat}). Cột cuối là số lần chạy, mỗi lần một cặp kịch bản và cấu hình.}
\label{bang:c4-taptu}
\footnotesize
\setlength{\tabcolsep}{2.5pt}
\begin{tabular}{@{}p{2.5cm}p{1.9cm}p{3.1cm}p{3.0cm}p{2.6cm}p{1.9cm}@{}}
\toprule
\textbf{Tập} & \textbf{Quy mô} & \textbf{Ngày, khung giờ} & \textbf{Cấu hình dùng ở chương này} &
\textbf{Bộ chạy} & \textbf{Lần chạy}\\
\midrule
Tập kiểm tra chính & 34 kịch bản, 3.672 yêu cầu & 14--18/11/2018; khung 7, 8, 17 giờ có 9, 14, 11
kịch bản & U, C-30, C-120, M$^{+}$-30, M$^{+}$-60, V-30 & Bản đã lưu, có luật phục hồi không tệ hơn &
170 (năm cấu hình), 34 (C-120); thêm 68 ưu tiên chi phí\\
% EV: B08; B09; B10 (170 = 34 × 5); L1a-01 (34 lần chạy C-120); A19; A28; X3-10 (R120:3-8)
Tập chỉnh trần buổi sáng & 8 kịch bản, 864 yêu cầu & 12--13/11/2018; khung 7 giờ & U; C và M$^{+}$ ở 30, 60,
90, 120, 150, 180, 240 và 300 giây & Như tập chính & 136 (trong 192 lần chạy đã lưu)\\
% EV: B15; B16; D01; D04; X3-10 (RW7:1-8; RW7X:6-8); X3-28; C4A-07 (17 cấu hình × 8 = 136; 88 + 104 = 192, RW7:33-34; RW7X:35-36)
Tập tắc đường cố định & 16 kịch bản, 1.728 yêu cầu & 12--13/11/2018; khung 8 và 17 giờ; hệ số không
đổi trong lần chạy & FREE, U, C-30, C-60, C-120 & Bản cũ, không có luật phục hồi & 16 mỗi cấu hình\\
% EV: B17; D08; D09; H13 (n = 16 mỗi cấu hình); LAD:2-6 (travelSnapshotVersion {1: 16}: một bảng giờ đi); C4A-04
Chiều tuyết 15/11/2018 & 8 kịch bản, 864 yêu cầu & 15/11/2018; khung 14 và 15 giờ & C-30 &
Như tập chính & 8\\
% EV: B20; N-26 (8 lần chạy C-30, 601 khách được phục vụ); C4A-06 (864 = 8 × 108, INV:14-15,17-22); N4B8-05 (CK8:65-72); R26 §Q8; agent-numeric-output-v2.txt (manifest, cấu hình phục hồi và DLL)
\bottomrule
\end{tabular}
\end{table}

Bản đã lưu là bản bộ chạy biên dịch và lưu lại, dùng chung cho tập kiểm tra chính, tập buổi sáng, các lần chạy trần
120 giây và chiều tuyết; bản cũ chưa có luật phục hồi (mục~\ref{muc:doichung}, \ref{muc:kiemtra}). Trừ ghép tự do, mọi
cấu hình trong từng thứ tự ưu tiên dùng cùng bộ sinh, bộ kiểm và bộ chọn với luật C; điều phối không trần chỉ
khác luật C trần 30 giây ở trần giờ trả. Các mục trước mục~\ref{muc:c4-vaitro}, cùng các phép so dự báo và đệm
sau đó, dùng thứ tự ưu tiên ổn định nếu không ghi khác. Ghép tự do không khóa gì và không qua bộ kiểm cam kết.
% EV: A28; X3-10 (RW7:1-8; RW7X:6-8; R120:3-8); N4B8-05 (CK8:65-72); R26 §Q8; B17; NEW:13 (bộ chạy cũ không có commitmentRecovery); A08; H01; A17; K50; X3-38 (HO:109, b1-rolling-cost)

Lần chạy hỏng vì lỗi hạ tầng (một lần ở tập kiểm tra chính) được thay bằng bản chạy lại hoàn tất (Phụ
lục~\ref{phuluc:C}). Bản kiểm kê ngày 08/10/2026 ghi 1.991 bản ghi lần chạy; kết quả của chương chỉ lấy từ bốn tập
trên và 68 lượt ưu tiên chi phí; các đợt chạy khác chỉ được nhắc khi nêu phạm vi kiểm tra (mục~\ref{muc:c4-congdon}).
% EV: C4A-17 (T38:12; ADD1:6: C-30-d20181114-s10-r3-w17-rerun1); ADDENDUM-1 mục 0.3; B19; C4A-20 (INV:1, bản kiểm kê ngày 08/10)

\subsection{Chỉ số}
\label{muc:c4-chiso}

Một lượt phục vụ là một yêu cầu, tức một khách, được nhận và chở tới nơi. Lượt trễ và lượt tới sớm được đếm trong nhóm
khách được phục vụ, so với lời hứa đầu (Bảng~\ref{bang:c4-chiso}). Số lượt không trễ bằng số lượt phục vụ trừ số lượt
trễ, tính trên toàn bộ yêu cầu, nên yêu cầu bị từ chối không được tính là đúng giờ. Ngưỡng chính là 60 giây; ngưỡng
120 giây được báo kèm vì có phép so đổi kết luận theo ngưỡng (mục~\ref{muc:c4-ketqua}).
% EV: C01; N-07 (khách được nhận = khách được phục vụ); C06; C08 (DIALPY:8-9: bị từ chối không tính là kịp; ở 120 s U hơn C-30 5 lượt)

Lời hứa đầu là giờ trả của kế hoạch lúc khách xác nhận đặt xe, tính với thời gian đi lúc đó. Tính từ lúc gọi xe, độ dài
trung bình của lời hứa là 14,15 phút trên 2.434 khách khi không trần và 14,12 phút trên 2.338 khách dưới luật C trần 30
giây, gộp theo khách được phục vụ. Hai giá trị trung bình gần nhau trên hai tập khách khác nhau; phép so này không nói
về độ dài lời hứa của từng khách.
% EV: A13; A05; B14 (BURD:5,11); R36-A03 (p0 trừ request.arrivalTimeMs, gộp theo khách); agent-numeric-output-v2.txt

Tỉ lệ \emph{gộp theo khách} chia tổng số khách thỏa điều kiện cho tổng số khách được phục vụ; \emph{trung bình theo
kịch bản} lấy trung bình của tỉ lệ từng kịch bản, mỗi kịch bản tính như nhau. Ví dụ dựng tay: với 5 trong 10 khách ở
một kịch bản và 5 trong 90 ở kịch bản kia, gộp theo khách cho 10\%, trung bình theo kịch bản cho 27,8\%. Chênh lệch là
cấu hình được thử trừ đối chứng; số dương nghĩa là nhiều hơn.
% EV: C27 (hai cách gộp cho hai số khác nhau); K24; ví dụ dựng tay (10/100 = 10%; (50% + 5,56%)/2 = 27,78%)

Kết quả ``có dự báo'' là phép phát lại trên các lần chạy đã ghi (mục~\ref{muc:dubao}): xe nào chở ai giữ nguyên, chỉ
các giờ báo cho khách được tính lại theo dự báo, và các chỉ số dùng $D_0$. Dự báo chỉ có ở 21 kịch bản ngày thường của
tập kiểm tra chính; 13 kịch bản cuối tuần và tập buổi sáng giữ lời hứa không dự báo.
% EV: C09; F03 (chỉ vế phát lại); X3-11 (D = max(p, S) áp cho mọi lời hứa đón và trả đã công bố); F02; N-28; D02 (12–13/11 không có dự báo)

\begin{table}[htbp]
\centering
\caption[Định nghĩa các chỉ số đánh giá]{Các chỉ số dùng trong chương. $a$ là giờ trả thực tế, $p_0$ là lời hứa đầu không dự báo, $D_0$
là lời hứa đầu theo dự báo (không có dự báo thì $D_0 = p_0$); $c$ là một lần dời giờ trả do điều phối, $z$ là phần
khách thấy của một lần cập nhật lời hứa (mục~\ref{muc:bamoc}). ``Quá'' nghĩa là lớn hơn hẳn: trễ đúng 60 giây
không tính là trễ quá 60 giây.}
\label{bang:c4-chiso}
\footnotesize
\setlength{\tabcolsep}{3pt}
\begin{tabular}{@{}p{2.7cm}p{5.3cm}p{3.5cm}p{3.5cm}@{}}
\toprule
\textbf{Chỉ số} & \textbf{Định nghĩa} & \textbf{Mẫu số, cách gộp} & \textbf{Vì sao đo}\\
\midrule
Lượt phục vụ & Yêu cầu được nhận và chở tới nơi & Số đếm; chia cho số yêu cầu thì ra tỉ lệ phục vụ & Giá của
luật tính bằng chuyến\\
Trễ quá $\theta$ & Khách được phục vụ có $a - D_0 > \theta$; $\theta$ = 30, 60, 120, 300 hoặc 600 giây & Đếm trong
khách được phục vụ & Lời hứa bị vượt ở phía muộn\\
Không trễ quá $\theta$ & Lượt phục vụ trừ lượt trễ quá $\theta$ & Chia cho số yêu cầu; yêu cầu bị từ chối không
tính là đúng giờ & Phục vụ và đúng giờ trên một mẫu số\\
Tới sớm quá 60 giây & Khách được phục vụ có $a - D_0 < -60$ giây & Đếm trong khách được phục vụ & Lời
hứa quá dài cũng sai\\
Lời hứa dài thêm & $D_0 - p_0$ & Trung bình gộp theo khách được phục vụ, kể cả khách không có dự báo & Giá của
dự báo hay đệm\\
Khách đi chung & Có khách khác trên cùng xe trong khoảng từ lúc được đón tới lúc được trả & Gộp theo
khách được phục vụ & Mức ghép chuyến còn lại\\
Khách bị điều phối dời giờ trả & Có ít nhất một lần $c \neq 0$; độ dời của khách là $\sum|c|$ &
Gộp theo khách được phục vụ & Đại lượng luật C giới hạn\\
Quyết định bị ép & Lượt quyết định có ít nhất một xe bị ép: luật lời hứa đã loại cả phương án giữ nguyên kế hoạch
của xe, và xe đi theo phương án được giữ nhờ luật phục hồi, có ít nhất một khách được miễn trừ
(mục~\ref{muc:doichung}) & Chia cho số lượt quyết định đã được bộ mô phỏng xác nhận áp dụng; gộp theo
lượt quyết định, mỗi lượt đếm một lần & Luật có tự giữ được nghĩa vụ không\\
Khách bị phá lời hứa & Khách mà luật đang dùng không giữ được điều nó hứa. Luật C: $\sum|c| > B$; luật hạn chót:
$a - p_0 > B$; luật trần độ dời khách thấy: $\sum|z| > B$ & Chia cho số khách được phục vụ; gộp theo
khách & Luật có giữ được điều nó hứa không\\
\bottomrule
\end{tabular}
\end{table}
% EV (phục vụ): C01; L1a-01 (2.387 hoàn tất + 1.285 bị từ chối = 3.672); N-07
% EV (trễ, không trễ, sớm): C06; C08 (DIALPY:8-9); N-04 (trễ > 60 s và sớm < −60 s rời nhau); C4A-01 (ngưỡng 30 s, cùng quy tắc sweep_analyze.py:107,128-132); C09; C10
% EV (dài thêm, đi chung, bị dời): F06; N-06 (gộp theo khách, ext_sum / served, ADD1:14,19,24); N-09 (định nghĩa đi chung); C20; N-07; A06; X3-05
% EV (bị ép): C4A-10 (RS:1012-1027; VALID:141-145; HVA:252-300; POL:443-463); X3-30; X3-09; C13 (trung bình theo kịch bản)
% EV (phá lời hứa): N-21 (định nghĩa của M⁺-30 và C-30; số gộp); A11; GLOSSARY-VI mục B "phá lời hứa"

\subsection{Mức trần 30 giây và tính chất thăm dò của kết quả}
\label{muc:c4-thamdo}

Trần chính $B$ = 30 giây là trần giờ trả của cấu hình có trần ở tập kịch bản tĩnh của đợt thử trước, nơi thời gian đi
không đổi (Phụ lục~\ref{phuluc:A}). Mức này được giữ nguyên cho tập kiểm tra chính, không chỉnh theo kết quả của tập
này. Nó không phải ngưỡng cảm nhận của khách hay chuẩn ngành, nên mỗi mức trần chỉ là một điểm trên đường đánh đổi
(mục~\ref{muc:c4-chontran}).
% EV: C4A-15 (TEX08:136-137: C1 = trần 30 giây cho tổng độ dời giờ trả; thay X3-18 đã bỏ); K29 (tập thời gian đi tĩnh); K33; B10 (năm cấu hình cố định trước lần chạy đầu, T38:3-6); D20; K52; K36

Các kết quả là thăm dò và dùng lại dữ liệu đã xem. Ngưỡng phân tích và cách chia tập được chọn sau khi đã xem tập buổi
sáng và các đợt chạy thử trước đó. Mức thử 120 giây được chọn sau khi đã xem kết quả của điều phối không trần và luật
C trần 30 giây trên chính tập kiểm tra chính, kể cả các phép tính thử ở 120 giây (mục~\ref{muc:c4-dieuphoi},
\ref{muc:c4-congdon}). Các con số vì vậy mô tả điều đã đo, không phải phép kiểm trên dữ liệu chưa từng xem.
% EV: I06; T38:8 (chọn sau khi xem Việc A, B, w07x); C4A-11 (BUD 07/10 = D14; CVP 10/10 = E06; C120R:10 chạy 10/10; T38:5 tập chính chạy 1–2/10); K42

% -----------------------------------------------------------------------------------------------------
\section{Kết quả chính khi chưa có dự báo}
\label{muc:c4-ketqua}

Với bộ chọn ưu tiên ổn định trên 3.672 yêu cầu của tập kiểm tra chính, luật C trần 30 giây không để phần dời giờ trả do điều phối của khách
nào vượt 30 giây: lớn nhất 29,2 giây, so với 590,1 giây (gần 10 phút) khi không trần. Không lượt nào của luật C trễ
quá 5 phút so với lời hứa đầu, so với 18 lượt khi không trần. Đổi lại, luật C phục vụ 2.338 lượt, ít hơn điều phối
không trần 96 lượt, và có ít hơn 242 khách đi chung. Số lượt không trễ quá một phút tăng 31, từ 2.231 lên 2.262, nhưng
số lượt không trễ quá hai phút giảm 5, từ 2.332 xuống 2.327.
% EV: C20; N-23; N-24; C07; C01; C02; N-20; C08

Mục~\ref{muc:c4-dem} cho thấy khi có dự báo, cộng thêm 30 giây vào mọi lời hứa của điều phối không trần cũng cho 2.328
lượt không trễ quá một phút, bằng luật C có dự báo. Vì vậy khóa luận không dùng chỉ số này làm lý do chọn luật C.
% EV: C10 (U + dự báo + 30 s: 2.328 = C-30 + dự báo: 2.328); K06; K51

Luật hạn chót chỉ mượn phần hạn chót của Schulz và Pfeiffer~\cite{schulz2026}: giờ trả mới không được muộn hơn lời hứa
đầu quá 30 hoặc 60 giây, thay vì 900 giây của họ. Nó chạy cùng bộ điều phối và luật phục hồi với luật C, nên đây không
phải phép so với hệ thống của hai tác giả.
% EV: J10; J11; A09; A19; H01; X3-28; X3-41

\begin{table}[htbp]
\centering
\caption[Kết quả chính khi chưa có dự báo]{Kết quả trên tập kiểm tra chính với lời hứa đầu không dự báo (34 kịch bản, 3.672 yêu cầu mỗi
cấu hình, với bộ chọn ưu tiên ổn định). Số đếm gộp trên mọi kịch bản; trễ tính so với lời hứa đầu. Hai cột ``không trễ quá'' là số lượt phục vụ
không trễ quá ngưỡng; chia cho 3.672 thì ra tỉ lệ trên yêu cầu. ``Chưa đếm'': chưa có số đếm chính xác.}
\label{bang:c4-chinh}
\small
\setlength{\tabcolsep}{4pt}
\begin{tabular}{@{}p{\dimexpr .12\linewidth-2\tabcolsep\relax}
>{\raggedleft\arraybackslash}p{\dimexpr .13\linewidth-2\tabcolsep\relax}
>{\raggedleft\arraybackslash}p{\dimexpr .16\linewidth-2\tabcolsep\relax}
>{\raggedleft\arraybackslash}p{\dimexpr .22\linewidth-2\tabcolsep\relax}
>{\raggedleft\arraybackslash}p{\dimexpr .23\linewidth-2\tabcolsep\relax}
>{\raggedleft\arraybackslash}p{\dimexpr .14\linewidth-2\tabcolsep\relax}@{}}
\toprule
\textbf{Cấu hình} & \textbf{Phục vụ} & \textbf{Trễ quá 60 s} & \textbf{Không trễ quá 60 s} &
\textbf{Không trễ quá 120 s} & \textbf{Khách đi chung}\\
\midrule
U & 2.434 & 203 & 2.231 & 2.332 & 555\\
C-30 & 2.338 & 76 & 2.262 & 2.327 & 313\\
M$^{+}$-30 & 2.338 & 74 & 2.264 & 2.327 & 307\\
% EV (U, C-30, M⁺-30): C01; C06; C08; N-09 (555 / 313 / 307, số đếm chính xác)
M$^{+}$-60 & 2.347 & 79 & 2.268 & 2.336 & chưa đếm\\
V-30 & 2.336 & 74 & 2.262 & 2.325 & chưa đếm\\
% EV (M⁺-60, V-30): C01; C4A-02 (CK4A check: trễ quá 60 / 120 s 79 / 11 và 74 / 11; EXP:5-6 trễ quá 60 s 79 và 74; đi chung chỉ có % làm tròn ở EXP)
\bottomrule
\end{tabular}
\end{table}

Cách đọc bảng: dưới trần 30 giây, 2.338 yêu cầu được phục vụ; 76 lượt trong số đó tới muộn hơn lời hứa đầu quá 60
giây, nên 2.338 $-$ 76 = 2.262 lượt không trễ quá 60 giây.
% EV: C01; C06; C08

Cái giá trước hết là số lượt phục vụ: bớt 96 lượt nghĩa là tỉ lệ phục vụ giảm 2,6 điểm phần trăm (96 trên 3.672 yêu
cầu), tức ít hơn 3,9\% số lượt của điều phối không trần (96 trên 2.434). Mức giảm nằm ở khung 8 giờ (1.014 xuống 975)
và 17 giờ (904 xuống 852); khung 7 giờ chỉ giảm 5. Số khách đi chung giảm từ 555 xuống 313, tức từ 22,8\% xuống 13,4\%
số khách được phục vụ (gộp theo khách). Khóa luận không có bằng chứng về giá trị của các chuyến đi chung này đối với
khách hay hãng xe.
% EV: C02; C03; N-09; N-20 (lưu ý: D20, I07, I10)

Số lượt không trễ quá một phút vẫn tăng vì phục vụ bớt 96 lượt nhưng trễ bớt 127 lượt: 127 $-$ 96 = 31. Hai cấu hình
cũng phục vụ những khách khác nhau: 2.207 khách chung, 227 khách chỉ được điều phối không trần và 131 khách chỉ được
luật C phục vụ. Vì vậy 31 là chênh lệch hai tổng, không phải 31 khách cụ thể được đưa tới đúng giờ.
% EV: C06; C08; C04 (câu về chờ, ngồi xe của khách chung, C05, cắt vì ngân sách trang; xem NOTES)

So với luật hạn chót 30 giây, luật C phục vụ đúng bằng số lượt (2.338) nhưng không hoàn toàn cùng những khách (2.332
khách chung, mỗi bên 6 khách riêng), và có 76 so với 74 lượt trễ quá một phút. Luật hạn chót 60 giây và luật trần độ
dời khách thấy cũng gần như vậy (Bảng~\ref{bang:c4-chinh}). Các luật lời hứa khác nhau ở chỗ mỗi luật có giữ được điều
chính nó hứa hay không (Bảng~\ref{bang:c4-nghiavu}).
% EV: C01; C04; C06; K05; K07; C32; C4A-02

\begin{table}[htbp]
\centering
\caption[Nghĩa vụ của chính từng luật lời hứa]{Nghĩa vụ của chính từng luật trên tập kiểm tra chính (34 kịch bản). Hai cột đầu ghi số đếm gộp: số quyết
định bị ép trên số quyết định được xác nhận, và số khách bị phá lời hứa trên số khách được phục vụ. Cột cuối ghi số
kịch bản có ít nhất một khách bị phá lời hứa trên 34 kịch bản. Định nghĩa ở Bảng~\ref{bang:c4-chiso}. Ở hàng luật C,
các số 0 là số đo trên tập đã chạy, phù hợp với tính chất có điều kiện ở Nhận xét~\ref{nx:giunguyen}.}
\label{bang:c4-nghiavu}
\small
\setlength{\tabcolsep}{4pt}
\begin{tabular}{@{}p{\dimexpr .12\linewidth-2\tabcolsep\relax}
>{\raggedleft\arraybackslash}p{\dimexpr .28\linewidth-2\tabcolsep\relax}
>{\raggedleft\arraybackslash}p{\dimexpr .31\linewidth-2\tabcolsep\relax}
>{\raggedleft\arraybackslash}p{\dimexpr .29\linewidth-2\tabcolsep\relax}@{}}
\toprule
\textbf{Cấu hình} & \textbf{Quyết định bị ép / được xác nhận} &
\textbf{Khách bị phá lời hứa / được phục vụ} &
\textbf{Kịch bản có khách bị phá / tổng kịch bản}\\
\midrule
C-30 & 0 / 20.101 & 0 / 2.338 & 0 / 34\\
M$^{+}$-30 & 7.710 / 20.103 & 289 / 2.338 & 28 / 34\\
M$^{+}$-60 & 4.031 / 20.163 & 79 / 2.347 & 22 / 34\\
V-30 & 10.442 / 20.090 & 456 / 2.336 & 34 / 34\\
\bottomrule
\end{tabular}
\end{table}
% EV: R36-A01 (sổ bổ sung trong CHANGES-v3.md; 170 tệp IVC và đối chiếu PAIRS); N-21; X3-30; A22

Cách đọc bảng: luật hạn chót 30 giây có 7.710 trong 20.103 lượt quyết định bị ép và 289 trong 2.338 khách tới muộn
hơn lời hứa đầu quá 30 giây; những khách này nằm ở 28 trong 34 kịch bản.
% EV: N-21; R36-A01

Trên cả 34 kịch bản đã chạy, luật C không có quyết định bị ép và không khách nào vượt trần. Theo Nhận
xét~\ref{nx:giunguyen}, nếu giả thiết (N) đúng ở lượt đang xét và mọi khách có lời hứa của xe có $0 \le L_r \le B$,
phương án giữ nguyên không thêm phần dời do điều phối nên vẫn qua kiểm tra trần. Khi (N) đúng ở mọi lượt, trần được
giữ từ lúc xác nhận. Các số đo ở đây phù hợp với tính chất đó; nhận xét không bảo đảm giờ trả thực tế. Điều phối không
trần cũng có 0 trong 20.760 quyết định bị ép. Luật hạn chót so giờ trả mới với lời hứa đầu, nên phần thay đổi khi giữ
nguyên kế hoạch vẫn có thể làm xe bị ép (mục~\ref{muc:doichung}).
% EV: A21; A22; N-21; X3-16 (U 0 / 20.760); X3-30; A10; K05; Ch3 v4-v3:371-390; R36-A01

Tính gộp, luật hạn chót 30 giây bị ép ở 7.710 trong 20.103 lượt quyết định (38,4\%) và phá lời hứa với 289 trong 2.338
khách (12,4\%). Gánh nặng dồn vào khung 7 giờ: cả 9 kịch bản có khách bị phá lời hứa, với 181 trong 509 khách được phục
vụ và 3.447 trong 4.564 lượt quyết định bị ép. Ở khung 17 giờ, 5 trong 11 kịch bản có khách bị phá, với 8 trong 852
khách và 328 trong 7.163 lượt quyết định bị ép. Các số này được gộp theo từng khung giờ. Ở mức 60 giây, các tỉ lệ toàn
tập giảm nhưng không về 0. Luật trần độ dời khách thấy bị ép nhiều nhất, phù hợp với định nghĩa: tổng độ lớn các lần
cập nhật mà nó kiểm không nhỏ hơn độ lùi so với lời hứa đầu (công thức~\eqref{eq:dongnhat}).
% EV: N-21 (gộp 7.710/20.103 và 289/2.338); R36-A02 (cặp đếm theo khung giờ, sổ bổ sung trong CHANGES-v3.md); C14; C13; R36-A01; A11; A12; X3-33 (eq:dongnhat: Σ|z| ≥ |Σz| = |p_n − p_0|)

Bảng~\ref{bang:c4-nghiavu} không so hai chuẩn đúng giờ: luật C hứa về phần dời do điều phối, còn luật hạn chót hứa về
toàn bộ độ trễ so với lời hứa đầu, một lời hứa lớn hơn. Đo bằng chính thước của luật hạn chót, tức tới muộn hơn lời
hứa đầu quá 30 giây, luật C có 292 trong 2.338 khách (12,5\%), so với 289 khách (12,4\%) của luật hạn chót. Cả hai
luật đều có khách tới muộn hơn lời hứa đầu quá 30 giây ở 28 trong 34 kịch bản. Vậy về độ đúng giờ, hai luật gần như
ngang nhau; luật C khác ở chỗ chỉ hứa phần mà bộ điều phối kiểm soát được, nên giữ được lời hứa đó.
% EV: C13; K05; C4A-01 (C-30 292/2.338 = 12,49%, 28/34); N-21; R36-A01

Bảng~\ref{bang:c4-phanbotre} cho thấy số lượt trễ giảm ở cả bốn ngưỡng đã đo, từ 1 tới 10 phút.
% EV: N-24

\begin{table}[htbp]
\centering
\caption[Phân bố trễ theo ngưỡng khi chưa có dự báo]{Phân bố trễ so với lời hứa đầu không dự báo trên tập kiểm tra chính: số lượt trễ quá mỗi ngưỡng
(trong số khách được phục vụ) và số lượt đó tính trên 100 yêu cầu (3.672 yêu cầu mỗi cấu hình).}
\label{bang:c4-phanbotre}
\small
\begin{tabular}{@{}lrrrr@{}}
\toprule
& \multicolumn{2}{c}{\textbf{U}} & \multicolumn{2}{c}{\textbf{C-30}}\\
\cmidrule(lr){2-3}\cmidrule(lr){4-5}
\textbf{Trễ quá} & \textbf{Lượt} & \textbf{Trên 100 yêu cầu} & \textbf{Lượt} & \textbf{Trên 100 yêu cầu}\\
\midrule
1 phút & 203 & 5,53 & 76 & 2,07\\
2 phút & 102 & 2,78 & 11 & 0,30\\
5 phút & 18 & 0,49 & 0 & 0\\
10 phút & 2 & 0,05 & 0 & 0\\
\midrule
Trễ lớn nhất (s) & \multicolumn{2}{c}{660,1} & \multicolumn{2}{c}{232,3}\\
\bottomrule
\end{tabular}
\end{table}
% EV: N-24 (203/102/18/2 và 76/11/0/0; 5,53/2,78/0,49/0,05 và 2,07/0,30/0/0); C07; C28; C4A-14 (CON:4,11: +660,1 s và +232,3 s; khớp C26 "11 phút" và "3,87 phút")

Cách đọc bảng: khi không trần, 18 lượt tới muộn hơn lời hứa đầu quá 5 phút, tức 0,49 lượt trên 100 yêu cầu.
% EV: N-24; C28

Khi không trần, lượt trễ nhất muộn 660,1 giây (khoảng 11 phút); dưới trần 30 giây là 232,3 giây (3,87 phút). Trong 76
khách còn trễ quá một phút dưới luật C, 72 người không có lần dời do điều phối nào được ghi và 61 người thuộc khung 7
giờ. Lời hứa theo dự báo (mục~\ref{muc:c4-dubao}) nhắm tới phần đoán được của loại trễ này.
% EV: N-24; C4A-14; C26; C21 (4/76 từng bị dời); K39 (không viết "trễ thuần vì kẹt xe"); N-11 (trễ quá 60 s của C-30 theo khung 61 / 15 / 0)

Tách theo kịch bản, luật C không có nhiều lượt trễ quá một phút hơn điều phối không trần ở cả 34 kịch bản và có số
lượt không trễ quá một phút bằng hoặc nhiều hơn ở 29 kịch bản. Ở khung 8 và 17 giờ, điều phối không trần nhiều lượt
trễ hơn ở cả 25 kịch bản; ở khung 7 giờ, luật C ít hơn ở 3 kịch bản và bằng ở 6. Vì 34 kịch bản thuộc 5 ngày và không
độc lập, đây là phép đếm mô tả, không phải kiểm định thống kê.
% EV: N-22 (34/34 trễ ≤ U; 29/34 không trễ ≥ U; bảng theo ngày × khung ở ADD1:173-186); C30; I05; B11

% -----------------------------------------------------------------------------------------------------
\section{Phần do điều phối: luật C chặn gì và mất gì}
\label{muc:c4-dieuphoi}

Dưới luật C trần 30 giây, 33 trong 2.338 khách được phục vụ bị điều phối dời giờ trả, mỗi khách nhiều nhất 29,2 giây;
khi không trần là 174 trong 2.434 khách, một khách bị dời tổng cộng 590,1 giây (Bảng~\ref{bang:c4-doidieuphoi}). Cái
giá là số yêu cầu bị từ chối tăng từ 1.238 lên 1.334; chỉ 148 trong số đó có phương án bị trần chặn được ghi lại.
% EV: C20; N-23; C01 (1.238 = 3.672 − 2.434; CON:3); D13 (1.334; 148)

\begin{table}[htbp]
\centering
\caption[Phần dời giờ trả do điều phối trên tập kiểm tra chính]{Phần dời giờ trả do điều phối trên tập kiểm tra chính, lời hứa đầu không dự báo. Mẫu số của
dòng đầu là số khách được phục vụ (2.434, 2.338 và 2.387); độ dời của một khách là tổng $|c|$ của khách
đó. Dòng cuối lấy tổng phần dời có dấu do điều phối chia cho tổng số giây trễ so với lời hứa đầu của chính các khách
trễ quá 60 s: 23.027 trên 31.353 giây, 82 trên 6.892 giây và 4.836 trên 12.450 giây. Cột trần 120 giây được dùng
lại ở mục~\ref{muc:c4-chontran}.}
\label{bang:c4-doidieuphoi}
\small
\begin{tabular}{@{}p{7.6cm}rrr@{}}
\toprule
\textbf{Chỉ số} & \textbf{U} & \textbf{C-30} & \textbf{C-120}\\
\midrule
Khách bị điều phối dời giờ trả & 174 (7,15\%) & 33 (1,41\%) & 103 (4,32\%)\\
Khách bị dời từ hai lần trở lên & 12 & 0 & 0\\
Tổng dời lớn nhất của một khách (s) & 590,1 & 29,2 & 118,9\\
Trung vị độ lớn mỗi lần dời (s) & 108,6 & 8,3 & 62,5\\
% EV: C20; C22 (trung vị theo lần: 108,6 s trên 187 lần, 8,3 s trên 33 lần); N-07 (C-120: 103, 0 khách ≥ 2 lần, nên trung vị theo lần = trung vị theo khách 62,54 s); N-23
Khách trễ quá 60 s từng bị dời & 134 trên 203 & 4 trên 76 & 60 trên 132\\
Phần tính cho điều phối trong tổng giây trễ của các khách trễ quá 60 s (\%) & 73,4 & 1,2 & 38,8\\
% EV: C21; C24; N-23 (cột C-120: 60/132, 38,8%); C4A-03 (tử số, mẫu số: 23.027,239 / 31.353,407; 81,742 / 6.891,526; 4.836,281 / 12.450,223 s)
\bottomrule
\end{tabular}
\end{table}

Cách đọc bảng: khi không trần, 174 khách bị điều phối dời giờ trả ít nhất một lần, 12 người trong số đó bị dời từ hai
lần, và tổng dời lớn nhất của một khách là 590,1 giây.
% EV: C20

Dưới trần 30 giây, mỗi khách bị dời chỉ một lần, trung vị 8,3 giây, so với 108,6 giây trên 187 lần dời khi không trần.
Trong 33 khách bị dời, 29 người bị dời muộn đi; 4 người bị kéo sớm hơn, mỗi người đúng 1 mili giây. Khi không trần,
134 trong 203 khách trễ quá một phút từng bị dời; con số chỉ cho biết hai nhóm trùng nhau, không cho biết việc bị dời
là nguyên nhân của độ trễ. Theo sổ theo từng khách (công thức~\eqref{eq:so}), phần tính cho điều phối bằng 73,4\% tổng
số giây trễ của 203 khách đó, còn dưới trần 30 giây chỉ 1,2\%: gần như mọi giây trễ còn lại nằm ở hai phần kia của sổ.
% EV: E02; C22; C23 (29 muộn, 4 sớm); N-05 (mỗi lần dời sớm đúng −1 ms; chỉ kiểm ở U, C-30, C-120, M⁺-30; nguyên nhân chưa kiểm, ADDENDUM-1 mục 0.2; C4A-19); C21 (giao hai tập); C24 (sổ mô hình, không nhân quả); K04; X3-12

Giới hạn này không dựa vào dự báo, vì dự báo không đi vào bộ điều phối. Ở 13 kịch bản cuối tuần không có dự báo, tổng
dời lớn nhất của một khách là 29,2 giây ở trần 30 giây và 116,5 giây ở trần 120 giây. Trên 8 kịch bản chiều tuyết,
8 trong 601 khách được phục vụ bị dời, lớn nhất 25,2 giây; không ai vượt trần và không quyết định nào bị ép trong
5.067 lượt.
% EV: F03; X3-11 (D chỉ tính sau khi chạy); N-28; N-26 (thay anchor F13 của RD §8 mục 4, ADDENDUM-1 mục 0.4); N4B8-05 (CK8:65-72); R26 §Q8

Một ví dụ trên tập buổi sáng cho thấy luật C và luật hạn chót khác nhau ở đâu. Trước một quyết định, giờ trả đang hứa
của một khách đã lùi 79,0 giây so với lời hứa đầu, toàn bộ là phần thay đổi khi giữ nguyên kế hoạch. Phương án ghép
thêm một khách mới sẽ dời giờ trả dự kiến của khách đó thêm 68,3 giây do điều phối. Ở mức 90 giây, luật hạn chót cộng
hai phần, khoảng 147 giây, nên loại phương án; luật C chỉ tính 68,3 giây nên cho qua. Trên tập này hai luật tách nhau
ở đúng ba chỗ cùng kiểu, ở mức 90 hoặc 120 giây; mỗi chỗ luật C nhận thêm một khách. Đây là phép so hai luật trên
giờ trả dự kiến tại cùng một trạng thái của mô hình, không phải phép đo nhân quả.
% EV: D05 (79,0 + 68,3 s; 60,9 + 61,2 s); C4A-12 (W07XM:6,11,16: "sum x 79.0, sum c 0.0, p_last − p_first 79.0", c = 68,3 s, x = 0,1 s; W07X:54-62: "khách đang có", giai đoạn chờ đón hay trên xe không ghi); K04

Khi luật C loại một phương án vì nó làm một khách đã có lời hứa vượt trần, bộ kiểm ghi một bản ghi chặn gồm khách đó,
phần trần đã dùng, lần dời mới và trần (Bảng~\ref{bang:kiemtran}). Dưới trần 30 giây, 148 trong 1.334 yêu cầu bị từ
chối có ít nhất một phương án mang bản ghi chặn; điều phối không trần phục vụ 116 trong số đó. 1.186 yêu cầu còn lại
không có bản ghi chặn nào được ghi, và điều phối không trần cũng từ chối 1.075 yêu cầu trong số này; điều đó không
chứng minh chúng bị từ chối chỉ vì ràng buộc dịch vụ. Bản ghi chặn cũng không phải khách: 845 bản ghi nằm ở 643 yêu
cầu, và 495 yêu cầu trong số đó cuối cùng vẫn được phục vụ.
Các số đếm theo yêu cầu trong đoạn này tính theo từng cặp kịch bản và yêu cầu.
% EV: X3-05 (nội dung bản ghi chặn); D13 (BUD:2-5: 148 = 116 + 32; 1.186 = 111 + 1.075); K10; C4A-09 (BUDPY:29-44: nhóm 1.186 là phần dư không có bản ghi chặn được ghi, không phải lý do cấu trúc đã chứng minh); D16 (SF:2: 643 cặp, 495 phục vụ, 148 từ chối); E01 (845)

Với mỗi yêu cầu trong 148 yêu cầu, phép phân tích chọn bản ghi có mức vượt nhỏ nhất, tức phần vượt quá 30 giây; trung
vị là 109,6 giây, lớn nhất 524,8 giây. Chỉ ở 2 trong 148 bản ghi được chọn, khách nêu trong bản ghi đã dùng một phần
trần trước đó; ở cả 148, riêng lần dời mới đã vượt 30 giây. Với trần 120 giây, mức vượt của 59 trong 148 yêu cầu nằm
trong phần nới thêm 90 giây. Không thể kết luận 59 yêu cầu đó sẽ được phục vụ: bộ kiểm dừng ở khách vi phạm gặp trước,
và khi trần đổi thì các quyết định trước cũng đổi. Mục~\ref{muc:c4-chontran} cho kết quả chạy thật.
% EV: D14 (BUD:6-8: trung vị 109,6 s; lớn nhất 524,8 s; 59/148 có mức vượt ≤ 90 s); D19 (BUD:10; FJC:4-6); C4A-09 (BUDPY:22,35-42,51-52: mức vượt = after − limit; 2/148 đếm theo bản ghi được chọn của từng yêu cầu, không phải số khách khác nhau); K11; A23; I08

% -----------------------------------------------------------------------------------------------------
\section{Chọn mức trần}
\label{muc:c4-chontran}

Trong ba điểm đã đo trên tập kiểm tra chính, không điểm nào tốt nhất ở mọi chỉ số. Trần 30 giây có ít lượt trễ quá một
phút nhất (76). Trần 120 giây có nhiều lượt không trễ quá hai phút nhất (2.366, so với 2.327 của trần 30 giây và 2.332
khi không trần) và phục vụ 2.387 lượt, nằm giữa hai điểm kia. Chọn mức nào tùy ngưỡng mà hãng xe coi trọng, điều thực
nghiệm không trả lời được.
% EV: N-06; L1a-02; C01; C06; C08; K12; K36

Mức trần được đo trên ba tập không cộng gộp (Hình~\ref{hinh:c4-tran}). Tập tắc đường cố định cho thấy trần làm gì khi
thời gian đi không đổi, tập buổi sáng dò tám mức trên cùng 8 kịch bản, tập kiểm tra chính thêm mức 120 giây.
% EV: B17; B15; B16; X3-28; D01; L1a-01; I12

\textbf{Tập tắc đường cố định.} Bảng~\ref{bang:c4-codinh} xếp các cấu hình từ lỏng nhất (ghép tự do) tới chặt nhất
(trần 30 giây), với hệ số tắc đường cố định suốt lần chạy.
% EV: B17; LAD:2-6

\begin{table}[htbp]
\centering
\caption[Các mức trần trên tập tắc đường cố định]{Các mức trần trên tập tắc đường cố định (16 kịch bản, 1.728 yêu cầu mỗi cấu hình; bộ chạy cũ,
không có luật phục hồi). Ba cột số đầu là số đếm; hai cột cuối là số khách được phục vụ thỏa điều kiện, trong ngoặc
là tỉ lệ gộp theo khách được phục vụ (\%).}
\label{bang:c4-codinh}
\small
\begin{tabular}{@{}lrrrrr@{}}
\toprule
\textbf{Cấu hình} & \textbf{Phục vụ} & \textbf{Trễ quá 60 s} & \textbf{Không trễ quá 120 s} &
\textbf{Đi chung} & \textbf{Bị dời giờ trả}\\
\midrule
FREE & 1.629 & 274 & 1.465 & 865 (53,1) & 388 (23,8)\\
U & 1.584 & 47 & 1.557 & 346 (21,8) & 62 (3,9)\\
C-120 & 1.563 & 23 & 1.563 & 303 (19,4) & 41 (2,6)\\
C-60 & 1.540 & 0 & 1.540 & 266 (17,3) & 20 (1,3)\\
C-30 & 1.534 & 0 & 1.534 & 241 (15,7) & 6 (0,4)\\
\bottomrule
\end{tabular}
\end{table}
% EV: D09 (số đếm, WP17); D08 (tỉ lệ đi chung, bị dời; gộp theo khách); C4A-04 (đi chung 865/346/303/266/241, CK4A fixed-shared; bị dời 388/62/41/20/6, WBU:4,7,10,13,16); B17

Cách đọc bảng: với trần 120 giây, 1.563 trong 1.728 yêu cầu được phục vụ, 23 lượt trễ quá 60 giây; 303 khách (19,4\%)
đi chung và 41 khách (2,6\%) bị dời giờ trả.
% EV: D08; D09; C4A-04

Ghép tự do khác điều phối không trần ở cả cách chọn phương án lẫn việc kiểm cam kết, nên chênh lệch giữa hai cấu hình
này không chỉ do trần. Trong nhóm cùng chính sách, siết từ không trần xuống trần 30 giây làm số lượt phục vụ giảm từ
1.584 xuống 1.534 và số lượt trễ quá 60 giây từ 47 xuống 0. Giữa hai mức đã chạy 60 và 120 giây, số lượt trễ tăng từ 0
lên 23. Các con số của bộ chạy cũ này không chuyển sang tập kiểm tra chính.
% EV: A17; K50; X3-38; NEW:13; D09; D11; B17; I12

\textbf{Tập chỉnh trần buổi sáng.} Trên 864 yêu cầu, số lượt phục vụ chỉ dao động trong 6 lượt, từ 660 ở trần 30 giây
tới 666 ở trần 120--180 giây (Bảng~\ref{bang:c4-sang}).
% EV: D01; B15; B16

\begin{table}[htbp]
\centering
\caption[Các mức trần trên tập chỉnh trần buổi sáng]{Tám mức trần trên tập chỉnh trần buổi sáng (8 kịch bản khung 7 giờ ngày 12--13/11/2018, 864 yêu cầu
mỗi cấu hình; lời hứa đầu không dự báo). Số đếm gộp; hai cột ``không trễ quá'' là số lượt trên 864 yêu cầu. Các mức
có cùng mọi số đếm được gộp vào một dòng.}
\label{bang:c4-sang}
\small
\begin{tabular}{@{}lrrrr@{}}
\toprule
\textbf{Cấu hình} & \textbf{Phục vụ} & \textbf{Trễ quá 60 s} & \textbf{Không trễ quá 60 s} &
\textbf{Không trễ quá 120 s}\\
\midrule
U & 665 & 50 & 615 & 656\\
C-30 & 660 & 47 & 613 & 658\\
C-60 & 661 & 47 & 614 & 659\\
C-90 & 664 & 50 & 614 & 660\\
C-120, C-150, C-180 & 666 & 51 & 615 & 659\\
C-240, C-300 & 665 & 50 & 615 & 656\\
\bottomrule
\end{tabular}
\end{table}
% EV: D01; D02 (U, C-30, C-90, C-120: CAP:3-7; DIAL17:3-7, hai bộ đọc độc lập); C4A-05 (C-60, C-150, C-180, C-240, C-300: CAP:5,8-11; DIAL17:5,8-11); D02 lưu ý (cột dự báo bằng ngây thơ ở 12–13/11)

Cách đọc bảng: với trần 120, 150 hoặc 180 giây, 666 yêu cầu được phục vụ, 51 lượt trễ quá 60 giây, 615 lượt không trễ
quá 60 giây và 659 lượt không trễ quá 120 giây.
% EV: D01; D02; C4A-05

Trần 120 giây phục vụ nhiều hơn điều phối không trần một lượt, nên điều phối không trần không phải mức phục vụ cao
nhất có thể đạt, chỉ là một điểm đo. Từ 30 lên 120 giây, trần thêm 6 lượt phục vụ, 4 lượt trễ quá 60 giây và 2 lượt
không trễ quá 60 giây. Luật C không có quyết định bị ép trong 64 lần chạy của tám mức trần này. Ở cả tám mức, hai
luật phục vụ chênh nhau không quá 2 lượt.
% EV: D01; K12; D03 (+6 / +4 / +2); D04; X3-28 (tám mức 30–300 s); E01 (64 lần C, 0 bị ép); agent-numeric-output-v2.txt (0 / 44.270 lượt quyết định)

\textbf{Tập kiểm tra chính.} Bảng~\ref{bang:c4-tran120} đặt trần 120 giây cạnh hai cấu hình của
mục~\ref{muc:c4-ketqua}, khi không có và khi có dự báo.

\begin{table}[htbp]
\centering
\caption[Trần 120 giây trên tập kiểm tra chính]{Trần 120 giây trên tập kiểm tra chính (34 kịch bản, 3.672 yêu cầu mỗi cấu hình). Số đếm gộp; các dòng
``không trễ quá'' là số lượt trên 3.672 yêu cầu. Khối dưới là phép phát lại lời hứa theo dự báo trên cùng các lần
chạy, bàn ở mục~\ref{muc:c4-dubao}; dự báo chỉ có ở 21 kịch bản ngày thường, và độ dài thêm của lời hứa là trung
bình gộp trên mọi khách được phục vụ, kể cả khách cuối tuần có độ dài thêm bằng 0.}
\label{bang:c4-tran120}
\small
\begin{tabular}{@{}p{8.4cm}rrr@{}}
\toprule
\textbf{Chỉ số} & \textbf{U} & \textbf{C-30} & \textbf{C-120}\\
\midrule
\multicolumn{4}{@{}l}{\textit{Lời hứa đầu không dự báo}}\\
Lượt phục vụ & 2.434 & 2.338 & 2.387\\
Lượt trễ quá 60 s & 203 & 76 & 132\\
Lượt trễ quá 120 s & 102 & 11 & 21\\
Lượt trễ quá 300 s & 18 & 0 & 0\\
Lượt tới sớm quá 60 s & 7 & 3 & 1\\
Không trễ quá 60 s & 2.231 & 2.262 & 2.255\\
Không trễ quá 120 s & 2.332 & 2.327 & 2.366\\
Khách đi chung & 555 & 313 & 428\\
% EV: L1a-02; L1a-04 (U, C-30 tái lập); N-06; C06; C07; C08; N-09 (L1a-07); C4A-13 (sớm quá 60 s 7 / 3 / 1, ADD1:12,17,22; C120R:14,16,18)
\midrule
\multicolumn{4}{@{}l}{\textit{Lời hứa đầu theo dự báo}}\\
Lượt trễ quá 60 s & 137 & 10 & 65\\
Lượt tới sớm quá 60 s & 8 & 4 & 2\\
Không trễ quá 60 s & 2.297 & 2.328 & 2.322\\
Không trễ quá 120 s & 2.347 & 2.338 & 2.379\\
Lời hứa dài thêm, trung bình gộp theo khách (s) & 7,18 & 7,37 & 7,27\\
% EV: L1a-03; N-03; N-06; N-16 (sớm quá 60 s 8 / 4 / 2); C09; F06
\bottomrule
\end{tabular}
\end{table}

Cách đọc bảng: với trần 120 giây và lời hứa đầu không dự báo, 2.387 yêu cầu được phục vụ, 132 lượt trễ quá 60 giây,
nên 2.255 lượt không trễ quá 60 giây trên 3.672 yêu cầu. Chỉ 1 lượt tới sớm quá 60 giây.
% EV: N-06; L1a-02; C4A-13

So với trần 30 giây, trần 120 giây phục vụ thêm 49 lượt và có thêm 56 lượt trễ quá một phút, nên ít hơn 7 lượt không
trễ quá một phút. Nhưng nó thêm 39 lượt không trễ quá hai phút, nhiều hơn cả trần 30 giây lẫn điều phối không trần.
Khi có dự báo, nó ít hơn trần 30 giây 6 lượt ở ngưỡng một phút và thêm 41 lượt ở ngưỡng hai phút. Độ dài thêm của lời
hứa gần như bằng nhau (7,27 so với 7,37 giây, trên hai tập khách khác nhau). Mục~\ref{muc:c4-dem} đặt điểm này cạnh
cách đệm lời hứa.
% EV: L1a-09 (+49 / +56 / −7 / +39; dự báo −6 / +41); N-06; L1a-02 (2.366 > U 2.332 và C-30 2.327); K12; L1a-03; N-16; L1a-08

Dưới trần 120 giây, 9 khách bị dời từ 108 giây trở lên (90\% của trần), không khách nào vượt trần và không quyết định
nào bị ép trong 20.435 lượt.
% EV: N-07 (≥ 0,9 B: 9); L1a-05; N-23

Hai tập khách khác nhau: 2.266 khách được cả hai mức trần phục vụ, 121 khách chỉ được trần 120 giây và 72 khách chỉ
được trần 30 giây phục vụ; 93 trong 121 khách đó cũng được phục vụ khi không trần. Vì vậy 49 lượt thêm là chênh lệch
hai tổng, không phải 49 trong 96 lượt mà trần 30 giây bỏ. Phần thêm dồn vào khung 8 và 17 giờ; so từng kịch bản, trần
120 giây phục vụ nhiều hơn ở 23, bằng ở 8 và ít hơn ở 3.
% EV: N-10 (lưu ý: không viết "lấy lại 49 trong 96 lượt"); L1a-08; N-11 (C-120 − C-30: phục vụ +4 / +26 / +19, trễ quá 60 s +2 / +34 / +20; 23 / 8 / 3 kịch bản); L1a-09

\begin{figure}[htbp]
\centering
\begin{tikzpicture}[font=\footnotesize,
  vuong/.style={draw,fill=white,rectangle,minimum size=4.2pt,inner sep=0pt},
  tamgiac/.style={draw,fill=white,regular polygon,regular polygon sides=3,minimum size=6pt,inner sep=0pt},
  tron/.style={draw,fill=white,circle,minimum size=4.6pt,inner sep=0pt},
  thoi/.style={draw,fill=black,diamond,minimum size=6.5pt,inner sep=0pt},
  nhan/.style={font=\scriptsize,inner sep=1.5pt}]
% (a) Tập tắc đường cố định. Trục ngang: phục vụ − 1.530, 0,06 cm mỗi lượt; trục dọc: trễ, 0,05 cm mỗi lượt.
\begin{scope}[shift={(0,0)}]
\draw[->] (0,0)--(3.95,0);
\draw[->] (0,0)--(0,3.35);
\foreach \vx/\so in {0/1.530,1.2/1.550,2.4/1.570,3.6/1.590}
  \draw (\vx,0)--(\vx,-0.07) node[below,nhan]{\so};
\foreach \vy/\so in {0/0,1/20,2/40,3/60}
  \draw (0,\vy)--(-0.07,\vy) node[left,nhan]{\so};
\draw[dashed] (0.24,0)--(0.60,0)--(1.98,1.15)--(3.24,2.35);
\node[vuong] at (0.24,0) {}; \node[nhan,above] at (0.24,0.05) {30};
\node[vuong] at (0.60,0) {}; \node[nhan,above] at (0.62,0.05) {60};
\node[vuong] at (1.98,1.15) {}; \node[nhan,above left] at (1.98,1.15) {120};
\node[vuong] at (3.24,2.35) {}; \node[nhan,above left] at (3.24,2.35) {U};
\node at (1.8,-0.72) {Lượt phục vụ};
\node[rotate=90] at (-0.85,1.6) {Lượt trễ quá 60 s};
\node[anchor=south] at (1.9,3.4) {(a) Tắc đường cố định};
\end{scope}
% (b) Tập chỉnh trần buổi sáng. Trục ngang: phục vụ − 658, 0,36 cm mỗi lượt; trục dọc: trễ, 0,05 cm mỗi lượt.
% C-120, C-150, C-180 trùng một điểm (666; 51); U, C-240, C-300 trùng một điểm (665; 50).
\begin{scope}[shift={(5.3,0)}]
\draw[->] (0,0)--(3.95,0);
\draw[->] (0,0)--(0,3.35);
\foreach \vx/\so in {0/658,0.72/660,1.44/662,2.16/664,2.88/666,3.6/668}
  \draw (\vx,0)--(\vx,-0.07) node[below,nhan]{\so};
\foreach \vy/\so in {0/0,1/20,2/40,3/60}
  \draw (0,\vy)--(-0.07,\vy) node[left,nhan]{\so};
\draw[densely dotted,thick] (0.72,2.35)--(1.08,2.35)--(2.16,2.50)--(2.88,2.55)--(2.52,2.50);
\node[tamgiac] at (0.72,2.35) {}; \node[nhan,below] at (0.72,2.25) {30};
\node[tamgiac] at (1.08,2.35) {}; \node[nhan,below] at (1.10,2.25) {60};
\node[tamgiac] at (2.16,2.50) {}; \node[nhan,below] at (2.16,2.40) {90};
\node[tamgiac] at (2.88,2.55) {}; \node[nhan,right] at (2.98,2.55) {120--180};
\node[tamgiac] at (2.52,2.50) {}; \node[nhan,above left] at (2.50,2.58) {U, 240, 300};
\node at (1.8,-0.72) {Lượt phục vụ};
\node[rotate=90] at (-0.85,1.6) {Lượt trễ quá 60 s};
\node[anchor=south] at (1.9,3.4) {(b) Buổi sáng};
\end{scope}
% (c) Tập kiểm tra chính. Trục ngang: phục vụ − 2.320, 0,03 cm mỗi lượt; trục dọc: trễ, 0,0125 cm mỗi lượt.
\begin{scope}[shift={(10.6,0)}]
\draw[->] (0,0)--(3.95,0);
\draw[->] (0,0)--(0,3.35);
\foreach \vx/\so in {0/2.320,1.2/2.360,2.4/2.400,3.6/2.440}
  \draw (\vx,0)--(\vx,-0.07) node[below,nhan]{\so};
\foreach \vy/\so in {0/0,1/80,2/160,3/240}
  \draw (0,\vy)--(-0.07,\vy) node[left,nhan]{\so};
\draw (0.54,0.95)--(2.01,1.65)--(3.42,2.54);
\node[tron] at (0.54,0.95) {}; \node[nhan,below right] at (0.54,0.95) {30};
\node[thoi] at (2.01,1.65) {}; \node[nhan,below right] at (2.06,1.62) {120};
\node[tron] at (3.42,2.54) {}; \node[nhan,above left] at (3.42,2.54) {U};
\node at (1.8,-0.72) {Lượt phục vụ};
\node[rotate=90] at (-0.85,1.6) {Lượt trễ quá 60 s};
\node[anchor=south] at (1.9,3.4) {(c) Tập kiểm tra chính};
\end{scope}
\end{tikzpicture}
\caption[Lượt trễ theo lượt phục vụ ở các mức trần]{Lượt trễ quá 60 giây so với lời hứa đầu (trục dọc) theo lượt phục vụ (trục ngang) ở các mức
trần, lời hứa đầu không dự báo. Mỗi khung là một tập với thang riêng: (a) tập tắc đường cố định, 1.728
yêu cầu; (b) tập chỉnh trần buổi sáng, 864 yêu cầu, đủ tám mức, các mức có cùng số đếm chung một điểm; (c) tập
kiểm tra chính, 3.672 yêu cầu. Số cạnh mỗi điểm là trần tính bằng giây, U là điều phối không trần. Đường nối các mức
của cùng một tập theo thứ tự trần tăng dần, chỉ giúp theo dõi các điểm đã đo và không biểu diễn kết quả ở các mức
chưa chạy; các tập không nối với nhau. Ở (a), ghép tự do (1.629 lượt phục vụ, 274 lượt trễ) nằm ngoài khung. Hình
thoi đen là trần 120 giây trên tập kiểm tra chính.}
\label{hinh:c4-tran}
\end{figure}
% EV (a): D09 (C-30 1.534/0; C-60 1.540/0; C-120 1.563/23; U 1.584/47; FREE 1.629/274)
% EV (b): D01; D02; C4A-05 (C-30 660/47; C-60 661/47; C-90 664/50; C-120 = C-150 = C-180 666/51; U = C-240 = C-300 665/50)
% EV (c): C01; C06; N-06 (C-30 2.338/76; C-120 2.387/132; U 2.434/203)
% Tọa độ: (a) x = (phục vụ − 1530) × 0,06, y = trễ × 0,05; (b) x = (phục vụ − 658) × 0,36, y = trễ × 0,05
% (C-60: x = 3 × 0,36 = 1,08; y = 47 × 0,05 = 2,35); (c) x = (phục vụ − 2320) × 0,03, y = trễ × 0,0125.

Trên cả ba tập, nới trần từ 30 lên 120 giây thêm lượt phục vụ và thêm lượt trễ quá một phút. Số lượt không trễ quá một
phút thêm 6, thêm 2 và bớt 7 lần lượt ở tập tắc đường cố định, tập buổi sáng và tập chính, còn ở ngưỡng hai phút trần
120 giây không kém trần 30 giây ở cả ba tập. Vì mỗi kịch bản chỉ chạy một lần và mức 120 giây được chọn sau khi đã xem
tập kiểm tra chính, trần 120 giây chỉ là một điểm đo thêm trên đường đánh đổi.
% EV: D10; D03; L1a-09; D09 (không trễ quá 120 s: 1.563 so với 1.534); D02 (659 so với 658); L1a-02; K36; D11; B11; I06; C4A-11; K42

% -----------------------------------------------------------------------------------------------------
\section{Phần cộng dồn của luật C: kiểm toàn bộ bản ghi chặn}
\label{muc:c4-congdon}

Luật C cộng mọi lần dời của một khách qua cả chuyến rồi mới so với trần. Trên các lần chạy đã kiểm, hiếm khi chỉ riêng
phép cộng ấy làm một phương án bị loại. Trong 98 lần chạy trần 30--300 giây của tập kiểm tra chính và tập buổi sáng,
cả 1.256 bản ghi chặn đều có một lần dời đơn lẻ đã vượt trần của cấu hình tương ứng. Trên 34 lần chạy trần 120 giây
của tập kiểm tra chính, 3 trong 524 bản ghi chỉ vượt trần khi cộng thêm phần trần đã dùng trước đó.
% EV: E01 (98 = 34 + 64; 1.256 = 845 + 411); K09; N-08; L1a-06

Một bản ghi chặn là \emph{chỉ do cộng dồn} khi lần dời mới không vượt $B$ nhưng phần đã dùng $L_r$ cộng lần dời mới
thì vượt (mục~\ref{muc:luatc}). Ví dụ, ở một bản ghi ngày 14/11 với trần 120 giây, khách đã dùng 25,844 giây và phương
án mới dời thêm 99,027 giây: mỗi phần dưới 120 giây, nhưng tổng 124,871 giây vượt trần. Bảng~\ref{bang:c4-congdon} đếm
các bản ghi này trên mọi bản ghi chặn của các lần chạy đã kiểm.
% EV: A06; E05; N-27 (25,844 + 99,027 = 124,871 s); N-08 (quy tắc riêng + cumulative_decisive_check)

\begin{table}[htbp]
\centering
\caption[Kiểm toàn bộ bản ghi chặn của luật C]{Kiểm toàn bộ bản ghi chặn do trần trên 140 lần chạy có luật C dùng hệ số tắc đường đo được đã được kiểm,
thuộc bốn nhóm dưới đây. Một bản ghi ứng với một phương án bị loại, không phải một khách.}
\label{bang:c4-congdon}
\small
\setlength{\tabcolsep}{4pt}
\begin{tabular}{@{}>{\raggedright\arraybackslash}p{\dimexpr .38\linewidth-2\tabcolsep\relax}
>{\raggedleft\arraybackslash}p{\dimexpr .11\linewidth-2\tabcolsep\relax}
>{\raggedleft\arraybackslash}p{\dimexpr .16\linewidth-2\tabcolsep\relax}
>{\raggedleft\arraybackslash}p{\dimexpr .16\linewidth-2\tabcolsep\relax}
>{\raggedleft\arraybackslash}p{\dimexpr .19\linewidth-2\tabcolsep\relax}@{}}
\toprule
\textbf{Lần chạy} & \textbf{Số lần chạy} & \textbf{Bản ghi chặn} & \textbf{Chỉ do cộng dồn} &
\textbf{Khách phục vụ bị dời từ hai lần}\\
\midrule
Tập kiểm tra chính, trần 30 giây & 34 & 845 & 0 & 0\\
Tập chỉnh trần buổi sáng, trần 30--300 giây & 64 & 411 & 0 & 0\\
% EV: E01 (98 = 34 + 64; 1.256 = 845 + 411; 0 chỉ do cộng dồn); E02 (repeated = 0 trên 98 lần chạy)
Tập kiểm tra chính, trần 120 giây & 34 & 524 & 3 & 0\\
% EV: N-08; L1a-06; N-07 (0 khách ≥ 2 lần)
Chiều tuyết 15/11/2018, trần 30 giây & 8 & 163 & 0 & 0\\
% EV: N-26 (163 bản ghi, 0 chỉ do cộng dồn, ≥ 2 lần 0); N4B8-05 (bộ chạy chiều tuyết đã kiểm, CK8:65-72); R26 §Q8
\midrule
Cộng & 140 & 1.943 & 3 & 0\\
% EV: N-15 (132 / 1.780 / 3; kể cả chiều tuyết 140 / 1.943 / 3); ADDENDUM-1 mục 0.3
\bottomrule
\end{tabular}
\end{table}

Cách đọc bảng: trên 34 lần chạy trần 120 giây có 524 phương án bị loại kèm bản ghi chặn do trần, 3 bản ghi trong số đó
chỉ do cộng dồn, và không khách được phục vụ nào bị dời từ hai lần.
% EV: N-08; N-07

Phạm vi kiểm là 132 lần chạy và 1.780 bản ghi (98 lần chạy với 64.371 lượt quyết định được xác nhận, cộng 34 lần chạy
trần 120 giây), hay 140 lần chạy và 1.943 bản ghi kể cả chiều tuyết. Chưa được kiểm: 15 lần chạy hiệu chỉnh, 16 lần
chạy trong đó bộ điều phối dùng dự báo khi quyết định (khác phép phát lại) và 12 lần chạy luật C dùng hệ số đo được
của một phép quét trước đó. Kết quả vì vậy không áp cho mọi lần chạy.
% EV: N-15 (SCOPE); ADDENDUM-1 mục 0.3 (15 lần chạy hiệu chỉnh, 16 thư mục C-30F); E01 (64.371); K09; F12; C4A-08 (SWEEP:13,70: C-x, x = 30…900 s, trên thế giới đo thật W5, 2 kịch bản mỗi mức; NEW:1 by_arm chỉ có t38 và w07)

Ba bản ghi chỉ do cộng dồn nằm ở ba kịch bản khung 8 giờ (14, 15 và 18/11), trong 3 lượt quyết định; khi không trần,
khách của bản ghi ngày 14/11 bị dời hai lần đúng bằng hai lượng ở ví dụ trên. Trong 3 yêu cầu mới liên quan, 2 cuối
cùng không được nhận dưới trần 120 giây (một trong hai cũng không dưới trần 30 giây). Vì bản ghi dừng ở lỗi gặp trước
và diễn biến của các cấu hình khác nhau, điều này không chứng minh phần cộng dồn làm mất hai yêu cầu đó.
% EV: L1a-06 (ba kịch bản w08 ngày 14, 15, 18/11); N-08 (3 quyết định; 2 không được nhận, 1 phục vụ); N-27; E07; A23; I08

Ngoài hệ số đo được, bản ghi chỉ do cộng dồn còn có ở hai nơi. Trên tập tắc đường cố định, nơi bộ kiểm ghi đủ mọi bản
ghi, có 2 bản ghi ở trần 60 giây và 4 ở trần 120 giây; mọi yêu cầu mới liên quan vẫn được phục vụ. Trên các kiểu tắc
đường tự tạo của phép quét trước đó, sau khi bỏ bản chạy trùng, có 4 bản ghi trong 3 lượt quyết định ở trần 120 giây;
1 yêu cầu mới liên quan cuối cùng bị từ chối.
% EV: E03; E04 (4/3 sau khi bỏ bản lặp; thô 6/5); A23 (NEW:13: tập tắc đường cố định bật chế độ ghi đủ); FJC:31-54; FJR:2-7; C4A-08 (thế giới tự tạo thuộc tier3-sweep-v1)

Về định nghĩa, luật C vẫn khác luật ``mỗi lần dời không quá 30 giây'': luật sau cho một khách bị dời ba lần, tổng 90
giây, còn luật C trần 30 giây thì không. Trên số liệu, khác biệt này hiếm: khi không trần chỉ 12 trong 2.434 khách bị
dời từ hai lần. Tính thử cả hai cách trên các lần chạy không trần đã ghi, chúng chỉ khác kết luận ở 4 khách với trần
120 giây và 0 khách với trần 30, 60 và 300 giây. Phép tính thử không cho biết kết quả nếu chạy thật với luật kia, vì
các quyết định sẽ khác; luật mỗi lần một trần chưa được cài đặt và chưa chạy.

Khi chuyển sang ưu tiên chi phí (mục~\ref{muc:c4-vaitro}), cấu hình có trần 30 giây có một bản ghi chỉ vượt
do cộng dồn: phần đã dùng 24,404 giây cộng lần dời 13,558 giây thành 37,962 giây. Yêu cầu mới liên quan cuối
cùng vẫn được phục vụ. Cấu hình này có hai khách bị sửa giờ trả từ hai lần. Các số 0 của bảng trên chỉ thuộc
phạm vi ưu tiên ổn định đã nêu, không áp cho mọi thứ tự chọn phương án.
% EV: S15-01 (5d--5e); S15-04. Một witness, không phải một yêu cầu cuối cùng bị từ chối.

Trên 2.502 lịch sử khách được phục vụ của cấu hình ưu tiên chi phí không trần, có 7 lịch sử mà mỗi lần dời
không quá 120 giây nhưng tổng độ lớn vượt 120 giây. Ví dụ minh họa, hai lần dời 70 và 60 giây đều dưới trần
từng lần 120 giây, nhưng tổng 130 giây vượt trần cộng dồn. Số 7 đếm lịch sử khách theo phép kiểm ở mức 120
giây, không đếm khách vi phạm trần 30 giây. Nhánh so trực tiếp với luật giới hạn từng lần không được thực
hiện vì phép sàng lọc chưa đạt mốc 10 lịch sử đã đặt để phân bổ nguồn lực. Mốc này không phải ngưỡng kiểm
định thống kê; việc dừng chưa chứng minh phép cộng dồn không có tác dụng.
% EV: S15-01 (H7); S15-02 (h7_served); STAGE1B-PLAN.md mục 6,9. Ví dụ 70+60 là dựng tay.

% EV: E08; C20; E06 (CVP:14-18: 0/0/4/0 trên 2.434 khách U); K42; E09; K08

Các mục trên so luật C với điều phối không trần khi giữ nguyên bộ chọn ưu tiên ổn định, tức bộ chọn xét mức dùng trần
và phần dời trước chi phí kế hoạch (Bảng~\ref{bang:rangbuoc}). Mục~\ref{muc:c4-vaitro} xét kết quả có phụ thuộc vào
thứ tự đó không.
% EV: A08; H01; H02; A14; X3-07

% =====================================================================================================
% Chương 4, PHẦN B (mục 4.6–4.11) — bản sửa v11 (2026-10-13), BẢN DÙNG ĐỂ GHÉP.
% Sửa có mục tiêu từ drafts/09-chuong4-phanB-v10.tex theo R37 và brief R39; các bản trước giữ nguyên.
% Quyết định trước: reviews/CH4B-REVISION-LOG.md; ghi chú cũ: drafts/09-chuong4-phanB-v8-NOTES.md.
% Từng thay đổi, phép kiểm số mới và việc bàn giao: 09-chuong4-phanB-v11-CHANGES-v2.md bên cạnh.
% Nối ngay sau 08-chuong4-phanA-v5-v3.tex trong cùng % =====================================================================================================
% Chương 4, PHẦN A (mở chương, mục 4.1–4.5) — bản sửa v5 (2026-10-13), BẢN DÙNG ĐỂ GHÉP.
% Sửa có mục tiêu từ drafts/08-chuong4-phanA-v4.tex theo R34-CH4A-V4-VERIFY-CODEX-v2.md và brief R36.
% Quyết định trước: reviews/CH4A-REVISION-LOG-v2.md; ghi chú cũ: drafts/08-chuong4-phanA-v2-NOTES.md.
% Bản v4 và mọi bản trước giữ nguyên. Tệp tiếp tục R36: 08-chuong4-phanA-v5-v3.tex;
% từng thay đổi, bằng chứng mới và việc bàn giao: 08-chuong4-phanA-v5-CHANGES-v3.md. Phần B v10 nối sau tệp này.
% Mỗi dòng "% EV:" trỏ tới EVIDENCE-REGISTER.md (A01…K52), các phụ lục 1 (N-xx, L1a-xx), 2 (X3-xx),
% 4 (C4A-xx), 6 (N4B8-xx), hoặc sổ bổ sung R36-A01…03 trong tệp CHANGES-v3.md bên cạnh.
% Không dùng gói mới: TikZ và array đã nạp trong main.tex; chỉ sửa độ rộng các bảng được brief chỉ ra.
% =====================================================================================================

\chapter{Thực nghiệm và đánh giá}
\label{chuong:thucnghiem}

Chương này đo luật C được gì và mất gì trên mô phỏng FleetPy, so với điều phối không trần, luật hạn chót và cách đệm
lời hứa, tức cộng thêm một khoảng cố định vào mọi lời hứa. Chương cũng xét lời hứa theo dự báo và những chỗ dự báo
chưa đủ. Các mục đi theo ba phần của độ lệch giờ trả so với lời hứa đầu (mục~\ref{muc:so}). Phần tính cho điều phối,
gắn với việc ghép thêm khách, ở mục~\ref{muc:c4-ketqua}--\ref{muc:c4-vaitro}; phần được hồ sơ tắc đường lịch sử giải
thích, cùng cách đệm lời hứa, ở mục~\ref{muc:c4-dubao} và~\ref{muc:c4-dem}; phần còn lại ở mục~\ref{muc:c4-tuyet}.
Bảng~\ref{bang:c4-lotrinh} nối ba câu hỏi nghiên cứu ở mục~\ref{muc:c1-cauhoi} với các mục trả lời chúng.
% EV: — (đoạn dẫn, không nêu số); ba phần theo eq:so của chương 3 (X3-12); "đệm" theo GLOSSARY-VI mục B

\begin{table}[htbp]
\centering
\caption[Câu hỏi nghiên cứu và các mục trả lời]{Ba câu hỏi nghiên cứu và các mục của Chương~\ref{chuong:thucnghiem} trả lời chúng}
\label{bang:c4-lotrinh}
\small
\begin{tabular}{@{}p{6.3cm}p{2.4cm}p{6.3cm}@{}}
\toprule
\textbf{Câu hỏi} & \textbf{Mục} & \textbf{Đại lượng được đo}\\
\midrule
CH1. Phần dời do điều phối được đo thế nào, và trần trên phần đó giữ được gì, không giữ được gì?
& \ref{muc:c4-dieuphoi}, \ref{muc:c4-congdon}, \ref{muc:c4-hople}
& Số khách bị dời giờ trả vì ghép thêm khách, độ dời lớn nhất; yêu cầu bị từ chối có phương án bị trần chặn;
phép cộng các lần dời qua cả chuyến; các bộ đếm kiểm tra\\
CH2. Giữ trần đổi số lượt phục vụ, số lượt không trễ quá một và hai phút, cùng độ dài lời hứa thế nào so
với điều phối không trần, luật hạn chót và cách đệm lời hứa?
& \ref{muc:c4-ketqua}, \ref{muc:c4-chontran}, \ref{muc:c4-vaitro}, \ref{muc:c4-dem}
& Lượt phục vụ; lượt trễ và không trễ so với lời hứa đầu; quyết định bị ép; các mức trần; thứ tự tiêu chí
của bộ chọn; độ dài thêm của lời hứa\\
CH3. Lời hứa theo hồ sơ tắc đường lịch sử đổi số lượt không trễ quá một phút thế nào, và phần nào còn
chưa được giải thích?
& \ref{muc:c4-dubao}, \ref{muc:c4-tuyet}
& Lượt trễ khi có và không có dự báo; độ dài thêm của lời hứa; ba phần của độ lệch trên ngày thường và chiều
tuyết\\
\bottomrule
\end{tabular}
\end{table}
% EV: — (lộ trình, không có khẳng định; câu hỏi theo OUTLINE-V2 cuối §1.2)

% -----------------------------------------------------------------------------------------------------
\section{Thiết lập thực nghiệm}
\label{muc:c4-thietlap}

\subsection{Kịch bản thử}
\label{muc:c4-kichban}

FleetPy 1.0.2 mô phỏng xe và khách trên bộ dữ liệu chuẩn Manhattan; ở mỗi lượt quyết định, bộ chuyển đổi gửi tình
trạng hiện tại cho bộ chạy RideBound và nhận lại quyết định (mục~\ref{muc:kientruc}). Nhu cầu là các chuyến taxi New
York tuần 11--18/11/2018 trong bộ dữ liệu này~\cite{engelhardt2026}. Một kịch bản gồm một ngày, một khung giờ hai
tiếng và một trong bốn mẫu nhu cầu của ngày đó, mỗi mẫu gồm 10\% số chuyến của ngày: 108 yêu cầu trong khung giờ, phục
vụ bởi 8 xe 4 chỗ (Bảng~\ref{bang:c4-thamso}).
% EV: B01; X3-15; X3-21 (FC:395-429); B02 (FIX: 2018-11-14_sample_10_3.csv); J18 (mẫu 10% của một ngày); B08 (mẫu r1–r4); B03

Tắc đường là một hệ số chung cho cả thành phố, đổi theo khung 15 phút và lấy từ dữ liệu chuyến taxi, không phải tốc độ
từng đoạn đường; dưới đây gọi là hệ số tắc đường đo được. Mô phỏng không có nhiễu ngẫu nhiên, nên mỗi cặp
kịch bản và cấu hình chạy một lần và cho một kết quả. Ba kịch bản cùng ngày, cùng mẫu có chung 3 đến 20 trong 108 yêu
cầu, nên các kịch bản không độc lập với nhau.
% EV: B06 (sigma = shockRate = 0); K49; B11 (T38:90, 93); I05

\begin{table}[htbp]
\centering
\caption[Tham số của một kịch bản thử]{Tham số của một kịch bản thử. Các giá trị là cấu hình của thực nghiệm, không phải chuẩn ngành. Tập tắc
đường cố định (mục~\ref{muc:c4-taptu}) khác ở dòng tắc đường và chạy bằng bộ chạy cũ, có giới hạn bộ sinh và bộ
giải chưa được kiểm lại.}
\label{bang:c4-thamso}
\small
\begin{tabular}{@{}p{4.0cm}p{11.0cm}@{}}
\toprule
\textbf{Tham số} & \textbf{Giá trị}\\
\midrule
Khung giờ & 7:00--9:00, 8:00--10:00 hoặc 17:00--19:00, gọi là khung 7, 8 và 17 giờ; riêng chiều tuyết
là 14:00--16:00 và 15:00--17:00\\
Yêu cầu & 108 mỗi kịch bản, lấy trong khung giờ từ một mẫu 10\% số chuyến taxi\\
Xe & 8 xe, mỗi xe 4 chỗ\\
Đồ thị & 96 nút là các điểm đi, điểm đến lấy từ dữ liệu Manhattan; mỗi cặp nút có sẵn một thời gian đi. Đây là
đồ thị rút gọn, không phải mạng đường phố\\
Ràng buộc dịch vụ & Đón trong vòng 600 giây kể từ lúc gọi; ngồi xe không quá 1,5 lần thời gian đi
thẳng (Bảng~\ref{bang:rangbuoc})\\
Thời gian dừng đón, trả & 0 giây: xe không mất thời gian khi khách lên, xuống (kiểm trên 3 lần chạy mẫu)\\
Tắc đường & Một hệ số chung cho cả thành phố theo khung 15 phút; tốc độ gốc của mọi kịch bản lấy theo hệ số lúc
00:00\\
Bộ sinh phương án & Tối đa 100 phương án mỗi xe; mỗi lượt quyết định xét tối đa 2 yêu cầu mới\\
Bộ giải & OR-Tools 9.15, hạt giống ngẫu nhiên cố định bằng 7\\
\bottomrule
\end{tabular}
\end{table}
% EV (khung giờ, yêu cầu): B08; B20; B02; C4A-06 (chiều tuyết cũng 108 yêu cầu mỗi kịch bản, INV:14-15,17-22)
% EV (xe, đồ thị, ràng buộc dịch vụ): B03; K48 (PREF:22-26: mỗi cặp điểm đi–đến là một cạnh trực tiếp)
% EV (dừng, tắc đường): B05 (chỉ 3 lần chạy mẫu, SERV:2-4); B06; B07 (tốc độ nền lấy lúc 00:00)
% EV (bộ sinh, bộ giải): X3-17; X3-32 (giới hạn 2 yêu cầu áp cho cả lô của lượt, GEN:54-96); C4A-18 (tập tắc đường cố định: bộ chạy khác, NEW:13; giới hạn chưa kiểm)

\subsection{Các tập kịch bản và cấu hình được so}
\label{muc:c4-taptu}

Chương này dùng bốn tập kịch bản (Bảng~\ref{bang:c4-taptu}). Trên tập kiểm tra chính còn có phép so hai
thứ tự chọn phương án ở mục~\ref{muc:c4-vaitro}; phép so này dùng lại cùng kịch bản, không bổ sung ngày dữ liệu. Từ các ngày 14--18/11/2018 dựng được 49 kịch bản, chia
theo cặp ngày và mẫu để ba khung giờ của một cặp nằm cùng một phía. 34 kịch bản làm tập kiểm tra chính; 15 kịch bản
còn lại không dùng để tính các kết quả phục vụ và đúng giờ của chương này. Các lần chạy trên 15 kịch bản đó chỉ tham
gia phép kiểm tính nhất quán của sổ ở mục~\ref{muc:kiemtra}. Ngày 14--16/11 là ngày thường, 17 và 18/11 là cuối tuần.
% EV: B08 (SPLIT: chia theo nhóm ngày, mẫu); C17 (15 kịch bản hiệu chỉnh khoảng dự phòng γ); B09; B21 (11/11/2018 là Chủ nhật); A26; RUN38; T38:99-106; chương 3 v4-v3:725-727 (185 lần chạy gồm 15 lần hiệu chỉnh)

\begin{table}[htbp]
\centering
\caption[Các tập kịch bản dùng trong thực nghiệm]{Bốn tập kịch bản dùng trong chương. U: điều phối không trần; C-$B$: luật C trần $B$ giây;
M$^{+}$-$B$: luật hạn chót $B$ giây; V-30: luật trần độ dời khách thấy 30 giây; FREE: ghép tự do
(Bảng~\ref{bang:luat}). Cột cuối là số lần chạy, mỗi lần một cặp kịch bản và cấu hình.}
\label{bang:c4-taptu}
\footnotesize
\setlength{\tabcolsep}{2.5pt}
\begin{tabular}{@{}p{2.5cm}p{1.9cm}p{3.1cm}p{3.0cm}p{2.6cm}p{1.9cm}@{}}
\toprule
\textbf{Tập} & \textbf{Quy mô} & \textbf{Ngày, khung giờ} & \textbf{Cấu hình dùng ở chương này} &
\textbf{Bộ chạy} & \textbf{Lần chạy}\\
\midrule
Tập kiểm tra chính & 34 kịch bản, 3.672 yêu cầu & 14--18/11/2018; khung 7, 8, 17 giờ có 9, 14, 11
kịch bản & U, C-30, C-120, M$^{+}$-30, M$^{+}$-60, V-30 & Bản đã lưu, có luật phục hồi không tệ hơn &
170 (năm cấu hình), 34 (C-120); thêm 68 ưu tiên chi phí\\
% EV: B08; B09; B10 (170 = 34 × 5); L1a-01 (34 lần chạy C-120); A19; A28; X3-10 (R120:3-8)
Tập chỉnh trần buổi sáng & 8 kịch bản, 864 yêu cầu & 12--13/11/2018; khung 7 giờ & U; C và M$^{+}$ ở 30, 60,
90, 120, 150, 180, 240 và 300 giây & Như tập chính & 136 (trong 192 lần chạy đã lưu)\\
% EV: B15; B16; D01; D04; X3-10 (RW7:1-8; RW7X:6-8); X3-28; C4A-07 (17 cấu hình × 8 = 136; 88 + 104 = 192, RW7:33-34; RW7X:35-36)
Tập tắc đường cố định & 16 kịch bản, 1.728 yêu cầu & 12--13/11/2018; khung 8 và 17 giờ; hệ số không
đổi trong lần chạy & FREE, U, C-30, C-60, C-120 & Bản cũ, không có luật phục hồi & 16 mỗi cấu hình\\
% EV: B17; D08; D09; H13 (n = 16 mỗi cấu hình); LAD:2-6 (travelSnapshotVersion {1: 16}: một bảng giờ đi); C4A-04
Chiều tuyết 15/11/2018 & 8 kịch bản, 864 yêu cầu & 15/11/2018; khung 14 và 15 giờ & C-30 &
Như tập chính & 8\\
% EV: B20; N-26 (8 lần chạy C-30, 601 khách được phục vụ); C4A-06 (864 = 8 × 108, INV:14-15,17-22); N4B8-05 (CK8:65-72); R26 §Q8; agent-numeric-output-v2.txt (manifest, cấu hình phục hồi và DLL)
\bottomrule
\end{tabular}
\end{table}

Bản đã lưu là bản bộ chạy biên dịch và lưu lại, dùng chung cho tập kiểm tra chính, tập buổi sáng, các lần chạy trần
120 giây và chiều tuyết; bản cũ chưa có luật phục hồi (mục~\ref{muc:doichung}, \ref{muc:kiemtra}). Trừ ghép tự do, mọi
cấu hình trong từng thứ tự ưu tiên dùng cùng bộ sinh, bộ kiểm và bộ chọn với luật C; điều phối không trần chỉ
khác luật C trần 30 giây ở trần giờ trả. Các mục trước mục~\ref{muc:c4-vaitro}, cùng các phép so dự báo và đệm
sau đó, dùng thứ tự ưu tiên ổn định nếu không ghi khác. Ghép tự do không khóa gì và không qua bộ kiểm cam kết.
% EV: A28; X3-10 (RW7:1-8; RW7X:6-8; R120:3-8); N4B8-05 (CK8:65-72); R26 §Q8; B17; NEW:13 (bộ chạy cũ không có commitmentRecovery); A08; H01; A17; K50; X3-38 (HO:109, b1-rolling-cost)

Lần chạy hỏng vì lỗi hạ tầng (một lần ở tập kiểm tra chính) được thay bằng bản chạy lại hoàn tất (Phụ
lục~\ref{phuluc:C}). Bản kiểm kê ngày 08/10/2026 ghi 1.991 bản ghi lần chạy; kết quả của chương chỉ lấy từ bốn tập
trên và 68 lượt ưu tiên chi phí; các đợt chạy khác chỉ được nhắc khi nêu phạm vi kiểm tra (mục~\ref{muc:c4-congdon}).
% EV: C4A-17 (T38:12; ADD1:6: C-30-d20181114-s10-r3-w17-rerun1); ADDENDUM-1 mục 0.3; B19; C4A-20 (INV:1, bản kiểm kê ngày 08/10)

\subsection{Chỉ số}
\label{muc:c4-chiso}

Một lượt phục vụ là một yêu cầu, tức một khách, được nhận và chở tới nơi. Lượt trễ và lượt tới sớm được đếm trong nhóm
khách được phục vụ, so với lời hứa đầu (Bảng~\ref{bang:c4-chiso}). Số lượt không trễ bằng số lượt phục vụ trừ số lượt
trễ, tính trên toàn bộ yêu cầu, nên yêu cầu bị từ chối không được tính là đúng giờ. Ngưỡng chính là 60 giây; ngưỡng
120 giây được báo kèm vì có phép so đổi kết luận theo ngưỡng (mục~\ref{muc:c4-ketqua}).
% EV: C01; N-07 (khách được nhận = khách được phục vụ); C06; C08 (DIALPY:8-9: bị từ chối không tính là kịp; ở 120 s U hơn C-30 5 lượt)

Lời hứa đầu là giờ trả của kế hoạch lúc khách xác nhận đặt xe, tính với thời gian đi lúc đó. Tính từ lúc gọi xe, độ dài
trung bình của lời hứa là 14,15 phút trên 2.434 khách khi không trần và 14,12 phút trên 2.338 khách dưới luật C trần 30
giây, gộp theo khách được phục vụ. Hai giá trị trung bình gần nhau trên hai tập khách khác nhau; phép so này không nói
về độ dài lời hứa của từng khách.
% EV: A13; A05; B14 (BURD:5,11); R36-A03 (p0 trừ request.arrivalTimeMs, gộp theo khách); agent-numeric-output-v2.txt

Tỉ lệ \emph{gộp theo khách} chia tổng số khách thỏa điều kiện cho tổng số khách được phục vụ; \emph{trung bình theo
kịch bản} lấy trung bình của tỉ lệ từng kịch bản, mỗi kịch bản tính như nhau. Ví dụ dựng tay: với 5 trong 10 khách ở
một kịch bản và 5 trong 90 ở kịch bản kia, gộp theo khách cho 10\%, trung bình theo kịch bản cho 27,8\%. Chênh lệch là
cấu hình được thử trừ đối chứng; số dương nghĩa là nhiều hơn.
% EV: C27 (hai cách gộp cho hai số khác nhau); K24; ví dụ dựng tay (10/100 = 10%; (50% + 5,56%)/2 = 27,78%)

Kết quả ``có dự báo'' là phép phát lại trên các lần chạy đã ghi (mục~\ref{muc:dubao}): xe nào chở ai giữ nguyên, chỉ
các giờ báo cho khách được tính lại theo dự báo, và các chỉ số dùng $D_0$. Dự báo chỉ có ở 21 kịch bản ngày thường của
tập kiểm tra chính; 13 kịch bản cuối tuần và tập buổi sáng giữ lời hứa không dự báo.
% EV: C09; F03 (chỉ vế phát lại); X3-11 (D = max(p, S) áp cho mọi lời hứa đón và trả đã công bố); F02; N-28; D02 (12–13/11 không có dự báo)

\begin{table}[htbp]
\centering
\caption[Định nghĩa các chỉ số đánh giá]{Các chỉ số dùng trong chương. $a$ là giờ trả thực tế, $p_0$ là lời hứa đầu không dự báo, $D_0$
là lời hứa đầu theo dự báo (không có dự báo thì $D_0 = p_0$); $c$ là một lần dời giờ trả do điều phối, $z$ là phần
khách thấy của một lần cập nhật lời hứa (mục~\ref{muc:bamoc}). ``Quá'' nghĩa là lớn hơn hẳn: trễ đúng 60 giây
không tính là trễ quá 60 giây.}
\label{bang:c4-chiso}
\footnotesize
\setlength{\tabcolsep}{3pt}
\begin{tabular}{@{}p{2.7cm}p{5.3cm}p{3.5cm}p{3.5cm}@{}}
\toprule
\textbf{Chỉ số} & \textbf{Định nghĩa} & \textbf{Mẫu số, cách gộp} & \textbf{Vì sao đo}\\
\midrule
Lượt phục vụ & Yêu cầu được nhận và chở tới nơi & Số đếm; chia cho số yêu cầu thì ra tỉ lệ phục vụ & Giá của
luật tính bằng chuyến\\
Trễ quá $\theta$ & Khách được phục vụ có $a - D_0 > \theta$; $\theta$ = 30, 60, 120, 300 hoặc 600 giây & Đếm trong
khách được phục vụ & Lời hứa bị vượt ở phía muộn\\
Không trễ quá $\theta$ & Lượt phục vụ trừ lượt trễ quá $\theta$ & Chia cho số yêu cầu; yêu cầu bị từ chối không
tính là đúng giờ & Phục vụ và đúng giờ trên một mẫu số\\
Tới sớm quá 60 giây & Khách được phục vụ có $a - D_0 < -60$ giây & Đếm trong khách được phục vụ & Lời
hứa quá dài cũng sai\\
Lời hứa dài thêm & $D_0 - p_0$ & Trung bình gộp theo khách được phục vụ, kể cả khách không có dự báo & Giá của
dự báo hay đệm\\
Khách đi chung & Có khách khác trên cùng xe trong khoảng từ lúc được đón tới lúc được trả & Gộp theo
khách được phục vụ & Mức ghép chuyến còn lại\\
Khách bị điều phối dời giờ trả & Có ít nhất một lần $c \neq 0$; độ dời của khách là $\sum|c|$ &
Gộp theo khách được phục vụ & Đại lượng luật C giới hạn\\
Quyết định bị ép & Lượt quyết định có ít nhất một xe bị ép: luật lời hứa đã loại cả phương án giữ nguyên kế hoạch
của xe, và xe đi theo phương án được giữ nhờ luật phục hồi, có ít nhất một khách được miễn trừ
(mục~\ref{muc:doichung}) & Chia cho số lượt quyết định đã được bộ mô phỏng xác nhận áp dụng; gộp theo
lượt quyết định, mỗi lượt đếm một lần & Luật có tự giữ được nghĩa vụ không\\
Khách bị phá lời hứa & Khách mà luật đang dùng không giữ được điều nó hứa. Luật C: $\sum|c| > B$; luật hạn chót:
$a - p_0 > B$; luật trần độ dời khách thấy: $\sum|z| > B$ & Chia cho số khách được phục vụ; gộp theo
khách & Luật có giữ được điều nó hứa không\\
\bottomrule
\end{tabular}
\end{table}
% EV (phục vụ): C01; L1a-01 (2.387 hoàn tất + 1.285 bị từ chối = 3.672); N-07
% EV (trễ, không trễ, sớm): C06; C08 (DIALPY:8-9); N-04 (trễ > 60 s và sớm < −60 s rời nhau); C4A-01 (ngưỡng 30 s, cùng quy tắc sweep_analyze.py:107,128-132); C09; C10
% EV (dài thêm, đi chung, bị dời): F06; N-06 (gộp theo khách, ext_sum / served, ADD1:14,19,24); N-09 (định nghĩa đi chung); C20; N-07; A06; X3-05
% EV (bị ép): C4A-10 (RS:1012-1027; VALID:141-145; HVA:252-300; POL:443-463); X3-30; X3-09; C13 (trung bình theo kịch bản)
% EV (phá lời hứa): N-21 (định nghĩa của M⁺-30 và C-30; số gộp); A11; GLOSSARY-VI mục B "phá lời hứa"

\subsection{Mức trần 30 giây và tính chất thăm dò của kết quả}
\label{muc:c4-thamdo}

Trần chính $B$ = 30 giây là trần giờ trả của cấu hình có trần ở tập kịch bản tĩnh của đợt thử trước, nơi thời gian đi
không đổi (Phụ lục~\ref{phuluc:A}). Mức này được giữ nguyên cho tập kiểm tra chính, không chỉnh theo kết quả của tập
này. Nó không phải ngưỡng cảm nhận của khách hay chuẩn ngành, nên mỗi mức trần chỉ là một điểm trên đường đánh đổi
(mục~\ref{muc:c4-chontran}).
% EV: C4A-15 (TEX08:136-137: C1 = trần 30 giây cho tổng độ dời giờ trả; thay X3-18 đã bỏ); K29 (tập thời gian đi tĩnh); K33; B10 (năm cấu hình cố định trước lần chạy đầu, T38:3-6); D20; K52; K36

Các kết quả là thăm dò và dùng lại dữ liệu đã xem. Ngưỡng phân tích và cách chia tập được chọn sau khi đã xem tập buổi
sáng và các đợt chạy thử trước đó. Mức thử 120 giây được chọn sau khi đã xem kết quả của điều phối không trần và luật
C trần 30 giây trên chính tập kiểm tra chính, kể cả các phép tính thử ở 120 giây (mục~\ref{muc:c4-dieuphoi},
\ref{muc:c4-congdon}). Các con số vì vậy mô tả điều đã đo, không phải phép kiểm trên dữ liệu chưa từng xem.
% EV: I06; T38:8 (chọn sau khi xem Việc A, B, w07x); C4A-11 (BUD 07/10 = D14; CVP 10/10 = E06; C120R:10 chạy 10/10; T38:5 tập chính chạy 1–2/10); K42

% -----------------------------------------------------------------------------------------------------
\section{Kết quả chính khi chưa có dự báo}
\label{muc:c4-ketqua}

Với bộ chọn ưu tiên ổn định trên 3.672 yêu cầu của tập kiểm tra chính, luật C trần 30 giây không để phần dời giờ trả do điều phối của khách
nào vượt 30 giây: lớn nhất 29,2 giây, so với 590,1 giây (gần 10 phút) khi không trần. Không lượt nào của luật C trễ
quá 5 phút so với lời hứa đầu, so với 18 lượt khi không trần. Đổi lại, luật C phục vụ 2.338 lượt, ít hơn điều phối
không trần 96 lượt, và có ít hơn 242 khách đi chung. Số lượt không trễ quá một phút tăng 31, từ 2.231 lên 2.262, nhưng
số lượt không trễ quá hai phút giảm 5, từ 2.332 xuống 2.327.
% EV: C20; N-23; N-24; C07; C01; C02; N-20; C08

Mục~\ref{muc:c4-dem} cho thấy khi có dự báo, cộng thêm 30 giây vào mọi lời hứa của điều phối không trần cũng cho 2.328
lượt không trễ quá một phút, bằng luật C có dự báo. Vì vậy khóa luận không dùng chỉ số này làm lý do chọn luật C.
% EV: C10 (U + dự báo + 30 s: 2.328 = C-30 + dự báo: 2.328); K06; K51

Luật hạn chót chỉ mượn phần hạn chót của Schulz và Pfeiffer~\cite{schulz2026}: giờ trả mới không được muộn hơn lời hứa
đầu quá 30 hoặc 60 giây, thay vì 900 giây của họ. Nó chạy cùng bộ điều phối và luật phục hồi với luật C, nên đây không
phải phép so với hệ thống của hai tác giả.
% EV: J10; J11; A09; A19; H01; X3-28; X3-41

\begin{table}[htbp]
\centering
\caption[Kết quả chính khi chưa có dự báo]{Kết quả trên tập kiểm tra chính với lời hứa đầu không dự báo (34 kịch bản, 3.672 yêu cầu mỗi
cấu hình, với bộ chọn ưu tiên ổn định). Số đếm gộp trên mọi kịch bản; trễ tính so với lời hứa đầu. Hai cột ``không trễ quá'' là số lượt phục vụ
không trễ quá ngưỡng; chia cho 3.672 thì ra tỉ lệ trên yêu cầu. ``Chưa đếm'': chưa có số đếm chính xác.}
\label{bang:c4-chinh}
\small
\setlength{\tabcolsep}{4pt}
\begin{tabular}{@{}p{\dimexpr .12\linewidth-2\tabcolsep\relax}
>{\raggedleft\arraybackslash}p{\dimexpr .13\linewidth-2\tabcolsep\relax}
>{\raggedleft\arraybackslash}p{\dimexpr .16\linewidth-2\tabcolsep\relax}
>{\raggedleft\arraybackslash}p{\dimexpr .22\linewidth-2\tabcolsep\relax}
>{\raggedleft\arraybackslash}p{\dimexpr .23\linewidth-2\tabcolsep\relax}
>{\raggedleft\arraybackslash}p{\dimexpr .14\linewidth-2\tabcolsep\relax}@{}}
\toprule
\textbf{Cấu hình} & \textbf{Phục vụ} & \textbf{Trễ quá 60 s} & \textbf{Không trễ quá 60 s} &
\textbf{Không trễ quá 120 s} & \textbf{Khách đi chung}\\
\midrule
U & 2.434 & 203 & 2.231 & 2.332 & 555\\
C-30 & 2.338 & 76 & 2.262 & 2.327 & 313\\
M$^{+}$-30 & 2.338 & 74 & 2.264 & 2.327 & 307\\
% EV (U, C-30, M⁺-30): C01; C06; C08; N-09 (555 / 313 / 307, số đếm chính xác)
M$^{+}$-60 & 2.347 & 79 & 2.268 & 2.336 & chưa đếm\\
V-30 & 2.336 & 74 & 2.262 & 2.325 & chưa đếm\\
% EV (M⁺-60, V-30): C01; C4A-02 (CK4A check: trễ quá 60 / 120 s 79 / 11 và 74 / 11; EXP:5-6 trễ quá 60 s 79 và 74; đi chung chỉ có % làm tròn ở EXP)
\bottomrule
\end{tabular}
\end{table}

Cách đọc bảng: dưới trần 30 giây, 2.338 yêu cầu được phục vụ; 76 lượt trong số đó tới muộn hơn lời hứa đầu quá 60
giây, nên 2.338 $-$ 76 = 2.262 lượt không trễ quá 60 giây.
% EV: C01; C06; C08

Cái giá trước hết là số lượt phục vụ: bớt 96 lượt nghĩa là tỉ lệ phục vụ giảm 2,6 điểm phần trăm (96 trên 3.672 yêu
cầu), tức ít hơn 3,9\% số lượt của điều phối không trần (96 trên 2.434). Mức giảm nằm ở khung 8 giờ (1.014 xuống 975)
và 17 giờ (904 xuống 852); khung 7 giờ chỉ giảm 5. Số khách đi chung giảm từ 555 xuống 313, tức từ 22,8\% xuống 13,4\%
số khách được phục vụ (gộp theo khách). Khóa luận không có bằng chứng về giá trị của các chuyến đi chung này đối với
khách hay hãng xe.
% EV: C02; C03; N-09; N-20 (lưu ý: D20, I07, I10)

Số lượt không trễ quá một phút vẫn tăng vì phục vụ bớt 96 lượt nhưng trễ bớt 127 lượt: 127 $-$ 96 = 31. Hai cấu hình
cũng phục vụ những khách khác nhau: 2.207 khách chung, 227 khách chỉ được điều phối không trần và 131 khách chỉ được
luật C phục vụ. Vì vậy 31 là chênh lệch hai tổng, không phải 31 khách cụ thể được đưa tới đúng giờ.
% EV: C06; C08; C04 (câu về chờ, ngồi xe của khách chung, C05, cắt vì ngân sách trang; xem NOTES)

So với luật hạn chót 30 giây, luật C phục vụ đúng bằng số lượt (2.338) nhưng không hoàn toàn cùng những khách (2.332
khách chung, mỗi bên 6 khách riêng), và có 76 so với 74 lượt trễ quá một phút. Luật hạn chót 60 giây và luật trần độ
dời khách thấy cũng gần như vậy (Bảng~\ref{bang:c4-chinh}). Các luật lời hứa khác nhau ở chỗ mỗi luật có giữ được điều
chính nó hứa hay không (Bảng~\ref{bang:c4-nghiavu}).
% EV: C01; C04; C06; K05; K07; C32; C4A-02

\begin{table}[htbp]
\centering
\caption[Nghĩa vụ của chính từng luật lời hứa]{Nghĩa vụ của chính từng luật trên tập kiểm tra chính (34 kịch bản). Hai cột đầu ghi số đếm gộp: số quyết
định bị ép trên số quyết định được xác nhận, và số khách bị phá lời hứa trên số khách được phục vụ. Cột cuối ghi số
kịch bản có ít nhất một khách bị phá lời hứa trên 34 kịch bản. Định nghĩa ở Bảng~\ref{bang:c4-chiso}. Ở hàng luật C,
các số 0 là số đo trên tập đã chạy, phù hợp với tính chất có điều kiện ở Nhận xét~\ref{nx:giunguyen}.}
\label{bang:c4-nghiavu}
\small
\setlength{\tabcolsep}{4pt}
\begin{tabular}{@{}p{\dimexpr .12\linewidth-2\tabcolsep\relax}
>{\raggedleft\arraybackslash}p{\dimexpr .28\linewidth-2\tabcolsep\relax}
>{\raggedleft\arraybackslash}p{\dimexpr .31\linewidth-2\tabcolsep\relax}
>{\raggedleft\arraybackslash}p{\dimexpr .29\linewidth-2\tabcolsep\relax}@{}}
\toprule
\textbf{Cấu hình} & \textbf{Quyết định bị ép / được xác nhận} &
\textbf{Khách bị phá lời hứa / được phục vụ} &
\textbf{Kịch bản có khách bị phá / tổng kịch bản}\\
\midrule
C-30 & 0 / 20.101 & 0 / 2.338 & 0 / 34\\
M$^{+}$-30 & 7.710 / 20.103 & 289 / 2.338 & 28 / 34\\
M$^{+}$-60 & 4.031 / 20.163 & 79 / 2.347 & 22 / 34\\
V-30 & 10.442 / 20.090 & 456 / 2.336 & 34 / 34\\
\bottomrule
\end{tabular}
\end{table}
% EV: R36-A01 (sổ bổ sung trong CHANGES-v3.md; 170 tệp IVC và đối chiếu PAIRS); N-21; X3-30; A22

Cách đọc bảng: luật hạn chót 30 giây có 7.710 trong 20.103 lượt quyết định bị ép và 289 trong 2.338 khách tới muộn
hơn lời hứa đầu quá 30 giây; những khách này nằm ở 28 trong 34 kịch bản.
% EV: N-21; R36-A01

Trên cả 34 kịch bản đã chạy, luật C không có quyết định bị ép và không khách nào vượt trần. Theo Nhận
xét~\ref{nx:giunguyen}, nếu giả thiết (N) đúng ở lượt đang xét và mọi khách có lời hứa của xe có $0 \le L_r \le B$,
phương án giữ nguyên không thêm phần dời do điều phối nên vẫn qua kiểm tra trần. Khi (N) đúng ở mọi lượt, trần được
giữ từ lúc xác nhận. Các số đo ở đây phù hợp với tính chất đó; nhận xét không bảo đảm giờ trả thực tế. Điều phối không
trần cũng có 0 trong 20.760 quyết định bị ép. Luật hạn chót so giờ trả mới với lời hứa đầu, nên phần thay đổi khi giữ
nguyên kế hoạch vẫn có thể làm xe bị ép (mục~\ref{muc:doichung}).
% EV: A21; A22; N-21; X3-16 (U 0 / 20.760); X3-30; A10; K05; Ch3 v4-v3:371-390; R36-A01

Tính gộp, luật hạn chót 30 giây bị ép ở 7.710 trong 20.103 lượt quyết định (38,4\%) và phá lời hứa với 289 trong 2.338
khách (12,4\%). Gánh nặng dồn vào khung 7 giờ: cả 9 kịch bản có khách bị phá lời hứa, với 181 trong 509 khách được phục
vụ và 3.447 trong 4.564 lượt quyết định bị ép. Ở khung 17 giờ, 5 trong 11 kịch bản có khách bị phá, với 8 trong 852
khách và 328 trong 7.163 lượt quyết định bị ép. Các số này được gộp theo từng khung giờ. Ở mức 60 giây, các tỉ lệ toàn
tập giảm nhưng không về 0. Luật trần độ dời khách thấy bị ép nhiều nhất, phù hợp với định nghĩa: tổng độ lớn các lần
cập nhật mà nó kiểm không nhỏ hơn độ lùi so với lời hứa đầu (công thức~\eqref{eq:dongnhat}).
% EV: N-21 (gộp 7.710/20.103 và 289/2.338); R36-A02 (cặp đếm theo khung giờ, sổ bổ sung trong CHANGES-v3.md); C14; C13; R36-A01; A11; A12; X3-33 (eq:dongnhat: Σ|z| ≥ |Σz| = |p_n − p_0|)

Bảng~\ref{bang:c4-nghiavu} không so hai chuẩn đúng giờ: luật C hứa về phần dời do điều phối, còn luật hạn chót hứa về
toàn bộ độ trễ so với lời hứa đầu, một lời hứa lớn hơn. Đo bằng chính thước của luật hạn chót, tức tới muộn hơn lời
hứa đầu quá 30 giây, luật C có 292 trong 2.338 khách (12,5\%), so với 289 khách (12,4\%) của luật hạn chót. Cả hai
luật đều có khách tới muộn hơn lời hứa đầu quá 30 giây ở 28 trong 34 kịch bản. Vậy về độ đúng giờ, hai luật gần như
ngang nhau; luật C khác ở chỗ chỉ hứa phần mà bộ điều phối kiểm soát được, nên giữ được lời hứa đó.
% EV: C13; K05; C4A-01 (C-30 292/2.338 = 12,49%, 28/34); N-21; R36-A01

Bảng~\ref{bang:c4-phanbotre} cho thấy số lượt trễ giảm ở cả bốn ngưỡng đã đo, từ 1 tới 10 phút.
% EV: N-24

\begin{table}[htbp]
\centering
\caption[Phân bố trễ theo ngưỡng khi chưa có dự báo]{Phân bố trễ so với lời hứa đầu không dự báo trên tập kiểm tra chính: số lượt trễ quá mỗi ngưỡng
(trong số khách được phục vụ) và số lượt đó tính trên 100 yêu cầu (3.672 yêu cầu mỗi cấu hình).}
\label{bang:c4-phanbotre}
\small
\begin{tabular}{@{}lrrrr@{}}
\toprule
& \multicolumn{2}{c}{\textbf{U}} & \multicolumn{2}{c}{\textbf{C-30}}\\
\cmidrule(lr){2-3}\cmidrule(lr){4-5}
\textbf{Trễ quá} & \textbf{Lượt} & \textbf{Trên 100 yêu cầu} & \textbf{Lượt} & \textbf{Trên 100 yêu cầu}\\
\midrule
1 phút & 203 & 5,53 & 76 & 2,07\\
2 phút & 102 & 2,78 & 11 & 0,30\\
5 phút & 18 & 0,49 & 0 & 0\\
10 phút & 2 & 0,05 & 0 & 0\\
\midrule
Trễ lớn nhất (s) & \multicolumn{2}{c}{660,1} & \multicolumn{2}{c}{232,3}\\
\bottomrule
\end{tabular}
\end{table}
% EV: N-24 (203/102/18/2 và 76/11/0/0; 5,53/2,78/0,49/0,05 và 2,07/0,30/0/0); C07; C28; C4A-14 (CON:4,11: +660,1 s và +232,3 s; khớp C26 "11 phút" và "3,87 phút")

Cách đọc bảng: khi không trần, 18 lượt tới muộn hơn lời hứa đầu quá 5 phút, tức 0,49 lượt trên 100 yêu cầu.
% EV: N-24; C28

Khi không trần, lượt trễ nhất muộn 660,1 giây (khoảng 11 phút); dưới trần 30 giây là 232,3 giây (3,87 phút). Trong 76
khách còn trễ quá một phút dưới luật C, 72 người không có lần dời do điều phối nào được ghi và 61 người thuộc khung 7
giờ. Lời hứa theo dự báo (mục~\ref{muc:c4-dubao}) nhắm tới phần đoán được của loại trễ này.
% EV: N-24; C4A-14; C26; C21 (4/76 từng bị dời); K39 (không viết "trễ thuần vì kẹt xe"); N-11 (trễ quá 60 s của C-30 theo khung 61 / 15 / 0)

Tách theo kịch bản, luật C không có nhiều lượt trễ quá một phút hơn điều phối không trần ở cả 34 kịch bản và có số
lượt không trễ quá một phút bằng hoặc nhiều hơn ở 29 kịch bản. Ở khung 8 và 17 giờ, điều phối không trần nhiều lượt
trễ hơn ở cả 25 kịch bản; ở khung 7 giờ, luật C ít hơn ở 3 kịch bản và bằng ở 6. Vì 34 kịch bản thuộc 5 ngày và không
độc lập, đây là phép đếm mô tả, không phải kiểm định thống kê.
% EV: N-22 (34/34 trễ ≤ U; 29/34 không trễ ≥ U; bảng theo ngày × khung ở ADD1:173-186); C30; I05; B11

% -----------------------------------------------------------------------------------------------------
\section{Phần do điều phối: luật C chặn gì và mất gì}
\label{muc:c4-dieuphoi}

Dưới luật C trần 30 giây, 33 trong 2.338 khách được phục vụ bị điều phối dời giờ trả, mỗi khách nhiều nhất 29,2 giây;
khi không trần là 174 trong 2.434 khách, một khách bị dời tổng cộng 590,1 giây (Bảng~\ref{bang:c4-doidieuphoi}). Cái
giá là số yêu cầu bị từ chối tăng từ 1.238 lên 1.334; chỉ 148 trong số đó có phương án bị trần chặn được ghi lại.
% EV: C20; N-23; C01 (1.238 = 3.672 − 2.434; CON:3); D13 (1.334; 148)

\begin{table}[htbp]
\centering
\caption[Phần dời giờ trả do điều phối trên tập kiểm tra chính]{Phần dời giờ trả do điều phối trên tập kiểm tra chính, lời hứa đầu không dự báo. Mẫu số của
dòng đầu là số khách được phục vụ (2.434, 2.338 và 2.387); độ dời của một khách là tổng $|c|$ của khách
đó. Dòng cuối lấy tổng phần dời có dấu do điều phối chia cho tổng số giây trễ so với lời hứa đầu của chính các khách
trễ quá 60 s: 23.027 trên 31.353 giây, 82 trên 6.892 giây và 4.836 trên 12.450 giây. Cột trần 120 giây được dùng
lại ở mục~\ref{muc:c4-chontran}.}
\label{bang:c4-doidieuphoi}
\small
\begin{tabular}{@{}p{7.6cm}rrr@{}}
\toprule
\textbf{Chỉ số} & \textbf{U} & \textbf{C-30} & \textbf{C-120}\\
\midrule
Khách bị điều phối dời giờ trả & 174 (7,15\%) & 33 (1,41\%) & 103 (4,32\%)\\
Khách bị dời từ hai lần trở lên & 12 & 0 & 0\\
Tổng dời lớn nhất của một khách (s) & 590,1 & 29,2 & 118,9\\
Trung vị độ lớn mỗi lần dời (s) & 108,6 & 8,3 & 62,5\\
% EV: C20; C22 (trung vị theo lần: 108,6 s trên 187 lần, 8,3 s trên 33 lần); N-07 (C-120: 103, 0 khách ≥ 2 lần, nên trung vị theo lần = trung vị theo khách 62,54 s); N-23
Khách trễ quá 60 s từng bị dời & 134 trên 203 & 4 trên 76 & 60 trên 132\\
Phần tính cho điều phối trong tổng giây trễ của các khách trễ quá 60 s (\%) & 73,4 & 1,2 & 38,8\\
% EV: C21; C24; N-23 (cột C-120: 60/132, 38,8%); C4A-03 (tử số, mẫu số: 23.027,239 / 31.353,407; 81,742 / 6.891,526; 4.836,281 / 12.450,223 s)
\bottomrule
\end{tabular}
\end{table}

Cách đọc bảng: khi không trần, 174 khách bị điều phối dời giờ trả ít nhất một lần, 12 người trong số đó bị dời từ hai
lần, và tổng dời lớn nhất của một khách là 590,1 giây.
% EV: C20

Dưới trần 30 giây, mỗi khách bị dời chỉ một lần, trung vị 8,3 giây, so với 108,6 giây trên 187 lần dời khi không trần.
Trong 33 khách bị dời, 29 người bị dời muộn đi; 4 người bị kéo sớm hơn, mỗi người đúng 1 mili giây. Khi không trần,
134 trong 203 khách trễ quá một phút từng bị dời; con số chỉ cho biết hai nhóm trùng nhau, không cho biết việc bị dời
là nguyên nhân của độ trễ. Theo sổ theo từng khách (công thức~\eqref{eq:so}), phần tính cho điều phối bằng 73,4\% tổng
số giây trễ của 203 khách đó, còn dưới trần 30 giây chỉ 1,2\%: gần như mọi giây trễ còn lại nằm ở hai phần kia của sổ.
% EV: E02; C22; C23 (29 muộn, 4 sớm); N-05 (mỗi lần dời sớm đúng −1 ms; chỉ kiểm ở U, C-30, C-120, M⁺-30; nguyên nhân chưa kiểm, ADDENDUM-1 mục 0.2; C4A-19); C21 (giao hai tập); C24 (sổ mô hình, không nhân quả); K04; X3-12

Giới hạn này không dựa vào dự báo, vì dự báo không đi vào bộ điều phối. Ở 13 kịch bản cuối tuần không có dự báo, tổng
dời lớn nhất của một khách là 29,2 giây ở trần 30 giây và 116,5 giây ở trần 120 giây. Trên 8 kịch bản chiều tuyết,
8 trong 601 khách được phục vụ bị dời, lớn nhất 25,2 giây; không ai vượt trần và không quyết định nào bị ép trong
5.067 lượt.
% EV: F03; X3-11 (D chỉ tính sau khi chạy); N-28; N-26 (thay anchor F13 của RD §8 mục 4, ADDENDUM-1 mục 0.4); N4B8-05 (CK8:65-72); R26 §Q8

Một ví dụ trên tập buổi sáng cho thấy luật C và luật hạn chót khác nhau ở đâu. Trước một quyết định, giờ trả đang hứa
của một khách đã lùi 79,0 giây so với lời hứa đầu, toàn bộ là phần thay đổi khi giữ nguyên kế hoạch. Phương án ghép
thêm một khách mới sẽ dời giờ trả dự kiến của khách đó thêm 68,3 giây do điều phối. Ở mức 90 giây, luật hạn chót cộng
hai phần, khoảng 147 giây, nên loại phương án; luật C chỉ tính 68,3 giây nên cho qua. Trên tập này hai luật tách nhau
ở đúng ba chỗ cùng kiểu, ở mức 90 hoặc 120 giây; mỗi chỗ luật C nhận thêm một khách. Đây là phép so hai luật trên
giờ trả dự kiến tại cùng một trạng thái của mô hình, không phải phép đo nhân quả.
% EV: D05 (79,0 + 68,3 s; 60,9 + 61,2 s); C4A-12 (W07XM:6,11,16: "sum x 79.0, sum c 0.0, p_last − p_first 79.0", c = 68,3 s, x = 0,1 s; W07X:54-62: "khách đang có", giai đoạn chờ đón hay trên xe không ghi); K04

Khi luật C loại một phương án vì nó làm một khách đã có lời hứa vượt trần, bộ kiểm ghi một bản ghi chặn gồm khách đó,
phần trần đã dùng, lần dời mới và trần (Bảng~\ref{bang:kiemtran}). Dưới trần 30 giây, 148 trong 1.334 yêu cầu bị từ
chối có ít nhất một phương án mang bản ghi chặn; điều phối không trần phục vụ 116 trong số đó. 1.186 yêu cầu còn lại
không có bản ghi chặn nào được ghi, và điều phối không trần cũng từ chối 1.075 yêu cầu trong số này; điều đó không
chứng minh chúng bị từ chối chỉ vì ràng buộc dịch vụ. Bản ghi chặn cũng không phải khách: 845 bản ghi nằm ở 643 yêu
cầu, và 495 yêu cầu trong số đó cuối cùng vẫn được phục vụ.
Các số đếm theo yêu cầu trong đoạn này tính theo từng cặp kịch bản và yêu cầu.
% EV: X3-05 (nội dung bản ghi chặn); D13 (BUD:2-5: 148 = 116 + 32; 1.186 = 111 + 1.075); K10; C4A-09 (BUDPY:29-44: nhóm 1.186 là phần dư không có bản ghi chặn được ghi, không phải lý do cấu trúc đã chứng minh); D16 (SF:2: 643 cặp, 495 phục vụ, 148 từ chối); E01 (845)

Với mỗi yêu cầu trong 148 yêu cầu, phép phân tích chọn bản ghi có mức vượt nhỏ nhất, tức phần vượt quá 30 giây; trung
vị là 109,6 giây, lớn nhất 524,8 giây. Chỉ ở 2 trong 148 bản ghi được chọn, khách nêu trong bản ghi đã dùng một phần
trần trước đó; ở cả 148, riêng lần dời mới đã vượt 30 giây. Với trần 120 giây, mức vượt của 59 trong 148 yêu cầu nằm
trong phần nới thêm 90 giây. Không thể kết luận 59 yêu cầu đó sẽ được phục vụ: bộ kiểm dừng ở khách vi phạm gặp trước,
và khi trần đổi thì các quyết định trước cũng đổi. Mục~\ref{muc:c4-chontran} cho kết quả chạy thật.
% EV: D14 (BUD:6-8: trung vị 109,6 s; lớn nhất 524,8 s; 59/148 có mức vượt ≤ 90 s); D19 (BUD:10; FJC:4-6); C4A-09 (BUDPY:22,35-42,51-52: mức vượt = after − limit; 2/148 đếm theo bản ghi được chọn của từng yêu cầu, không phải số khách khác nhau); K11; A23; I08

% -----------------------------------------------------------------------------------------------------
\section{Chọn mức trần}
\label{muc:c4-chontran}

Trong ba điểm đã đo trên tập kiểm tra chính, không điểm nào tốt nhất ở mọi chỉ số. Trần 30 giây có ít lượt trễ quá một
phút nhất (76). Trần 120 giây có nhiều lượt không trễ quá hai phút nhất (2.366, so với 2.327 của trần 30 giây và 2.332
khi không trần) và phục vụ 2.387 lượt, nằm giữa hai điểm kia. Chọn mức nào tùy ngưỡng mà hãng xe coi trọng, điều thực
nghiệm không trả lời được.
% EV: N-06; L1a-02; C01; C06; C08; K12; K36

Mức trần được đo trên ba tập không cộng gộp (Hình~\ref{hinh:c4-tran}). Tập tắc đường cố định cho thấy trần làm gì khi
thời gian đi không đổi, tập buổi sáng dò tám mức trên cùng 8 kịch bản, tập kiểm tra chính thêm mức 120 giây.
% EV: B17; B15; B16; X3-28; D01; L1a-01; I12

\textbf{Tập tắc đường cố định.} Bảng~\ref{bang:c4-codinh} xếp các cấu hình từ lỏng nhất (ghép tự do) tới chặt nhất
(trần 30 giây), với hệ số tắc đường cố định suốt lần chạy.
% EV: B17; LAD:2-6

\begin{table}[htbp]
\centering
\caption[Các mức trần trên tập tắc đường cố định]{Các mức trần trên tập tắc đường cố định (16 kịch bản, 1.728 yêu cầu mỗi cấu hình; bộ chạy cũ,
không có luật phục hồi). Ba cột số đầu là số đếm; hai cột cuối là số khách được phục vụ thỏa điều kiện, trong ngoặc
là tỉ lệ gộp theo khách được phục vụ (\%).}
\label{bang:c4-codinh}
\small
\begin{tabular}{@{}lrrrrr@{}}
\toprule
\textbf{Cấu hình} & \textbf{Phục vụ} & \textbf{Trễ quá 60 s} & \textbf{Không trễ quá 120 s} &
\textbf{Đi chung} & \textbf{Bị dời giờ trả}\\
\midrule
FREE & 1.629 & 274 & 1.465 & 865 (53,1) & 388 (23,8)\\
U & 1.584 & 47 & 1.557 & 346 (21,8) & 62 (3,9)\\
C-120 & 1.563 & 23 & 1.563 & 303 (19,4) & 41 (2,6)\\
C-60 & 1.540 & 0 & 1.540 & 266 (17,3) & 20 (1,3)\\
C-30 & 1.534 & 0 & 1.534 & 241 (15,7) & 6 (0,4)\\
\bottomrule
\end{tabular}
\end{table}
% EV: D09 (số đếm, WP17); D08 (tỉ lệ đi chung, bị dời; gộp theo khách); C4A-04 (đi chung 865/346/303/266/241, CK4A fixed-shared; bị dời 388/62/41/20/6, WBU:4,7,10,13,16); B17

Cách đọc bảng: với trần 120 giây, 1.563 trong 1.728 yêu cầu được phục vụ, 23 lượt trễ quá 60 giây; 303 khách (19,4\%)
đi chung và 41 khách (2,6\%) bị dời giờ trả.
% EV: D08; D09; C4A-04

Ghép tự do khác điều phối không trần ở cả cách chọn phương án lẫn việc kiểm cam kết, nên chênh lệch giữa hai cấu hình
này không chỉ do trần. Trong nhóm cùng chính sách, siết từ không trần xuống trần 30 giây làm số lượt phục vụ giảm từ
1.584 xuống 1.534 và số lượt trễ quá 60 giây từ 47 xuống 0. Giữa hai mức đã chạy 60 và 120 giây, số lượt trễ tăng từ 0
lên 23. Các con số của bộ chạy cũ này không chuyển sang tập kiểm tra chính.
% EV: A17; K50; X3-38; NEW:13; D09; D11; B17; I12

\textbf{Tập chỉnh trần buổi sáng.} Trên 864 yêu cầu, số lượt phục vụ chỉ dao động trong 6 lượt, từ 660 ở trần 30 giây
tới 666 ở trần 120--180 giây (Bảng~\ref{bang:c4-sang}).
% EV: D01; B15; B16

\begin{table}[htbp]
\centering
\caption[Các mức trần trên tập chỉnh trần buổi sáng]{Tám mức trần trên tập chỉnh trần buổi sáng (8 kịch bản khung 7 giờ ngày 12--13/11/2018, 864 yêu cầu
mỗi cấu hình; lời hứa đầu không dự báo). Số đếm gộp; hai cột ``không trễ quá'' là số lượt trên 864 yêu cầu. Các mức
có cùng mọi số đếm được gộp vào một dòng.}
\label{bang:c4-sang}
\small
\begin{tabular}{@{}lrrrr@{}}
\toprule
\textbf{Cấu hình} & \textbf{Phục vụ} & \textbf{Trễ quá 60 s} & \textbf{Không trễ quá 60 s} &
\textbf{Không trễ quá 120 s}\\
\midrule
U & 665 & 50 & 615 & 656\\
C-30 & 660 & 47 & 613 & 658\\
C-60 & 661 & 47 & 614 & 659\\
C-90 & 664 & 50 & 614 & 660\\
C-120, C-150, C-180 & 666 & 51 & 615 & 659\\
C-240, C-300 & 665 & 50 & 615 & 656\\
\bottomrule
\end{tabular}
\end{table}
% EV: D01; D02 (U, C-30, C-90, C-120: CAP:3-7; DIAL17:3-7, hai bộ đọc độc lập); C4A-05 (C-60, C-150, C-180, C-240, C-300: CAP:5,8-11; DIAL17:5,8-11); D02 lưu ý (cột dự báo bằng ngây thơ ở 12–13/11)

Cách đọc bảng: với trần 120, 150 hoặc 180 giây, 666 yêu cầu được phục vụ, 51 lượt trễ quá 60 giây, 615 lượt không trễ
quá 60 giây và 659 lượt không trễ quá 120 giây.
% EV: D01; D02; C4A-05

Trần 120 giây phục vụ nhiều hơn điều phối không trần một lượt, nên điều phối không trần không phải mức phục vụ cao
nhất có thể đạt, chỉ là một điểm đo. Từ 30 lên 120 giây, trần thêm 6 lượt phục vụ, 4 lượt trễ quá 60 giây và 2 lượt
không trễ quá 60 giây. Luật C không có quyết định bị ép trong 64 lần chạy của tám mức trần này. Ở cả tám mức, hai
luật phục vụ chênh nhau không quá 2 lượt.
% EV: D01; K12; D03 (+6 / +4 / +2); D04; X3-28 (tám mức 30–300 s); E01 (64 lần C, 0 bị ép); agent-numeric-output-v2.txt (0 / 44.270 lượt quyết định)

\textbf{Tập kiểm tra chính.} Bảng~\ref{bang:c4-tran120} đặt trần 120 giây cạnh hai cấu hình của
mục~\ref{muc:c4-ketqua}, khi không có và khi có dự báo.

\begin{table}[htbp]
\centering
\caption[Trần 120 giây trên tập kiểm tra chính]{Trần 120 giây trên tập kiểm tra chính (34 kịch bản, 3.672 yêu cầu mỗi cấu hình). Số đếm gộp; các dòng
``không trễ quá'' là số lượt trên 3.672 yêu cầu. Khối dưới là phép phát lại lời hứa theo dự báo trên cùng các lần
chạy, bàn ở mục~\ref{muc:c4-dubao}; dự báo chỉ có ở 21 kịch bản ngày thường, và độ dài thêm của lời hứa là trung
bình gộp trên mọi khách được phục vụ, kể cả khách cuối tuần có độ dài thêm bằng 0.}
\label{bang:c4-tran120}
\small
\begin{tabular}{@{}p{8.4cm}rrr@{}}
\toprule
\textbf{Chỉ số} & \textbf{U} & \textbf{C-30} & \textbf{C-120}\\
\midrule
\multicolumn{4}{@{}l}{\textit{Lời hứa đầu không dự báo}}\\
Lượt phục vụ & 2.434 & 2.338 & 2.387\\
Lượt trễ quá 60 s & 203 & 76 & 132\\
Lượt trễ quá 120 s & 102 & 11 & 21\\
Lượt trễ quá 300 s & 18 & 0 & 0\\
Lượt tới sớm quá 60 s & 7 & 3 & 1\\
Không trễ quá 60 s & 2.231 & 2.262 & 2.255\\
Không trễ quá 120 s & 2.332 & 2.327 & 2.366\\
Khách đi chung & 555 & 313 & 428\\
% EV: L1a-02; L1a-04 (U, C-30 tái lập); N-06; C06; C07; C08; N-09 (L1a-07); C4A-13 (sớm quá 60 s 7 / 3 / 1, ADD1:12,17,22; C120R:14,16,18)
\midrule
\multicolumn{4}{@{}l}{\textit{Lời hứa đầu theo dự báo}}\\
Lượt trễ quá 60 s & 137 & 10 & 65\\
Lượt tới sớm quá 60 s & 8 & 4 & 2\\
Không trễ quá 60 s & 2.297 & 2.328 & 2.322\\
Không trễ quá 120 s & 2.347 & 2.338 & 2.379\\
Lời hứa dài thêm, trung bình gộp theo khách (s) & 7,18 & 7,37 & 7,27\\
% EV: L1a-03; N-03; N-06; N-16 (sớm quá 60 s 8 / 4 / 2); C09; F06
\bottomrule
\end{tabular}
\end{table}

Cách đọc bảng: với trần 120 giây và lời hứa đầu không dự báo, 2.387 yêu cầu được phục vụ, 132 lượt trễ quá 60 giây,
nên 2.255 lượt không trễ quá 60 giây trên 3.672 yêu cầu. Chỉ 1 lượt tới sớm quá 60 giây.
% EV: N-06; L1a-02; C4A-13

So với trần 30 giây, trần 120 giây phục vụ thêm 49 lượt và có thêm 56 lượt trễ quá một phút, nên ít hơn 7 lượt không
trễ quá một phút. Nhưng nó thêm 39 lượt không trễ quá hai phút, nhiều hơn cả trần 30 giây lẫn điều phối không trần.
Khi có dự báo, nó ít hơn trần 30 giây 6 lượt ở ngưỡng một phút và thêm 41 lượt ở ngưỡng hai phút. Độ dài thêm của lời
hứa gần như bằng nhau (7,27 so với 7,37 giây, trên hai tập khách khác nhau). Mục~\ref{muc:c4-dem} đặt điểm này cạnh
cách đệm lời hứa.
% EV: L1a-09 (+49 / +56 / −7 / +39; dự báo −6 / +41); N-06; L1a-02 (2.366 > U 2.332 và C-30 2.327); K12; L1a-03; N-16; L1a-08

Dưới trần 120 giây, 9 khách bị dời từ 108 giây trở lên (90\% của trần), không khách nào vượt trần và không quyết định
nào bị ép trong 20.435 lượt.
% EV: N-07 (≥ 0,9 B: 9); L1a-05; N-23

Hai tập khách khác nhau: 2.266 khách được cả hai mức trần phục vụ, 121 khách chỉ được trần 120 giây và 72 khách chỉ
được trần 30 giây phục vụ; 93 trong 121 khách đó cũng được phục vụ khi không trần. Vì vậy 49 lượt thêm là chênh lệch
hai tổng, không phải 49 trong 96 lượt mà trần 30 giây bỏ. Phần thêm dồn vào khung 8 và 17 giờ; so từng kịch bản, trần
120 giây phục vụ nhiều hơn ở 23, bằng ở 8 và ít hơn ở 3.
% EV: N-10 (lưu ý: không viết "lấy lại 49 trong 96 lượt"); L1a-08; N-11 (C-120 − C-30: phục vụ +4 / +26 / +19, trễ quá 60 s +2 / +34 / +20; 23 / 8 / 3 kịch bản); L1a-09

\begin{figure}[htbp]
\centering
\begin{tikzpicture}[font=\footnotesize,
  vuong/.style={draw,fill=white,rectangle,minimum size=4.2pt,inner sep=0pt},
  tamgiac/.style={draw,fill=white,regular polygon,regular polygon sides=3,minimum size=6pt,inner sep=0pt},
  tron/.style={draw,fill=white,circle,minimum size=4.6pt,inner sep=0pt},
  thoi/.style={draw,fill=black,diamond,minimum size=6.5pt,inner sep=0pt},
  nhan/.style={font=\scriptsize,inner sep=1.5pt}]
% (a) Tập tắc đường cố định. Trục ngang: phục vụ − 1.530, 0,06 cm mỗi lượt; trục dọc: trễ, 0,05 cm mỗi lượt.
\begin{scope}[shift={(0,0)}]
\draw[->] (0,0)--(3.95,0);
\draw[->] (0,0)--(0,3.35);
\foreach \vx/\so in {0/1.530,1.2/1.550,2.4/1.570,3.6/1.590}
  \draw (\vx,0)--(\vx,-0.07) node[below,nhan]{\so};
\foreach \vy/\so in {0/0,1/20,2/40,3/60}
  \draw (0,\vy)--(-0.07,\vy) node[left,nhan]{\so};
\draw[dashed] (0.24,0)--(0.60,0)--(1.98,1.15)--(3.24,2.35);
\node[vuong] at (0.24,0) {}; \node[nhan,above] at (0.24,0.05) {30};
\node[vuong] at (0.60,0) {}; \node[nhan,above] at (0.62,0.05) {60};
\node[vuong] at (1.98,1.15) {}; \node[nhan,above left] at (1.98,1.15) {120};
\node[vuong] at (3.24,2.35) {}; \node[nhan,above left] at (3.24,2.35) {U};
\node at (1.8,-0.72) {Lượt phục vụ};
\node[rotate=90] at (-0.85,1.6) {Lượt trễ quá 60 s};
\node[anchor=south] at (1.9,3.4) {(a) Tắc đường cố định};
\end{scope}
% (b) Tập chỉnh trần buổi sáng. Trục ngang: phục vụ − 658, 0,36 cm mỗi lượt; trục dọc: trễ, 0,05 cm mỗi lượt.
% C-120, C-150, C-180 trùng một điểm (666; 51); U, C-240, C-300 trùng một điểm (665; 50).
\begin{scope}[shift={(5.3,0)}]
\draw[->] (0,0)--(3.95,0);
\draw[->] (0,0)--(0,3.35);
\foreach \vx/\so in {0/658,0.72/660,1.44/662,2.16/664,2.88/666,3.6/668}
  \draw (\vx,0)--(\vx,-0.07) node[below,nhan]{\so};
\foreach \vy/\so in {0/0,1/20,2/40,3/60}
  \draw (0,\vy)--(-0.07,\vy) node[left,nhan]{\so};
\draw[densely dotted,thick] (0.72,2.35)--(1.08,2.35)--(2.16,2.50)--(2.88,2.55)--(2.52,2.50);
\node[tamgiac] at (0.72,2.35) {}; \node[nhan,below] at (0.72,2.25) {30};
\node[tamgiac] at (1.08,2.35) {}; \node[nhan,below] at (1.10,2.25) {60};
\node[tamgiac] at (2.16,2.50) {}; \node[nhan,below] at (2.16,2.40) {90};
\node[tamgiac] at (2.88,2.55) {}; \node[nhan,right] at (2.98,2.55) {120--180};
\node[tamgiac] at (2.52,2.50) {}; \node[nhan,above left] at (2.50,2.58) {U, 240, 300};
\node at (1.8,-0.72) {Lượt phục vụ};
\node[rotate=90] at (-0.85,1.6) {Lượt trễ quá 60 s};
\node[anchor=south] at (1.9,3.4) {(b) Buổi sáng};
\end{scope}
% (c) Tập kiểm tra chính. Trục ngang: phục vụ − 2.320, 0,03 cm mỗi lượt; trục dọc: trễ, 0,0125 cm mỗi lượt.
\begin{scope}[shift={(10.6,0)}]
\draw[->] (0,0)--(3.95,0);
\draw[->] (0,0)--(0,3.35);
\foreach \vx/\so in {0/2.320,1.2/2.360,2.4/2.400,3.6/2.440}
  \draw (\vx,0)--(\vx,-0.07) node[below,nhan]{\so};
\foreach \vy/\so in {0/0,1/80,2/160,3/240}
  \draw (0,\vy)--(-0.07,\vy) node[left,nhan]{\so};
\draw (0.54,0.95)--(2.01,1.65)--(3.42,2.54);
\node[tron] at (0.54,0.95) {}; \node[nhan,below right] at (0.54,0.95) {30};
\node[thoi] at (2.01,1.65) {}; \node[nhan,below right] at (2.06,1.62) {120};
\node[tron] at (3.42,2.54) {}; \node[nhan,above left] at (3.42,2.54) {U};
\node at (1.8,-0.72) {Lượt phục vụ};
\node[rotate=90] at (-0.85,1.6) {Lượt trễ quá 60 s};
\node[anchor=south] at (1.9,3.4) {(c) Tập kiểm tra chính};
\end{scope}
\end{tikzpicture}
\caption[Lượt trễ theo lượt phục vụ ở các mức trần]{Lượt trễ quá 60 giây so với lời hứa đầu (trục dọc) theo lượt phục vụ (trục ngang) ở các mức
trần, lời hứa đầu không dự báo. Mỗi khung là một tập với thang riêng: (a) tập tắc đường cố định, 1.728
yêu cầu; (b) tập chỉnh trần buổi sáng, 864 yêu cầu, đủ tám mức, các mức có cùng số đếm chung một điểm; (c) tập
kiểm tra chính, 3.672 yêu cầu. Số cạnh mỗi điểm là trần tính bằng giây, U là điều phối không trần. Đường nối các mức
của cùng một tập theo thứ tự trần tăng dần, chỉ giúp theo dõi các điểm đã đo và không biểu diễn kết quả ở các mức
chưa chạy; các tập không nối với nhau. Ở (a), ghép tự do (1.629 lượt phục vụ, 274 lượt trễ) nằm ngoài khung. Hình
thoi đen là trần 120 giây trên tập kiểm tra chính.}
\label{hinh:c4-tran}
\end{figure}
% EV (a): D09 (C-30 1.534/0; C-60 1.540/0; C-120 1.563/23; U 1.584/47; FREE 1.629/274)
% EV (b): D01; D02; C4A-05 (C-30 660/47; C-60 661/47; C-90 664/50; C-120 = C-150 = C-180 666/51; U = C-240 = C-300 665/50)
% EV (c): C01; C06; N-06 (C-30 2.338/76; C-120 2.387/132; U 2.434/203)
% Tọa độ: (a) x = (phục vụ − 1530) × 0,06, y = trễ × 0,05; (b) x = (phục vụ − 658) × 0,36, y = trễ × 0,05
% (C-60: x = 3 × 0,36 = 1,08; y = 47 × 0,05 = 2,35); (c) x = (phục vụ − 2320) × 0,03, y = trễ × 0,0125.

Trên cả ba tập, nới trần từ 30 lên 120 giây thêm lượt phục vụ và thêm lượt trễ quá một phút. Số lượt không trễ quá một
phút thêm 6, thêm 2 và bớt 7 lần lượt ở tập tắc đường cố định, tập buổi sáng và tập chính, còn ở ngưỡng hai phút trần
120 giây không kém trần 30 giây ở cả ba tập. Vì mỗi kịch bản chỉ chạy một lần và mức 120 giây được chọn sau khi đã xem
tập kiểm tra chính, trần 120 giây chỉ là một điểm đo thêm trên đường đánh đổi.
% EV: D10; D03; L1a-09; D09 (không trễ quá 120 s: 1.563 so với 1.534); D02 (659 so với 658); L1a-02; K36; D11; B11; I06; C4A-11; K42

% -----------------------------------------------------------------------------------------------------
\section{Phần cộng dồn của luật C: kiểm toàn bộ bản ghi chặn}
\label{muc:c4-congdon}

Luật C cộng mọi lần dời của một khách qua cả chuyến rồi mới so với trần. Trên các lần chạy đã kiểm, hiếm khi chỉ riêng
phép cộng ấy làm một phương án bị loại. Trong 98 lần chạy trần 30--300 giây của tập kiểm tra chính và tập buổi sáng,
cả 1.256 bản ghi chặn đều có một lần dời đơn lẻ đã vượt trần của cấu hình tương ứng. Trên 34 lần chạy trần 120 giây
của tập kiểm tra chính, 3 trong 524 bản ghi chỉ vượt trần khi cộng thêm phần trần đã dùng trước đó.
% EV: E01 (98 = 34 + 64; 1.256 = 845 + 411); K09; N-08; L1a-06

Một bản ghi chặn là \emph{chỉ do cộng dồn} khi lần dời mới không vượt $B$ nhưng phần đã dùng $L_r$ cộng lần dời mới
thì vượt (mục~\ref{muc:luatc}). Ví dụ, ở một bản ghi ngày 14/11 với trần 120 giây, khách đã dùng 25,844 giây và phương
án mới dời thêm 99,027 giây: mỗi phần dưới 120 giây, nhưng tổng 124,871 giây vượt trần. Bảng~\ref{bang:c4-congdon} đếm
các bản ghi này trên mọi bản ghi chặn của các lần chạy đã kiểm.
% EV: A06; E05; N-27 (25,844 + 99,027 = 124,871 s); N-08 (quy tắc riêng + cumulative_decisive_check)

\begin{table}[htbp]
\centering
\caption[Kiểm toàn bộ bản ghi chặn của luật C]{Kiểm toàn bộ bản ghi chặn do trần trên 140 lần chạy có luật C dùng hệ số tắc đường đo được đã được kiểm,
thuộc bốn nhóm dưới đây. Một bản ghi ứng với một phương án bị loại, không phải một khách.}
\label{bang:c4-congdon}
\small
\setlength{\tabcolsep}{4pt}
\begin{tabular}{@{}>{\raggedright\arraybackslash}p{\dimexpr .38\linewidth-2\tabcolsep\relax}
>{\raggedleft\arraybackslash}p{\dimexpr .11\linewidth-2\tabcolsep\relax}
>{\raggedleft\arraybackslash}p{\dimexpr .16\linewidth-2\tabcolsep\relax}
>{\raggedleft\arraybackslash}p{\dimexpr .16\linewidth-2\tabcolsep\relax}
>{\raggedleft\arraybackslash}p{\dimexpr .19\linewidth-2\tabcolsep\relax}@{}}
\toprule
\textbf{Lần chạy} & \textbf{Số lần chạy} & \textbf{Bản ghi chặn} & \textbf{Chỉ do cộng dồn} &
\textbf{Khách phục vụ bị dời từ hai lần}\\
\midrule
Tập kiểm tra chính, trần 30 giây & 34 & 845 & 0 & 0\\
Tập chỉnh trần buổi sáng, trần 30--300 giây & 64 & 411 & 0 & 0\\
% EV: E01 (98 = 34 + 64; 1.256 = 845 + 411; 0 chỉ do cộng dồn); E02 (repeated = 0 trên 98 lần chạy)
Tập kiểm tra chính, trần 120 giây & 34 & 524 & 3 & 0\\
% EV: N-08; L1a-06; N-07 (0 khách ≥ 2 lần)
Chiều tuyết 15/11/2018, trần 30 giây & 8 & 163 & 0 & 0\\
% EV: N-26 (163 bản ghi, 0 chỉ do cộng dồn, ≥ 2 lần 0); N4B8-05 (bộ chạy chiều tuyết đã kiểm, CK8:65-72); R26 §Q8
\midrule
Cộng & 140 & 1.943 & 3 & 0\\
% EV: N-15 (132 / 1.780 / 3; kể cả chiều tuyết 140 / 1.943 / 3); ADDENDUM-1 mục 0.3
\bottomrule
\end{tabular}
\end{table}

Cách đọc bảng: trên 34 lần chạy trần 120 giây có 524 phương án bị loại kèm bản ghi chặn do trần, 3 bản ghi trong số đó
chỉ do cộng dồn, và không khách được phục vụ nào bị dời từ hai lần.
% EV: N-08; N-07

Phạm vi kiểm là 132 lần chạy và 1.780 bản ghi (98 lần chạy với 64.371 lượt quyết định được xác nhận, cộng 34 lần chạy
trần 120 giây), hay 140 lần chạy và 1.943 bản ghi kể cả chiều tuyết. Chưa được kiểm: 15 lần chạy hiệu chỉnh, 16 lần
chạy trong đó bộ điều phối dùng dự báo khi quyết định (khác phép phát lại) và 12 lần chạy luật C dùng hệ số đo được
của một phép quét trước đó. Kết quả vì vậy không áp cho mọi lần chạy.
% EV: N-15 (SCOPE); ADDENDUM-1 mục 0.3 (15 lần chạy hiệu chỉnh, 16 thư mục C-30F); E01 (64.371); K09; F12; C4A-08 (SWEEP:13,70: C-x, x = 30…900 s, trên thế giới đo thật W5, 2 kịch bản mỗi mức; NEW:1 by_arm chỉ có t38 và w07)

Ba bản ghi chỉ do cộng dồn nằm ở ba kịch bản khung 8 giờ (14, 15 và 18/11), trong 3 lượt quyết định; khi không trần,
khách của bản ghi ngày 14/11 bị dời hai lần đúng bằng hai lượng ở ví dụ trên. Trong 3 yêu cầu mới liên quan, 2 cuối
cùng không được nhận dưới trần 120 giây (một trong hai cũng không dưới trần 30 giây). Vì bản ghi dừng ở lỗi gặp trước
và diễn biến của các cấu hình khác nhau, điều này không chứng minh phần cộng dồn làm mất hai yêu cầu đó.
% EV: L1a-06 (ba kịch bản w08 ngày 14, 15, 18/11); N-08 (3 quyết định; 2 không được nhận, 1 phục vụ); N-27; E07; A23; I08

Ngoài hệ số đo được, bản ghi chỉ do cộng dồn còn có ở hai nơi. Trên tập tắc đường cố định, nơi bộ kiểm ghi đủ mọi bản
ghi, có 2 bản ghi ở trần 60 giây và 4 ở trần 120 giây; mọi yêu cầu mới liên quan vẫn được phục vụ. Trên các kiểu tắc
đường tự tạo của phép quét trước đó, sau khi bỏ bản chạy trùng, có 4 bản ghi trong 3 lượt quyết định ở trần 120 giây;
1 yêu cầu mới liên quan cuối cùng bị từ chối.
% EV: E03; E04 (4/3 sau khi bỏ bản lặp; thô 6/5); A23 (NEW:13: tập tắc đường cố định bật chế độ ghi đủ); FJC:31-54; FJR:2-7; C4A-08 (thế giới tự tạo thuộc tier3-sweep-v1)

Về định nghĩa, luật C vẫn khác luật ``mỗi lần dời không quá 30 giây'': luật sau cho một khách bị dời ba lần, tổng 90
giây, còn luật C trần 30 giây thì không. Trên số liệu, khác biệt này hiếm: khi không trần chỉ 12 trong 2.434 khách bị
dời từ hai lần. Tính thử cả hai cách trên các lần chạy không trần đã ghi, chúng chỉ khác kết luận ở 4 khách với trần
120 giây và 0 khách với trần 30, 60 và 300 giây. Phép tính thử không cho biết kết quả nếu chạy thật với luật kia, vì
các quyết định sẽ khác; luật mỗi lần một trần chưa được cài đặt và chưa chạy.

Khi chuyển sang ưu tiên chi phí (mục~\ref{muc:c4-vaitro}), cấu hình có trần 30 giây có một bản ghi chỉ vượt
do cộng dồn: phần đã dùng 24,404 giây cộng lần dời 13,558 giây thành 37,962 giây. Yêu cầu mới liên quan cuối
cùng vẫn được phục vụ. Cấu hình này có hai khách bị sửa giờ trả từ hai lần. Các số 0 của bảng trên chỉ thuộc
phạm vi ưu tiên ổn định đã nêu, không áp cho mọi thứ tự chọn phương án.
% EV: S15-01 (5d--5e); S15-04. Một witness, không phải một yêu cầu cuối cùng bị từ chối.

Trên 2.502 lịch sử khách được phục vụ của cấu hình ưu tiên chi phí không trần, có 7 lịch sử mà mỗi lần dời
không quá 120 giây nhưng tổng độ lớn vượt 120 giây. Ví dụ minh họa, hai lần dời 70 và 60 giây đều dưới trần
từng lần 120 giây, nhưng tổng 130 giây vượt trần cộng dồn. Số 7 đếm lịch sử khách theo phép kiểm ở mức 120
giây, không đếm khách vi phạm trần 30 giây. Nhánh so trực tiếp với luật giới hạn từng lần không được thực
hiện vì phép sàng lọc chưa đạt mốc 10 lịch sử đã đặt để phân bổ nguồn lực. Mốc này không phải ngưỡng kiểm
định thống kê; việc dừng chưa chứng minh phép cộng dồn không có tác dụng.
% EV: S15-01 (H7); S15-02 (h7_served); STAGE1B-PLAN.md mục 6,9. Ví dụ 70+60 là dựng tay.

% EV: E08; C20; E06 (CVP:14-18: 0/0/4/0 trên 2.434 khách U); K42; E09; K08

Các mục trên so luật C với điều phối không trần khi giữ nguyên bộ chọn ưu tiên ổn định, tức bộ chọn xét mức dùng trần
và phần dời trước chi phí kế hoạch (Bảng~\ref{bang:rangbuoc}). Mục~\ref{muc:c4-vaitro} xét kết quả có phụ thuộc vào
thứ tự đó không.
% EV: A08; H01; H02; A14; X3-07

% =====================================================================================================
% Chương 4, PHẦN B (mục 4.6–4.11) — bản sửa v11 (2026-10-13), BẢN DÙNG ĐỂ GHÉP.
% Sửa có mục tiêu từ drafts/09-chuong4-phanB-v10.tex theo R37 và brief R39; các bản trước giữ nguyên.
% Quyết định trước: reviews/CH4B-REVISION-LOG.md; ghi chú cũ: drafts/09-chuong4-phanB-v8-NOTES.md.
% Từng thay đổi, phép kiểm số mới và việc bàn giao: 09-chuong4-phanB-v11-CHANGES-v2.md bên cạnh.
% Nối ngay sau 08-chuong4-phanA-v5-v3.tex trong cùng \input{noi-dung/08-chuong4}; không có \chapter mới.
% Mỗi dòng "% EV:" trỏ tới EVIDENCE-REGISTER.md (A01…K52), EVIDENCE-REGISTER-ADDENDUM-1.md (N-xx, L1a-xx, mục 0.x),
% EVIDENCE-REGISTER-ADDENDUM-2-CH3.md (X3-xx), EVIDENCE-REGISTER-ADDENDUM-3-CH4B.md (N4B-xx),
% EVIDENCE-REGISTER-ADDENDUM-4-CH4A.md (C4A-xx) hoặc EVIDENCE-REGISTER-ADDENDUM-6-CH4B-v8.md (N4B8-xx; script chỉ đọc
% evidence/checks-ch4b-v8/ch4b_v8_check.py, output ch4b-v8-check-2026-10-12.txt, phần A/B/C gọi là "check A/B/C").
% Không dùng gói mới: TikZ, enumitem và graphicx đã nạp trong main.tex; thu bảng/hình bằng resizebox.
% =====================================================================================================

% -----------------------------------------------------------------------------------------------------
\section{Vai trò của trần và thứ tự tiêu chí của bộ chọn}
\label{muc:c4-vaitro}

\subsection{Vì sao cần đổi thứ tự ưu tiên?}

Các kết quả trước mục này dùng bộ chọn ưu tiên ổn định: sau khi xét số yêu cầu nhận được, bộ chọn so mức dùng
trần và phần dời giờ rồi mới xét chi phí kế hoạch. Như vậy, cả việc xếp hạng phương án lẫn trần cứng đều có thể
góp phần hạn chế thay đổi lời hứa. Chỉ so có trần với không trần dưới một thứ tự chưa cho biết kết quả sẽ ra sao
khi bộ điều phối đặt chi phí kế hoạch lên trước phần dời giờ.

Để kiểm tra điều đó, khóa luận giữ ưu tiên nhận yêu cầu và các phép kiểm khả thi, nhưng cho bộ chọn xét chi phí
kế hoạch ngay sau đó. Các mức xếp hạng theo ổn định được bỏ, còn tiêu chí phân xử theo mã định danh giữ nguyên.
Trần cứng của luật C vẫn được kiểm trước khi chọn nếu cấu hình có trần. Hai cách gọi dưới đây là \emph{ưu tiên
ổn định} và \emph{ưu tiên chi phí kế hoạch}; chúng biểu thị mục tiêu khác nhau, không phải hai mức chất lượng
mặc định. Chi phí kế hoạch là tổng thời gian còn lại của tuyến dự kiến, chưa phải tiền hay chi phí vận hành thực.
% EV: S15-03; A14; X3-07; K15. Cùng ưu tiên nhận khách trong từng lượt, không bảo đảm cùng số phục vụ cuối.

Phép so gồm bốn cấu hình trên cùng 34 kịch bản và 3.672 yêu cầu mỗi cấu hình: hai thứ tự ưu tiên, mỗi thứ tự có
trần 30 giây và không trần. Hai cấu hình ưu tiên chi phí có 68 lượt mô phỏng hoàn tất, được đối chiếu với 68 lượt
ưu tiên ổn định đã lưu. Tập kịch bản và tập mã yêu cầu khớp ở cả 34 kịch bản. Các diễn biến sau khi chọn có thể
khác nhau, nên nhóm khách được phục vụ cũng khác nhau.
% EV: S15-01 (mục 0--3), main-final-admission.json; S15-02. Không có 34 lần lặp độc lập.

\subsection{Đánh đổi giữa phục vụ, ghép khách và sửa lời hứa}

\begin{table}[htbp]
\centering
\caption[Đánh đổi dưới hai thứ tự chọn phương án]{Hai thứ tự chọn phương án trên tập kiểm tra chính, mỗi cấu hình
có 3.672 yêu cầu. Trễ được so với lời hứa trả đầu chưa tính dự báo, trừ dòng ghi riêng. Khách có đi chung được
đếm một lần nếu hoàn thành chuyến và có khoảng thời gian cùng xe với khách khác.}
\label{bang:c4-doi-uutien}
\small
\setlength{\tabcolsep}{3pt}
\begin{tabular}{@{}>{\raggedright\arraybackslash}p{6.5cm}rrrr@{}}
\toprule
& \multicolumn{2}{c}{\textbf{Ưu tiên ổn định}} & \multicolumn{2}{c}{\textbf{Ưu tiên chi phí}}\\
\cmidrule(lr){2-3}\cmidrule(l){4-5}
\textbf{Chỉ số} & \textbf{Không trần} & \textbf{Trần 30} & \textbf{Không trần} & \textbf{Trần 30}\\
\midrule
Lượt phục vụ & 2.434 & 2.338 & 2.502 & 2.364\\
Muộn quá một phút & 203 & 76 & 379 & 72\\
Không muộn quá một phút & 2.231 & 2.262 & 2.123 & 2.292\\
Không muộn quá hai phút & 2.332 & 2.327 & 2.280 & 2.353\\
Không muộn quá năm phút & 2.416 & 2.338 & 2.466 & 2.364\\
Khách hoàn thành có đi chung & 555 & 313 & 1.027 & 432\\
Khách bị sửa giờ trả do điều phối & 174 & 33 & 434 & 106\\
Trong đó, bị sửa từ hai lần & 12 & 0 & 50 & 2\\
Đón muộn theo giới hạn dịch vụ & 51 & 43 & 39 & 44\\
Muộn quá một phút khi lời hứa có dự báo & 137 & 10 & 318 & 11\\
Tổng dời lớn nhất của một khách (s) & 590,1 & 29,2 & 709,1 & 29,7\\
\bottomrule
\end{tabular}
\end{table}
% EV: S15-01 (mục 2, 5a--5f); S15-02 (pooled); S15-04 (kiểm bảng và các chênh lệch từ TSV).

Khi cùng giữ trần 30 giây, ưu tiên chi phí phục vụ thêm 26 lượt và có thêm 119 khách hoàn thành có đi chung
so với ưu tiên ổn định (Bảng~\ref{bang:c4-doi-uutien}). Số lượt không muộn quá một phút tăng 30. Đổi lại, số khách
bị sửa giờ trả tăng từ 33 trong 2.338 lên 106 trong 2.364, tương ứng 1,41\% và 4,48\%; số khách bị sửa từ hai lần
tăng từ 0 lên 2. Số lượt đón muộn tăng từ 43 lên 44. Như vậy, cách ưu tiên ổn định vẫn cho ít khách bị sửa lời
hứa hơn, còn ưu tiên chi phí nhận và ghép được nhiều khách hơn trong thiết lập này.

Đúng giờ đến và ít sửa lời hứa là hai chỉ số khác nhau. Ưu tiên chi phí có trần có 72 lượt muộn quá một phút,
ít hơn 4 lượt so với ưu tiên ổn định có trần, nhưng nhiều khách bị sửa giờ hơn. Khi dùng lời hứa có dự báo, số
lượt muộn quá một phút là 11 và 10 trên hai nhóm khách tương ứng. Các chênh lệch nhỏ này không đủ để chọn một
cách làm tốt hơn về mọi mặt. Trung vị độ dài lời hứa đầu chưa dự báo là 809,642 giây với ưu tiên chi phí và
815,060 giây với ưu tiên ổn định, trên hai tập khách khác nhau.
% EV: S15-01 (5b--5e); S15-02. Độ dài dùng p0-t0; không quy thành chờ/đi vòng của từng khách.

Không trần cho thấy rõ cái giá của đổi thứ tự. Ưu tiên chi phí phục vụ thêm 68 lượt và có thêm 472 khách hoàn
thành có đi chung so với ưu tiên ổn định, nhưng có thêm 176 lượt muộn quá một phút và thêm 260 khách bị sửa giờ
trả. Trần vì vậy vẫn có vai trò khi bộ chọn chuyển sang ưu tiên chi phí. Trong thứ tự này, thêm trần 30 giây làm
giảm 138 lượt phục vụ và 595 khách hoàn thành có đi chung; đồng thời có ít hơn 307 lượt muộn quá một phút và
thêm 169 lượt được phục vụ, không muộn quá một phút.
% EV: S15-01 (mục 2, 5f--5g); S15-04. Các chênh lệch là tổng trên tập yêu cầu, không cùng tập khách phục vụ.

Phần giảm số lượt trễ phải đọc cùng phần không phục vụ. Với ưu tiên ổn định, 127 lượt trễ ít hơn đi cùng 96 lượt
phục vụ ít hơn và 31 lượt không muộn quá một phút nhiều hơn. Với ưu tiên chi phí, các số tương ứng là 307, 138
và 169. Chênh lệch giữa hai mức giảm trễ là 180, nhưng không có nghĩa 180 khách cụ thể được bảo vệ chỉ nhờ đổi
bộ chọn. Chênh lệch này cùng chiều ở 29 trong 34 kịch bản và cả năm ngày; các kịch bản phụ thuộc nhau nên đây
chỉ là mô tả trên dữ liệu đã dùng.
% EV: S15-01 (mục 2, 4); S15-02 (interactions); 180 = (138-96) + (169-31).

Ở ưu tiên chi phí có trần, tổng dời lớn nhất do điều phối là 29,657 giây và không khách nào vượt 30 giây.
Kết quả cho thấy cấu hình có trần giữ được giới hạn đã đặt khi thay thứ tự xếp hạng. Nó chưa chứng minh rằng
riêng trần đủ bảo đảm giờ đến thực tế, vì khóa, ràng buộc dịch vụ, bộ sinh và luật phục hồi vẫn cùng tham gia.

Mức trần thích hợp còn phụ thuộc tiêu chí dịch vụ. Với ưu tiên chi phí, trần 30 giây tăng số lượt được phục vụ
và không muộn quá một phút từ 2.123 lên 2.292; nếu chấp nhận muộn tới năm phút thì số lượt đáp ứng lại giảm
từ 2.466 xuống 2.364. Đây là hai cách đọc cùng kết quả đã có, không phải khảo sát mức khách chấp nhận chờ.
Trần 30 giây là điểm thực nghiệm, không được chọn làm ngưỡng tối ưu hay ngưỡng cảm nhận của hành khách.
% EV: S15-01 (5a); S15-02 (timely60, timely300); K36; K52.

\subsection{Lời hứa có dự báo và giới hạn của phép so}

Với ưu tiên chi phí có trần, tính lại lời hứa đầu theo dự báo cho 11 lượt muộn quá một phút thay vì 72,
trên cùng 2.364 hành trình; lời hứa dài thêm trung bình 7,30 giây. Với cấu hình không trần, các số lượt muộn
là 318 thay vì 379, và lời hứa dài thêm trung bình 7,02 giây. Giờ trả thực tế và việc xe chở ai không đổi
trong phép tính này. Dự báo điều chỉnh phần đường dự kiến chậm đi, còn luật C giới hạn phần dời do điều phối;
hai thành phần không thay thế nhau và không tạo thành bảo đảm vô điều kiện về giờ đến.
% EV: S15-01 (5b); F01--F04. 21 ô bật causal-no15; 13 ô còn lại giữ p0.

Các lượt ưu tiên chi phí dùng bản chạy riêng. Phép đối chiếu hành vi trên bốn cặp mẫu và việc tính lại các
chỉ số ưu tiên ổn định giúp kiểm tra khả năng so sánh, nhưng không chứng minh mọi khác biệt giữa hai bản
chạy đã được loại bỏ. Các con số là phép so thăm dò theo cấu hình trên một bộ mô phỏng và năm ngày đã xem,
không phải ước lượng tác động nhân quả riêng của một tiêu chí. Phép sàng lọc trên các phương án đã ghi của
tập tắc đường cố định được giữ ở Phụ lục~\ref{phuluc:D}; nó bổ trợ việc hiểu bộ chọn, không thay kết quả
của những lần mô phỏng đầy đủ ở đây.
% EV: STAGE1B-PLAN.md mục 1.1 và amendment mục 3; S15-01 (mục 0, 3); G0 chỉ behavioural-only.


\section{Lời hứa theo dự báo}
\label{muc:c4-dubao}

Trên 21 kịch bản ngày thường có dự báo, số khách của luật C trần 30 giây tới muộn hơn lời hứa đầu quá một phút giảm từ
68 xuống 2 trong 1.471 khách được phục vụ (Bảng~\ref{bang:c4-dubao}). Lời hứa theo dự báo (mục~\ref{muc:dubao}) được
tính lại trên các lần chạy đã ghi; mỗi cách hứa được so với lời hứa đầu của chính nó.
% EV: F17; F04; N4B8-01 (check A: 68 → 2 trên 1.471); F03 (chỉ vế phát lại); FFA:8-13

Đối chứng thứ hai là đệm tỉ lệ: mỗi lời hứa đầu được kéo dài một tỉ lệ cố định của khoảng từ lúc hứa tới giờ trả đã
hứa. Tỉ lệ được chọn theo từng khung giờ, trên chính các khách này, để tổng phần kéo dài bằng của lời hứa theo dự báo.
% EV: N4B-03 (FFA:15-17: lời hứa đệm = P0 + k_w(P0 − t0), k_w = Σ(D0 − P0)/Σ(P0 − t0) trên mọi khách C-30 của khung w ở 21 kịch bản, tính từ lời hứa; FA:2)

\begin{table}[htbp]
\centering
\caption[Lời hứa theo dự báo và đệm tỉ lệ trên ngày thường]{Khách của luật C trần 30 giây trên 21 kịch bản ngày thường có dự báo, theo ba cách hứa: lời hứa đầu không dự
báo, đệm tỉ lệ và lời hứa đầu theo dự báo; mỗi cách so với lời hứa đầu của chính nó. Khối trên là số đếm gộp theo khách;
dòng cuối là trung bình theo kịch bản của tỉ lệ trễ, mỗi kịch bản tính như nhau.}
\label{bang:c4-dubao}
\small
\begin{tabular}{@{}lrrrr@{}}
\toprule
\textbf{Nhóm kịch bản} & \textbf{Phục vụ} & \textbf{Không dự báo} & \textbf{Đệm tỉ lệ} & \textbf{Dự báo}\\
\midrule
\multicolumn{5}{@{}l}{\textit{Số khách tới muộn hơn lời hứa đầu quá 60 s}}\\
Cả 21 kịch bản & 1.471 & 68 & 22 & 2\\
Khung 7 giờ (8 kịch bản) & 444 & 60 & 17 & 1\\
Khung 8 giờ (8 kịch bản) & 609 & 8 & 5 & 1\\
Khung 17 giờ (5 kịch bản) & 418 & 0 & 0 & 0\\
\multicolumn{5}{@{}l}{\textit{Số khách tới sớm hơn lời hứa đầu quá 60 s}}\\
Cả 21 kịch bản & 1.471 & 0 & 0 & 1\\
% EV: N4B8-01 (check A: CELLS, share × served từng kịch bản, phần dư 0); N4B8-08 (FFA:139-141 tính E_pad; sửa N4B-03: đệm tỉ lệ có đo tới sớm, bằng 0); F17
\midrule
\multicolumn{2}{@{}l}{Tỉ lệ trễ quá 60 s, trung bình theo kịch bản (\%)} & 6,01 & 1,94 & 0,17\\
% EV: F04 (FA:8); N4B-03 (FA:8: pad 0,0194)
\bottomrule
\end{tabular}
\end{table}

Cách đọc bảng: ở khung 7 giờ, 60 trong 444 khách trễ quá một phút khi không dự báo, 17 khi đệm tỉ lệ và 1 khi có dự
báo.
% EV: N4B8-01

Tính trung bình theo kịch bản, tỉ lệ trễ giảm 5,84 điểm phần trăm. Gom 21 kịch bản thành 8 nhóm theo ngày và khung giờ,
khoảng tin cậy 95\% theo phân phối Student là từ 1,15 đến 10,54 điểm phần trăm. So với đệm tỉ lệ, chênh lệch là
$-1{,}77$ điểm phần trăm, khoảng từ $-3{,}62$ đến $+0{,}07$, tức có chứa 0: chưa thể nói dự báo tốt hơn cách đệm này.
Các khoảng chỉ dựa trên ba buổi sáng đã được xem trước, và mọi hồ sơ dự báo đều có ngày 12/11 trong các ngày nguồn.
% EV: F04; F05 (FA:33,36: dY −5,844 [−10,542; −1,146]; dPad −1,775 [−3,620; +0,071]; FFA cluster_t: tầng = khung giờ, cụm = ngày–khung, df = 8 − 3 = 5); F11 (FR:97-98; factor_nowcast_screen.py:21-22: 14 ← 12, 13; 15, 16 ← 12, 13, 14); FR:170-171; I05; K24

Lợi ích không riêng của luật C. Trên cùng 21 kịch bản, khi có dự báo, luật hạn chót 30 giây giảm từ 66 xuống 2 khách
trễ trong 1.471 khách được phục vụ; luật trần độ dời khách thấy 30 giây cũng giảm từ 66 xuống 2 trong 1.470 khách.
Điều phối không trần giảm từ 129 xuống 63 trong 1.512 khách. Riêng khung 17 giờ, 28 khách trễ trong 437 khách của điều
phối không trần không đổi, còn luật C có 0 trong 418 khách. Theo sổ, 86,1\% số giây trễ của 68 khách trễ khi không dự
báo thuộc phần được hồ sơ lịch sử giải thích (Hình~\ref{hinh:c4-so}); dự báo và trần nhắm vào hai phần khác nhau của sổ.
% EV: F08; K37; N4B8-02 (check A: M⁺-30 66 → 2 trên 1.471; V-30 66 → 2 trên 1.470; U 129 → 63 trên 1.512; U khung 17 giờ 28 → 28 trên 437; C-30 khung 17 giờ 0 trên 418); F09 (FA:30; RX:3); eq:so; K04

Trên 21 kịch bản này, giá của dự báo là lời hứa giờ trả dài thêm 13,7 giây (khung 7 giờ 27,8 giây) và lời hứa giờ đón 3,2 giây, trung bình
theo kịch bản. Trên cả 34 kịch bản, khách trễ của luật C giảm từ 76 xuống 10 trong 2.338 (Bảng~\ref{bang:c4-tran120}).
% EV: F06 (FR:67-69; FA:8-9); C06 (76); C09 (10 trên 34 kịch bản); N-16; F02; N-28 (13 kịch bản cuối tuần không có dự báo; độ dài thêm gộp 7,37 s, N4B-10, nằm ở bang:c4-tran120 và bang:c4-dem)

Kết quả 2 trong 1.471 khách đạt được trong điều kiện thuận lợi cho dự báo theo giờ trong ngày, không phải hiệu quả ngoài
đời. Tắc đường là một hệ số chung cho cả thành phố. Trên cả 34 kịch bản, với luật C, kéo dài lời hứa theo đúng hệ số
sẽ xảy ra cho giờ trả khớp trong vòng 2 giây ở cả 2.305 khách không bị điều phối dời. Phép kiểm giờ đón cũng chưa dùng
lời hứa này (mục~\ref{muc:c4-thaoluan}).
% EV: I02 (FR:167-169: mức tốt nhất cho loại dự báo theo giờ trong ngày trong mô phỏng gần như tất định); F10 (FA:28; RX:8: 34 kịch bản); B06; F07; X3-11

\begin{figure}[htbp]
\centering
\resizebox{\linewidth}{!}{%
\begin{tikzpicture}[x=1cm,y=1cm,font=\footnotesize,
  nhan/.style={font=\scriptsize,inner sep=1.5pt}]
% Thanh trên: ngày thường. Độ dài 11 cm = 6.217,4 s; ranh giới 0,1445 và 9,6145 cm (FA:30).
\draw[fill=black!75] (0,1.6) rectangle (0.1445,2.2);
\draw[fill=black!20] (0.1445,1.6) rectangle (9.6145,2.2);
\draw[fill=white] (9.6145,1.6) rectangle (11,2.2);
\node at (4.88,1.9) {5.353 s (86,1\%)};
\node[align=center,nhan] at (10.307,1.9) {783 s\\(12,6\%)};
\draw (0.072,2.2)--(0.072,2.45);
\node[nhan,anchor=south west] at (0.0,2.45) {82 s (1,3\%)};
\node[anchor=east,align=right] at (-0.15,1.9) {Ngày thường, 21 kịch bản\\68 khách, 6.217 s};
% Thanh dưới: chiều tuyết. Độ dài 11 cm = 3.445,9 s; ranh giới 0,0273 và 0,8302 cm (NEW:5).
\draw[fill=black!75] (0,0) rectangle (0.0273,0.6);
\draw[fill=black!20] (0.0273,0) rectangle (0.8302,0.6);
\draw[fill=white] (0.8302,0) rectangle (11,0.6);
\node at (5.915,0.3) {3.185,9 s (92,5\%)};
% v8: nhãn 8,6 s đặt phía trên thanh tuyết (khoảng trống giữa hai thanh) để đường dẫn của nhãn 251,5 s không cắt chữ.
\draw (0.014,0.6)--(0.014,0.85);
\node[nhan,anchor=south west] at (-0.02,0.85) {8,6 s (0,25\%)};
\draw (0.43,0)--(0.43,-0.2);
\node[nhan,anchor=north west] at (0.39,-0.2) {251,5 s (7,3\%)};
\node[anchor=east,align=right] at (-0.15,0.3) {Chiều tuyết, 8 kịch bản\\41 khách, 3.445,9 s};
% Chú giải (ba ô cách nhau đủ để chữ không chồng lên nhau)
\draw[fill=black!75] (0,-1.25) rectangle (0.3,-1.0);
\node[anchor=west] at (0.35,-1.125) {phần tính cho điều phối};
\draw[fill=black!20] (4.0,-1.25) rectangle (4.3,-1.0);
\node[anchor=west] at (4.35,-1.125) {phần được hồ sơ lịch sử giải thích};
\draw[fill=white] (9.6,-1.25) rectangle (9.9,-1.0);
\node[anchor=west] at (9.95,-1.125) {phần còn lại};
\end{tikzpicture}
}
\caption[Ba phần của số giây trễ theo sổ của từng khách]{Số giây trễ chia thành ba phần theo sổ của từng khách (công thức~\eqref{eq:so}), với luật C trần 30 giây và các
khách tới muộn hơn lời hứa đầu không dự báo quá 60 giây: 68 khách của 21 kịch bản ngày thường và 41 khách của 8 kịch
bản chiều tuyết 15/11/2018. Mỗi thanh dài bằng 100\% số giây trễ của nhóm đó. Phần được hồ sơ lịch sử giải thích phụ
thuộc dự báo đang dùng; phần còn lại là phần dự báo này chưa giải thích được, không phải phần không thể biết trước.}
\label{hinh:c4-so}
\end{figure}
% EV (ngày thường): F09; N4B8-11 (FA:30 và CELLS: 81,742 + 5.352,646 + 783,018 = 6.217,406 s; các nhãn làm tròn đến giây độc lập, nên tổng ba nhãn là 6.218 s, lệch 1 s so với nhãn tổng 6.217 s); RX:3 (1,3 / 86,1 / 12,6%)
% EV (chiều tuyết): F13; N4B-05 (NEW:5: 8,566 / 251,496 / 3.445,933 s; RP1:24-26: 7,3%; SNOW:10: 92,5%; FR:104,186); N4B8-12 ([A] 8,566 / 3.445,933 = 0,25%)
% Tọa độ: x = giây / tổng × 11 cm (tính tay; xem NOTES v7 mục 4).

\textbf{Đưa dự báo vào bộ điều phối.} Phép thử chạy luật C và luật hạn chót 30 giây trên 16 kịch bản khung 7 và 8 giờ.
Thời gian đi của bộ điều phối được nhân thêm mức tăng dự báo trung bình của khung 15 phút hiện tại và khung kế tiếp. Có
3 trong 32 lần chạy dừng vì FleetPy tính rằng xe sẽ tới điểm đón muộn hơn giờ đón muộn nhất chưa tới 3 giây. Luật C
không có quyết định bị ép ở 14 lần chạy hoàn tất. Luật hạn chót bị ép ở 47,6\% lượt quyết định ở 15 lần chạy hoàn tất,
so với 62,8\% ở các lần chạy không đưa dự báo vào trên cùng 15 kịch bản. Tính gộp trên cùng 15 kịch bản, đó là 3.849
trong 8.496 lượt quyết định khi đưa dự báo vào, so với 5.121 trong 8.468 lượt khi không đưa vào; các tỉ lệ 47,6\% và
62,8\% ở trên là trung bình theo kịch bản. Tỉ lệ khách bị luật hạn chót phá lời hứa ở khung 7 giờ giảm từ 39,0\% xuống
11,0\%, trung bình trên cùng tám kịch bản. Tính gộp theo khách trên tám kịch bản này, đó là 170 trong 442 khách khi
không đưa dự báo vào, so với 45 trong 439 khách khi đưa vào. Vì vậy, các số của Bảng~\ref{bang:c4-nghiavu} phụ thuộc cách
hứa. Khóa luận chỉ dùng dự báo ở lời hứa báo cho khách; kết luận này chỉ áp cho cách đã thử.
% EV: F12; X3-29 (W13:85-89); N4B8-06 (K = 2: run_forecast.py:9,72; configs/forecast-k2-2018-11-1{4,5,6}.json:1; preflight_world_v1_3.py:85-89; FR:117; R26 Q9); N4B-07 (RX:9,11-15; FR:120-121,124,138,157: C hoàn tất 14, luật hạn chót 15; 47,6 và 62,8% trên cùng 15 kịch bản); C14; FR:141-146; R37 verify-numbers-v2.json:dispatcher; R39 agent-numbers-v1.json:dispatcher.paired, morning_promise_breaks (cùng15 hoàn tất; khung7 cùng8; số gộp tách khỏi trung bình theo kịch bản)

\textbf{Học máy.} Ba mô hình học một lần rồi giữ nguyên ước lượng phân vị 90 của độ trễ so với lời hứa theo dự báo (mức
mà 90\% khách giống vậy không vượt) từ thông tin đã có ở thời điểm hứa. Đó là mô hình cây quyết định tăng cường
(gradient boosting), mô hình ấy khi bỏ hai biến trạng thái xe (số khách khác đã gán và đang trên xe), và hồi quy phân vị
tuyến tính. Mỗi khách được cộng vào lời hứa theo dự báo một phần đệm tỉ lệ với mức ước lượng; đối chứng cộng cùng một
khoảng cho mọi khách. Cả hai được chỉnh trên ngày hiệu chỉnh để lời hứa dài thêm trung bình 30 giây so với lời hứa không
dự báo, nên trên ngày kiểm mức này có thể khác. Phép thử dùng hồ sơ dự báo cố định của ngày
12--13/11 và kiểm trên 15, 16/11 khung 7 giờ và 13/11 của tập tắc đường cố định, các ngày đã được xem trước. Mô hình
chỉ được coi là tốt hơn khi, so với đối chứng, thỏa cả bốn điều kiện:
\begin{itemize}[leftmargin=1.2cm,itemsep=0pt,topsep=2pt]
\item tỉ lệ khách trễ quá một phút thấp hơn;
\item độ dài thêm trung bình không hơn quá 1 giây;
\item phân vị 90 của độ dài thêm không hơn quá 15 giây;
\item tỉ lệ khách tới sớm quá một phút không hơn quá 1 điểm phần trăm.
\end{itemize}
% EV: G01; G02; G03; N4B-08; N4B8-04 (MLS:9-10 các ngày đã xem; 19-23, 67-68 hai biến trạng thái; 37-39 hồ sơ cố định 12–13/11; 41-61 chia ngày; 63-69 ba mô hình; 71-82 chỉnh trên ngày hiệu chỉnh, F = D0 + ρ·w, đối chứng D0 + (E − mean δ), không chỉnh lại khi kiểm; 91-95 bốn điều kiện; MLT: đối chứng của U dài thêm 37,219 / 31,772 s ở 15/11 / 16/11)

Không mô hình nào thỏa đủ bốn điều kiện. Các số đếm sau được suy từ tỉ lệ gộp theo khách đã làm tròn đến ba chữ số
thập phân trong bảng kết quả; với mẫu số tương ứng, mỗi tỉ lệ chỉ ứng với một số nguyên. Với luật C, đối chứng và mô hình đều không có khách trễ (147 và 197 khách).
Với điều phối không trần, mô hình cây tăng cường chỉ bớt 1 trong 2 khách trễ ở ngày 15/11 (145 khách), không bớt ai ở
ngày 16/11 (200 khách), và không thỏa điều kiện thứ hai, thứ ba. Với ghép tự do, khách trễ giảm từ 104 xuống 94 trong 807
nhưng 57 khách tới sớm quá một phút, so với 0 khi đệm đều. Bỏ hai biến trạng thái xe không đổi số khách trễ ở hai phép
so này.
% EV: G05–G10; N4B8-03 (check B: C-30 0/147, 0/197; U 2 → 1 trên 145, 1 → 1 trên 200; FREE 104 → 94 trên 807, sớm 57 so với 0; bỏ trạng thái xe 1/145 và 94/807); ML:6-7,14-15,21,45-53,67-69 (độ dài thêm +4,4 / +3,8 s; phân vị 90 +21,8 / +21,7 s); K40; R37 verify-numbers-v2.json:ml; R39 agent-numbers-v1.json:ml (xét mọi số nguyên, sai số 0,0005 điểm phần trăm); ml_pilot_run.py:37-42,175-176

% -----------------------------------------------------------------------------------------------------
\section{So với cách đệm lời hứa}
\label{muc:c4-dem}

Khi có dự báo, cộng thêm 30 giây vào mọi lời hứa của điều phối không trần cũng cho 2.328 lượt không trễ quá một phút
trên 3.672 yêu cầu, bằng luật C trần 30 giây có dự báo. Ở ngưỡng hai phút, cách đệm này cho 2.366 lượt, so với 2.338
của luật C. Hai cách này cho hai thứ khác nhau. Đệm kéo dài lời hứa, trung bình 37,2 thay vì 7,4 giây, nên ít khách bị tính là trễ. Nhưng đệm không đổi
quyết định điều phối: vẫn có 174 khách bị điều phối dời giờ trả, với tổng độ dời lớn nhất là 590,1 giây (gần 10 phút).
Luật C trần 30 giây giới hạn tổng độ dời do điều phối của mỗi khách ở 30 giây dưới giả thiết (N); trên 34 kịch bản của
tập kiểm tra chính, mức lớn nhất đo được là 29,2 giây, kể cả ở các kịch bản cuối tuần không có dự báo. Đổi lại, luật C
phục vụ ít hơn 96 lượt.
% EV: C10; N-16; N-18; K06; K51; C20; N-23; N-28 (cuối tuần 29,2 s); C02; A06 (dưới giả thiết (N); không bảo đảm giờ đến: A22, K05); RD §8; khớp câu trỏ của phần A v5-v3 mục 4.2 (C10); R37 verify-numbers-v2.json:padding; R39 agent-numbers-v1.json:padding,cohort_totals

Cách đệm là phát lại trên các lần chạy của điều phối không trần, nên hai bên được so trên hai nhóm khách khác nhau, chỉ
chung 2.207 khách. Bảng~\ref{bang:c4-dem} xếp các cách hứa thành ba khối có lời hứa dài thêm gần bằng nhau.
% EV: C10 (phát lại, không phải đối chứng ghép cặp); ADDENDUM-1 phần đầu (quy tắc đệm D = D0 + k); C04

\begin{table}[htbp]
\centering
\caption[Trần và cách đệm lời hứa khi có dự báo]{Các cách hứa khi có dự báo trên tập kiểm tra chính (34 kịch bản, 3.672 yêu cầu mỗi cấu hình; dự báo chỉ có ở
21 kịch bản ngày thường). U: điều phối không trần; C-$B$: luật C trần $B$ giây; trong ngoặc là số khách được phục vụ.
Số đếm gộp; trễ và tới sớm tính trong số khách được phục vụ, so với lời hứa của chính cách hứa; hai cột ``không trễ
quá'' tính trên 3.672 yêu cầu; độ dài thêm là trung bình gộp theo khách được phục vụ. Đệm không đổi quyết định điều
phối, nên hàng của U và hàng của C là các nhóm khách khác nhau.}
\label{bang:c4-dem}
\small
\begin{tabular}{@{}p{4.6cm}rrrrr@{}}
\toprule
\textbf{Cách hứa} &
\textbf{\begin{tabular}[b]{@{}r@{}}Trễ\\ quá 60 s\end{tabular}} &
\textbf{\begin{tabular}[b]{@{}r@{}}Không trễ\\ quá 60 s\end{tabular}} &
\textbf{\begin{tabular}[b]{@{}r@{}}Không trễ\\ quá 120 s\end{tabular}} &
\textbf{\begin{tabular}[b]{@{}r@{}}Tới sớm\\ quá 60 s\end{tabular}} &
\textbf{\begin{tabular}[b]{@{}r@{}}Dài thêm\\ (s)\end{tabular}}\\
\midrule
\multicolumn{6}{@{}l}{\textit{Không đệm (dài thêm khoảng 7 giây)}}\\
U, dự báo (2.434) & 137 & 2.297 & 2.347 & 8 & 7,18\\
C-30, dự báo (2.338) & 10 & 2.328 & 2.338 & 4 & 7,37\\
C-120, dự báo (2.387) & 65 & 2.322 & 2.379 & 2 & 7,27\\
% EV: N-16; N-03; C09; N-06; L1a-03; N4B-09 (sớm 8 / 4 / 2)
\midrule
\multicolumn{6}{@{}l}{\textit{Đệm 30 giây (dài thêm khoảng 37 giây)}}\\
U, dự báo, đệm 30 s (2.434) & 106 & 2.328 & 2.366 & 44 & 37,18\\
C-30, dự báo, đệm 30 s (2.338) & 2 & 2.336 & 2.338 & 44 & 37,37\\
% EV: C10; N-02; N-16
\midrule
\multicolumn{6}{@{}l}{\textit{Đệm 60 giây (dài thêm khoảng 67 giây)}}\\
U, dự báo, đệm 60 s (2.434) & 87 & 2.347 & 2.386 & 1.033 & 67,18\\
% EV: C11; N-16
\bottomrule
\end{tabular}
\end{table}

Cách đọc bảng: với đệm 30 giây trên điều phối không trần, 106 trong 2.434 khách tới muộn hơn lời hứa đã đệm quá 60
giây, nên có 2.328 lượt không trễ quá 60 giây trên 3.672 yêu cầu; 44 khách tới sớm quá 60 giây.
% EV: C10; N-16

\textbf{Vai trò của từng thành phần.} Thứ nhất, ở mức kéo dài lời hứa gần bằng nhau, luật C trần 30 giây hơn cách đệm
ở ngưỡng một phút và kém ở ngưỡng hai phút: khi cùng đệm 30 giây, hai bên có 2.336 so với 2.328 lượt và 2.338 so với
2.366 lượt. Ở hàng đệm 30 giây, điều phối không trần đạt cùng 2.328 lượt không trễ quá một phút như luật C có dự báo
không đệm, với lời hứa dài thêm trung bình 37,18 so với 7,37 giây, tức khoảng năm lần.
% EV: N-17; N-18 (37,1833 / 7,3671 = 5,05); N-16; R37 verify-numbers-v2.json:padding; R39 agent-numbers-v1.json:padding,cohort_totals

Thứ hai, lời hứa dài thì ít trễ nhưng nhiều khách tới sớm: 44 khách với cách đệm 30 giây và 1.033 khách với cách đệm 60
giây, so với 4 khách với luật C có dự báo. Số lượt tới trong khoảng $\pm$60 giây quanh lời hứa là 2.284 với cách đệm 30
giây, so với 2.324 với luật C, trên 3.672 yêu cầu.
% EV: C10; C11; N-04 (±60 s = phục vụ − trễ60 − sớm60, ADD1:92); N4B-09 (ADD1:84-89)

Thứ ba, cách đệm 30 giây vẫn có 68 khách trễ quá hai phút, 11 khách quá năm phút và 1 khách quá mười phút; luật C có
dự báo không có ai trễ quá hai phút. Cách đệm hơn 28 lượt ở ngưỡng hai phút vì phục vụ nhiều hơn 96 lượt, không phải vì
ít khách trễ hơn.
% EV: N-01 (68 / 11 / 1 so với 0 / 0 / 0; ADD1:15,19); N-19 (2.366 − 2.338 = 28 = 96 − 68)

Khi phát lại cùng lời hứa theo dự báo, luật hạn chót 30 giây cũng có 2.338 lượt phục vụ và 2.328 lượt không trễ quá
một phút. Vì phép phát lại không đổi các lần chạy, nghĩa vụ của luật này vẫn tính từ lời hứa đầu không dự báo. Theo
nghĩa vụ đó, tính trung bình theo kịch bản, nó bị ép ở 41,3\% lượt quyết định; tính gộp là 7.710 trong 20.103 lượt.
Luật hạn chót có khách tới muộn hơn lời hứa đầu không dự báo quá 30 giây ở 28 trong 34 kịch bản, còn luật C không có
quyết định bị ép hay khách vượt trần (Bảng~\ref{bang:c4-nghiavu}).
% EV: N-16 (M⁺-30 + dự báo 2.338 / 10 / 2.328 / 2.338, ADD1:29); C13 (41,3% trung bình theo kịch bản); N-21 (phá hứa của M⁺-30 = A − p₀ > 30 s với lời hứa không dự báo; C-30 0/34, 0 quyết định bị ép); X3-16; R36-A01 (bảng A-v5-v3 dùng số gộp); R39 agent-numbers-v1.json:main_original_obligation

Trần 120 giây có dự báo, với lời hứa chỉ dài thêm 7,27 giây, cho 2.379 lượt không trễ quá hai phút, nhiều hơn 2.366
của cách đệm 30 giây, nhưng 2.322 lượt ở ngưỡng một phút, ít hơn 2.328. Đây là một điểm đo thêm trên một tập khách khác,
không phải mức trần tốt nhất.
% EV: N-06; N-16; N-18; N-10; K36

\textbf{Khi nào chọn trần, khi nào chọn đệm.} Nếu hãng xe chỉ tính số lượt không trễ quá hai phút và chấp nhận lời hứa
dài hơn cùng nhiều khách tới sớm hơn, điều phối không trần có dự báo và đệm 30 giây đạt 2.366 lượt, so với 2.338 lượt
của luật C trần 30 giây có dự báo, trên 3.672 yêu cầu; điều phối không trần cũng phục vụ thêm 96 lượt. Nếu cần giới hạn
tổng độ dời do điều phối của từng khách, kể cả khi không có dự báo, luật C trần 30 giây làm được điều này dưới giả thiết
(N), đổi lại 96 lượt phục vụ và 242 khách đi chung. Hai cách dùng chung được; chọn cách nào là quyết định của hãng xe,
không phải kết luận của thực nghiệm.
% EV: N-19; C02; N-20; N-02 (hàng C-30 có dự báo và đệm 30 s); D20; I07; I10; K38; K51; A22 (giả thiết N); R37 verify-numbers-v2.json:padding; R39 agent-numbers-v1.json:padding,cohort_totals

% -----------------------------------------------------------------------------------------------------
\section{Khi dự báo không đủ: chiều tuyết 15/11/2018}
\label{muc:c4-tuyet}

Trên 601 khách được phục vụ của 8 kịch bản chiều tuyết 15/11/2018 (luật C trần 30 giây), lời hứa theo dự báo chỉ giảm
số khách trễ quá một phút từ 41 xuống 34. Tính trung bình theo kịch bản, tỉ lệ này giảm từ 7,19\% xuống 5,94\%. Theo
sổ, trong 3.445,9 giây trễ của 41 khách, 8,6 giây tính cho điều phối và 251,5 giây được hồ sơ lịch sử giải thích
(Hình~\ref{hinh:c4-so}).
% EV: F13; K25; B20; N4B-05 (NEW:5)

Bảng~\ref{bang:c4-tuyet} cho kết quả phát lại vài cách hứa dùng thêm hệ số tắc đường của chính ngày tuyết. Cách ``trộn'' lấy trung
bình nhân của mức tăng theo hồ sơ lịch sử và mức tăng của hệ số cùng ngày từ khung 15 phút trước sang khung hiện tại.
% EV: F14; K20; K21; N4B-06 (factor_nowcast_screen.py:14-20 quy tắc thông tin, 28-29: MOM = (f(k0)/f(k0−1))^(b−b0), BLEND = √(HIST·MOM); RP1:7-20; NEW:5)

\begin{table}[htbp]
\centering
\caption[Phát lại các cách hứa trong chiều tuyết]{Phát lại các cách hứa trên 8 kịch bản chiều tuyết 15/11/2018 (luật C trần 30 giây, 601 khách được phục vụ).
Cột ``khách trễ'' đếm gộp số khách tới muộn hơn lời hứa đầu quá 60 giây; cột cuối là độ dài thêm trung bình theo kịch
bản. Không cách hứa nào trong bảng có khách tới sớm quá 60 giây. Mọi cách hứa đều tính lại sau khi chạy, trên một ngày
đã được xem trước khi thiết kế các cách này; hai hàng ``có buổi sáng'' chỉ để tham khảo.}
\label{bang:c4-tuyet}
\small
\begin{tabular}{@{}p{3.4cm}p{6.6cm}rr@{}}
\toprule
\textbf{Cách hứa} & \textbf{Thông tin dùng lúc hứa} &
\textbf{\begin{tabular}[b]{@{}r@{}}Khách\\ trễ\end{tabular}} &
\textbf{\begin{tabular}[b]{@{}r@{}}Dài thêm\\ (s)\end{tabular}}\\
\midrule
Không dự báo & Thời gian đi hiện tại & 41 & 0\\
Theo dự báo & Hồ sơ tắc đường của các ngày nguồn & 34 & 2,0\\
Trộn trong khung & Thêm hệ số cùng ngày đã có trong khung giờ & 21 & 8,5\\
Trộn có buổi sáng & Thêm hệ số cùng ngày từ 6 giờ; thiết kế thử không cho khung chiều dùng dữ liệu này & 19 & 9,5\\
Trộn có buổi sáng, trễ một khung & Như hàng trên, nhưng mức tăng gần nhất chỉ lấy từ các khung đã kết thúc; hệ số của
khung hiện tại vẫn dùng giá trị đúng & 29 & 9,0\\
Đệm tỉ lệ & Tỉ lệ chọn để tổng phần kéo dài bằng của hàng ``trộn trong khung'', trên chính ngày tuyết & 27 & 8,3\\
Hệ số thật & Hệ số sau lúc hứa, không có ở thời điểm hứa & 0 & 18,8\\
\bottomrule
\end{tabular}
\end{table}
% EV: F14 (RP1:7-20; NEW:5); N4B-06 (định nghĩa; cột dài thêm là trung bình theo kịch bản, RP1:6, rider_nowcast_replay.py:96-106); K20; K21; factor_nowcast_screen.py:17-20 (hàng trễ một khung cũng dùng thông tin từ 06:00); N4B8-09 (RP1:7-20: E60 0,00% ở mọi hàng ⇒ 0 khách tới sớm ở mọi kịch bản [A])

Cách đọc bảng: khi trộn hồ sơ lịch sử với xu hướng trong khung giờ, 21 trong 601 khách tới muộn hơn lời hứa quá một
phút, với lời hứa dài thêm trung bình 8,5 giây.
% EV: F14; N4B-06

Phép phát lại này tách số giây trễ theo cách cộng có giới hạn, khác phần có dấu của công thức~\eqref{eq:so}. Phần kéo
dài của mỗi khách được chặn trong khoảng từ 0 đến phần trễ còn lại sau khi trừ phần tính cho điều phối. Theo cách cộng
này, cách trộn
có buổi sáng giải thích thêm 886,9 giây, tức 25,7\% của 3.445,9 giây trễ; 2.299,0 giây (66,7\%) chưa được giải thích
bằng các cách đã thử. Theo số liệu trạm Central Park của Cơ quan Quản lý Khí quyển và Đại dương Quốc gia Hoa Kỳ
(National Oceanic and Atmospheric Administration, NOAA), từ 9/2018 đến 12/2019 chỉ có 5 đợt tuyết từ 25,4 mm trở lên,
kể cả đợt này; muốn kiểm cách hứa cho ngày bất thường cần thêm nhiều đợt khác.
% EV: F15 (NEW:5; RP1:24-35); N4B8-07 (rider_nowcast_replay.py:14-15, 113: explained = Σ clip(D − P0, 0, A − P0 − dispatch); RP1:29: phần vượt 6 s; với hồ sơ lịch sử phần vượt 0 nên hai cách cộng trùng nhau, N4B-05); K19; I11; B26 (NEW:9; WX:1-3: tệp NOAA GHCN-D trạm Central Park; WX17:1-4; ghép ngày cách nhau ≤ 2)

% -----------------------------------------------------------------------------------------------------
\section{Kiểm tra tính hợp lệ của kết quả}
\label{muc:c4-hople}

Trên 34 kịch bản của tập kiểm tra chính, không khách nào vượt trần: tổng dời lớn nhất là 29,2 giây với trần 30 giây và
118,9 giây với trần 120 giây. Không lượt quyết định nào bị ép, 0 trong 20.101 và 0 trong 20.435 lượt. Đây là phép kiểm cài
đặt: với sổ hợp lệ $0 \le L_r \le B$ ở lượt đang xét và giả thiết (N), Nhận xét~\ref{nx:giunguyen} đã cho trước các
số 0 này; kết luận xuyên chuyến cần giả thiết (N) đúng ở mọi lượt. Điều phối không trần cũng có 0 trong 20.760 lượt.
% EV: C20; N-07; X3-16 (0/20.101; 0/20.435; U 0/20.760); N-21; A22; chương 3, nx:giunguyen và đoạn sau nó (0 ≤ L_r ≤ B ở lượt đang xét; (N) đúng ở mọi lượt cho kết luận xuyên chuyến); R39 agent-numbers-v1.json:cohort_totals

Bộ đếm của bộ sinh được đọc ở điều phối không trần, luật C trần 30, 120 giây và luật hạn chót 30 giây. Ở cả bốn cấu
hình, không phương án nào bị bỏ vì giới hạn 100 phương án mỗi xe hay giới hạn số bước tìm cách chèn. Giới hạn 2 yêu cầu
mới mỗi lượt chỉ để một yêu cầu sang lượt sau ở mỗi cấu hình. Các số này không đo mọi giới hạn tính toán (mục~\ref{muc:gioihan}): giới hạn
của bước kiểm, của bộ giải, một bộ đếm khác của bộ sinh và các bộ đếm ở tập tắc đường cố định chưa được đọc.
% EV: N-25 (ADD1:119-122: 34 lần chạy mỗi cấu hình); N4B-11 (GEN:54-95, 236-246, 298-345; CandidateModels.cs:68-82, 100-118: điều kiện đầy đủ còn gồm OmittedRepairRequestCount, chưa đọc); X3-17; X3-03; X3-32; C4A-18 (tập tắc đường cố định: giới hạn chưa kiểm); K34 thu hẹp (N4B8-13)

Như mục~\ref{muc:doichung} đã nêu, mức dùng trần lớn nhất của tiêu chí 3 bằng 100\% ở mọi lượt có ghi tiêu chí này:
3.503, 3.506 và 3.504 lượt khi không trần, với trần 30 và 120 giây.
% EV: A15; N-25 (ADD1:119-121); A07; X3-07 (chương 3 mục doichung nêu hệ quả: tiêu chí 3 chỉ tác dụng qua tổng phần dời)

Các số đếm chính của tập kiểm tra chính được ít nhất hai chương trình đọc dữ liệu viết riêng tính lại và khớp nhau
(mục~\ref{muc:kiemtra}). Phép so gồm cả từng khách ở 7.159 khách được phục vụ của điều phối không trần và luật C trần
30, 120 giây; riêng số khách đi chung chỉ được một quy tắc đếm. Theo bản ghi khởi tạo, 8 lần chạy chiều tuyết dùng cùng
bản bộ chạy đã lưu và cùng cấu hình với các lần chạy luật C trần 30 giây của tập kiểm tra chính, kể cả luật phục hồi
không tệ hơn.
% EV: N4B-14 (ADD1:66-78); K41; C19; N-09 (đi chung: Codex không đếm lại); N4B8-05 (check C: 8 lần chạy chiều tuyết và 51 lần chạy C-30 của T3.8 cùng binarySha256 a631038b…, coreCommitSha 1ce15e92…, policyConfigurationHash a2f11a3a…; sha256 runner-tier3-v2/RideBound.Runner.dll khớp; wp4-tier3-v2.json:7 no-worse-than-reference-v1; R26 Q8)

Phép kiểm số đếm học máy chỉ đối chiếu số học của bảng kết quả; nó không kiểm lại quá trình học hay dự đoán của mô hình.
% EV: R37 §1.8; R39 agent-numbers-v1.py phần C; agent-numbers-v1.json:ml; ml_pilot_run.py:37-42,175-176

Chạy lại trên đúng mã nguồn đã dựng bộ chạy, cả 1.025 trong 1.025 kiểm thử tự động đạt trong một lần chạy đầy đủ ngày
13/10/2026, với kiểu xuống dòng LF và một tệp chưa ghi nhận trong cây làm việc (Phụ lục~\ref{phuluc:C}). Các phép kiểm
trên không cho biết mô phỏng giống thực tế đến đâu.
% EV: TESTRUN-45f8026-2026-10-12.md:18 (attempt 4: LF và tệp chưa ghi nhận, 1.025/1.025 trong một lần chạy ngày 13/10/2026); C:/RideBoundData/research/testrun-45f8026-eollf-dirty-fullrun-2026-10-13.log; R39 agent-numbers-v1.json:test_attempt4; I13; A28; I01; K41

% -----------------------------------------------------------------------------------------------------
\section{Thảo luận và hạn chế}
\label{muc:c4-thaoluan}

Với một hãng xe, các kết quả trên là các điểm trên một đường đánh đổi, không phải một khuyến nghị. So với điều phối
không trần, luật C trần 30 giây phục vụ ít hơn 96 lượt và có ít hơn 242 khách đi chung; đổi lại, phần dời do điều phối
của mỗi khách không quá 30 giây và số lượt trễ thấp hơn ở cả bốn ngưỡng đã đo, trên hai nhóm khách khác nhau
(Bảng~\ref{bang:c4-phanbotre}). Khóa luận không đo giá trị của các điều này với khách hay hãng xe.
% EV: C02; N-20; C20; N-24; C04; D20; I10; K38 (không kể 0 quyết định bị ép là phần luật C đổi được: X3-16, U cũng 0/20.760)

Một phép tính hòa vốn giả định cho thấy cách đọc các điểm này. Gọi $\lambda$ là chi phí của hãng xe cho mỗi lượt trễ quá
một phút và $\mu$ là biên lợi nhuận mỗi chuyến. Với 96 lượt phục vụ ít hơn và 127 lượt trễ ít hơn, luật C chỉ có lợi hơn
khi $\lambda > 0{,}756\,\mu$, tức $96/127$. Phép tính bỏ qua mọi yếu tố khác và so hai nhóm khách khác nhau; khóa luận
không có giá trị thật nào của $\lambda$ hay $\mu$.
% EV: C29 (R15:231-235; 96/127 = 0,756); K35; C28; C04; I10

Các hạn chế dưới đây chia theo bốn nhóm; mỗi ý nêu điều chưa làm và kết luận chịu ảnh hưởng.

\textbf{Mô hình mô phỏng.}
\begin{itemize}[leftmargin=1.2cm,itemsep=2pt,topsep=3pt]
\item Mọi kết quả đến từ FleetPy với một hệ số tắc đường chung cho cả thành phố, mức nền lấy lúc 00:00, nên mô phỏng gần
như không có yếu tố ngẫu nhiên. Các kết quả, kể cả mức trễ còn lại sau dự báo, không cho biết hiệu quả khi vận hành thật.
\item Thời gian dừng đón, trả bằng 0 giây (kiểm trên 3 lần chạy mẫu), so với 30 giây mỗi khách trong nghiên cứu điển hình
của FleetPy~\cite{engelhardt2026}, và đồ thị rút gọn chỉ có 96 nút. Phần dời do điều phối và cái giá 96 lượt vì vậy chỉ áp cho
thiết lập này.
\item Không có phản ứng của khách và không gán lại xe cho khách đã nhận, nên số lượt phục vụ không tính khách bỏ đi, và
chưa biết cái giá của trần khi được gán lại xe.
\end{itemize}
% EV: I01; B01; B06; B07; I02; F10; B05 (SERV:2-4: kiểm trên 3 lần chạy mẫu; FLEET:423-424); J18; B03; K48; B12; I07; X3-39

\textbf{Dữ liệu và thiết kế thực nghiệm.}
\begin{itemize}[leftmargin=1.2cm,itemsep=2pt,topsep=3pt]
\item 34 kịch bản của tập kiểm tra chính phụ thuộc nhau, thuộc 5 ngày, mỗi kịch bản chạy một lần; các phép đếm theo
kịch bản là mô tả, không phải kiểm định.
\item Chưa kiểm độ nhạy theo hạt giống của bộ giải, nên chưa biết các chênh lệch có giữ nguyên khi đổi hạt giống.
% EV: R37 §1.8/L-26; B09; B11; chương 3 v4-v3 muc:gioihan (hạt giống cố định 7, một lần chạy mỗi kịch bản)
\item Các kết quả dùng lại dữ liệu đã xem. Ngưỡng phân tích và cách chia tập được chọn sau khi xem tập buổi sáng và các
đợt chạy thử trước đó. Mức 120 giây được chọn sau khi xem kết quả của điều phối không trần và luật C trần 30 giây trên
chính tập kiểm tra chính, kể cả các phép tính thử ở 120 giây (mục~\ref{muc:c4-thamdo}). Chiều tuyết và các ngày kiểm học máy cũng đã được xem, nên mọi kết quả là thăm dò.
\item Tải thấp: 0,535 khách trên xe kể cả lúc xe rỗi khi không trần, so với khoảng 1,33 trong nghiên cứu điển hình của
FleetPy. Ở 49 kịch bản dựng từ dữ liệu thật, chuyến được chọn ngắn hơn mọi chuyến đủ điều kiện (trung vị thời gian đi thẳng 165
so với 196 giây). Kết quả không suy ra cho tải cao hay chuyến dài hơn.
\end{itemize}
% EV: B09; I05; B11; I06; C4A-11 (mức 120 giây, như phần A v5-v3 mục c4-thamdo); MLS:9-10; I11; B12 (LOAD:2-3; FLEET:447-449); J18; B13 (SEL6:2-4: 49/49 kịch bản); K42 (kể cả các phép tính thử ở 120 giây); chương 3 v4-v3 muc:gioihan (hạt giống7; chưa có kiểm độ nhạy)

\textbf{Tích hợp.}
\begin{itemize}[leftmargin=1.2cm,itemsep=2pt,topsep=3pt]
\item Bộ điều phối vẫn kiểm giờ đón bằng giờ không dự báo, nên 34 trong 1.471 khách được báo giờ đón muộn hơn giờ đón
muộn nhất. Các số có dự báo ở mục~\ref{muc:c4-dubao} và~\ref{muc:c4-dem} chưa tính việc sửa phép kiểm này.
\item Bản chạy ưu tiên chi phí khác bản lưu của ưu tiên ổn định. Bốn cặp đối chiếu hành vi và việc tái lập các
chỉ số nền không chứng minh tương đương toàn bộ hai bản. Kết quả ở mục~\ref{muc:c4-vaitro} vì vậy vẫn là phép so
thăm dò theo cấu hình, không loại trừ mọi khác biệt thời gian chạy.
\end{itemize}
% EV: F07 (RX:2; FR:70); I13; A28; R18 B24 (qua REG A28)

\textbf{Phạm vi kết luận.}
\begin{itemize}[leftmargin=1.2cm,itemsep=2pt,topsep=3pt]
\item Phép so hai thứ tự chọn phương án đã có kết quả ở trần 30 giây và không trần, nhưng chưa tạo thành
đường đánh đổi nhiều mức cho ưu tiên chi phí. Luật giới hạn từng lần chưa được cài đặt; các bản ghi chỉ vượt do
cộng dồn và bảy lịch sử sàng lọc không thay cho phép so hai luật trên diễn biến mô phỏng đầy đủ.
\item Luật hạn chót chỉ mượn thành phần hạn chót của Schulz và Pfeiffer~\cite{schulz2026}, với 30 và 60 giây trên tập
kiểm tra chính và 30--300 giây trên tập buổi sáng; mức mặc định của bài đó là 900 giây. Hạn chót nới kiểu Hörl và
Zwick~\cite{horl2022} (mục~\ref{muc:doichung}) trùng với cấu hình M$^{+}$ có luật phục hồi ``không tệ hơn'' ở trên,
nên mới được chạy ở 30 và 60 giây trên tập kiểm tra chính. Các phép so không nói về phương pháp của các tác giả đó.
\item Mức 30 giây không phải ngưỡng cảm nhận hay chuẩn ngành, và không có dữ liệu về mức khách sẵn lòng trả cho lời
hứa ổn định hơn; vì vậy không mức trần nào được chọn là tốt nhất.
\end{itemize}
% EV: H02; H12; K18; E09; K08; J10; J11; X3-28 (luật hạn chót 30–300 s ở tập buổi sáng; chương 3 Bảng luat); K27; D20; K52; K36
; không có \chapter mới.
% Mỗi dòng "% EV:" trỏ tới EVIDENCE-REGISTER.md (A01…K52), EVIDENCE-REGISTER-ADDENDUM-1.md (N-xx, L1a-xx, mục 0.x),
% EVIDENCE-REGISTER-ADDENDUM-2-CH3.md (X3-xx), EVIDENCE-REGISTER-ADDENDUM-3-CH4B.md (N4B-xx),
% EVIDENCE-REGISTER-ADDENDUM-4-CH4A.md (C4A-xx) hoặc EVIDENCE-REGISTER-ADDENDUM-6-CH4B-v8.md (N4B8-xx; script chỉ đọc
% evidence/checks-ch4b-v8/ch4b_v8_check.py, output ch4b-v8-check-2026-10-12.txt, phần A/B/C gọi là "check A/B/C").
% Không dùng gói mới: TikZ, enumitem và graphicx đã nạp trong main.tex; thu bảng/hình bằng resizebox.
% =====================================================================================================

% -----------------------------------------------------------------------------------------------------
\section{Vai trò của trần và thứ tự tiêu chí của bộ chọn}
\label{muc:c4-vaitro}

\subsection{Vì sao cần đổi thứ tự ưu tiên?}

Các kết quả trước mục này dùng bộ chọn ưu tiên ổn định: sau khi xét số yêu cầu nhận được, bộ chọn so mức dùng
trần và phần dời giờ rồi mới xét chi phí kế hoạch. Như vậy, cả việc xếp hạng phương án lẫn trần cứng đều có thể
góp phần hạn chế thay đổi lời hứa. Chỉ so có trần với không trần dưới một thứ tự chưa cho biết kết quả sẽ ra sao
khi bộ điều phối đặt chi phí kế hoạch lên trước phần dời giờ.

Để kiểm tra điều đó, khóa luận giữ ưu tiên nhận yêu cầu và các phép kiểm khả thi, nhưng cho bộ chọn xét chi phí
kế hoạch ngay sau đó. Các mức xếp hạng theo ổn định được bỏ, còn tiêu chí phân xử theo mã định danh giữ nguyên.
Trần cứng của luật C vẫn được kiểm trước khi chọn nếu cấu hình có trần. Hai cách gọi dưới đây là \emph{ưu tiên
ổn định} và \emph{ưu tiên chi phí kế hoạch}; chúng biểu thị mục tiêu khác nhau, không phải hai mức chất lượng
mặc định. Chi phí kế hoạch là tổng thời gian còn lại của tuyến dự kiến, chưa phải tiền hay chi phí vận hành thực.
% EV: S15-03; A14; X3-07; K15. Cùng ưu tiên nhận khách trong từng lượt, không bảo đảm cùng số phục vụ cuối.

Phép so gồm bốn cấu hình trên cùng 34 kịch bản và 3.672 yêu cầu mỗi cấu hình: hai thứ tự ưu tiên, mỗi thứ tự có
trần 30 giây và không trần. Hai cấu hình ưu tiên chi phí có 68 lượt mô phỏng hoàn tất, được đối chiếu với 68 lượt
ưu tiên ổn định đã lưu. Tập kịch bản và tập mã yêu cầu khớp ở cả 34 kịch bản. Các diễn biến sau khi chọn có thể
khác nhau, nên nhóm khách được phục vụ cũng khác nhau.
% EV: S15-01 (mục 0--3), main-final-admission.json; S15-02. Không có 34 lần lặp độc lập.

\subsection{Đánh đổi giữa phục vụ, ghép khách và sửa lời hứa}

\begin{table}[htbp]
\centering
\caption[Đánh đổi dưới hai thứ tự chọn phương án]{Hai thứ tự chọn phương án trên tập kiểm tra chính, mỗi cấu hình
có 3.672 yêu cầu. Trễ được so với lời hứa trả đầu chưa tính dự báo, trừ dòng ghi riêng. Khách có đi chung được
đếm một lần nếu hoàn thành chuyến và có khoảng thời gian cùng xe với khách khác.}
\label{bang:c4-doi-uutien}
\small
\setlength{\tabcolsep}{3pt}
\begin{tabular}{@{}>{\raggedright\arraybackslash}p{6.5cm}rrrr@{}}
\toprule
& \multicolumn{2}{c}{\textbf{Ưu tiên ổn định}} & \multicolumn{2}{c}{\textbf{Ưu tiên chi phí}}\\
\cmidrule(lr){2-3}\cmidrule(l){4-5}
\textbf{Chỉ số} & \textbf{Không trần} & \textbf{Trần 30} & \textbf{Không trần} & \textbf{Trần 30}\\
\midrule
Lượt phục vụ & 2.434 & 2.338 & 2.502 & 2.364\\
Muộn quá một phút & 203 & 76 & 379 & 72\\
Không muộn quá một phút & 2.231 & 2.262 & 2.123 & 2.292\\
Không muộn quá hai phút & 2.332 & 2.327 & 2.280 & 2.353\\
Không muộn quá năm phút & 2.416 & 2.338 & 2.466 & 2.364\\
Khách hoàn thành có đi chung & 555 & 313 & 1.027 & 432\\
Khách bị sửa giờ trả do điều phối & 174 & 33 & 434 & 106\\
Trong đó, bị sửa từ hai lần & 12 & 0 & 50 & 2\\
Đón muộn theo giới hạn dịch vụ & 51 & 43 & 39 & 44\\
Muộn quá một phút khi lời hứa có dự báo & 137 & 10 & 318 & 11\\
Tổng dời lớn nhất của một khách (s) & 590,1 & 29,2 & 709,1 & 29,7\\
\bottomrule
\end{tabular}
\end{table}
% EV: S15-01 (mục 2, 5a--5f); S15-02 (pooled); S15-04 (kiểm bảng và các chênh lệch từ TSV).

Khi cùng giữ trần 30 giây, ưu tiên chi phí phục vụ thêm 26 lượt và có thêm 119 khách hoàn thành có đi chung
so với ưu tiên ổn định (Bảng~\ref{bang:c4-doi-uutien}). Số lượt không muộn quá một phút tăng 30. Đổi lại, số khách
bị sửa giờ trả tăng từ 33 trong 2.338 lên 106 trong 2.364, tương ứng 1,41\% và 4,48\%; số khách bị sửa từ hai lần
tăng từ 0 lên 2. Số lượt đón muộn tăng từ 43 lên 44. Như vậy, cách ưu tiên ổn định vẫn cho ít khách bị sửa lời
hứa hơn, còn ưu tiên chi phí nhận và ghép được nhiều khách hơn trong thiết lập này.

Đúng giờ đến và ít sửa lời hứa là hai chỉ số khác nhau. Ưu tiên chi phí có trần có 72 lượt muộn quá một phút,
ít hơn 4 lượt so với ưu tiên ổn định có trần, nhưng nhiều khách bị sửa giờ hơn. Khi dùng lời hứa có dự báo, số
lượt muộn quá một phút là 11 và 10 trên hai nhóm khách tương ứng. Các chênh lệch nhỏ này không đủ để chọn một
cách làm tốt hơn về mọi mặt. Trung vị độ dài lời hứa đầu chưa dự báo là 809,642 giây với ưu tiên chi phí và
815,060 giây với ưu tiên ổn định, trên hai tập khách khác nhau.
% EV: S15-01 (5b--5e); S15-02. Độ dài dùng p0-t0; không quy thành chờ/đi vòng của từng khách.

Không trần cho thấy rõ cái giá của đổi thứ tự. Ưu tiên chi phí phục vụ thêm 68 lượt và có thêm 472 khách hoàn
thành có đi chung so với ưu tiên ổn định, nhưng có thêm 176 lượt muộn quá một phút và thêm 260 khách bị sửa giờ
trả. Trần vì vậy vẫn có vai trò khi bộ chọn chuyển sang ưu tiên chi phí. Trong thứ tự này, thêm trần 30 giây làm
giảm 138 lượt phục vụ và 595 khách hoàn thành có đi chung; đồng thời có ít hơn 307 lượt muộn quá một phút và
thêm 169 lượt được phục vụ, không muộn quá một phút.
% EV: S15-01 (mục 2, 5f--5g); S15-04. Các chênh lệch là tổng trên tập yêu cầu, không cùng tập khách phục vụ.

Phần giảm số lượt trễ phải đọc cùng phần không phục vụ. Với ưu tiên ổn định, 127 lượt trễ ít hơn đi cùng 96 lượt
phục vụ ít hơn và 31 lượt không muộn quá một phút nhiều hơn. Với ưu tiên chi phí, các số tương ứng là 307, 138
và 169. Chênh lệch giữa hai mức giảm trễ là 180, nhưng không có nghĩa 180 khách cụ thể được bảo vệ chỉ nhờ đổi
bộ chọn. Chênh lệch này cùng chiều ở 29 trong 34 kịch bản và cả năm ngày; các kịch bản phụ thuộc nhau nên đây
chỉ là mô tả trên dữ liệu đã dùng.
% EV: S15-01 (mục 2, 4); S15-02 (interactions); 180 = (138-96) + (169-31).

Ở ưu tiên chi phí có trần, tổng dời lớn nhất do điều phối là 29,657 giây và không khách nào vượt 30 giây.
Kết quả cho thấy cấu hình có trần giữ được giới hạn đã đặt khi thay thứ tự xếp hạng. Nó chưa chứng minh rằng
riêng trần đủ bảo đảm giờ đến thực tế, vì khóa, ràng buộc dịch vụ, bộ sinh và luật phục hồi vẫn cùng tham gia.

Mức trần thích hợp còn phụ thuộc tiêu chí dịch vụ. Với ưu tiên chi phí, trần 30 giây tăng số lượt được phục vụ
và không muộn quá một phút từ 2.123 lên 2.292; nếu chấp nhận muộn tới năm phút thì số lượt đáp ứng lại giảm
từ 2.466 xuống 2.364. Đây là hai cách đọc cùng kết quả đã có, không phải khảo sát mức khách chấp nhận chờ.
Trần 30 giây là điểm thực nghiệm, không được chọn làm ngưỡng tối ưu hay ngưỡng cảm nhận của hành khách.
% EV: S15-01 (5a); S15-02 (timely60, timely300); K36; K52.

\subsection{Lời hứa có dự báo và giới hạn của phép so}

Với ưu tiên chi phí có trần, tính lại lời hứa đầu theo dự báo cho 11 lượt muộn quá một phút thay vì 72,
trên cùng 2.364 hành trình; lời hứa dài thêm trung bình 7,30 giây. Với cấu hình không trần, các số lượt muộn
là 318 thay vì 379, và lời hứa dài thêm trung bình 7,02 giây. Giờ trả thực tế và việc xe chở ai không đổi
trong phép tính này. Dự báo điều chỉnh phần đường dự kiến chậm đi, còn luật C giới hạn phần dời do điều phối;
hai thành phần không thay thế nhau và không tạo thành bảo đảm vô điều kiện về giờ đến.
% EV: S15-01 (5b); F01--F04. 21 ô bật causal-no15; 13 ô còn lại giữ p0.

Các lượt ưu tiên chi phí dùng bản chạy riêng. Phép đối chiếu hành vi trên bốn cặp mẫu và việc tính lại các
chỉ số ưu tiên ổn định giúp kiểm tra khả năng so sánh, nhưng không chứng minh mọi khác biệt giữa hai bản
chạy đã được loại bỏ. Các con số là phép so thăm dò theo cấu hình trên một bộ mô phỏng và năm ngày đã xem,
không phải ước lượng tác động nhân quả riêng của một tiêu chí. Phép sàng lọc trên các phương án đã ghi của
tập tắc đường cố định được giữ ở Phụ lục~\ref{phuluc:D}; nó bổ trợ việc hiểu bộ chọn, không thay kết quả
của những lần mô phỏng đầy đủ ở đây.
% EV: STAGE1B-PLAN.md mục 1.1 và amendment mục 3; S15-01 (mục 0, 3); G0 chỉ behavioural-only.


\section{Lời hứa theo dự báo}
\label{muc:c4-dubao}

Trên 21 kịch bản ngày thường có dự báo, số khách của luật C trần 30 giây tới muộn hơn lời hứa đầu quá một phút giảm từ
68 xuống 2 trong 1.471 khách được phục vụ (Bảng~\ref{bang:c4-dubao}). Lời hứa theo dự báo (mục~\ref{muc:dubao}) được
tính lại trên các lần chạy đã ghi; mỗi cách hứa được so với lời hứa đầu của chính nó.
% EV: F17; F04; N4B8-01 (check A: 68 → 2 trên 1.471); F03 (chỉ vế phát lại); FFA:8-13

Đối chứng thứ hai là đệm tỉ lệ: mỗi lời hứa đầu được kéo dài một tỉ lệ cố định của khoảng từ lúc hứa tới giờ trả đã
hứa. Tỉ lệ được chọn theo từng khung giờ, trên chính các khách này, để tổng phần kéo dài bằng của lời hứa theo dự báo.
% EV: N4B-03 (FFA:15-17: lời hứa đệm = P0 + k_w(P0 − t0), k_w = Σ(D0 − P0)/Σ(P0 − t0) trên mọi khách C-30 của khung w ở 21 kịch bản, tính từ lời hứa; FA:2)

\begin{table}[htbp]
\centering
\caption[Lời hứa theo dự báo và đệm tỉ lệ trên ngày thường]{Khách của luật C trần 30 giây trên 21 kịch bản ngày thường có dự báo, theo ba cách hứa: lời hứa đầu không dự
báo, đệm tỉ lệ và lời hứa đầu theo dự báo; mỗi cách so với lời hứa đầu của chính nó. Khối trên là số đếm gộp theo khách;
dòng cuối là trung bình theo kịch bản của tỉ lệ trễ, mỗi kịch bản tính như nhau.}
\label{bang:c4-dubao}
\small
\begin{tabular}{@{}lrrrr@{}}
\toprule
\textbf{Nhóm kịch bản} & \textbf{Phục vụ} & \textbf{Không dự báo} & \textbf{Đệm tỉ lệ} & \textbf{Dự báo}\\
\midrule
\multicolumn{5}{@{}l}{\textit{Số khách tới muộn hơn lời hứa đầu quá 60 s}}\\
Cả 21 kịch bản & 1.471 & 68 & 22 & 2\\
Khung 7 giờ (8 kịch bản) & 444 & 60 & 17 & 1\\
Khung 8 giờ (8 kịch bản) & 609 & 8 & 5 & 1\\
Khung 17 giờ (5 kịch bản) & 418 & 0 & 0 & 0\\
\multicolumn{5}{@{}l}{\textit{Số khách tới sớm hơn lời hứa đầu quá 60 s}}\\
Cả 21 kịch bản & 1.471 & 0 & 0 & 1\\
% EV: N4B8-01 (check A: CELLS, share × served từng kịch bản, phần dư 0); N4B8-08 (FFA:139-141 tính E_pad; sửa N4B-03: đệm tỉ lệ có đo tới sớm, bằng 0); F17
\midrule
\multicolumn{2}{@{}l}{Tỉ lệ trễ quá 60 s, trung bình theo kịch bản (\%)} & 6,01 & 1,94 & 0,17\\
% EV: F04 (FA:8); N4B-03 (FA:8: pad 0,0194)
\bottomrule
\end{tabular}
\end{table}

Cách đọc bảng: ở khung 7 giờ, 60 trong 444 khách trễ quá một phút khi không dự báo, 17 khi đệm tỉ lệ và 1 khi có dự
báo.
% EV: N4B8-01

Tính trung bình theo kịch bản, tỉ lệ trễ giảm 5,84 điểm phần trăm. Gom 21 kịch bản thành 8 nhóm theo ngày và khung giờ,
khoảng tin cậy 95\% theo phân phối Student là từ 1,15 đến 10,54 điểm phần trăm. So với đệm tỉ lệ, chênh lệch là
$-1{,}77$ điểm phần trăm, khoảng từ $-3{,}62$ đến $+0{,}07$, tức có chứa 0: chưa thể nói dự báo tốt hơn cách đệm này.
Các khoảng chỉ dựa trên ba buổi sáng đã được xem trước, và mọi hồ sơ dự báo đều có ngày 12/11 trong các ngày nguồn.
% EV: F04; F05 (FA:33,36: dY −5,844 [−10,542; −1,146]; dPad −1,775 [−3,620; +0,071]; FFA cluster_t: tầng = khung giờ, cụm = ngày–khung, df = 8 − 3 = 5); F11 (FR:97-98; factor_nowcast_screen.py:21-22: 14 ← 12, 13; 15, 16 ← 12, 13, 14); FR:170-171; I05; K24

Lợi ích không riêng của luật C. Trên cùng 21 kịch bản, khi có dự báo, luật hạn chót 30 giây giảm từ 66 xuống 2 khách
trễ trong 1.471 khách được phục vụ; luật trần độ dời khách thấy 30 giây cũng giảm từ 66 xuống 2 trong 1.470 khách.
Điều phối không trần giảm từ 129 xuống 63 trong 1.512 khách. Riêng khung 17 giờ, 28 khách trễ trong 437 khách của điều
phối không trần không đổi, còn luật C có 0 trong 418 khách. Theo sổ, 86,1\% số giây trễ của 68 khách trễ khi không dự
báo thuộc phần được hồ sơ lịch sử giải thích (Hình~\ref{hinh:c4-so}); dự báo và trần nhắm vào hai phần khác nhau của sổ.
% EV: F08; K37; N4B8-02 (check A: M⁺-30 66 → 2 trên 1.471; V-30 66 → 2 trên 1.470; U 129 → 63 trên 1.512; U khung 17 giờ 28 → 28 trên 437; C-30 khung 17 giờ 0 trên 418); F09 (FA:30; RX:3); eq:so; K04

Trên 21 kịch bản này, giá của dự báo là lời hứa giờ trả dài thêm 13,7 giây (khung 7 giờ 27,8 giây) và lời hứa giờ đón 3,2 giây, trung bình
theo kịch bản. Trên cả 34 kịch bản, khách trễ của luật C giảm từ 76 xuống 10 trong 2.338 (Bảng~\ref{bang:c4-tran120}).
% EV: F06 (FR:67-69; FA:8-9); C06 (76); C09 (10 trên 34 kịch bản); N-16; F02; N-28 (13 kịch bản cuối tuần không có dự báo; độ dài thêm gộp 7,37 s, N4B-10, nằm ở bang:c4-tran120 và bang:c4-dem)

Kết quả 2 trong 1.471 khách đạt được trong điều kiện thuận lợi cho dự báo theo giờ trong ngày, không phải hiệu quả ngoài
đời. Tắc đường là một hệ số chung cho cả thành phố. Trên cả 34 kịch bản, với luật C, kéo dài lời hứa theo đúng hệ số
sẽ xảy ra cho giờ trả khớp trong vòng 2 giây ở cả 2.305 khách không bị điều phối dời. Phép kiểm giờ đón cũng chưa dùng
lời hứa này (mục~\ref{muc:c4-thaoluan}).
% EV: I02 (FR:167-169: mức tốt nhất cho loại dự báo theo giờ trong ngày trong mô phỏng gần như tất định); F10 (FA:28; RX:8: 34 kịch bản); B06; F07; X3-11

\begin{figure}[htbp]
\centering
\resizebox{\linewidth}{!}{%
\begin{tikzpicture}[x=1cm,y=1cm,font=\footnotesize,
  nhan/.style={font=\scriptsize,inner sep=1.5pt}]
% Thanh trên: ngày thường. Độ dài 11 cm = 6.217,4 s; ranh giới 0,1445 và 9,6145 cm (FA:30).
\draw[fill=black!75] (0,1.6) rectangle (0.1445,2.2);
\draw[fill=black!20] (0.1445,1.6) rectangle (9.6145,2.2);
\draw[fill=white] (9.6145,1.6) rectangle (11,2.2);
\node at (4.88,1.9) {5.353 s (86,1\%)};
\node[align=center,nhan] at (10.307,1.9) {783 s\\(12,6\%)};
\draw (0.072,2.2)--(0.072,2.45);
\node[nhan,anchor=south west] at (0.0,2.45) {82 s (1,3\%)};
\node[anchor=east,align=right] at (-0.15,1.9) {Ngày thường, 21 kịch bản\\68 khách, 6.217 s};
% Thanh dưới: chiều tuyết. Độ dài 11 cm = 3.445,9 s; ranh giới 0,0273 và 0,8302 cm (NEW:5).
\draw[fill=black!75] (0,0) rectangle (0.0273,0.6);
\draw[fill=black!20] (0.0273,0) rectangle (0.8302,0.6);
\draw[fill=white] (0.8302,0) rectangle (11,0.6);
\node at (5.915,0.3) {3.185,9 s (92,5\%)};
% v8: nhãn 8,6 s đặt phía trên thanh tuyết (khoảng trống giữa hai thanh) để đường dẫn của nhãn 251,5 s không cắt chữ.
\draw (0.014,0.6)--(0.014,0.85);
\node[nhan,anchor=south west] at (-0.02,0.85) {8,6 s (0,25\%)};
\draw (0.43,0)--(0.43,-0.2);
\node[nhan,anchor=north west] at (0.39,-0.2) {251,5 s (7,3\%)};
\node[anchor=east,align=right] at (-0.15,0.3) {Chiều tuyết, 8 kịch bản\\41 khách, 3.445,9 s};
% Chú giải (ba ô cách nhau đủ để chữ không chồng lên nhau)
\draw[fill=black!75] (0,-1.25) rectangle (0.3,-1.0);
\node[anchor=west] at (0.35,-1.125) {phần tính cho điều phối};
\draw[fill=black!20] (4.0,-1.25) rectangle (4.3,-1.0);
\node[anchor=west] at (4.35,-1.125) {phần được hồ sơ lịch sử giải thích};
\draw[fill=white] (9.6,-1.25) rectangle (9.9,-1.0);
\node[anchor=west] at (9.95,-1.125) {phần còn lại};
\end{tikzpicture}
}
\caption[Ba phần của số giây trễ theo sổ của từng khách]{Số giây trễ chia thành ba phần theo sổ của từng khách (công thức~\eqref{eq:so}), với luật C trần 30 giây và các
khách tới muộn hơn lời hứa đầu không dự báo quá 60 giây: 68 khách của 21 kịch bản ngày thường và 41 khách của 8 kịch
bản chiều tuyết 15/11/2018. Mỗi thanh dài bằng 100\% số giây trễ của nhóm đó. Phần được hồ sơ lịch sử giải thích phụ
thuộc dự báo đang dùng; phần còn lại là phần dự báo này chưa giải thích được, không phải phần không thể biết trước.}
\label{hinh:c4-so}
\end{figure}
% EV (ngày thường): F09; N4B8-11 (FA:30 và CELLS: 81,742 + 5.352,646 + 783,018 = 6.217,406 s; các nhãn làm tròn đến giây độc lập, nên tổng ba nhãn là 6.218 s, lệch 1 s so với nhãn tổng 6.217 s); RX:3 (1,3 / 86,1 / 12,6%)
% EV (chiều tuyết): F13; N4B-05 (NEW:5: 8,566 / 251,496 / 3.445,933 s; RP1:24-26: 7,3%; SNOW:10: 92,5%; FR:104,186); N4B8-12 ([A] 8,566 / 3.445,933 = 0,25%)
% Tọa độ: x = giây / tổng × 11 cm (tính tay; xem NOTES v7 mục 4).

\textbf{Đưa dự báo vào bộ điều phối.} Phép thử chạy luật C và luật hạn chót 30 giây trên 16 kịch bản khung 7 và 8 giờ.
Thời gian đi của bộ điều phối được nhân thêm mức tăng dự báo trung bình của khung 15 phút hiện tại và khung kế tiếp. Có
3 trong 32 lần chạy dừng vì FleetPy tính rằng xe sẽ tới điểm đón muộn hơn giờ đón muộn nhất chưa tới 3 giây. Luật C
không có quyết định bị ép ở 14 lần chạy hoàn tất. Luật hạn chót bị ép ở 47,6\% lượt quyết định ở 15 lần chạy hoàn tất,
so với 62,8\% ở các lần chạy không đưa dự báo vào trên cùng 15 kịch bản. Tính gộp trên cùng 15 kịch bản, đó là 3.849
trong 8.496 lượt quyết định khi đưa dự báo vào, so với 5.121 trong 8.468 lượt khi không đưa vào; các tỉ lệ 47,6\% và
62,8\% ở trên là trung bình theo kịch bản. Tỉ lệ khách bị luật hạn chót phá lời hứa ở khung 7 giờ giảm từ 39,0\% xuống
11,0\%, trung bình trên cùng tám kịch bản. Tính gộp theo khách trên tám kịch bản này, đó là 170 trong 442 khách khi
không đưa dự báo vào, so với 45 trong 439 khách khi đưa vào. Vì vậy, các số của Bảng~\ref{bang:c4-nghiavu} phụ thuộc cách
hứa. Khóa luận chỉ dùng dự báo ở lời hứa báo cho khách; kết luận này chỉ áp cho cách đã thử.
% EV: F12; X3-29 (W13:85-89); N4B8-06 (K = 2: run_forecast.py:9,72; configs/forecast-k2-2018-11-1{4,5,6}.json:1; preflight_world_v1_3.py:85-89; FR:117; R26 Q9); N4B-07 (RX:9,11-15; FR:120-121,124,138,157: C hoàn tất 14, luật hạn chót 15; 47,6 và 62,8% trên cùng 15 kịch bản); C14; FR:141-146; R37 verify-numbers-v2.json:dispatcher; R39 agent-numbers-v1.json:dispatcher.paired, morning_promise_breaks (cùng15 hoàn tất; khung7 cùng8; số gộp tách khỏi trung bình theo kịch bản)

\textbf{Học máy.} Ba mô hình học một lần rồi giữ nguyên ước lượng phân vị 90 của độ trễ so với lời hứa theo dự báo (mức
mà 90\% khách giống vậy không vượt) từ thông tin đã có ở thời điểm hứa. Đó là mô hình cây quyết định tăng cường
(gradient boosting), mô hình ấy khi bỏ hai biến trạng thái xe (số khách khác đã gán và đang trên xe), và hồi quy phân vị
tuyến tính. Mỗi khách được cộng vào lời hứa theo dự báo một phần đệm tỉ lệ với mức ước lượng; đối chứng cộng cùng một
khoảng cho mọi khách. Cả hai được chỉnh trên ngày hiệu chỉnh để lời hứa dài thêm trung bình 30 giây so với lời hứa không
dự báo, nên trên ngày kiểm mức này có thể khác. Phép thử dùng hồ sơ dự báo cố định của ngày
12--13/11 và kiểm trên 15, 16/11 khung 7 giờ và 13/11 của tập tắc đường cố định, các ngày đã được xem trước. Mô hình
chỉ được coi là tốt hơn khi, so với đối chứng, thỏa cả bốn điều kiện:
\begin{itemize}[leftmargin=1.2cm,itemsep=0pt,topsep=2pt]
\item tỉ lệ khách trễ quá một phút thấp hơn;
\item độ dài thêm trung bình không hơn quá 1 giây;
\item phân vị 90 của độ dài thêm không hơn quá 15 giây;
\item tỉ lệ khách tới sớm quá một phút không hơn quá 1 điểm phần trăm.
\end{itemize}
% EV: G01; G02; G03; N4B-08; N4B8-04 (MLS:9-10 các ngày đã xem; 19-23, 67-68 hai biến trạng thái; 37-39 hồ sơ cố định 12–13/11; 41-61 chia ngày; 63-69 ba mô hình; 71-82 chỉnh trên ngày hiệu chỉnh, F = D0 + ρ·w, đối chứng D0 + (E − mean δ), không chỉnh lại khi kiểm; 91-95 bốn điều kiện; MLT: đối chứng của U dài thêm 37,219 / 31,772 s ở 15/11 / 16/11)

Không mô hình nào thỏa đủ bốn điều kiện. Các số đếm sau được suy từ tỉ lệ gộp theo khách đã làm tròn đến ba chữ số
thập phân trong bảng kết quả; với mẫu số tương ứng, mỗi tỉ lệ chỉ ứng với một số nguyên. Với luật C, đối chứng và mô hình đều không có khách trễ (147 và 197 khách).
Với điều phối không trần, mô hình cây tăng cường chỉ bớt 1 trong 2 khách trễ ở ngày 15/11 (145 khách), không bớt ai ở
ngày 16/11 (200 khách), và không thỏa điều kiện thứ hai, thứ ba. Với ghép tự do, khách trễ giảm từ 104 xuống 94 trong 807
nhưng 57 khách tới sớm quá một phút, so với 0 khi đệm đều. Bỏ hai biến trạng thái xe không đổi số khách trễ ở hai phép
so này.
% EV: G05–G10; N4B8-03 (check B: C-30 0/147, 0/197; U 2 → 1 trên 145, 1 → 1 trên 200; FREE 104 → 94 trên 807, sớm 57 so với 0; bỏ trạng thái xe 1/145 và 94/807); ML:6-7,14-15,21,45-53,67-69 (độ dài thêm +4,4 / +3,8 s; phân vị 90 +21,8 / +21,7 s); K40; R37 verify-numbers-v2.json:ml; R39 agent-numbers-v1.json:ml (xét mọi số nguyên, sai số 0,0005 điểm phần trăm); ml_pilot_run.py:37-42,175-176

% -----------------------------------------------------------------------------------------------------
\section{So với cách đệm lời hứa}
\label{muc:c4-dem}

Khi có dự báo, cộng thêm 30 giây vào mọi lời hứa của điều phối không trần cũng cho 2.328 lượt không trễ quá một phút
trên 3.672 yêu cầu, bằng luật C trần 30 giây có dự báo. Ở ngưỡng hai phút, cách đệm này cho 2.366 lượt, so với 2.338
của luật C. Hai cách này cho hai thứ khác nhau. Đệm kéo dài lời hứa, trung bình 37,2 thay vì 7,4 giây, nên ít khách bị tính là trễ. Nhưng đệm không đổi
quyết định điều phối: vẫn có 174 khách bị điều phối dời giờ trả, với tổng độ dời lớn nhất là 590,1 giây (gần 10 phút).
Luật C trần 30 giây giới hạn tổng độ dời do điều phối của mỗi khách ở 30 giây dưới giả thiết (N); trên 34 kịch bản của
tập kiểm tra chính, mức lớn nhất đo được là 29,2 giây, kể cả ở các kịch bản cuối tuần không có dự báo. Đổi lại, luật C
phục vụ ít hơn 96 lượt.
% EV: C10; N-16; N-18; K06; K51; C20; N-23; N-28 (cuối tuần 29,2 s); C02; A06 (dưới giả thiết (N); không bảo đảm giờ đến: A22, K05); RD §8; khớp câu trỏ của phần A v5-v3 mục 4.2 (C10); R37 verify-numbers-v2.json:padding; R39 agent-numbers-v1.json:padding,cohort_totals

Cách đệm là phát lại trên các lần chạy của điều phối không trần, nên hai bên được so trên hai nhóm khách khác nhau, chỉ
chung 2.207 khách. Bảng~\ref{bang:c4-dem} xếp các cách hứa thành ba khối có lời hứa dài thêm gần bằng nhau.
% EV: C10 (phát lại, không phải đối chứng ghép cặp); ADDENDUM-1 phần đầu (quy tắc đệm D = D0 + k); C04

\begin{table}[htbp]
\centering
\caption[Trần và cách đệm lời hứa khi có dự báo]{Các cách hứa khi có dự báo trên tập kiểm tra chính (34 kịch bản, 3.672 yêu cầu mỗi cấu hình; dự báo chỉ có ở
21 kịch bản ngày thường). U: điều phối không trần; C-$B$: luật C trần $B$ giây; trong ngoặc là số khách được phục vụ.
Số đếm gộp; trễ và tới sớm tính trong số khách được phục vụ, so với lời hứa của chính cách hứa; hai cột ``không trễ
quá'' tính trên 3.672 yêu cầu; độ dài thêm là trung bình gộp theo khách được phục vụ. Đệm không đổi quyết định điều
phối, nên hàng của U và hàng của C là các nhóm khách khác nhau.}
\label{bang:c4-dem}
\small
\begin{tabular}{@{}p{4.6cm}rrrrr@{}}
\toprule
\textbf{Cách hứa} &
\textbf{\begin{tabular}[b]{@{}r@{}}Trễ\\ quá 60 s\end{tabular}} &
\textbf{\begin{tabular}[b]{@{}r@{}}Không trễ\\ quá 60 s\end{tabular}} &
\textbf{\begin{tabular}[b]{@{}r@{}}Không trễ\\ quá 120 s\end{tabular}} &
\textbf{\begin{tabular}[b]{@{}r@{}}Tới sớm\\ quá 60 s\end{tabular}} &
\textbf{\begin{tabular}[b]{@{}r@{}}Dài thêm\\ (s)\end{tabular}}\\
\midrule
\multicolumn{6}{@{}l}{\textit{Không đệm (dài thêm khoảng 7 giây)}}\\
U, dự báo (2.434) & 137 & 2.297 & 2.347 & 8 & 7,18\\
C-30, dự báo (2.338) & 10 & 2.328 & 2.338 & 4 & 7,37\\
C-120, dự báo (2.387) & 65 & 2.322 & 2.379 & 2 & 7,27\\
% EV: N-16; N-03; C09; N-06; L1a-03; N4B-09 (sớm 8 / 4 / 2)
\midrule
\multicolumn{6}{@{}l}{\textit{Đệm 30 giây (dài thêm khoảng 37 giây)}}\\
U, dự báo, đệm 30 s (2.434) & 106 & 2.328 & 2.366 & 44 & 37,18\\
C-30, dự báo, đệm 30 s (2.338) & 2 & 2.336 & 2.338 & 44 & 37,37\\
% EV: C10; N-02; N-16
\midrule
\multicolumn{6}{@{}l}{\textit{Đệm 60 giây (dài thêm khoảng 67 giây)}}\\
U, dự báo, đệm 60 s (2.434) & 87 & 2.347 & 2.386 & 1.033 & 67,18\\
% EV: C11; N-16
\bottomrule
\end{tabular}
\end{table}

Cách đọc bảng: với đệm 30 giây trên điều phối không trần, 106 trong 2.434 khách tới muộn hơn lời hứa đã đệm quá 60
giây, nên có 2.328 lượt không trễ quá 60 giây trên 3.672 yêu cầu; 44 khách tới sớm quá 60 giây.
% EV: C10; N-16

\textbf{Vai trò của từng thành phần.} Thứ nhất, ở mức kéo dài lời hứa gần bằng nhau, luật C trần 30 giây hơn cách đệm
ở ngưỡng một phút và kém ở ngưỡng hai phút: khi cùng đệm 30 giây, hai bên có 2.336 so với 2.328 lượt và 2.338 so với
2.366 lượt. Ở hàng đệm 30 giây, điều phối không trần đạt cùng 2.328 lượt không trễ quá một phút như luật C có dự báo
không đệm, với lời hứa dài thêm trung bình 37,18 so với 7,37 giây, tức khoảng năm lần.
% EV: N-17; N-18 (37,1833 / 7,3671 = 5,05); N-16; R37 verify-numbers-v2.json:padding; R39 agent-numbers-v1.json:padding,cohort_totals

Thứ hai, lời hứa dài thì ít trễ nhưng nhiều khách tới sớm: 44 khách với cách đệm 30 giây và 1.033 khách với cách đệm 60
giây, so với 4 khách với luật C có dự báo. Số lượt tới trong khoảng $\pm$60 giây quanh lời hứa là 2.284 với cách đệm 30
giây, so với 2.324 với luật C, trên 3.672 yêu cầu.
% EV: C10; C11; N-04 (±60 s = phục vụ − trễ60 − sớm60, ADD1:92); N4B-09 (ADD1:84-89)

Thứ ba, cách đệm 30 giây vẫn có 68 khách trễ quá hai phút, 11 khách quá năm phút và 1 khách quá mười phút; luật C có
dự báo không có ai trễ quá hai phút. Cách đệm hơn 28 lượt ở ngưỡng hai phút vì phục vụ nhiều hơn 96 lượt, không phải vì
ít khách trễ hơn.
% EV: N-01 (68 / 11 / 1 so với 0 / 0 / 0; ADD1:15,19); N-19 (2.366 − 2.338 = 28 = 96 − 68)

Khi phát lại cùng lời hứa theo dự báo, luật hạn chót 30 giây cũng có 2.338 lượt phục vụ và 2.328 lượt không trễ quá
một phút. Vì phép phát lại không đổi các lần chạy, nghĩa vụ của luật này vẫn tính từ lời hứa đầu không dự báo. Theo
nghĩa vụ đó, tính trung bình theo kịch bản, nó bị ép ở 41,3\% lượt quyết định; tính gộp là 7.710 trong 20.103 lượt.
Luật hạn chót có khách tới muộn hơn lời hứa đầu không dự báo quá 30 giây ở 28 trong 34 kịch bản, còn luật C không có
quyết định bị ép hay khách vượt trần (Bảng~\ref{bang:c4-nghiavu}).
% EV: N-16 (M⁺-30 + dự báo 2.338 / 10 / 2.328 / 2.338, ADD1:29); C13 (41,3% trung bình theo kịch bản); N-21 (phá hứa của M⁺-30 = A − p₀ > 30 s với lời hứa không dự báo; C-30 0/34, 0 quyết định bị ép); X3-16; R36-A01 (bảng A-v5-v3 dùng số gộp); R39 agent-numbers-v1.json:main_original_obligation

Trần 120 giây có dự báo, với lời hứa chỉ dài thêm 7,27 giây, cho 2.379 lượt không trễ quá hai phút, nhiều hơn 2.366
của cách đệm 30 giây, nhưng 2.322 lượt ở ngưỡng một phút, ít hơn 2.328. Đây là một điểm đo thêm trên một tập khách khác,
không phải mức trần tốt nhất.
% EV: N-06; N-16; N-18; N-10; K36

\textbf{Khi nào chọn trần, khi nào chọn đệm.} Nếu hãng xe chỉ tính số lượt không trễ quá hai phút và chấp nhận lời hứa
dài hơn cùng nhiều khách tới sớm hơn, điều phối không trần có dự báo và đệm 30 giây đạt 2.366 lượt, so với 2.338 lượt
của luật C trần 30 giây có dự báo, trên 3.672 yêu cầu; điều phối không trần cũng phục vụ thêm 96 lượt. Nếu cần giới hạn
tổng độ dời do điều phối của từng khách, kể cả khi không có dự báo, luật C trần 30 giây làm được điều này dưới giả thiết
(N), đổi lại 96 lượt phục vụ và 242 khách đi chung. Hai cách dùng chung được; chọn cách nào là quyết định của hãng xe,
không phải kết luận của thực nghiệm.
% EV: N-19; C02; N-20; N-02 (hàng C-30 có dự báo và đệm 30 s); D20; I07; I10; K38; K51; A22 (giả thiết N); R37 verify-numbers-v2.json:padding; R39 agent-numbers-v1.json:padding,cohort_totals

% -----------------------------------------------------------------------------------------------------
\section{Khi dự báo không đủ: chiều tuyết 15/11/2018}
\label{muc:c4-tuyet}

Trên 601 khách được phục vụ của 8 kịch bản chiều tuyết 15/11/2018 (luật C trần 30 giây), lời hứa theo dự báo chỉ giảm
số khách trễ quá một phút từ 41 xuống 34. Tính trung bình theo kịch bản, tỉ lệ này giảm từ 7,19\% xuống 5,94\%. Theo
sổ, trong 3.445,9 giây trễ của 41 khách, 8,6 giây tính cho điều phối và 251,5 giây được hồ sơ lịch sử giải thích
(Hình~\ref{hinh:c4-so}).
% EV: F13; K25; B20; N4B-05 (NEW:5)

Bảng~\ref{bang:c4-tuyet} cho kết quả phát lại vài cách hứa dùng thêm hệ số tắc đường của chính ngày tuyết. Cách ``trộn'' lấy trung
bình nhân của mức tăng theo hồ sơ lịch sử và mức tăng của hệ số cùng ngày từ khung 15 phút trước sang khung hiện tại.
% EV: F14; K20; K21; N4B-06 (factor_nowcast_screen.py:14-20 quy tắc thông tin, 28-29: MOM = (f(k0)/f(k0−1))^(b−b0), BLEND = √(HIST·MOM); RP1:7-20; NEW:5)

\begin{table}[htbp]
\centering
\caption[Phát lại các cách hứa trong chiều tuyết]{Phát lại các cách hứa trên 8 kịch bản chiều tuyết 15/11/2018 (luật C trần 30 giây, 601 khách được phục vụ).
Cột ``khách trễ'' đếm gộp số khách tới muộn hơn lời hứa đầu quá 60 giây; cột cuối là độ dài thêm trung bình theo kịch
bản. Không cách hứa nào trong bảng có khách tới sớm quá 60 giây. Mọi cách hứa đều tính lại sau khi chạy, trên một ngày
đã được xem trước khi thiết kế các cách này; hai hàng ``có buổi sáng'' chỉ để tham khảo.}
\label{bang:c4-tuyet}
\small
\begin{tabular}{@{}p{3.4cm}p{6.6cm}rr@{}}
\toprule
\textbf{Cách hứa} & \textbf{Thông tin dùng lúc hứa} &
\textbf{\begin{tabular}[b]{@{}r@{}}Khách\\ trễ\end{tabular}} &
\textbf{\begin{tabular}[b]{@{}r@{}}Dài thêm\\ (s)\end{tabular}}\\
\midrule
Không dự báo & Thời gian đi hiện tại & 41 & 0\\
Theo dự báo & Hồ sơ tắc đường của các ngày nguồn & 34 & 2,0\\
Trộn trong khung & Thêm hệ số cùng ngày đã có trong khung giờ & 21 & 8,5\\
Trộn có buổi sáng & Thêm hệ số cùng ngày từ 6 giờ; thiết kế thử không cho khung chiều dùng dữ liệu này & 19 & 9,5\\
Trộn có buổi sáng, trễ một khung & Như hàng trên, nhưng mức tăng gần nhất chỉ lấy từ các khung đã kết thúc; hệ số của
khung hiện tại vẫn dùng giá trị đúng & 29 & 9,0\\
Đệm tỉ lệ & Tỉ lệ chọn để tổng phần kéo dài bằng của hàng ``trộn trong khung'', trên chính ngày tuyết & 27 & 8,3\\
Hệ số thật & Hệ số sau lúc hứa, không có ở thời điểm hứa & 0 & 18,8\\
\bottomrule
\end{tabular}
\end{table}
% EV: F14 (RP1:7-20; NEW:5); N4B-06 (định nghĩa; cột dài thêm là trung bình theo kịch bản, RP1:6, rider_nowcast_replay.py:96-106); K20; K21; factor_nowcast_screen.py:17-20 (hàng trễ một khung cũng dùng thông tin từ 06:00); N4B8-09 (RP1:7-20: E60 0,00% ở mọi hàng ⇒ 0 khách tới sớm ở mọi kịch bản [A])

Cách đọc bảng: khi trộn hồ sơ lịch sử với xu hướng trong khung giờ, 21 trong 601 khách tới muộn hơn lời hứa quá một
phút, với lời hứa dài thêm trung bình 8,5 giây.
% EV: F14; N4B-06

Phép phát lại này tách số giây trễ theo cách cộng có giới hạn, khác phần có dấu của công thức~\eqref{eq:so}. Phần kéo
dài của mỗi khách được chặn trong khoảng từ 0 đến phần trễ còn lại sau khi trừ phần tính cho điều phối. Theo cách cộng
này, cách trộn
có buổi sáng giải thích thêm 886,9 giây, tức 25,7\% của 3.445,9 giây trễ; 2.299,0 giây (66,7\%) chưa được giải thích
bằng các cách đã thử. Theo số liệu trạm Central Park của Cơ quan Quản lý Khí quyển và Đại dương Quốc gia Hoa Kỳ
(National Oceanic and Atmospheric Administration, NOAA), từ 9/2018 đến 12/2019 chỉ có 5 đợt tuyết từ 25,4 mm trở lên,
kể cả đợt này; muốn kiểm cách hứa cho ngày bất thường cần thêm nhiều đợt khác.
% EV: F15 (NEW:5; RP1:24-35); N4B8-07 (rider_nowcast_replay.py:14-15, 113: explained = Σ clip(D − P0, 0, A − P0 − dispatch); RP1:29: phần vượt 6 s; với hồ sơ lịch sử phần vượt 0 nên hai cách cộng trùng nhau, N4B-05); K19; I11; B26 (NEW:9; WX:1-3: tệp NOAA GHCN-D trạm Central Park; WX17:1-4; ghép ngày cách nhau ≤ 2)

% -----------------------------------------------------------------------------------------------------
\section{Kiểm tra tính hợp lệ của kết quả}
\label{muc:c4-hople}

Trên 34 kịch bản của tập kiểm tra chính, không khách nào vượt trần: tổng dời lớn nhất là 29,2 giây với trần 30 giây và
118,9 giây với trần 120 giây. Không lượt quyết định nào bị ép, 0 trong 20.101 và 0 trong 20.435 lượt. Đây là phép kiểm cài
đặt: với sổ hợp lệ $0 \le L_r \le B$ ở lượt đang xét và giả thiết (N), Nhận xét~\ref{nx:giunguyen} đã cho trước các
số 0 này; kết luận xuyên chuyến cần giả thiết (N) đúng ở mọi lượt. Điều phối không trần cũng có 0 trong 20.760 lượt.
% EV: C20; N-07; X3-16 (0/20.101; 0/20.435; U 0/20.760); N-21; A22; chương 3, nx:giunguyen và đoạn sau nó (0 ≤ L_r ≤ B ở lượt đang xét; (N) đúng ở mọi lượt cho kết luận xuyên chuyến); R39 agent-numbers-v1.json:cohort_totals

Bộ đếm của bộ sinh được đọc ở điều phối không trần, luật C trần 30, 120 giây và luật hạn chót 30 giây. Ở cả bốn cấu
hình, không phương án nào bị bỏ vì giới hạn 100 phương án mỗi xe hay giới hạn số bước tìm cách chèn. Giới hạn 2 yêu cầu
mới mỗi lượt chỉ để một yêu cầu sang lượt sau ở mỗi cấu hình. Các số này không đo mọi giới hạn tính toán (mục~\ref{muc:gioihan}): giới hạn
của bước kiểm, của bộ giải, một bộ đếm khác của bộ sinh và các bộ đếm ở tập tắc đường cố định chưa được đọc.
% EV: N-25 (ADD1:119-122: 34 lần chạy mỗi cấu hình); N4B-11 (GEN:54-95, 236-246, 298-345; CandidateModels.cs:68-82, 100-118: điều kiện đầy đủ còn gồm OmittedRepairRequestCount, chưa đọc); X3-17; X3-03; X3-32; C4A-18 (tập tắc đường cố định: giới hạn chưa kiểm); K34 thu hẹp (N4B8-13)

Như mục~\ref{muc:doichung} đã nêu, mức dùng trần lớn nhất của tiêu chí 3 bằng 100\% ở mọi lượt có ghi tiêu chí này:
3.503, 3.506 và 3.504 lượt khi không trần, với trần 30 và 120 giây.
% EV: A15; N-25 (ADD1:119-121); A07; X3-07 (chương 3 mục doichung nêu hệ quả: tiêu chí 3 chỉ tác dụng qua tổng phần dời)

Các số đếm chính của tập kiểm tra chính được ít nhất hai chương trình đọc dữ liệu viết riêng tính lại và khớp nhau
(mục~\ref{muc:kiemtra}). Phép so gồm cả từng khách ở 7.159 khách được phục vụ của điều phối không trần và luật C trần
30, 120 giây; riêng số khách đi chung chỉ được một quy tắc đếm. Theo bản ghi khởi tạo, 8 lần chạy chiều tuyết dùng cùng
bản bộ chạy đã lưu và cùng cấu hình với các lần chạy luật C trần 30 giây của tập kiểm tra chính, kể cả luật phục hồi
không tệ hơn.
% EV: N4B-14 (ADD1:66-78); K41; C19; N-09 (đi chung: Codex không đếm lại); N4B8-05 (check C: 8 lần chạy chiều tuyết và 51 lần chạy C-30 của T3.8 cùng binarySha256 a631038b…, coreCommitSha 1ce15e92…, policyConfigurationHash a2f11a3a…; sha256 runner-tier3-v2/RideBound.Runner.dll khớp; wp4-tier3-v2.json:7 no-worse-than-reference-v1; R26 Q8)

Phép kiểm số đếm học máy chỉ đối chiếu số học của bảng kết quả; nó không kiểm lại quá trình học hay dự đoán của mô hình.
% EV: R37 §1.8; R39 agent-numbers-v1.py phần C; agent-numbers-v1.json:ml; ml_pilot_run.py:37-42,175-176

Chạy lại trên đúng mã nguồn đã dựng bộ chạy, cả 1.025 trong 1.025 kiểm thử tự động đạt trong một lần chạy đầy đủ ngày
13/10/2026, với kiểu xuống dòng LF và một tệp chưa ghi nhận trong cây làm việc (Phụ lục~\ref{phuluc:C}). Các phép kiểm
trên không cho biết mô phỏng giống thực tế đến đâu.
% EV: TESTRUN-45f8026-2026-10-12.md:18 (attempt 4: LF và tệp chưa ghi nhận, 1.025/1.025 trong một lần chạy ngày 13/10/2026); C:/RideBoundData/research/testrun-45f8026-eollf-dirty-fullrun-2026-10-13.log; R39 agent-numbers-v1.json:test_attempt4; I13; A28; I01; K41

% -----------------------------------------------------------------------------------------------------
\section{Thảo luận và hạn chế}
\label{muc:c4-thaoluan}

Với một hãng xe, các kết quả trên là các điểm trên một đường đánh đổi, không phải một khuyến nghị. So với điều phối
không trần, luật C trần 30 giây phục vụ ít hơn 96 lượt và có ít hơn 242 khách đi chung; đổi lại, phần dời do điều phối
của mỗi khách không quá 30 giây và số lượt trễ thấp hơn ở cả bốn ngưỡng đã đo, trên hai nhóm khách khác nhau
(Bảng~\ref{bang:c4-phanbotre}). Khóa luận không đo giá trị của các điều này với khách hay hãng xe.
% EV: C02; N-20; C20; N-24; C04; D20; I10; K38 (không kể 0 quyết định bị ép là phần luật C đổi được: X3-16, U cũng 0/20.760)

Một phép tính hòa vốn giả định cho thấy cách đọc các điểm này. Gọi $\lambda$ là chi phí của hãng xe cho mỗi lượt trễ quá
một phút và $\mu$ là biên lợi nhuận mỗi chuyến. Với 96 lượt phục vụ ít hơn và 127 lượt trễ ít hơn, luật C chỉ có lợi hơn
khi $\lambda > 0{,}756\,\mu$, tức $96/127$. Phép tính bỏ qua mọi yếu tố khác và so hai nhóm khách khác nhau; khóa luận
không có giá trị thật nào của $\lambda$ hay $\mu$.
% EV: C29 (R15:231-235; 96/127 = 0,756); K35; C28; C04; I10

Các hạn chế dưới đây chia theo bốn nhóm; mỗi ý nêu điều chưa làm và kết luận chịu ảnh hưởng.

\textbf{Mô hình mô phỏng.}
\begin{itemize}[leftmargin=1.2cm,itemsep=2pt,topsep=3pt]
\item Mọi kết quả đến từ FleetPy với một hệ số tắc đường chung cho cả thành phố, mức nền lấy lúc 00:00, nên mô phỏng gần
như không có yếu tố ngẫu nhiên. Các kết quả, kể cả mức trễ còn lại sau dự báo, không cho biết hiệu quả khi vận hành thật.
\item Thời gian dừng đón, trả bằng 0 giây (kiểm trên 3 lần chạy mẫu), so với 30 giây mỗi khách trong nghiên cứu điển hình
của FleetPy~\cite{engelhardt2026}, và đồ thị rút gọn chỉ có 96 nút. Phần dời do điều phối và cái giá 96 lượt vì vậy chỉ áp cho
thiết lập này.
\item Không có phản ứng của khách và không gán lại xe cho khách đã nhận, nên số lượt phục vụ không tính khách bỏ đi, và
chưa biết cái giá của trần khi được gán lại xe.
\end{itemize}
% EV: I01; B01; B06; B07; I02; F10; B05 (SERV:2-4: kiểm trên 3 lần chạy mẫu; FLEET:423-424); J18; B03; K48; B12; I07; X3-39

\textbf{Dữ liệu và thiết kế thực nghiệm.}
\begin{itemize}[leftmargin=1.2cm,itemsep=2pt,topsep=3pt]
\item 34 kịch bản của tập kiểm tra chính phụ thuộc nhau, thuộc 5 ngày, mỗi kịch bản chạy một lần; các phép đếm theo
kịch bản là mô tả, không phải kiểm định.
\item Chưa kiểm độ nhạy theo hạt giống của bộ giải, nên chưa biết các chênh lệch có giữ nguyên khi đổi hạt giống.
% EV: R37 §1.8/L-26; B09; B11; chương 3 v4-v3 muc:gioihan (hạt giống cố định 7, một lần chạy mỗi kịch bản)
\item Các kết quả dùng lại dữ liệu đã xem. Ngưỡng phân tích và cách chia tập được chọn sau khi xem tập buổi sáng và các
đợt chạy thử trước đó. Mức 120 giây được chọn sau khi xem kết quả của điều phối không trần và luật C trần 30 giây trên
chính tập kiểm tra chính, kể cả các phép tính thử ở 120 giây (mục~\ref{muc:c4-thamdo}). Chiều tuyết và các ngày kiểm học máy cũng đã được xem, nên mọi kết quả là thăm dò.
\item Tải thấp: 0,535 khách trên xe kể cả lúc xe rỗi khi không trần, so với khoảng 1,33 trong nghiên cứu điển hình của
FleetPy. Ở 49 kịch bản dựng từ dữ liệu thật, chuyến được chọn ngắn hơn mọi chuyến đủ điều kiện (trung vị thời gian đi thẳng 165
so với 196 giây). Kết quả không suy ra cho tải cao hay chuyến dài hơn.
\end{itemize}
% EV: B09; I05; B11; I06; C4A-11 (mức 120 giây, như phần A v5-v3 mục c4-thamdo); MLS:9-10; I11; B12 (LOAD:2-3; FLEET:447-449); J18; B13 (SEL6:2-4: 49/49 kịch bản); K42 (kể cả các phép tính thử ở 120 giây); chương 3 v4-v3 muc:gioihan (hạt giống7; chưa có kiểm độ nhạy)

\textbf{Tích hợp.}
\begin{itemize}[leftmargin=1.2cm,itemsep=2pt,topsep=3pt]
\item Bộ điều phối vẫn kiểm giờ đón bằng giờ không dự báo, nên 34 trong 1.471 khách được báo giờ đón muộn hơn giờ đón
muộn nhất. Các số có dự báo ở mục~\ref{muc:c4-dubao} và~\ref{muc:c4-dem} chưa tính việc sửa phép kiểm này.
\item Bản chạy ưu tiên chi phí khác bản lưu của ưu tiên ổn định. Bốn cặp đối chiếu hành vi và việc tái lập các
chỉ số nền không chứng minh tương đương toàn bộ hai bản. Kết quả ở mục~\ref{muc:c4-vaitro} vì vậy vẫn là phép so
thăm dò theo cấu hình, không loại trừ mọi khác biệt thời gian chạy.
\end{itemize}
% EV: F07 (RX:2; FR:70); I13; A28; R18 B24 (qua REG A28)

\textbf{Phạm vi kết luận.}
\begin{itemize}[leftmargin=1.2cm,itemsep=2pt,topsep=3pt]
\item Phép so hai thứ tự chọn phương án đã có kết quả ở trần 30 giây và không trần, nhưng chưa tạo thành
đường đánh đổi nhiều mức cho ưu tiên chi phí. Luật giới hạn từng lần chưa được cài đặt; các bản ghi chỉ vượt do
cộng dồn và bảy lịch sử sàng lọc không thay cho phép so hai luật trên diễn biến mô phỏng đầy đủ.
\item Luật hạn chót chỉ mượn thành phần hạn chót của Schulz và Pfeiffer~\cite{schulz2026}, với 30 và 60 giây trên tập
kiểm tra chính và 30--300 giây trên tập buổi sáng; mức mặc định của bài đó là 900 giây. Hạn chót nới kiểu Hörl và
Zwick~\cite{horl2022} (mục~\ref{muc:doichung}) trùng với cấu hình M$^{+}$ có luật phục hồi ``không tệ hơn'' ở trên,
nên mới được chạy ở 30 và 60 giây trên tập kiểm tra chính. Các phép so không nói về phương pháp của các tác giả đó.
\item Mức 30 giây không phải ngưỡng cảm nhận hay chuẩn ngành, và không có dữ liệu về mức khách sẵn lòng trả cho lời
hứa ổn định hơn; vì vậy không mức trần nào được chọn là tốt nhất.
\end{itemize}
% EV: H02; H12; K18; E09; K08; J10; J11; X3-28 (luật hạn chót 30–300 s ở tập buổi sáng; chương 3 Bảng luat); K27; D20; K52; K36
; không có \chapter mới.
% Mỗi dòng "% EV:" trỏ tới EVIDENCE-REGISTER.md (A01…K52), EVIDENCE-REGISTER-ADDENDUM-1.md (N-xx, L1a-xx, mục 0.x),
% EVIDENCE-REGISTER-ADDENDUM-2-CH3.md (X3-xx), EVIDENCE-REGISTER-ADDENDUM-3-CH4B.md (N4B-xx),
% EVIDENCE-REGISTER-ADDENDUM-4-CH4A.md (C4A-xx) hoặc EVIDENCE-REGISTER-ADDENDUM-6-CH4B-v8.md (N4B8-xx; script chỉ đọc
% evidence/checks-ch4b-v8/ch4b_v8_check.py, output ch4b-v8-check-2026-10-12.txt, phần A/B/C gọi là "check A/B/C").
% Không dùng gói mới: TikZ, enumitem và graphicx đã nạp trong main.tex; thu bảng/hình bằng resizebox.
% =====================================================================================================

% -----------------------------------------------------------------------------------------------------
\section{Vai trò của trần và thứ tự tiêu chí của bộ chọn}
\label{muc:c4-vaitro}

\subsection{Vì sao cần đổi thứ tự ưu tiên?}

Các kết quả trước mục này dùng bộ chọn ưu tiên ổn định: sau khi xét số yêu cầu nhận được, bộ chọn so mức dùng
trần và phần dời giờ rồi mới xét chi phí kế hoạch. Như vậy, cả việc xếp hạng phương án lẫn trần cứng đều có thể
góp phần hạn chế thay đổi lời hứa. Chỉ so có trần với không trần dưới một thứ tự chưa cho biết kết quả sẽ ra sao
khi bộ điều phối đặt chi phí kế hoạch lên trước phần dời giờ.

Để kiểm tra điều đó, khóa luận giữ ưu tiên nhận yêu cầu và các phép kiểm khả thi, nhưng cho bộ chọn xét chi phí
kế hoạch ngay sau đó. Các mức xếp hạng theo ổn định được bỏ, còn tiêu chí phân xử theo mã định danh giữ nguyên.
Trần cứng của luật C vẫn được kiểm trước khi chọn nếu cấu hình có trần. Hai cách gọi dưới đây là \emph{ưu tiên
ổn định} và \emph{ưu tiên chi phí kế hoạch}; chúng biểu thị mục tiêu khác nhau, không phải hai mức chất lượng
mặc định. Chi phí kế hoạch là tổng thời gian còn lại của tuyến dự kiến, chưa phải tiền hay chi phí vận hành thực.
% EV: S15-03; A14; X3-07; K15. Cùng ưu tiên nhận khách trong từng lượt, không bảo đảm cùng số phục vụ cuối.

Phép so gồm bốn cấu hình trên cùng 34 kịch bản và 3.672 yêu cầu mỗi cấu hình: hai thứ tự ưu tiên, mỗi thứ tự có
trần 30 giây và không trần. Hai cấu hình ưu tiên chi phí có 68 lượt mô phỏng hoàn tất, được đối chiếu với 68 lượt
ưu tiên ổn định đã lưu. Tập kịch bản và tập mã yêu cầu khớp ở cả 34 kịch bản. Các diễn biến sau khi chọn có thể
khác nhau, nên nhóm khách được phục vụ cũng khác nhau.
% EV: S15-01 (mục 0--3), main-final-admission.json; S15-02. Không có 34 lần lặp độc lập.

\subsection{Đánh đổi giữa phục vụ, ghép khách và sửa lời hứa}

\begin{table}[htbp]
\centering
\caption[Đánh đổi dưới hai thứ tự chọn phương án]{Hai thứ tự chọn phương án trên tập kiểm tra chính, mỗi cấu hình
có 3.672 yêu cầu. Trễ được so với lời hứa trả đầu chưa tính dự báo, trừ dòng ghi riêng. Khách có đi chung được
đếm một lần nếu hoàn thành chuyến và có khoảng thời gian cùng xe với khách khác.}
\label{bang:c4-doi-uutien}
\small
\setlength{\tabcolsep}{3pt}
\begin{tabular}{@{}>{\raggedright\arraybackslash}p{6.5cm}rrrr@{}}
\toprule
& \multicolumn{2}{c}{\textbf{Ưu tiên ổn định}} & \multicolumn{2}{c}{\textbf{Ưu tiên chi phí}}\\
\cmidrule(lr){2-3}\cmidrule(l){4-5}
\textbf{Chỉ số} & \textbf{Không trần} & \textbf{Trần 30} & \textbf{Không trần} & \textbf{Trần 30}\\
\midrule
Lượt phục vụ & 2.434 & 2.338 & 2.502 & 2.364\\
Muộn quá một phút & 203 & 76 & 379 & 72\\
Không muộn quá một phút & 2.231 & 2.262 & 2.123 & 2.292\\
Không muộn quá hai phút & 2.332 & 2.327 & 2.280 & 2.353\\
Không muộn quá năm phút & 2.416 & 2.338 & 2.466 & 2.364\\
Khách hoàn thành có đi chung & 555 & 313 & 1.027 & 432\\
Khách bị sửa giờ trả do điều phối & 174 & 33 & 434 & 106\\
Trong đó, bị sửa từ hai lần & 12 & 0 & 50 & 2\\
Đón muộn theo giới hạn dịch vụ & 51 & 43 & 39 & 44\\
Muộn quá một phút khi lời hứa có dự báo & 137 & 10 & 318 & 11\\
Tổng dời lớn nhất của một khách (s) & 590,1 & 29,2 & 709,1 & 29,7\\
\bottomrule
\end{tabular}
\end{table}
% EV: S15-01 (mục 2, 5a--5f); S15-02 (pooled); S15-04 (kiểm bảng và các chênh lệch từ TSV).

Khi cùng giữ trần 30 giây, ưu tiên chi phí phục vụ thêm 26 lượt và có thêm 119 khách hoàn thành có đi chung
so với ưu tiên ổn định (Bảng~\ref{bang:c4-doi-uutien}). Số lượt không muộn quá một phút tăng 30. Đổi lại, số khách
bị sửa giờ trả tăng từ 33 trong 2.338 lên 106 trong 2.364, tương ứng 1,41\% và 4,48\%; số khách bị sửa từ hai lần
tăng từ 0 lên 2. Số lượt đón muộn tăng từ 43 lên 44. Như vậy, cách ưu tiên ổn định vẫn cho ít khách bị sửa lời
hứa hơn, còn ưu tiên chi phí nhận và ghép được nhiều khách hơn trong thiết lập này.

Đúng giờ đến và ít sửa lời hứa là hai chỉ số khác nhau. Ưu tiên chi phí có trần có 72 lượt muộn quá một phút,
ít hơn 4 lượt so với ưu tiên ổn định có trần, nhưng nhiều khách bị sửa giờ hơn. Khi dùng lời hứa có dự báo, số
lượt muộn quá một phút là 11 và 10 trên hai nhóm khách tương ứng. Các chênh lệch nhỏ này không đủ để chọn một
cách làm tốt hơn về mọi mặt. Trung vị độ dài lời hứa đầu chưa dự báo là 809,642 giây với ưu tiên chi phí và
815,060 giây với ưu tiên ổn định, trên hai tập khách khác nhau.
% EV: S15-01 (5b--5e); S15-02. Độ dài dùng p0-t0; không quy thành chờ/đi vòng của từng khách.

Không trần cho thấy rõ cái giá của đổi thứ tự. Ưu tiên chi phí phục vụ thêm 68 lượt và có thêm 472 khách hoàn
thành có đi chung so với ưu tiên ổn định, nhưng có thêm 176 lượt muộn quá một phút và thêm 260 khách bị sửa giờ
trả. Trần vì vậy vẫn có vai trò khi bộ chọn chuyển sang ưu tiên chi phí. Trong thứ tự này, thêm trần 30 giây làm
giảm 138 lượt phục vụ và 595 khách hoàn thành có đi chung; đồng thời có ít hơn 307 lượt muộn quá một phút và
thêm 169 lượt được phục vụ, không muộn quá một phút.
% EV: S15-01 (mục 2, 5f--5g); S15-04. Các chênh lệch là tổng trên tập yêu cầu, không cùng tập khách phục vụ.

Phần giảm số lượt trễ phải đọc cùng phần không phục vụ. Với ưu tiên ổn định, 127 lượt trễ ít hơn đi cùng 96 lượt
phục vụ ít hơn và 31 lượt không muộn quá một phút nhiều hơn. Với ưu tiên chi phí, các số tương ứng là 307, 138
và 169. Chênh lệch giữa hai mức giảm trễ là 180, nhưng không có nghĩa 180 khách cụ thể được bảo vệ chỉ nhờ đổi
bộ chọn. Chênh lệch này cùng chiều ở 29 trong 34 kịch bản và cả năm ngày; các kịch bản phụ thuộc nhau nên đây
chỉ là mô tả trên dữ liệu đã dùng.
% EV: S15-01 (mục 2, 4); S15-02 (interactions); 180 = (138-96) + (169-31).

Ở ưu tiên chi phí có trần, tổng dời lớn nhất do điều phối là 29,657 giây và không khách nào vượt 30 giây.
Kết quả cho thấy cấu hình có trần giữ được giới hạn đã đặt khi thay thứ tự xếp hạng. Nó chưa chứng minh rằng
riêng trần đủ bảo đảm giờ đến thực tế, vì khóa, ràng buộc dịch vụ, bộ sinh và luật phục hồi vẫn cùng tham gia.

Mức trần thích hợp còn phụ thuộc tiêu chí dịch vụ. Với ưu tiên chi phí, trần 30 giây tăng số lượt được phục vụ
và không muộn quá một phút từ 2.123 lên 2.292; nếu chấp nhận muộn tới năm phút thì số lượt đáp ứng lại giảm
từ 2.466 xuống 2.364. Đây là hai cách đọc cùng kết quả đã có, không phải khảo sát mức khách chấp nhận chờ.
Trần 30 giây là điểm thực nghiệm, không được chọn làm ngưỡng tối ưu hay ngưỡng cảm nhận của hành khách.
% EV: S15-01 (5a); S15-02 (timely60, timely300); K36; K52.

\subsection{Lời hứa có dự báo và giới hạn của phép so}

Với ưu tiên chi phí có trần, tính lại lời hứa đầu theo dự báo cho 11 lượt muộn quá một phút thay vì 72,
trên cùng 2.364 hành trình; lời hứa dài thêm trung bình 7,30 giây. Với cấu hình không trần, các số lượt muộn
là 318 thay vì 379, và lời hứa dài thêm trung bình 7,02 giây. Giờ trả thực tế và việc xe chở ai không đổi
trong phép tính này. Dự báo điều chỉnh phần đường dự kiến chậm đi, còn luật C giới hạn phần dời do điều phối;
hai thành phần không thay thế nhau và không tạo thành bảo đảm vô điều kiện về giờ đến.
% EV: S15-01 (5b); F01--F04. 21 ô bật causal-no15; 13 ô còn lại giữ p0.

Các lượt ưu tiên chi phí dùng bản chạy riêng. Phép đối chiếu hành vi trên bốn cặp mẫu và việc tính lại các
chỉ số ưu tiên ổn định giúp kiểm tra khả năng so sánh, nhưng không chứng minh mọi khác biệt giữa hai bản
chạy đã được loại bỏ. Các con số là phép so thăm dò theo cấu hình trên một bộ mô phỏng và năm ngày đã xem,
không phải ước lượng tác động nhân quả riêng của một tiêu chí. Phép sàng lọc trên các phương án đã ghi của
tập tắc đường cố định được giữ ở Phụ lục~\ref{phuluc:D}; nó bổ trợ việc hiểu bộ chọn, không thay kết quả
của những lần mô phỏng đầy đủ ở đây.
% EV: STAGE1B-PLAN.md mục 1.1 và amendment mục 3; S15-01 (mục 0, 3); G0 chỉ behavioural-only.


\section{Lời hứa theo dự báo}
\label{muc:c4-dubao}

Trên 21 kịch bản ngày thường có dự báo, số khách của luật C trần 30 giây tới muộn hơn lời hứa đầu quá một phút giảm từ
68 xuống 2 trong 1.471 khách được phục vụ (Bảng~\ref{bang:c4-dubao}). Lời hứa theo dự báo (mục~\ref{muc:dubao}) được
tính lại trên các lần chạy đã ghi; mỗi cách hứa được so với lời hứa đầu của chính nó.
% EV: F17; F04; N4B8-01 (check A: 68 → 2 trên 1.471); F03 (chỉ vế phát lại); FFA:8-13

Đối chứng thứ hai là đệm tỉ lệ: mỗi lời hứa đầu được kéo dài một tỉ lệ cố định của khoảng từ lúc hứa tới giờ trả đã
hứa. Tỉ lệ được chọn theo từng khung giờ, trên chính các khách này, để tổng phần kéo dài bằng của lời hứa theo dự báo.
% EV: N4B-03 (FFA:15-17: lời hứa đệm = P0 + k_w(P0 − t0), k_w = Σ(D0 − P0)/Σ(P0 − t0) trên mọi khách C-30 của khung w ở 21 kịch bản, tính từ lời hứa; FA:2)

\begin{table}[htbp]
\centering
\caption[Lời hứa theo dự báo và đệm tỉ lệ trên ngày thường]{Khách của luật C trần 30 giây trên 21 kịch bản ngày thường có dự báo, theo ba cách hứa: lời hứa đầu không dự
báo, đệm tỉ lệ và lời hứa đầu theo dự báo; mỗi cách so với lời hứa đầu của chính nó. Khối trên là số đếm gộp theo khách;
dòng cuối là trung bình theo kịch bản của tỉ lệ trễ, mỗi kịch bản tính như nhau.}
\label{bang:c4-dubao}
\small
\begin{tabular}{@{}lrrrr@{}}
\toprule
\textbf{Nhóm kịch bản} & \textbf{Phục vụ} & \textbf{Không dự báo} & \textbf{Đệm tỉ lệ} & \textbf{Dự báo}\\
\midrule
\multicolumn{5}{@{}l}{\textit{Số khách tới muộn hơn lời hứa đầu quá 60 s}}\\
Cả 21 kịch bản & 1.471 & 68 & 22 & 2\\
Khung 7 giờ (8 kịch bản) & 444 & 60 & 17 & 1\\
Khung 8 giờ (8 kịch bản) & 609 & 8 & 5 & 1\\
Khung 17 giờ (5 kịch bản) & 418 & 0 & 0 & 0\\
\multicolumn{5}{@{}l}{\textit{Số khách tới sớm hơn lời hứa đầu quá 60 s}}\\
Cả 21 kịch bản & 1.471 & 0 & 0 & 1\\
% EV: N4B8-01 (check A: CELLS, share × served từng kịch bản, phần dư 0); N4B8-08 (FFA:139-141 tính E_pad; sửa N4B-03: đệm tỉ lệ có đo tới sớm, bằng 0); F17
\midrule
\multicolumn{2}{@{}l}{Tỉ lệ trễ quá 60 s, trung bình theo kịch bản (\%)} & 6,01 & 1,94 & 0,17\\
% EV: F04 (FA:8); N4B-03 (FA:8: pad 0,0194)
\bottomrule
\end{tabular}
\end{table}

Cách đọc bảng: ở khung 7 giờ, 60 trong 444 khách trễ quá một phút khi không dự báo, 17 khi đệm tỉ lệ và 1 khi có dự
báo.
% EV: N4B8-01

Tính trung bình theo kịch bản, tỉ lệ trễ giảm 5,84 điểm phần trăm. Gom 21 kịch bản thành 8 nhóm theo ngày và khung giờ,
khoảng tin cậy 95\% theo phân phối Student là từ 1,15 đến 10,54 điểm phần trăm. So với đệm tỉ lệ, chênh lệch là
$-1{,}77$ điểm phần trăm, khoảng từ $-3{,}62$ đến $+0{,}07$, tức có chứa 0: chưa thể nói dự báo tốt hơn cách đệm này.
Các khoảng chỉ dựa trên ba buổi sáng đã được xem trước, và mọi hồ sơ dự báo đều có ngày 12/11 trong các ngày nguồn.
% EV: F04; F05 (FA:33,36: dY −5,844 [−10,542; −1,146]; dPad −1,775 [−3,620; +0,071]; FFA cluster_t: tầng = khung giờ, cụm = ngày–khung, df = 8 − 3 = 5); F11 (FR:97-98; factor_nowcast_screen.py:21-22: 14 ← 12, 13; 15, 16 ← 12, 13, 14); FR:170-171; I05; K24

Lợi ích không riêng của luật C. Trên cùng 21 kịch bản, khi có dự báo, luật hạn chót 30 giây giảm từ 66 xuống 2 khách
trễ trong 1.471 khách được phục vụ; luật trần độ dời khách thấy 30 giây cũng giảm từ 66 xuống 2 trong 1.470 khách.
Điều phối không trần giảm từ 129 xuống 63 trong 1.512 khách. Riêng khung 17 giờ, 28 khách trễ trong 437 khách của điều
phối không trần không đổi, còn luật C có 0 trong 418 khách. Theo sổ, 86,1\% số giây trễ của 68 khách trễ khi không dự
báo thuộc phần được hồ sơ lịch sử giải thích (Hình~\ref{hinh:c4-so}); dự báo và trần nhắm vào hai phần khác nhau của sổ.
% EV: F08; K37; N4B8-02 (check A: M⁺-30 66 → 2 trên 1.471; V-30 66 → 2 trên 1.470; U 129 → 63 trên 1.512; U khung 17 giờ 28 → 28 trên 437; C-30 khung 17 giờ 0 trên 418); F09 (FA:30; RX:3); eq:so; K04

Trên 21 kịch bản này, giá của dự báo là lời hứa giờ trả dài thêm 13,7 giây (khung 7 giờ 27,8 giây) và lời hứa giờ đón 3,2 giây, trung bình
theo kịch bản. Trên cả 34 kịch bản, khách trễ của luật C giảm từ 76 xuống 10 trong 2.338 (Bảng~\ref{bang:c4-tran120}).
% EV: F06 (FR:67-69; FA:8-9); C06 (76); C09 (10 trên 34 kịch bản); N-16; F02; N-28 (13 kịch bản cuối tuần không có dự báo; độ dài thêm gộp 7,37 s, N4B-10, nằm ở bang:c4-tran120 và bang:c4-dem)

Kết quả 2 trong 1.471 khách đạt được trong điều kiện thuận lợi cho dự báo theo giờ trong ngày, không phải hiệu quả ngoài
đời. Tắc đường là một hệ số chung cho cả thành phố. Trên cả 34 kịch bản, với luật C, kéo dài lời hứa theo đúng hệ số
sẽ xảy ra cho giờ trả khớp trong vòng 2 giây ở cả 2.305 khách không bị điều phối dời. Phép kiểm giờ đón cũng chưa dùng
lời hứa này (mục~\ref{muc:c4-thaoluan}).
% EV: I02 (FR:167-169: mức tốt nhất cho loại dự báo theo giờ trong ngày trong mô phỏng gần như tất định); F10 (FA:28; RX:8: 34 kịch bản); B06; F07; X3-11

\begin{figure}[htbp]
\centering
\resizebox{\linewidth}{!}{%
\begin{tikzpicture}[x=1cm,y=1cm,font=\footnotesize,
  nhan/.style={font=\scriptsize,inner sep=1.5pt}]
% Thanh trên: ngày thường. Độ dài 11 cm = 6.217,4 s; ranh giới 0,1445 và 9,6145 cm (FA:30).
\draw[fill=black!75] (0,1.6) rectangle (0.1445,2.2);
\draw[fill=black!20] (0.1445,1.6) rectangle (9.6145,2.2);
\draw[fill=white] (9.6145,1.6) rectangle (11,2.2);
\node at (4.88,1.9) {5.353 s (86,1\%)};
\node[align=center,nhan] at (10.307,1.9) {783 s\\(12,6\%)};
\draw (0.072,2.2)--(0.072,2.45);
\node[nhan,anchor=south west] at (0.0,2.45) {82 s (1,3\%)};
\node[anchor=east,align=right] at (-0.15,1.9) {Ngày thường, 21 kịch bản\\68 khách, 6.217 s};
% Thanh dưới: chiều tuyết. Độ dài 11 cm = 3.445,9 s; ranh giới 0,0273 và 0,8302 cm (NEW:5).
\draw[fill=black!75] (0,0) rectangle (0.0273,0.6);
\draw[fill=black!20] (0.0273,0) rectangle (0.8302,0.6);
\draw[fill=white] (0.8302,0) rectangle (11,0.6);
\node at (5.915,0.3) {3.185,9 s (92,5\%)};
% v8: nhãn 8,6 s đặt phía trên thanh tuyết (khoảng trống giữa hai thanh) để đường dẫn của nhãn 251,5 s không cắt chữ.
\draw (0.014,0.6)--(0.014,0.85);
\node[nhan,anchor=south west] at (-0.02,0.85) {8,6 s (0,25\%)};
\draw (0.43,0)--(0.43,-0.2);
\node[nhan,anchor=north west] at (0.39,-0.2) {251,5 s (7,3\%)};
\node[anchor=east,align=right] at (-0.15,0.3) {Chiều tuyết, 8 kịch bản\\41 khách, 3.445,9 s};
% Chú giải (ba ô cách nhau đủ để chữ không chồng lên nhau)
\draw[fill=black!75] (0,-1.25) rectangle (0.3,-1.0);
\node[anchor=west] at (0.35,-1.125) {phần tính cho điều phối};
\draw[fill=black!20] (4.0,-1.25) rectangle (4.3,-1.0);
\node[anchor=west] at (4.35,-1.125) {phần được hồ sơ lịch sử giải thích};
\draw[fill=white] (9.6,-1.25) rectangle (9.9,-1.0);
\node[anchor=west] at (9.95,-1.125) {phần còn lại};
\end{tikzpicture}
}
\caption[Ba phần của số giây trễ theo sổ của từng khách]{Số giây trễ chia thành ba phần theo sổ của từng khách (công thức~\eqref{eq:so}), với luật C trần 30 giây và các
khách tới muộn hơn lời hứa đầu không dự báo quá 60 giây: 68 khách của 21 kịch bản ngày thường và 41 khách của 8 kịch
bản chiều tuyết 15/11/2018. Mỗi thanh dài bằng 100\% số giây trễ của nhóm đó. Phần được hồ sơ lịch sử giải thích phụ
thuộc dự báo đang dùng; phần còn lại là phần dự báo này chưa giải thích được, không phải phần không thể biết trước.}
\label{hinh:c4-so}
\end{figure}
% EV (ngày thường): F09; N4B8-11 (FA:30 và CELLS: 81,742 + 5.352,646 + 783,018 = 6.217,406 s; các nhãn làm tròn đến giây độc lập, nên tổng ba nhãn là 6.218 s, lệch 1 s so với nhãn tổng 6.217 s); RX:3 (1,3 / 86,1 / 12,6%)
% EV (chiều tuyết): F13; N4B-05 (NEW:5: 8,566 / 251,496 / 3.445,933 s; RP1:24-26: 7,3%; SNOW:10: 92,5%; FR:104,186); N4B8-12 ([A] 8,566 / 3.445,933 = 0,25%)
% Tọa độ: x = giây / tổng × 11 cm (tính tay; xem NOTES v7 mục 4).

\textbf{Đưa dự báo vào bộ điều phối.} Phép thử chạy luật C và luật hạn chót 30 giây trên 16 kịch bản khung 7 và 8 giờ.
Thời gian đi của bộ điều phối được nhân thêm mức tăng dự báo trung bình của khung 15 phút hiện tại và khung kế tiếp. Có
3 trong 32 lần chạy dừng vì FleetPy tính rằng xe sẽ tới điểm đón muộn hơn giờ đón muộn nhất chưa tới 3 giây. Luật C
không có quyết định bị ép ở 14 lần chạy hoàn tất. Luật hạn chót bị ép ở 47,6\% lượt quyết định ở 15 lần chạy hoàn tất,
so với 62,8\% ở các lần chạy không đưa dự báo vào trên cùng 15 kịch bản. Tính gộp trên cùng 15 kịch bản, đó là 3.849
trong 8.496 lượt quyết định khi đưa dự báo vào, so với 5.121 trong 8.468 lượt khi không đưa vào; các tỉ lệ 47,6\% và
62,8\% ở trên là trung bình theo kịch bản. Tỉ lệ khách bị luật hạn chót phá lời hứa ở khung 7 giờ giảm từ 39,0\% xuống
11,0\%, trung bình trên cùng tám kịch bản. Tính gộp theo khách trên tám kịch bản này, đó là 170 trong 442 khách khi
không đưa dự báo vào, so với 45 trong 439 khách khi đưa vào. Vì vậy, các số của Bảng~\ref{bang:c4-nghiavu} phụ thuộc cách
hứa. Khóa luận chỉ dùng dự báo ở lời hứa báo cho khách; kết luận này chỉ áp cho cách đã thử.
% EV: F12; X3-29 (W13:85-89); N4B8-06 (K = 2: run_forecast.py:9,72; configs/forecast-k2-2018-11-1{4,5,6}.json:1; preflight_world_v1_3.py:85-89; FR:117; R26 Q9); N4B-07 (RX:9,11-15; FR:120-121,124,138,157: C hoàn tất 14, luật hạn chót 15; 47,6 và 62,8% trên cùng 15 kịch bản); C14; FR:141-146; R37 verify-numbers-v2.json:dispatcher; R39 agent-numbers-v1.json:dispatcher.paired, morning_promise_breaks (cùng15 hoàn tất; khung7 cùng8; số gộp tách khỏi trung bình theo kịch bản)

\textbf{Học máy.} Ba mô hình học một lần rồi giữ nguyên ước lượng phân vị 90 của độ trễ so với lời hứa theo dự báo (mức
mà 90\% khách giống vậy không vượt) từ thông tin đã có ở thời điểm hứa. Đó là mô hình cây quyết định tăng cường
(gradient boosting), mô hình ấy khi bỏ hai biến trạng thái xe (số khách khác đã gán và đang trên xe), và hồi quy phân vị
tuyến tính. Mỗi khách được cộng vào lời hứa theo dự báo một phần đệm tỉ lệ với mức ước lượng; đối chứng cộng cùng một
khoảng cho mọi khách. Cả hai được chỉnh trên ngày hiệu chỉnh để lời hứa dài thêm trung bình 30 giây so với lời hứa không
dự báo, nên trên ngày kiểm mức này có thể khác. Phép thử dùng hồ sơ dự báo cố định của ngày
12--13/11 và kiểm trên 15, 16/11 khung 7 giờ và 13/11 của tập tắc đường cố định, các ngày đã được xem trước. Mô hình
chỉ được coi là tốt hơn khi, so với đối chứng, thỏa cả bốn điều kiện:
\begin{itemize}[leftmargin=1.2cm,itemsep=0pt,topsep=2pt]
\item tỉ lệ khách trễ quá một phút thấp hơn;
\item độ dài thêm trung bình không hơn quá 1 giây;
\item phân vị 90 của độ dài thêm không hơn quá 15 giây;
\item tỉ lệ khách tới sớm quá một phút không hơn quá 1 điểm phần trăm.
\end{itemize}
% EV: G01; G02; G03; N4B-08; N4B8-04 (MLS:9-10 các ngày đã xem; 19-23, 67-68 hai biến trạng thái; 37-39 hồ sơ cố định 12–13/11; 41-61 chia ngày; 63-69 ba mô hình; 71-82 chỉnh trên ngày hiệu chỉnh, F = D0 + ρ·w, đối chứng D0 + (E − mean δ), không chỉnh lại khi kiểm; 91-95 bốn điều kiện; MLT: đối chứng của U dài thêm 37,219 / 31,772 s ở 15/11 / 16/11)

Không mô hình nào thỏa đủ bốn điều kiện. Các số đếm sau được suy từ tỉ lệ gộp theo khách đã làm tròn đến ba chữ số
thập phân trong bảng kết quả; với mẫu số tương ứng, mỗi tỉ lệ chỉ ứng với một số nguyên. Với luật C, đối chứng và mô hình đều không có khách trễ (147 và 197 khách).
Với điều phối không trần, mô hình cây tăng cường chỉ bớt 1 trong 2 khách trễ ở ngày 15/11 (145 khách), không bớt ai ở
ngày 16/11 (200 khách), và không thỏa điều kiện thứ hai, thứ ba. Với ghép tự do, khách trễ giảm từ 104 xuống 94 trong 807
nhưng 57 khách tới sớm quá một phút, so với 0 khi đệm đều. Bỏ hai biến trạng thái xe không đổi số khách trễ ở hai phép
so này.
% EV: G05–G10; N4B8-03 (check B: C-30 0/147, 0/197; U 2 → 1 trên 145, 1 → 1 trên 200; FREE 104 → 94 trên 807, sớm 57 so với 0; bỏ trạng thái xe 1/145 và 94/807); ML:6-7,14-15,21,45-53,67-69 (độ dài thêm +4,4 / +3,8 s; phân vị 90 +21,8 / +21,7 s); K40; R37 verify-numbers-v2.json:ml; R39 agent-numbers-v1.json:ml (xét mọi số nguyên, sai số 0,0005 điểm phần trăm); ml_pilot_run.py:37-42,175-176

% -----------------------------------------------------------------------------------------------------
\section{So với cách đệm lời hứa}
\label{muc:c4-dem}

Khi có dự báo, cộng thêm 30 giây vào mọi lời hứa của điều phối không trần cũng cho 2.328 lượt không trễ quá một phút
trên 3.672 yêu cầu, bằng luật C trần 30 giây có dự báo. Ở ngưỡng hai phút, cách đệm này cho 2.366 lượt, so với 2.338
của luật C. Hai cách này cho hai thứ khác nhau. Đệm kéo dài lời hứa, trung bình 37,2 thay vì 7,4 giây, nên ít khách bị tính là trễ. Nhưng đệm không đổi
quyết định điều phối: vẫn có 174 khách bị điều phối dời giờ trả, với tổng độ dời lớn nhất là 590,1 giây (gần 10 phút).
Luật C trần 30 giây giới hạn tổng độ dời do điều phối của mỗi khách ở 30 giây dưới giả thiết (N); trên 34 kịch bản của
tập kiểm tra chính, mức lớn nhất đo được là 29,2 giây, kể cả ở các kịch bản cuối tuần không có dự báo. Đổi lại, luật C
phục vụ ít hơn 96 lượt.
% EV: C10; N-16; N-18; K06; K51; C20; N-23; N-28 (cuối tuần 29,2 s); C02; A06 (dưới giả thiết (N); không bảo đảm giờ đến: A22, K05); RD §8; khớp câu trỏ của phần A v5-v3 mục 4.2 (C10); R37 verify-numbers-v2.json:padding; R39 agent-numbers-v1.json:padding,cohort_totals

Cách đệm là phát lại trên các lần chạy của điều phối không trần, nên hai bên được so trên hai nhóm khách khác nhau, chỉ
chung 2.207 khách. Bảng~\ref{bang:c4-dem} xếp các cách hứa thành ba khối có lời hứa dài thêm gần bằng nhau.
% EV: C10 (phát lại, không phải đối chứng ghép cặp); ADDENDUM-1 phần đầu (quy tắc đệm D = D0 + k); C04

\begin{table}[htbp]
\centering
\caption[Trần và cách đệm lời hứa khi có dự báo]{Các cách hứa khi có dự báo trên tập kiểm tra chính (34 kịch bản, 3.672 yêu cầu mỗi cấu hình; dự báo chỉ có ở
21 kịch bản ngày thường). U: điều phối không trần; C-$B$: luật C trần $B$ giây; trong ngoặc là số khách được phục vụ.
Số đếm gộp; trễ và tới sớm tính trong số khách được phục vụ, so với lời hứa của chính cách hứa; hai cột ``không trễ
quá'' tính trên 3.672 yêu cầu; độ dài thêm là trung bình gộp theo khách được phục vụ. Đệm không đổi quyết định điều
phối, nên hàng của U và hàng của C là các nhóm khách khác nhau.}
\label{bang:c4-dem}
\small
\begin{tabular}{@{}p{4.6cm}rrrrr@{}}
\toprule
\textbf{Cách hứa} &
\textbf{\begin{tabular}[b]{@{}r@{}}Trễ\\ quá 60 s\end{tabular}} &
\textbf{\begin{tabular}[b]{@{}r@{}}Không trễ\\ quá 60 s\end{tabular}} &
\textbf{\begin{tabular}[b]{@{}r@{}}Không trễ\\ quá 120 s\end{tabular}} &
\textbf{\begin{tabular}[b]{@{}r@{}}Tới sớm\\ quá 60 s\end{tabular}} &
\textbf{\begin{tabular}[b]{@{}r@{}}Dài thêm\\ (s)\end{tabular}}\\
\midrule
\multicolumn{6}{@{}l}{\textit{Không đệm (dài thêm khoảng 7 giây)}}\\
U, dự báo (2.434) & 137 & 2.297 & 2.347 & 8 & 7,18\\
C-30, dự báo (2.338) & 10 & 2.328 & 2.338 & 4 & 7,37\\
C-120, dự báo (2.387) & 65 & 2.322 & 2.379 & 2 & 7,27\\
% EV: N-16; N-03; C09; N-06; L1a-03; N4B-09 (sớm 8 / 4 / 2)
\midrule
\multicolumn{6}{@{}l}{\textit{Đệm 30 giây (dài thêm khoảng 37 giây)}}\\
U, dự báo, đệm 30 s (2.434) & 106 & 2.328 & 2.366 & 44 & 37,18\\
C-30, dự báo, đệm 30 s (2.338) & 2 & 2.336 & 2.338 & 44 & 37,37\\
% EV: C10; N-02; N-16
\midrule
\multicolumn{6}{@{}l}{\textit{Đệm 60 giây (dài thêm khoảng 67 giây)}}\\
U, dự báo, đệm 60 s (2.434) & 87 & 2.347 & 2.386 & 1.033 & 67,18\\
% EV: C11; N-16
\bottomrule
\end{tabular}
\end{table}

Cách đọc bảng: với đệm 30 giây trên điều phối không trần, 106 trong 2.434 khách tới muộn hơn lời hứa đã đệm quá 60
giây, nên có 2.328 lượt không trễ quá 60 giây trên 3.672 yêu cầu; 44 khách tới sớm quá 60 giây.
% EV: C10; N-16

\textbf{Vai trò của từng thành phần.} Thứ nhất, ở mức kéo dài lời hứa gần bằng nhau, luật C trần 30 giây hơn cách đệm
ở ngưỡng một phút và kém ở ngưỡng hai phút: khi cùng đệm 30 giây, hai bên có 2.336 so với 2.328 lượt và 2.338 so với
2.366 lượt. Ở hàng đệm 30 giây, điều phối không trần đạt cùng 2.328 lượt không trễ quá một phút như luật C có dự báo
không đệm, với lời hứa dài thêm trung bình 37,18 so với 7,37 giây, tức khoảng năm lần.
% EV: N-17; N-18 (37,1833 / 7,3671 = 5,05); N-16; R37 verify-numbers-v2.json:padding; R39 agent-numbers-v1.json:padding,cohort_totals

Thứ hai, lời hứa dài thì ít trễ nhưng nhiều khách tới sớm: 44 khách với cách đệm 30 giây và 1.033 khách với cách đệm 60
giây, so với 4 khách với luật C có dự báo. Số lượt tới trong khoảng $\pm$60 giây quanh lời hứa là 2.284 với cách đệm 30
giây, so với 2.324 với luật C, trên 3.672 yêu cầu.
% EV: C10; C11; N-04 (±60 s = phục vụ − trễ60 − sớm60, ADD1:92); N4B-09 (ADD1:84-89)

Thứ ba, cách đệm 30 giây vẫn có 68 khách trễ quá hai phút, 11 khách quá năm phút và 1 khách quá mười phút; luật C có
dự báo không có ai trễ quá hai phút. Cách đệm hơn 28 lượt ở ngưỡng hai phút vì phục vụ nhiều hơn 96 lượt, không phải vì
ít khách trễ hơn.
% EV: N-01 (68 / 11 / 1 so với 0 / 0 / 0; ADD1:15,19); N-19 (2.366 − 2.338 = 28 = 96 − 68)

Khi phát lại cùng lời hứa theo dự báo, luật hạn chót 30 giây cũng có 2.338 lượt phục vụ và 2.328 lượt không trễ quá
một phút. Vì phép phát lại không đổi các lần chạy, nghĩa vụ của luật này vẫn tính từ lời hứa đầu không dự báo. Theo
nghĩa vụ đó, tính trung bình theo kịch bản, nó bị ép ở 41,3\% lượt quyết định; tính gộp là 7.710 trong 20.103 lượt.
Luật hạn chót có khách tới muộn hơn lời hứa đầu không dự báo quá 30 giây ở 28 trong 34 kịch bản, còn luật C không có
quyết định bị ép hay khách vượt trần (Bảng~\ref{bang:c4-nghiavu}).
% EV: N-16 (M⁺-30 + dự báo 2.338 / 10 / 2.328 / 2.338, ADD1:29); C13 (41,3% trung bình theo kịch bản); N-21 (phá hứa của M⁺-30 = A − p₀ > 30 s với lời hứa không dự báo; C-30 0/34, 0 quyết định bị ép); X3-16; R36-A01 (bảng A-v5-v3 dùng số gộp); R39 agent-numbers-v1.json:main_original_obligation

Trần 120 giây có dự báo, với lời hứa chỉ dài thêm 7,27 giây, cho 2.379 lượt không trễ quá hai phút, nhiều hơn 2.366
của cách đệm 30 giây, nhưng 2.322 lượt ở ngưỡng một phút, ít hơn 2.328. Đây là một điểm đo thêm trên một tập khách khác,
không phải mức trần tốt nhất.
% EV: N-06; N-16; N-18; N-10; K36

\textbf{Khi nào chọn trần, khi nào chọn đệm.} Nếu hãng xe chỉ tính số lượt không trễ quá hai phút và chấp nhận lời hứa
dài hơn cùng nhiều khách tới sớm hơn, điều phối không trần có dự báo và đệm 30 giây đạt 2.366 lượt, so với 2.338 lượt
của luật C trần 30 giây có dự báo, trên 3.672 yêu cầu; điều phối không trần cũng phục vụ thêm 96 lượt. Nếu cần giới hạn
tổng độ dời do điều phối của từng khách, kể cả khi không có dự báo, luật C trần 30 giây làm được điều này dưới giả thiết
(N), đổi lại 96 lượt phục vụ và 242 khách đi chung. Hai cách dùng chung được; chọn cách nào là quyết định của hãng xe,
không phải kết luận của thực nghiệm.
% EV: N-19; C02; N-20; N-02 (hàng C-30 có dự báo và đệm 30 s); D20; I07; I10; K38; K51; A22 (giả thiết N); R37 verify-numbers-v2.json:padding; R39 agent-numbers-v1.json:padding,cohort_totals

% -----------------------------------------------------------------------------------------------------
\section{Khi dự báo không đủ: chiều tuyết 15/11/2018}
\label{muc:c4-tuyet}

Trên 601 khách được phục vụ của 8 kịch bản chiều tuyết 15/11/2018 (luật C trần 30 giây), lời hứa theo dự báo chỉ giảm
số khách trễ quá một phút từ 41 xuống 34. Tính trung bình theo kịch bản, tỉ lệ này giảm từ 7,19\% xuống 5,94\%. Theo
sổ, trong 3.445,9 giây trễ của 41 khách, 8,6 giây tính cho điều phối và 251,5 giây được hồ sơ lịch sử giải thích
(Hình~\ref{hinh:c4-so}).
% EV: F13; K25; B20; N4B-05 (NEW:5)

Bảng~\ref{bang:c4-tuyet} cho kết quả phát lại vài cách hứa dùng thêm hệ số tắc đường của chính ngày tuyết. Cách ``trộn'' lấy trung
bình nhân của mức tăng theo hồ sơ lịch sử và mức tăng của hệ số cùng ngày từ khung 15 phút trước sang khung hiện tại.
% EV: F14; K20; K21; N4B-06 (factor_nowcast_screen.py:14-20 quy tắc thông tin, 28-29: MOM = (f(k0)/f(k0−1))^(b−b0), BLEND = √(HIST·MOM); RP1:7-20; NEW:5)

\begin{table}[htbp]
\centering
\caption[Phát lại các cách hứa trong chiều tuyết]{Phát lại các cách hứa trên 8 kịch bản chiều tuyết 15/11/2018 (luật C trần 30 giây, 601 khách được phục vụ).
Cột ``khách trễ'' đếm gộp số khách tới muộn hơn lời hứa đầu quá 60 giây; cột cuối là độ dài thêm trung bình theo kịch
bản. Không cách hứa nào trong bảng có khách tới sớm quá 60 giây. Mọi cách hứa đều tính lại sau khi chạy, trên một ngày
đã được xem trước khi thiết kế các cách này; hai hàng ``có buổi sáng'' chỉ để tham khảo.}
\label{bang:c4-tuyet}
\small
\begin{tabular}{@{}p{3.4cm}p{6.6cm}rr@{}}
\toprule
\textbf{Cách hứa} & \textbf{Thông tin dùng lúc hứa} &
\textbf{\begin{tabular}[b]{@{}r@{}}Khách\\ trễ\end{tabular}} &
\textbf{\begin{tabular}[b]{@{}r@{}}Dài thêm\\ (s)\end{tabular}}\\
\midrule
Không dự báo & Thời gian đi hiện tại & 41 & 0\\
Theo dự báo & Hồ sơ tắc đường của các ngày nguồn & 34 & 2,0\\
Trộn trong khung & Thêm hệ số cùng ngày đã có trong khung giờ & 21 & 8,5\\
Trộn có buổi sáng & Thêm hệ số cùng ngày từ 6 giờ; thiết kế thử không cho khung chiều dùng dữ liệu này & 19 & 9,5\\
Trộn có buổi sáng, trễ một khung & Như hàng trên, nhưng mức tăng gần nhất chỉ lấy từ các khung đã kết thúc; hệ số của
khung hiện tại vẫn dùng giá trị đúng & 29 & 9,0\\
Đệm tỉ lệ & Tỉ lệ chọn để tổng phần kéo dài bằng của hàng ``trộn trong khung'', trên chính ngày tuyết & 27 & 8,3\\
Hệ số thật & Hệ số sau lúc hứa, không có ở thời điểm hứa & 0 & 18,8\\
\bottomrule
\end{tabular}
\end{table}
% EV: F14 (RP1:7-20; NEW:5); N4B-06 (định nghĩa; cột dài thêm là trung bình theo kịch bản, RP1:6, rider_nowcast_replay.py:96-106); K20; K21; factor_nowcast_screen.py:17-20 (hàng trễ một khung cũng dùng thông tin từ 06:00); N4B8-09 (RP1:7-20: E60 0,00% ở mọi hàng ⇒ 0 khách tới sớm ở mọi kịch bản [A])

Cách đọc bảng: khi trộn hồ sơ lịch sử với xu hướng trong khung giờ, 21 trong 601 khách tới muộn hơn lời hứa quá một
phút, với lời hứa dài thêm trung bình 8,5 giây.
% EV: F14; N4B-06

Phép phát lại này tách số giây trễ theo cách cộng có giới hạn, khác phần có dấu của công thức~\eqref{eq:so}. Phần kéo
dài của mỗi khách được chặn trong khoảng từ 0 đến phần trễ còn lại sau khi trừ phần tính cho điều phối. Theo cách cộng
này, cách trộn
có buổi sáng giải thích thêm 886,9 giây, tức 25,7\% của 3.445,9 giây trễ; 2.299,0 giây (66,7\%) chưa được giải thích
bằng các cách đã thử. Theo số liệu trạm Central Park của Cơ quan Quản lý Khí quyển và Đại dương Quốc gia Hoa Kỳ
(National Oceanic and Atmospheric Administration, NOAA), từ 9/2018 đến 12/2019 chỉ có 5 đợt tuyết từ 25,4 mm trở lên,
kể cả đợt này; muốn kiểm cách hứa cho ngày bất thường cần thêm nhiều đợt khác.
% EV: F15 (NEW:5; RP1:24-35); N4B8-07 (rider_nowcast_replay.py:14-15, 113: explained = Σ clip(D − P0, 0, A − P0 − dispatch); RP1:29: phần vượt 6 s; với hồ sơ lịch sử phần vượt 0 nên hai cách cộng trùng nhau, N4B-05); K19; I11; B26 (NEW:9; WX:1-3: tệp NOAA GHCN-D trạm Central Park; WX17:1-4; ghép ngày cách nhau ≤ 2)

% -----------------------------------------------------------------------------------------------------
\section{Kiểm tra tính hợp lệ của kết quả}
\label{muc:c4-hople}

Trên 34 kịch bản của tập kiểm tra chính, không khách nào vượt trần: tổng dời lớn nhất là 29,2 giây với trần 30 giây và
118,9 giây với trần 120 giây. Không lượt quyết định nào bị ép, 0 trong 20.101 và 0 trong 20.435 lượt. Đây là phép kiểm cài
đặt: với sổ hợp lệ $0 \le L_r \le B$ ở lượt đang xét và giả thiết (N), Nhận xét~\ref{nx:giunguyen} đã cho trước các
số 0 này; kết luận xuyên chuyến cần giả thiết (N) đúng ở mọi lượt. Điều phối không trần cũng có 0 trong 20.760 lượt.
% EV: C20; N-07; X3-16 (0/20.101; 0/20.435; U 0/20.760); N-21; A22; chương 3, nx:giunguyen và đoạn sau nó (0 ≤ L_r ≤ B ở lượt đang xét; (N) đúng ở mọi lượt cho kết luận xuyên chuyến); R39 agent-numbers-v1.json:cohort_totals

Bộ đếm của bộ sinh được đọc ở điều phối không trần, luật C trần 30, 120 giây và luật hạn chót 30 giây. Ở cả bốn cấu
hình, không phương án nào bị bỏ vì giới hạn 100 phương án mỗi xe hay giới hạn số bước tìm cách chèn. Giới hạn 2 yêu cầu
mới mỗi lượt chỉ để một yêu cầu sang lượt sau ở mỗi cấu hình. Các số này không đo mọi giới hạn tính toán (mục~\ref{muc:gioihan}): giới hạn
của bước kiểm, của bộ giải, một bộ đếm khác của bộ sinh và các bộ đếm ở tập tắc đường cố định chưa được đọc.
% EV: N-25 (ADD1:119-122: 34 lần chạy mỗi cấu hình); N4B-11 (GEN:54-95, 236-246, 298-345; CandidateModels.cs:68-82, 100-118: điều kiện đầy đủ còn gồm OmittedRepairRequestCount, chưa đọc); X3-17; X3-03; X3-32; C4A-18 (tập tắc đường cố định: giới hạn chưa kiểm); K34 thu hẹp (N4B8-13)

Như mục~\ref{muc:doichung} đã nêu, mức dùng trần lớn nhất của tiêu chí 3 bằng 100\% ở mọi lượt có ghi tiêu chí này:
3.503, 3.506 và 3.504 lượt khi không trần, với trần 30 và 120 giây.
% EV: A15; N-25 (ADD1:119-121); A07; X3-07 (chương 3 mục doichung nêu hệ quả: tiêu chí 3 chỉ tác dụng qua tổng phần dời)

Các số đếm chính của tập kiểm tra chính được ít nhất hai chương trình đọc dữ liệu viết riêng tính lại và khớp nhau
(mục~\ref{muc:kiemtra}). Phép so gồm cả từng khách ở 7.159 khách được phục vụ của điều phối không trần và luật C trần
30, 120 giây; riêng số khách đi chung chỉ được một quy tắc đếm. Theo bản ghi khởi tạo, 8 lần chạy chiều tuyết dùng cùng
bản bộ chạy đã lưu và cùng cấu hình với các lần chạy luật C trần 30 giây của tập kiểm tra chính, kể cả luật phục hồi
không tệ hơn.
% EV: N4B-14 (ADD1:66-78); K41; C19; N-09 (đi chung: Codex không đếm lại); N4B8-05 (check C: 8 lần chạy chiều tuyết và 51 lần chạy C-30 của T3.8 cùng binarySha256 a631038b…, coreCommitSha 1ce15e92…, policyConfigurationHash a2f11a3a…; sha256 runner-tier3-v2/RideBound.Runner.dll khớp; wp4-tier3-v2.json:7 no-worse-than-reference-v1; R26 Q8)

Phép kiểm số đếm học máy chỉ đối chiếu số học của bảng kết quả; nó không kiểm lại quá trình học hay dự đoán của mô hình.
% EV: R37 §1.8; R39 agent-numbers-v1.py phần C; agent-numbers-v1.json:ml; ml_pilot_run.py:37-42,175-176

Chạy lại trên đúng mã nguồn đã dựng bộ chạy, cả 1.025 trong 1.025 kiểm thử tự động đạt trong một lần chạy đầy đủ ngày
13/10/2026, với kiểu xuống dòng LF và một tệp chưa ghi nhận trong cây làm việc (Phụ lục~\ref{phuluc:C}). Các phép kiểm
trên không cho biết mô phỏng giống thực tế đến đâu.
% EV: TESTRUN-45f8026-2026-10-12.md:18 (attempt 4: LF và tệp chưa ghi nhận, 1.025/1.025 trong một lần chạy ngày 13/10/2026); C:/RideBoundData/research/testrun-45f8026-eollf-dirty-fullrun-2026-10-13.log; R39 agent-numbers-v1.json:test_attempt4; I13; A28; I01; K41

% -----------------------------------------------------------------------------------------------------
\section{Thảo luận và hạn chế}
\label{muc:c4-thaoluan}

Với một hãng xe, các kết quả trên là các điểm trên một đường đánh đổi, không phải một khuyến nghị. So với điều phối
không trần, luật C trần 30 giây phục vụ ít hơn 96 lượt và có ít hơn 242 khách đi chung; đổi lại, phần dời do điều phối
của mỗi khách không quá 30 giây và số lượt trễ thấp hơn ở cả bốn ngưỡng đã đo, trên hai nhóm khách khác nhau
(Bảng~\ref{bang:c4-phanbotre}). Khóa luận không đo giá trị của các điều này với khách hay hãng xe.
% EV: C02; N-20; C20; N-24; C04; D20; I10; K38 (không kể 0 quyết định bị ép là phần luật C đổi được: X3-16, U cũng 0/20.760)

Một phép tính hòa vốn giả định cho thấy cách đọc các điểm này. Gọi $\lambda$ là chi phí của hãng xe cho mỗi lượt trễ quá
một phút và $\mu$ là biên lợi nhuận mỗi chuyến. Với 96 lượt phục vụ ít hơn và 127 lượt trễ ít hơn, luật C chỉ có lợi hơn
khi $\lambda > 0{,}756\,\mu$, tức $96/127$. Phép tính bỏ qua mọi yếu tố khác và so hai nhóm khách khác nhau; khóa luận
không có giá trị thật nào của $\lambda$ hay $\mu$.
% EV: C29 (R15:231-235; 96/127 = 0,756); K35; C28; C04; I10

Các hạn chế dưới đây chia theo bốn nhóm; mỗi ý nêu điều chưa làm và kết luận chịu ảnh hưởng.

\textbf{Mô hình mô phỏng.}
\begin{itemize}[leftmargin=1.2cm,itemsep=2pt,topsep=3pt]
\item Mọi kết quả đến từ FleetPy với một hệ số tắc đường chung cho cả thành phố, mức nền lấy lúc 00:00, nên mô phỏng gần
như không có yếu tố ngẫu nhiên. Các kết quả, kể cả mức trễ còn lại sau dự báo, không cho biết hiệu quả khi vận hành thật.
\item Thời gian dừng đón, trả bằng 0 giây (kiểm trên 3 lần chạy mẫu), so với 30 giây mỗi khách trong nghiên cứu điển hình
của FleetPy~\cite{engelhardt2026}, và đồ thị rút gọn chỉ có 96 nút. Phần dời do điều phối và cái giá 96 lượt vì vậy chỉ áp cho
thiết lập này.
\item Không có phản ứng của khách và không gán lại xe cho khách đã nhận, nên số lượt phục vụ không tính khách bỏ đi, và
chưa biết cái giá của trần khi được gán lại xe.
\end{itemize}
% EV: I01; B01; B06; B07; I02; F10; B05 (SERV:2-4: kiểm trên 3 lần chạy mẫu; FLEET:423-424); J18; B03; K48; B12; I07; X3-39

\textbf{Dữ liệu và thiết kế thực nghiệm.}
\begin{itemize}[leftmargin=1.2cm,itemsep=2pt,topsep=3pt]
\item 34 kịch bản của tập kiểm tra chính phụ thuộc nhau, thuộc 5 ngày, mỗi kịch bản chạy một lần; các phép đếm theo
kịch bản là mô tả, không phải kiểm định.
\item Chưa kiểm độ nhạy theo hạt giống của bộ giải, nên chưa biết các chênh lệch có giữ nguyên khi đổi hạt giống.
% EV: R37 §1.8/L-26; B09; B11; chương 3 v4-v3 muc:gioihan (hạt giống cố định 7, một lần chạy mỗi kịch bản)
\item Các kết quả dùng lại dữ liệu đã xem. Ngưỡng phân tích và cách chia tập được chọn sau khi xem tập buổi sáng và các
đợt chạy thử trước đó. Mức 120 giây được chọn sau khi xem kết quả của điều phối không trần và luật C trần 30 giây trên
chính tập kiểm tra chính, kể cả các phép tính thử ở 120 giây (mục~\ref{muc:c4-thamdo}). Chiều tuyết và các ngày kiểm học máy cũng đã được xem, nên mọi kết quả là thăm dò.
\item Tải thấp: 0,535 khách trên xe kể cả lúc xe rỗi khi không trần, so với khoảng 1,33 trong nghiên cứu điển hình của
FleetPy. Ở 49 kịch bản dựng từ dữ liệu thật, chuyến được chọn ngắn hơn mọi chuyến đủ điều kiện (trung vị thời gian đi thẳng 165
so với 196 giây). Kết quả không suy ra cho tải cao hay chuyến dài hơn.
\end{itemize}
% EV: B09; I05; B11; I06; C4A-11 (mức 120 giây, như phần A v5-v3 mục c4-thamdo); MLS:9-10; I11; B12 (LOAD:2-3; FLEET:447-449); J18; B13 (SEL6:2-4: 49/49 kịch bản); K42 (kể cả các phép tính thử ở 120 giây); chương 3 v4-v3 muc:gioihan (hạt giống7; chưa có kiểm độ nhạy)

\textbf{Tích hợp.}
\begin{itemize}[leftmargin=1.2cm,itemsep=2pt,topsep=3pt]
\item Bộ điều phối vẫn kiểm giờ đón bằng giờ không dự báo, nên 34 trong 1.471 khách được báo giờ đón muộn hơn giờ đón
muộn nhất. Các số có dự báo ở mục~\ref{muc:c4-dubao} và~\ref{muc:c4-dem} chưa tính việc sửa phép kiểm này.
\item Bản chạy ưu tiên chi phí khác bản lưu của ưu tiên ổn định. Bốn cặp đối chiếu hành vi và việc tái lập các
chỉ số nền không chứng minh tương đương toàn bộ hai bản. Kết quả ở mục~\ref{muc:c4-vaitro} vì vậy vẫn là phép so
thăm dò theo cấu hình, không loại trừ mọi khác biệt thời gian chạy.
\end{itemize}
% EV: F07 (RX:2; FR:70); I13; A28; R18 B24 (qua REG A28)

\textbf{Phạm vi kết luận.}
\begin{itemize}[leftmargin=1.2cm,itemsep=2pt,topsep=3pt]
\item Phép so hai thứ tự chọn phương án đã có kết quả ở trần 30 giây và không trần, nhưng chưa tạo thành
đường đánh đổi nhiều mức cho ưu tiên chi phí. Luật giới hạn từng lần chưa được cài đặt; các bản ghi chỉ vượt do
cộng dồn và bảy lịch sử sàng lọc không thay cho phép so hai luật trên diễn biến mô phỏng đầy đủ.
\item Luật hạn chót chỉ mượn thành phần hạn chót của Schulz và Pfeiffer~\cite{schulz2026}, với 30 và 60 giây trên tập
kiểm tra chính và 30--300 giây trên tập buổi sáng; mức mặc định của bài đó là 900 giây. Hạn chót nới kiểu Hörl và
Zwick~\cite{horl2022} (mục~\ref{muc:doichung}) trùng với cấu hình M$^{+}$ có luật phục hồi ``không tệ hơn'' ở trên,
nên mới được chạy ở 30 và 60 giây trên tập kiểm tra chính. Các phép so không nói về phương pháp của các tác giả đó.
\item Mức 30 giây không phải ngưỡng cảm nhận hay chuẩn ngành, và không có dữ liệu về mức khách sẵn lòng trả cho lời
hứa ổn định hơn; vì vậy không mức trần nào được chọn là tốt nhất.
\end{itemize}
% EV: H02; H12; K18; E09; K08; J10; J11; X3-28 (luật hạn chót 30–300 s ở tập buổi sáng; chương 3 Bảng luat); K27; D20; K52; K36
; không có \chapter mới.
% Mỗi dòng "% EV:" trỏ tới EVIDENCE-REGISTER.md (A01…K52), EVIDENCE-REGISTER-ADDENDUM-1.md (N-xx, L1a-xx, mục 0.x),
% EVIDENCE-REGISTER-ADDENDUM-2-CH3.md (X3-xx), EVIDENCE-REGISTER-ADDENDUM-3-CH4B.md (N4B-xx),
% EVIDENCE-REGISTER-ADDENDUM-4-CH4A.md (C4A-xx) hoặc EVIDENCE-REGISTER-ADDENDUM-6-CH4B-v8.md (N4B8-xx; script chỉ đọc
% evidence/checks-ch4b-v8/ch4b_v8_check.py, output ch4b-v8-check-2026-10-12.txt, phần A/B/C gọi là "check A/B/C").
% Không dùng gói mới: TikZ, enumitem và graphicx đã nạp trong main.tex; thu bảng/hình bằng resizebox.
% =====================================================================================================

% -----------------------------------------------------------------------------------------------------
\section{Vai trò của trần và thứ tự tiêu chí của bộ chọn}
\label{muc:c4-vaitro}

\subsection{Vì sao cần đổi thứ tự ưu tiên?}

Các kết quả trước mục này dùng bộ chọn ưu tiên ổn định: sau khi xét số yêu cầu nhận được, bộ chọn so mức dùng
trần và phần dời giờ rồi mới xét chi phí kế hoạch. Như vậy, cả việc xếp hạng phương án lẫn trần cứng đều có thể
góp phần hạn chế thay đổi lời hứa. Chỉ so có trần với không trần dưới một thứ tự chưa cho biết kết quả sẽ ra sao
khi bộ điều phối đặt chi phí kế hoạch lên trước phần dời giờ.

Để kiểm tra điều đó, khóa luận giữ ưu tiên nhận yêu cầu và các phép kiểm khả thi, nhưng cho bộ chọn xét chi phí
kế hoạch ngay sau đó. Các mức xếp hạng theo ổn định được bỏ, còn tiêu chí phân xử theo mã định danh giữ nguyên.
Trần cứng của luật C vẫn được kiểm trước khi chọn nếu cấu hình có trần. Hai cách gọi dưới đây là \emph{ưu tiên
ổn định} và \emph{ưu tiên chi phí kế hoạch}; chúng biểu thị mục tiêu khác nhau, không phải hai mức chất lượng
mặc định. Chi phí kế hoạch là tổng thời gian còn lại của tuyến dự kiến, chưa phải tiền hay chi phí vận hành thực.
% EV: S15-03; A14; X3-07; K15. Cùng ưu tiên nhận khách trong từng lượt, không bảo đảm cùng số phục vụ cuối.

Phép so gồm bốn cấu hình trên cùng 34 kịch bản và 3.672 yêu cầu mỗi cấu hình: hai thứ tự ưu tiên, mỗi thứ tự có
trần 30 giây và không trần. Hai cấu hình ưu tiên chi phí có 68 lượt mô phỏng hoàn tất, được đối chiếu với 68 lượt
ưu tiên ổn định đã lưu. Tập kịch bản và tập mã yêu cầu khớp ở cả 34 kịch bản. Các diễn biến sau khi chọn có thể
khác nhau, nên nhóm khách được phục vụ cũng khác nhau.
% EV: S15-01 (mục 0--3), main-final-admission.json; S15-02. Không có 34 lần lặp độc lập.

\subsection{Đánh đổi giữa phục vụ, ghép khách và sửa lời hứa}

\begin{table}[htbp]
\centering
\caption[Đánh đổi dưới hai thứ tự chọn phương án]{Hai thứ tự chọn phương án trên tập kiểm tra chính, mỗi cấu hình
có 3.672 yêu cầu. Trễ được so với lời hứa trả đầu chưa tính dự báo, trừ dòng ghi riêng. Khách có đi chung được
đếm một lần nếu hoàn thành chuyến và có khoảng thời gian cùng xe với khách khác.}
\label{bang:c4-doi-uutien}
\small
\setlength{\tabcolsep}{3pt}
\begin{tabular}{@{}>{\raggedright\arraybackslash}p{6.5cm}rrrr@{}}
\toprule
& \multicolumn{2}{c}{\textbf{Ưu tiên ổn định}} & \multicolumn{2}{c}{\textbf{Ưu tiên chi phí}}\\
\cmidrule(lr){2-3}\cmidrule(l){4-5}
\textbf{Chỉ số} & \textbf{Không trần} & \textbf{Trần 30} & \textbf{Không trần} & \textbf{Trần 30}\\
\midrule
Lượt phục vụ & 2.434 & 2.338 & 2.502 & 2.364\\
Muộn quá một phút & 203 & 76 & 379 & 72\\
Không muộn quá một phút & 2.231 & 2.262 & 2.123 & 2.292\\
Không muộn quá hai phút & 2.332 & 2.327 & 2.280 & 2.353\\
Không muộn quá năm phút & 2.416 & 2.338 & 2.466 & 2.364\\
Khách hoàn thành có đi chung & 555 & 313 & 1.027 & 432\\
Khách bị sửa giờ trả do điều phối & 174 & 33 & 434 & 106\\
Trong đó, bị sửa từ hai lần & 12 & 0 & 50 & 2\\
Đón muộn theo giới hạn dịch vụ & 51 & 43 & 39 & 44\\
Muộn quá một phút khi lời hứa có dự báo & 137 & 10 & 318 & 11\\
Tổng dời lớn nhất của một khách (s) & 590,1 & 29,2 & 709,1 & 29,7\\
\bottomrule
\end{tabular}
\end{table}
% EV: S15-01 (mục 2, 5a--5f); S15-02 (pooled); S15-04 (kiểm bảng và các chênh lệch từ TSV).

Khi cùng giữ trần 30 giây, ưu tiên chi phí phục vụ thêm 26 lượt và có thêm 119 khách hoàn thành có đi chung
so với ưu tiên ổn định (Bảng~\ref{bang:c4-doi-uutien}). Số lượt không muộn quá một phút tăng 30. Đổi lại, số khách
bị sửa giờ trả tăng từ 33 trong 2.338 lên 106 trong 2.364, tương ứng 1,41\% và 4,48\%; số khách bị sửa từ hai lần
tăng từ 0 lên 2. Số lượt đón muộn tăng từ 43 lên 44. Như vậy, cách ưu tiên ổn định vẫn cho ít khách bị sửa lời
hứa hơn, còn ưu tiên chi phí nhận và ghép được nhiều khách hơn trong thiết lập này.

Đúng giờ đến và ít sửa lời hứa là hai chỉ số khác nhau. Ưu tiên chi phí có trần có 72 lượt muộn quá một phút,
ít hơn 4 lượt so với ưu tiên ổn định có trần, nhưng nhiều khách bị sửa giờ hơn. Khi dùng lời hứa có dự báo, số
lượt muộn quá một phút là 11 và 10 trên hai nhóm khách tương ứng. Các chênh lệch nhỏ này không đủ để chọn một
cách làm tốt hơn về mọi mặt. Trung vị độ dài lời hứa đầu chưa dự báo là 809,642 giây với ưu tiên chi phí và
815,060 giây với ưu tiên ổn định, trên hai tập khách khác nhau.
% EV: S15-01 (5b--5e); S15-02. Độ dài dùng p0-t0; không quy thành chờ/đi vòng của từng khách.

Không trần cho thấy rõ cái giá của đổi thứ tự. Ưu tiên chi phí phục vụ thêm 68 lượt và có thêm 472 khách hoàn
thành có đi chung so với ưu tiên ổn định, nhưng có thêm 176 lượt muộn quá một phút và thêm 260 khách bị sửa giờ
trả. Trần vì vậy vẫn có vai trò khi bộ chọn chuyển sang ưu tiên chi phí. Trong thứ tự này, thêm trần 30 giây làm
giảm 138 lượt phục vụ và 595 khách hoàn thành có đi chung; đồng thời có ít hơn 307 lượt muộn quá một phút và
thêm 169 lượt được phục vụ, không muộn quá một phút.
% EV: S15-01 (mục 2, 5f--5g); S15-04. Các chênh lệch là tổng trên tập yêu cầu, không cùng tập khách phục vụ.

Phần giảm số lượt trễ phải đọc cùng phần không phục vụ. Với ưu tiên ổn định, 127 lượt trễ ít hơn đi cùng 96 lượt
phục vụ ít hơn và 31 lượt không muộn quá một phút nhiều hơn. Với ưu tiên chi phí, các số tương ứng là 307, 138
và 169. Chênh lệch giữa hai mức giảm trễ là 180, nhưng không có nghĩa 180 khách cụ thể được bảo vệ chỉ nhờ đổi
bộ chọn. Chênh lệch này cùng chiều ở 29 trong 34 kịch bản và cả năm ngày; các kịch bản phụ thuộc nhau nên đây
chỉ là mô tả trên dữ liệu đã dùng.
% EV: S15-01 (mục 2, 4); S15-02 (interactions); 180 = (138-96) + (169-31).

Ở ưu tiên chi phí có trần, tổng dời lớn nhất do điều phối là 29,657 giây và không khách nào vượt 30 giây.
Kết quả cho thấy cấu hình có trần giữ được giới hạn đã đặt khi thay thứ tự xếp hạng. Nó chưa chứng minh rằng
riêng trần đủ bảo đảm giờ đến thực tế, vì khóa, ràng buộc dịch vụ, bộ sinh và luật phục hồi vẫn cùng tham gia.

Mức trần thích hợp còn phụ thuộc tiêu chí dịch vụ. Với ưu tiên chi phí, trần 30 giây tăng số lượt được phục vụ
và không muộn quá một phút từ 2.123 lên 2.292; nếu chấp nhận muộn tới năm phút thì số lượt đáp ứng lại giảm
từ 2.466 xuống 2.364. Đây là hai cách đọc cùng kết quả đã có, không phải khảo sát mức khách chấp nhận chờ.
Trần 30 giây là điểm thực nghiệm, không được chọn làm ngưỡng tối ưu hay ngưỡng cảm nhận của hành khách.
% EV: S15-01 (5a); S15-02 (timely60, timely300); K36; K52.

\subsection{Lời hứa có dự báo và giới hạn của phép so}

Với ưu tiên chi phí có trần, tính lại lời hứa đầu theo dự báo cho 11 lượt muộn quá một phút thay vì 72,
trên cùng 2.364 hành trình; lời hứa dài thêm trung bình 7,30 giây. Với cấu hình không trần, các số lượt muộn
là 318 thay vì 379, và lời hứa dài thêm trung bình 7,02 giây. Giờ trả thực tế và việc xe chở ai không đổi
trong phép tính này. Dự báo điều chỉnh phần đường dự kiến chậm đi, còn luật C giới hạn phần dời do điều phối;
hai thành phần không thay thế nhau và không tạo thành bảo đảm vô điều kiện về giờ đến.
% EV: S15-01 (5b); F01--F04. 21 ô bật causal-no15; 13 ô còn lại giữ p0.

Các lượt ưu tiên chi phí dùng bản chạy riêng. Phép đối chiếu hành vi trên bốn cặp mẫu và việc tính lại các
chỉ số ưu tiên ổn định giúp kiểm tra khả năng so sánh, nhưng không chứng minh mọi khác biệt giữa hai bản
chạy đã được loại bỏ. Các con số là phép so thăm dò theo cấu hình trên một bộ mô phỏng và năm ngày đã xem,
không phải ước lượng tác động nhân quả riêng của một tiêu chí. Phép sàng lọc trên các phương án đã ghi của
tập tắc đường cố định được giữ ở Phụ lục~\ref{phuluc:D}; nó bổ trợ việc hiểu bộ chọn, không thay kết quả
của những lần mô phỏng đầy đủ ở đây.
% EV: STAGE1B-PLAN.md mục 1.1 và amendment mục 3; S15-01 (mục 0, 3); G0 chỉ behavioural-only.


\section{Lời hứa theo dự báo}
\label{muc:c4-dubao}

Trên 21 kịch bản ngày thường có dự báo, số khách của luật C trần 30 giây tới muộn hơn lời hứa đầu quá một phút giảm từ
68 xuống 2 trong 1.471 khách được phục vụ (Bảng~\ref{bang:c4-dubao}). Lời hứa theo dự báo (mục~\ref{muc:dubao}) được
tính lại trên các lần chạy đã ghi; mỗi cách hứa được so với lời hứa đầu của chính nó.
% EV: F17; F04; N4B8-01 (check A: 68 → 2 trên 1.471); F03 (chỉ vế phát lại); FFA:8-13

Đối chứng thứ hai là đệm tỉ lệ: mỗi lời hứa đầu được kéo dài một tỉ lệ cố định của khoảng từ lúc hứa tới giờ trả đã
hứa. Tỉ lệ được chọn theo từng khung giờ, trên chính các khách này, để tổng phần kéo dài bằng của lời hứa theo dự báo.
% EV: N4B-03 (FFA:15-17: lời hứa đệm = P0 + k_w(P0 − t0), k_w = Σ(D0 − P0)/Σ(P0 − t0) trên mọi khách C-30 của khung w ở 21 kịch bản, tính từ lời hứa; FA:2)

\begin{table}[htbp]
\centering
\caption[Lời hứa theo dự báo và đệm tỉ lệ trên ngày thường]{Khách của luật C trần 30 giây trên 21 kịch bản ngày thường có dự báo, theo ba cách hứa: lời hứa đầu không dự
báo, đệm tỉ lệ và lời hứa đầu theo dự báo; mỗi cách so với lời hứa đầu của chính nó. Khối trên là số đếm gộp theo khách;
dòng cuối là trung bình theo kịch bản của tỉ lệ trễ, mỗi kịch bản tính như nhau.}
\label{bang:c4-dubao}
\small
\begin{tabular}{@{}lrrrr@{}}
\toprule
\textbf{Nhóm kịch bản} & \textbf{Phục vụ} & \textbf{Không dự báo} & \textbf{Đệm tỉ lệ} & \textbf{Dự báo}\\
\midrule
\multicolumn{5}{@{}l}{\textit{Số khách tới muộn hơn lời hứa đầu quá 60 s}}\\
Cả 21 kịch bản & 1.471 & 68 & 22 & 2\\
Khung 7 giờ (8 kịch bản) & 444 & 60 & 17 & 1\\
Khung 8 giờ (8 kịch bản) & 609 & 8 & 5 & 1\\
Khung 17 giờ (5 kịch bản) & 418 & 0 & 0 & 0\\
\multicolumn{5}{@{}l}{\textit{Số khách tới sớm hơn lời hứa đầu quá 60 s}}\\
Cả 21 kịch bản & 1.471 & 0 & 0 & 1\\
% EV: N4B8-01 (check A: CELLS, share × served từng kịch bản, phần dư 0); N4B8-08 (FFA:139-141 tính E_pad; sửa N4B-03: đệm tỉ lệ có đo tới sớm, bằng 0); F17
\midrule
\multicolumn{2}{@{}l}{Tỉ lệ trễ quá 60 s, trung bình theo kịch bản (\%)} & 6,01 & 1,94 & 0,17\\
% EV: F04 (FA:8); N4B-03 (FA:8: pad 0,0194)
\bottomrule
\end{tabular}
\end{table}

Cách đọc bảng: ở khung 7 giờ, 60 trong 444 khách trễ quá một phút khi không dự báo, 17 khi đệm tỉ lệ và 1 khi có dự
báo.
% EV: N4B8-01

Tính trung bình theo kịch bản, tỉ lệ trễ giảm 5,84 điểm phần trăm. Gom 21 kịch bản thành 8 nhóm theo ngày và khung giờ,
khoảng tin cậy 95\% theo phân phối Student là từ 1,15 đến 10,54 điểm phần trăm. So với đệm tỉ lệ, chênh lệch là
$-1{,}77$ điểm phần trăm, khoảng từ $-3{,}62$ đến $+0{,}07$, tức có chứa 0: chưa thể nói dự báo tốt hơn cách đệm này.
Các khoảng chỉ dựa trên ba buổi sáng đã được xem trước, và mọi hồ sơ dự báo đều có ngày 12/11 trong các ngày nguồn.
% EV: F04; F05 (FA:33,36: dY −5,844 [−10,542; −1,146]; dPad −1,775 [−3,620; +0,071]; FFA cluster_t: tầng = khung giờ, cụm = ngày–khung, df = 8 − 3 = 5); F11 (FR:97-98; factor_nowcast_screen.py:21-22: 14 ← 12, 13; 15, 16 ← 12, 13, 14); FR:170-171; I05; K24

Lợi ích không riêng của luật C. Trên cùng 21 kịch bản, khi có dự báo, luật hạn chót 30 giây giảm từ 66 xuống 2 khách
trễ trong 1.471 khách được phục vụ; luật trần độ dời khách thấy 30 giây cũng giảm từ 66 xuống 2 trong 1.470 khách.
Điều phối không trần giảm từ 129 xuống 63 trong 1.512 khách. Riêng khung 17 giờ, 28 khách trễ trong 437 khách của điều
phối không trần không đổi, còn luật C có 0 trong 418 khách. Theo sổ, 86,1\% số giây trễ của 68 khách trễ khi không dự
báo thuộc phần được hồ sơ lịch sử giải thích (Hình~\ref{hinh:c4-so}); dự báo và trần nhắm vào hai phần khác nhau của sổ.
% EV: F08; K37; N4B8-02 (check A: M⁺-30 66 → 2 trên 1.471; V-30 66 → 2 trên 1.470; U 129 → 63 trên 1.512; U khung 17 giờ 28 → 28 trên 437; C-30 khung 17 giờ 0 trên 418); F09 (FA:30; RX:3); eq:so; K04

Trên 21 kịch bản này, giá của dự báo là lời hứa giờ trả dài thêm 13,7 giây (khung 7 giờ 27,8 giây) và lời hứa giờ đón 3,2 giây, trung bình
theo kịch bản. Trên cả 34 kịch bản, khách trễ của luật C giảm từ 76 xuống 10 trong 2.338 (Bảng~\ref{bang:c4-tran120}).
% EV: F06 (FR:67-69; FA:8-9); C06 (76); C09 (10 trên 34 kịch bản); N-16; F02; N-28 (13 kịch bản cuối tuần không có dự báo; độ dài thêm gộp 7,37 s, N4B-10, nằm ở bang:c4-tran120 và bang:c4-dem)

Kết quả 2 trong 1.471 khách đạt được trong điều kiện thuận lợi cho dự báo theo giờ trong ngày, không phải hiệu quả ngoài
đời. Tắc đường là một hệ số chung cho cả thành phố. Trên cả 34 kịch bản, với luật C, kéo dài lời hứa theo đúng hệ số
sẽ xảy ra cho giờ trả khớp trong vòng 2 giây ở cả 2.305 khách không bị điều phối dời. Phép kiểm giờ đón cũng chưa dùng
lời hứa này (mục~\ref{muc:c4-thaoluan}).
% EV: I02 (FR:167-169: mức tốt nhất cho loại dự báo theo giờ trong ngày trong mô phỏng gần như tất định); F10 (FA:28; RX:8: 34 kịch bản); B06; F07; X3-11

\begin{figure}[htbp]
\centering
\resizebox{\linewidth}{!}{%
\begin{tikzpicture}[x=1cm,y=1cm,font=\footnotesize,
  nhan/.style={font=\scriptsize,inner sep=1.5pt}]
% Thanh trên: ngày thường. Độ dài 11 cm = 6.217,4 s; ranh giới 0,1445 và 9,6145 cm (FA:30).
\draw[fill=black!75] (0,1.6) rectangle (0.1445,2.2);
\draw[fill=black!20] (0.1445,1.6) rectangle (9.6145,2.2);
\draw[fill=white] (9.6145,1.6) rectangle (11,2.2);
\node at (4.88,1.9) {5.353 s (86,1\%)};
\node[align=center,nhan] at (10.307,1.9) {783 s\\(12,6\%)};
\draw (0.072,2.2)--(0.072,2.45);
\node[nhan,anchor=south west] at (0.0,2.45) {82 s (1,3\%)};
\node[anchor=east,align=right] at (-0.15,1.9) {Ngày thường, 21 kịch bản\\68 khách, 6.217 s};
% Thanh dưới: chiều tuyết. Độ dài 11 cm = 3.445,9 s; ranh giới 0,0273 và 0,8302 cm (NEW:5).
\draw[fill=black!75] (0,0) rectangle (0.0273,0.6);
\draw[fill=black!20] (0.0273,0) rectangle (0.8302,0.6);
\draw[fill=white] (0.8302,0) rectangle (11,0.6);
\node at (5.915,0.3) {3.185,9 s (92,5\%)};
% v8: nhãn 8,6 s đặt phía trên thanh tuyết (khoảng trống giữa hai thanh) để đường dẫn của nhãn 251,5 s không cắt chữ.
\draw (0.014,0.6)--(0.014,0.85);
\node[nhan,anchor=south west] at (-0.02,0.85) {8,6 s (0,25\%)};
\draw (0.43,0)--(0.43,-0.2);
\node[nhan,anchor=north west] at (0.39,-0.2) {251,5 s (7,3\%)};
\node[anchor=east,align=right] at (-0.15,0.3) {Chiều tuyết, 8 kịch bản\\41 khách, 3.445,9 s};
% Chú giải (ba ô cách nhau đủ để chữ không chồng lên nhau)
\draw[fill=black!75] (0,-1.25) rectangle (0.3,-1.0);
\node[anchor=west] at (0.35,-1.125) {phần tính cho điều phối};
\draw[fill=black!20] (4.0,-1.25) rectangle (4.3,-1.0);
\node[anchor=west] at (4.35,-1.125) {phần được hồ sơ lịch sử giải thích};
\draw[fill=white] (9.6,-1.25) rectangle (9.9,-1.0);
\node[anchor=west] at (9.95,-1.125) {phần còn lại};
\end{tikzpicture}
}
\caption[Ba phần của số giây trễ theo sổ của từng khách]{Số giây trễ chia thành ba phần theo sổ của từng khách (công thức~\eqref{eq:so}), với luật C trần 30 giây và các
khách tới muộn hơn lời hứa đầu không dự báo quá 60 giây: 68 khách của 21 kịch bản ngày thường và 41 khách của 8 kịch
bản chiều tuyết 15/11/2018. Mỗi thanh dài bằng 100\% số giây trễ của nhóm đó. Phần được hồ sơ lịch sử giải thích phụ
thuộc dự báo đang dùng; phần còn lại là phần dự báo này chưa giải thích được, không phải phần không thể biết trước.}
\label{hinh:c4-so}
\end{figure}
% EV (ngày thường): F09; N4B8-11 (FA:30 và CELLS: 81,742 + 5.352,646 + 783,018 = 6.217,406 s; các nhãn làm tròn đến giây độc lập, nên tổng ba nhãn là 6.218 s, lệch 1 s so với nhãn tổng 6.217 s); RX:3 (1,3 / 86,1 / 12,6%)
% EV (chiều tuyết): F13; N4B-05 (NEW:5: 8,566 / 251,496 / 3.445,933 s; RP1:24-26: 7,3%; SNOW:10: 92,5%; FR:104,186); N4B8-12 ([A] 8,566 / 3.445,933 = 0,25%)
% Tọa độ: x = giây / tổng × 11 cm (tính tay; xem NOTES v7 mục 4).

\textbf{Đưa dự báo vào bộ điều phối.} Phép thử chạy luật C và luật hạn chót 30 giây trên 16 kịch bản khung 7 và 8 giờ.
Thời gian đi của bộ điều phối được nhân thêm mức tăng dự báo trung bình của khung 15 phút hiện tại và khung kế tiếp. Có
3 trong 32 lần chạy dừng vì FleetPy tính rằng xe sẽ tới điểm đón muộn hơn giờ đón muộn nhất chưa tới 3 giây. Luật C
không có quyết định bị ép ở 14 lần chạy hoàn tất. Luật hạn chót bị ép ở 47,6\% lượt quyết định ở 15 lần chạy hoàn tất,
so với 62,8\% ở các lần chạy không đưa dự báo vào trên cùng 15 kịch bản. Tính gộp trên cùng 15 kịch bản, đó là 3.849
trong 8.496 lượt quyết định khi đưa dự báo vào, so với 5.121 trong 8.468 lượt khi không đưa vào; các tỉ lệ 47,6\% và
62,8\% ở trên là trung bình theo kịch bản. Tỉ lệ khách bị luật hạn chót phá lời hứa ở khung 7 giờ giảm từ 39,0\% xuống
11,0\%, trung bình trên cùng tám kịch bản. Tính gộp theo khách trên tám kịch bản này, đó là 170 trong 442 khách khi
không đưa dự báo vào, so với 45 trong 439 khách khi đưa vào. Vì vậy, các số của Bảng~\ref{bang:c4-nghiavu} phụ thuộc cách
hứa. Khóa luận chỉ dùng dự báo ở lời hứa báo cho khách; kết luận này chỉ áp cho cách đã thử.
% EV: F12; X3-29 (W13:85-89); N4B8-06 (K = 2: run_forecast.py:9,72; configs/forecast-k2-2018-11-1{4,5,6}.json:1; preflight_world_v1_3.py:85-89; FR:117; R26 Q9); N4B-07 (RX:9,11-15; FR:120-121,124,138,157: C hoàn tất 14, luật hạn chót 15; 47,6 và 62,8% trên cùng 15 kịch bản); C14; FR:141-146; R37 verify-numbers-v2.json:dispatcher; R39 agent-numbers-v1.json:dispatcher.paired, morning_promise_breaks (cùng15 hoàn tất; khung7 cùng8; số gộp tách khỏi trung bình theo kịch bản)

\textbf{Học máy.} Ba mô hình học một lần rồi giữ nguyên ước lượng phân vị 90 của độ trễ so với lời hứa theo dự báo (mức
mà 90\% khách giống vậy không vượt) từ thông tin đã có ở thời điểm hứa. Đó là mô hình cây quyết định tăng cường
(gradient boosting), mô hình ấy khi bỏ hai biến trạng thái xe (số khách khác đã gán và đang trên xe), và hồi quy phân vị
tuyến tính. Mỗi khách được cộng vào lời hứa theo dự báo một phần đệm tỉ lệ với mức ước lượng; đối chứng cộng cùng một
khoảng cho mọi khách. Cả hai được chỉnh trên ngày hiệu chỉnh để lời hứa dài thêm trung bình 30 giây so với lời hứa không
dự báo, nên trên ngày kiểm mức này có thể khác. Phép thử dùng hồ sơ dự báo cố định của ngày
12--13/11 và kiểm trên 15, 16/11 khung 7 giờ và 13/11 của tập tắc đường cố định, các ngày đã được xem trước. Mô hình
chỉ được coi là tốt hơn khi, so với đối chứng, thỏa cả bốn điều kiện:
\begin{itemize}[leftmargin=1.2cm,itemsep=0pt,topsep=2pt]
\item tỉ lệ khách trễ quá một phút thấp hơn;
\item độ dài thêm trung bình không hơn quá 1 giây;
\item phân vị 90 của độ dài thêm không hơn quá 15 giây;
\item tỉ lệ khách tới sớm quá một phút không hơn quá 1 điểm phần trăm.
\end{itemize}
% EV: G01; G02; G03; N4B-08; N4B8-04 (MLS:9-10 các ngày đã xem; 19-23, 67-68 hai biến trạng thái; 37-39 hồ sơ cố định 12–13/11; 41-61 chia ngày; 63-69 ba mô hình; 71-82 chỉnh trên ngày hiệu chỉnh, F = D0 + ρ·w, đối chứng D0 + (E − mean δ), không chỉnh lại khi kiểm; 91-95 bốn điều kiện; MLT: đối chứng của U dài thêm 37,219 / 31,772 s ở 15/11 / 16/11)

Không mô hình nào thỏa đủ bốn điều kiện. Các số đếm sau được suy từ tỉ lệ gộp theo khách đã làm tròn đến ba chữ số
thập phân trong bảng kết quả; với mẫu số tương ứng, mỗi tỉ lệ chỉ ứng với một số nguyên. Với luật C, đối chứng và mô hình đều không có khách trễ (147 và 197 khách).
Với điều phối không trần, mô hình cây tăng cường chỉ bớt 1 trong 2 khách trễ ở ngày 15/11 (145 khách), không bớt ai ở
ngày 16/11 (200 khách), và không thỏa điều kiện thứ hai, thứ ba. Với ghép tự do, khách trễ giảm từ 104 xuống 94 trong 807
nhưng 57 khách tới sớm quá một phút, so với 0 khi đệm đều. Bỏ hai biến trạng thái xe không đổi số khách trễ ở hai phép
so này.
% EV: G05–G10; N4B8-03 (check B: C-30 0/147, 0/197; U 2 → 1 trên 145, 1 → 1 trên 200; FREE 104 → 94 trên 807, sớm 57 so với 0; bỏ trạng thái xe 1/145 và 94/807); ML:6-7,14-15,21,45-53,67-69 (độ dài thêm +4,4 / +3,8 s; phân vị 90 +21,8 / +21,7 s); K40; R37 verify-numbers-v2.json:ml; R39 agent-numbers-v1.json:ml (xét mọi số nguyên, sai số 0,0005 điểm phần trăm); ml_pilot_run.py:37-42,175-176

% -----------------------------------------------------------------------------------------------------
\section{So với cách đệm lời hứa}
\label{muc:c4-dem}

Khi có dự báo, cộng thêm 30 giây vào mọi lời hứa của điều phối không trần cũng cho 2.328 lượt không trễ quá một phút
trên 3.672 yêu cầu, bằng luật C trần 30 giây có dự báo. Ở ngưỡng hai phút, cách đệm này cho 2.366 lượt, so với 2.338
của luật C. Hai cách này cho hai thứ khác nhau. Đệm kéo dài lời hứa, trung bình 37,2 thay vì 7,4 giây, nên ít khách bị tính là trễ. Nhưng đệm không đổi
quyết định điều phối: vẫn có 174 khách bị điều phối dời giờ trả, với tổng độ dời lớn nhất là 590,1 giây (gần 10 phút).
Luật C trần 30 giây giới hạn tổng độ dời do điều phối của mỗi khách ở 30 giây dưới giả thiết (N); trên 34 kịch bản của
tập kiểm tra chính, mức lớn nhất đo được là 29,2 giây, kể cả ở các kịch bản cuối tuần không có dự báo. Đổi lại, luật C
phục vụ ít hơn 96 lượt.
% EV: C10; N-16; N-18; K06; K51; C20; N-23; N-28 (cuối tuần 29,2 s); C02; A06 (dưới giả thiết (N); không bảo đảm giờ đến: A22, K05); RD §8; khớp câu trỏ của phần A v5-v3 mục 4.2 (C10); R37 verify-numbers-v2.json:padding; R39 agent-numbers-v1.json:padding,cohort_totals

Cách đệm là phát lại trên các lần chạy của điều phối không trần, nên hai bên được so trên hai nhóm khách khác nhau, chỉ
chung 2.207 khách. Bảng~\ref{bang:c4-dem} xếp các cách hứa thành ba khối có lời hứa dài thêm gần bằng nhau.
% EV: C10 (phát lại, không phải đối chứng ghép cặp); ADDENDUM-1 phần đầu (quy tắc đệm D = D0 + k); C04

\begin{table}[htbp]
\centering
\caption[Trần và cách đệm lời hứa khi có dự báo]{Các cách hứa khi có dự báo trên tập kiểm tra chính (34 kịch bản, 3.672 yêu cầu mỗi cấu hình; dự báo chỉ có ở
21 kịch bản ngày thường). U: điều phối không trần; C-$B$: luật C trần $B$ giây; trong ngoặc là số khách được phục vụ.
Số đếm gộp; trễ và tới sớm tính trong số khách được phục vụ, so với lời hứa của chính cách hứa; hai cột ``không trễ
quá'' tính trên 3.672 yêu cầu; độ dài thêm là trung bình gộp theo khách được phục vụ. Đệm không đổi quyết định điều
phối, nên hàng của U và hàng của C là các nhóm khách khác nhau.}
\label{bang:c4-dem}
\small
\begin{tabular}{@{}p{4.6cm}rrrrr@{}}
\toprule
\textbf{Cách hứa} &
\textbf{\begin{tabular}[b]{@{}r@{}}Trễ\\ quá 60 s\end{tabular}} &
\textbf{\begin{tabular}[b]{@{}r@{}}Không trễ\\ quá 60 s\end{tabular}} &
\textbf{\begin{tabular}[b]{@{}r@{}}Không trễ\\ quá 120 s\end{tabular}} &
\textbf{\begin{tabular}[b]{@{}r@{}}Tới sớm\\ quá 60 s\end{tabular}} &
\textbf{\begin{tabular}[b]{@{}r@{}}Dài thêm\\ (s)\end{tabular}}\\
\midrule
\multicolumn{6}{@{}l}{\textit{Không đệm (dài thêm khoảng 7 giây)}}\\
U, dự báo (2.434) & 137 & 2.297 & 2.347 & 8 & 7,18\\
C-30, dự báo (2.338) & 10 & 2.328 & 2.338 & 4 & 7,37\\
C-120, dự báo (2.387) & 65 & 2.322 & 2.379 & 2 & 7,27\\
% EV: N-16; N-03; C09; N-06; L1a-03; N4B-09 (sớm 8 / 4 / 2)
\midrule
\multicolumn{6}{@{}l}{\textit{Đệm 30 giây (dài thêm khoảng 37 giây)}}\\
U, dự báo, đệm 30 s (2.434) & 106 & 2.328 & 2.366 & 44 & 37,18\\
C-30, dự báo, đệm 30 s (2.338) & 2 & 2.336 & 2.338 & 44 & 37,37\\
% EV: C10; N-02; N-16
\midrule
\multicolumn{6}{@{}l}{\textit{Đệm 60 giây (dài thêm khoảng 67 giây)}}\\
U, dự báo, đệm 60 s (2.434) & 87 & 2.347 & 2.386 & 1.033 & 67,18\\
% EV: C11; N-16
\bottomrule
\end{tabular}
\end{table}

Cách đọc bảng: với đệm 30 giây trên điều phối không trần, 106 trong 2.434 khách tới muộn hơn lời hứa đã đệm quá 60
giây, nên có 2.328 lượt không trễ quá 60 giây trên 3.672 yêu cầu; 44 khách tới sớm quá 60 giây.
% EV: C10; N-16

\textbf{Vai trò của từng thành phần.} Thứ nhất, ở mức kéo dài lời hứa gần bằng nhau, luật C trần 30 giây hơn cách đệm
ở ngưỡng một phút và kém ở ngưỡng hai phút: khi cùng đệm 30 giây, hai bên có 2.336 so với 2.328 lượt và 2.338 so với
2.366 lượt. Ở hàng đệm 30 giây, điều phối không trần đạt cùng 2.328 lượt không trễ quá một phút như luật C có dự báo
không đệm, với lời hứa dài thêm trung bình 37,18 so với 7,37 giây, tức khoảng năm lần.
% EV: N-17; N-18 (37,1833 / 7,3671 = 5,05); N-16; R37 verify-numbers-v2.json:padding; R39 agent-numbers-v1.json:padding,cohort_totals

Thứ hai, lời hứa dài thì ít trễ nhưng nhiều khách tới sớm: 44 khách với cách đệm 30 giây và 1.033 khách với cách đệm 60
giây, so với 4 khách với luật C có dự báo. Số lượt tới trong khoảng $\pm$60 giây quanh lời hứa là 2.284 với cách đệm 30
giây, so với 2.324 với luật C, trên 3.672 yêu cầu.
% EV: C10; C11; N-04 (±60 s = phục vụ − trễ60 − sớm60, ADD1:92); N4B-09 (ADD1:84-89)

Thứ ba, cách đệm 30 giây vẫn có 68 khách trễ quá hai phút, 11 khách quá năm phút và 1 khách quá mười phút; luật C có
dự báo không có ai trễ quá hai phút. Cách đệm hơn 28 lượt ở ngưỡng hai phút vì phục vụ nhiều hơn 96 lượt, không phải vì
ít khách trễ hơn.
% EV: N-01 (68 / 11 / 1 so với 0 / 0 / 0; ADD1:15,19); N-19 (2.366 − 2.338 = 28 = 96 − 68)

Khi phát lại cùng lời hứa theo dự báo, luật hạn chót 30 giây cũng có 2.338 lượt phục vụ và 2.328 lượt không trễ quá
một phút. Vì phép phát lại không đổi các lần chạy, nghĩa vụ của luật này vẫn tính từ lời hứa đầu không dự báo. Theo
nghĩa vụ đó, tính trung bình theo kịch bản, nó bị ép ở 41,3\% lượt quyết định; tính gộp là 7.710 trong 20.103 lượt.
Luật hạn chót có khách tới muộn hơn lời hứa đầu không dự báo quá 30 giây ở 28 trong 34 kịch bản, còn luật C không có
quyết định bị ép hay khách vượt trần (Bảng~\ref{bang:c4-nghiavu}).
% EV: N-16 (M⁺-30 + dự báo 2.338 / 10 / 2.328 / 2.338, ADD1:29); C13 (41,3% trung bình theo kịch bản); N-21 (phá hứa của M⁺-30 = A − p₀ > 30 s với lời hứa không dự báo; C-30 0/34, 0 quyết định bị ép); X3-16; R36-A01 (bảng A-v5-v3 dùng số gộp); R39 agent-numbers-v1.json:main_original_obligation

Trần 120 giây có dự báo, với lời hứa chỉ dài thêm 7,27 giây, cho 2.379 lượt không trễ quá hai phút, nhiều hơn 2.366
của cách đệm 30 giây, nhưng 2.322 lượt ở ngưỡng một phút, ít hơn 2.328. Đây là một điểm đo thêm trên một tập khách khác,
không phải mức trần tốt nhất.
% EV: N-06; N-16; N-18; N-10; K36

\textbf{Khi nào chọn trần, khi nào chọn đệm.} Nếu hãng xe chỉ tính số lượt không trễ quá hai phút và chấp nhận lời hứa
dài hơn cùng nhiều khách tới sớm hơn, điều phối không trần có dự báo và đệm 30 giây đạt 2.366 lượt, so với 2.338 lượt
của luật C trần 30 giây có dự báo, trên 3.672 yêu cầu; điều phối không trần cũng phục vụ thêm 96 lượt. Nếu cần giới hạn
tổng độ dời do điều phối của từng khách, kể cả khi không có dự báo, luật C trần 30 giây làm được điều này dưới giả thiết
(N), đổi lại 96 lượt phục vụ và 242 khách đi chung. Hai cách dùng chung được; chọn cách nào là quyết định của hãng xe,
không phải kết luận của thực nghiệm.
% EV: N-19; C02; N-20; N-02 (hàng C-30 có dự báo và đệm 30 s); D20; I07; I10; K38; K51; A22 (giả thiết N); R37 verify-numbers-v2.json:padding; R39 agent-numbers-v1.json:padding,cohort_totals

% -----------------------------------------------------------------------------------------------------
\section{Khi dự báo không đủ: chiều tuyết 15/11/2018}
\label{muc:c4-tuyet}

Trên 601 khách được phục vụ của 8 kịch bản chiều tuyết 15/11/2018 (luật C trần 30 giây), lời hứa theo dự báo chỉ giảm
số khách trễ quá một phút từ 41 xuống 34. Tính trung bình theo kịch bản, tỉ lệ này giảm từ 7,19\% xuống 5,94\%. Theo
sổ, trong 3.445,9 giây trễ của 41 khách, 8,6 giây tính cho điều phối và 251,5 giây được hồ sơ lịch sử giải thích
(Hình~\ref{hinh:c4-so}).
% EV: F13; K25; B20; N4B-05 (NEW:5)

Bảng~\ref{bang:c4-tuyet} cho kết quả phát lại vài cách hứa dùng thêm hệ số tắc đường của chính ngày tuyết. Cách ``trộn'' lấy trung
bình nhân của mức tăng theo hồ sơ lịch sử và mức tăng của hệ số cùng ngày từ khung 15 phút trước sang khung hiện tại.
% EV: F14; K20; K21; N4B-06 (factor_nowcast_screen.py:14-20 quy tắc thông tin, 28-29: MOM = (f(k0)/f(k0−1))^(b−b0), BLEND = √(HIST·MOM); RP1:7-20; NEW:5)

\begin{table}[htbp]
\centering
\caption[Phát lại các cách hứa trong chiều tuyết]{Phát lại các cách hứa trên 8 kịch bản chiều tuyết 15/11/2018 (luật C trần 30 giây, 601 khách được phục vụ).
Cột ``khách trễ'' đếm gộp số khách tới muộn hơn lời hứa đầu quá 60 giây; cột cuối là độ dài thêm trung bình theo kịch
bản. Không cách hứa nào trong bảng có khách tới sớm quá 60 giây. Mọi cách hứa đều tính lại sau khi chạy, trên một ngày
đã được xem trước khi thiết kế các cách này; hai hàng ``có buổi sáng'' chỉ để tham khảo.}
\label{bang:c4-tuyet}
\small
\begin{tabular}{@{}p{3.4cm}p{6.6cm}rr@{}}
\toprule
\textbf{Cách hứa} & \textbf{Thông tin dùng lúc hứa} &
\textbf{\begin{tabular}[b]{@{}r@{}}Khách\\ trễ\end{tabular}} &
\textbf{\begin{tabular}[b]{@{}r@{}}Dài thêm\\ (s)\end{tabular}}\\
\midrule
Không dự báo & Thời gian đi hiện tại & 41 & 0\\
Theo dự báo & Hồ sơ tắc đường của các ngày nguồn & 34 & 2,0\\
Trộn trong khung & Thêm hệ số cùng ngày đã có trong khung giờ & 21 & 8,5\\
Trộn có buổi sáng & Thêm hệ số cùng ngày từ 6 giờ; thiết kế thử không cho khung chiều dùng dữ liệu này & 19 & 9,5\\
Trộn có buổi sáng, trễ một khung & Như hàng trên, nhưng mức tăng gần nhất chỉ lấy từ các khung đã kết thúc; hệ số của
khung hiện tại vẫn dùng giá trị đúng & 29 & 9,0\\
Đệm tỉ lệ & Tỉ lệ chọn để tổng phần kéo dài bằng của hàng ``trộn trong khung'', trên chính ngày tuyết & 27 & 8,3\\
Hệ số thật & Hệ số sau lúc hứa, không có ở thời điểm hứa & 0 & 18,8\\
\bottomrule
\end{tabular}
\end{table}
% EV: F14 (RP1:7-20; NEW:5); N4B-06 (định nghĩa; cột dài thêm là trung bình theo kịch bản, RP1:6, rider_nowcast_replay.py:96-106); K20; K21; factor_nowcast_screen.py:17-20 (hàng trễ một khung cũng dùng thông tin từ 06:00); N4B8-09 (RP1:7-20: E60 0,00% ở mọi hàng ⇒ 0 khách tới sớm ở mọi kịch bản [A])

Cách đọc bảng: khi trộn hồ sơ lịch sử với xu hướng trong khung giờ, 21 trong 601 khách tới muộn hơn lời hứa quá một
phút, với lời hứa dài thêm trung bình 8,5 giây.
% EV: F14; N4B-06

Phép phát lại này tách số giây trễ theo cách cộng có giới hạn, khác phần có dấu của công thức~\eqref{eq:so}. Phần kéo
dài của mỗi khách được chặn trong khoảng từ 0 đến phần trễ còn lại sau khi trừ phần tính cho điều phối. Theo cách cộng
này, cách trộn
có buổi sáng giải thích thêm 886,9 giây, tức 25,7\% của 3.445,9 giây trễ; 2.299,0 giây (66,7\%) chưa được giải thích
bằng các cách đã thử. Theo số liệu trạm Central Park của Cơ quan Quản lý Khí quyển và Đại dương Quốc gia Hoa Kỳ
(National Oceanic and Atmospheric Administration, NOAA), từ 9/2018 đến 12/2019 chỉ có 5 đợt tuyết từ 25,4 mm trở lên,
kể cả đợt này; muốn kiểm cách hứa cho ngày bất thường cần thêm nhiều đợt khác.
% EV: F15 (NEW:5; RP1:24-35); N4B8-07 (rider_nowcast_replay.py:14-15, 113: explained = Σ clip(D − P0, 0, A − P0 − dispatch); RP1:29: phần vượt 6 s; với hồ sơ lịch sử phần vượt 0 nên hai cách cộng trùng nhau, N4B-05); K19; I11; B26 (NEW:9; WX:1-3: tệp NOAA GHCN-D trạm Central Park; WX17:1-4; ghép ngày cách nhau ≤ 2)

% -----------------------------------------------------------------------------------------------------
\section{Kiểm tra tính hợp lệ của kết quả}
\label{muc:c4-hople}

Trên 34 kịch bản của tập kiểm tra chính, không khách nào vượt trần: tổng dời lớn nhất là 29,2 giây với trần 30 giây và
118,9 giây với trần 120 giây. Không lượt quyết định nào bị ép, 0 trong 20.101 và 0 trong 20.435 lượt. Đây là phép kiểm cài
đặt: với sổ hợp lệ $0 \le L_r \le B$ ở lượt đang xét và giả thiết (N), Nhận xét~\ref{nx:giunguyen} đã cho trước các
số 0 này; kết luận xuyên chuyến cần giả thiết (N) đúng ở mọi lượt. Điều phối không trần cũng có 0 trong 20.760 lượt.
% EV: C20; N-07; X3-16 (0/20.101; 0/20.435; U 0/20.760); N-21; A22; chương 3, nx:giunguyen và đoạn sau nó (0 ≤ L_r ≤ B ở lượt đang xét; (N) đúng ở mọi lượt cho kết luận xuyên chuyến); R39 agent-numbers-v1.json:cohort_totals

Bộ đếm của bộ sinh được đọc ở điều phối không trần, luật C trần 30, 120 giây và luật hạn chót 30 giây. Ở cả bốn cấu
hình, không phương án nào bị bỏ vì giới hạn 100 phương án mỗi xe hay giới hạn số bước tìm cách chèn. Giới hạn 2 yêu cầu
mới mỗi lượt chỉ để một yêu cầu sang lượt sau ở mỗi cấu hình. Các số này không đo mọi giới hạn tính toán (mục~\ref{muc:gioihan}): giới hạn
của bước kiểm, của bộ giải, một bộ đếm khác của bộ sinh và các bộ đếm ở tập tắc đường cố định chưa được đọc.
% EV: N-25 (ADD1:119-122: 34 lần chạy mỗi cấu hình); N4B-11 (GEN:54-95, 236-246, 298-345; CandidateModels.cs:68-82, 100-118: điều kiện đầy đủ còn gồm OmittedRepairRequestCount, chưa đọc); X3-17; X3-03; X3-32; C4A-18 (tập tắc đường cố định: giới hạn chưa kiểm); K34 thu hẹp (N4B8-13)

Như mục~\ref{muc:doichung} đã nêu, mức dùng trần lớn nhất của tiêu chí 3 bằng 100\% ở mọi lượt có ghi tiêu chí này:
3.503, 3.506 và 3.504 lượt khi không trần, với trần 30 và 120 giây.
% EV: A15; N-25 (ADD1:119-121); A07; X3-07 (chương 3 mục doichung nêu hệ quả: tiêu chí 3 chỉ tác dụng qua tổng phần dời)

Các số đếm chính của tập kiểm tra chính được ít nhất hai chương trình đọc dữ liệu viết riêng tính lại và khớp nhau
(mục~\ref{muc:kiemtra}). Phép so gồm cả từng khách ở 7.159 khách được phục vụ của điều phối không trần và luật C trần
30, 120 giây; riêng số khách đi chung chỉ được một quy tắc đếm. Theo bản ghi khởi tạo, 8 lần chạy chiều tuyết dùng cùng
bản bộ chạy đã lưu và cùng cấu hình với các lần chạy luật C trần 30 giây của tập kiểm tra chính, kể cả luật phục hồi
không tệ hơn.
% EV: N4B-14 (ADD1:66-78); K41; C19; N-09 (đi chung: Codex không đếm lại); N4B8-05 (check C: 8 lần chạy chiều tuyết và 51 lần chạy C-30 của T3.8 cùng binarySha256 a631038b…, coreCommitSha 1ce15e92…, policyConfigurationHash a2f11a3a…; sha256 runner-tier3-v2/RideBound.Runner.dll khớp; wp4-tier3-v2.json:7 no-worse-than-reference-v1; R26 Q8)

Phép kiểm số đếm học máy chỉ đối chiếu số học của bảng kết quả; nó không kiểm lại quá trình học hay dự đoán của mô hình.
% EV: R37 §1.8; R39 agent-numbers-v1.py phần C; agent-numbers-v1.json:ml; ml_pilot_run.py:37-42,175-176

Chạy lại trên đúng mã nguồn đã dựng bộ chạy, cả 1.025 trong 1.025 kiểm thử tự động đạt trong một lần chạy đầy đủ ngày
13/10/2026, với kiểu xuống dòng LF và một tệp chưa ghi nhận trong cây làm việc (Phụ lục~\ref{phuluc:C}). Các phép kiểm
trên không cho biết mô phỏng giống thực tế đến đâu.
% EV: TESTRUN-45f8026-2026-10-12.md:18 (attempt 4: LF và tệp chưa ghi nhận, 1.025/1.025 trong một lần chạy ngày 13/10/2026); C:/RideBoundData/research/testrun-45f8026-eollf-dirty-fullrun-2026-10-13.log; R39 agent-numbers-v1.json:test_attempt4; I13; A28; I01; K41

% -----------------------------------------------------------------------------------------------------
\section{Thảo luận và hạn chế}
\label{muc:c4-thaoluan}

Với một hãng xe, các kết quả trên là các điểm trên một đường đánh đổi, không phải một khuyến nghị. So với điều phối
không trần, luật C trần 30 giây phục vụ ít hơn 96 lượt và có ít hơn 242 khách đi chung; đổi lại, phần dời do điều phối
của mỗi khách không quá 30 giây và số lượt trễ thấp hơn ở cả bốn ngưỡng đã đo, trên hai nhóm khách khác nhau
(Bảng~\ref{bang:c4-phanbotre}). Khóa luận không đo giá trị của các điều này với khách hay hãng xe.
% EV: C02; N-20; C20; N-24; C04; D20; I10; K38 (không kể 0 quyết định bị ép là phần luật C đổi được: X3-16, U cũng 0/20.760)

Một phép tính hòa vốn giả định cho thấy cách đọc các điểm này. Gọi $\lambda$ là chi phí của hãng xe cho mỗi lượt trễ quá
một phút và $\mu$ là biên lợi nhuận mỗi chuyến. Với 96 lượt phục vụ ít hơn và 127 lượt trễ ít hơn, luật C chỉ có lợi hơn
khi $\lambda > 0{,}756\,\mu$, tức $96/127$. Phép tính bỏ qua mọi yếu tố khác và so hai nhóm khách khác nhau; khóa luận
không có giá trị thật nào của $\lambda$ hay $\mu$.
% EV: C29 (R15:231-235; 96/127 = 0,756); K35; C28; C04; I10

Các hạn chế dưới đây chia theo bốn nhóm; mỗi ý nêu điều chưa làm và kết luận chịu ảnh hưởng.

\textbf{Mô hình mô phỏng.}
\begin{itemize}[leftmargin=1.2cm,itemsep=2pt,topsep=3pt]
\item Mọi kết quả đến từ FleetPy với một hệ số tắc đường chung cho cả thành phố, mức nền lấy lúc 00:00, nên mô phỏng gần
như không có yếu tố ngẫu nhiên. Các kết quả, kể cả mức trễ còn lại sau dự báo, không cho biết hiệu quả khi vận hành thật.
\item Thời gian dừng đón, trả bằng 0 giây (kiểm trên 3 lần chạy mẫu), so với 30 giây mỗi khách trong nghiên cứu điển hình
của FleetPy~\cite{engelhardt2026}, và đồ thị rút gọn chỉ có 96 nút. Phần dời do điều phối và cái giá 96 lượt vì vậy chỉ áp cho
thiết lập này.
\item Không có phản ứng của khách và không gán lại xe cho khách đã nhận, nên số lượt phục vụ không tính khách bỏ đi, và
chưa biết cái giá của trần khi được gán lại xe.
\end{itemize}
% EV: I01; B01; B06; B07; I02; F10; B05 (SERV:2-4: kiểm trên 3 lần chạy mẫu; FLEET:423-424); J18; B03; K48; B12; I07; X3-39

\textbf{Dữ liệu và thiết kế thực nghiệm.}
\begin{itemize}[leftmargin=1.2cm,itemsep=2pt,topsep=3pt]
\item 34 kịch bản của tập kiểm tra chính phụ thuộc nhau, thuộc 5 ngày, mỗi kịch bản chạy một lần; các phép đếm theo
kịch bản là mô tả, không phải kiểm định.
\item Chưa kiểm độ nhạy theo hạt giống của bộ giải, nên chưa biết các chênh lệch có giữ nguyên khi đổi hạt giống.
% EV: R37 §1.8/L-26; B09; B11; chương 3 v4-v3 muc:gioihan (hạt giống cố định 7, một lần chạy mỗi kịch bản)
\item Các kết quả dùng lại dữ liệu đã xem. Ngưỡng phân tích và cách chia tập được chọn sau khi xem tập buổi sáng và các
đợt chạy thử trước đó. Mức 120 giây được chọn sau khi xem kết quả của điều phối không trần và luật C trần 30 giây trên
chính tập kiểm tra chính, kể cả các phép tính thử ở 120 giây (mục~\ref{muc:c4-thamdo}). Chiều tuyết và các ngày kiểm học máy cũng đã được xem, nên mọi kết quả là thăm dò.
\item Tải thấp: 0,535 khách trên xe kể cả lúc xe rỗi khi không trần, so với khoảng 1,33 trong nghiên cứu điển hình của
FleetPy. Ở 49 kịch bản dựng từ dữ liệu thật, chuyến được chọn ngắn hơn mọi chuyến đủ điều kiện (trung vị thời gian đi thẳng 165
so với 196 giây). Kết quả không suy ra cho tải cao hay chuyến dài hơn.
\end{itemize}
% EV: B09; I05; B11; I06; C4A-11 (mức 120 giây, như phần A v5-v3 mục c4-thamdo); MLS:9-10; I11; B12 (LOAD:2-3; FLEET:447-449); J18; B13 (SEL6:2-4: 49/49 kịch bản); K42 (kể cả các phép tính thử ở 120 giây); chương 3 v4-v3 muc:gioihan (hạt giống7; chưa có kiểm độ nhạy)

\textbf{Tích hợp.}
\begin{itemize}[leftmargin=1.2cm,itemsep=2pt,topsep=3pt]
\item Bộ điều phối vẫn kiểm giờ đón bằng giờ không dự báo, nên 34 trong 1.471 khách được báo giờ đón muộn hơn giờ đón
muộn nhất. Các số có dự báo ở mục~\ref{muc:c4-dubao} và~\ref{muc:c4-dem} chưa tính việc sửa phép kiểm này.
\item Bản chạy ưu tiên chi phí khác bản lưu của ưu tiên ổn định. Bốn cặp đối chiếu hành vi và việc tái lập các
chỉ số nền không chứng minh tương đương toàn bộ hai bản. Kết quả ở mục~\ref{muc:c4-vaitro} vì vậy vẫn là phép so
thăm dò theo cấu hình, không loại trừ mọi khác biệt thời gian chạy.
\end{itemize}
% EV: F07 (RX:2; FR:70); I13; A28; R18 B24 (qua REG A28)

\textbf{Phạm vi kết luận.}
\begin{itemize}[leftmargin=1.2cm,itemsep=2pt,topsep=3pt]
\item Phép so hai thứ tự chọn phương án đã có kết quả ở trần 30 giây và không trần, nhưng chưa tạo thành
đường đánh đổi nhiều mức cho ưu tiên chi phí. Luật giới hạn từng lần chưa được cài đặt; các bản ghi chỉ vượt do
cộng dồn và bảy lịch sử sàng lọc không thay cho phép so hai luật trên diễn biến mô phỏng đầy đủ.
\item Luật hạn chót chỉ mượn thành phần hạn chót của Schulz và Pfeiffer~\cite{schulz2026}, với 30 và 60 giây trên tập
kiểm tra chính và 30--300 giây trên tập buổi sáng; mức mặc định của bài đó là 900 giây. Hạn chót nới kiểu Hörl và
Zwick~\cite{horl2022} (mục~\ref{muc:doichung}) trùng với cấu hình M$^{+}$ có luật phục hồi ``không tệ hơn'' ở trên,
nên mới được chạy ở 30 và 60 giây trên tập kiểm tra chính. Các phép so không nói về phương pháp của các tác giả đó.
\item Mức 30 giây không phải ngưỡng cảm nhận hay chuẩn ngành, và không có dữ liệu về mức khách sẵn lòng trả cho lời
hứa ổn định hơn; vì vậy không mức trần nào được chọn là tốt nhất.
\end{itemize}
% EV: H02; H12; K18; E09; K08; J10; J11; X3-28 (luật hạn chót 30–300 s ở tập buổi sáng; chương 3 Bảng luat); K27; D20; K52; K36
